\maketitle

\paperprovenance

\begin{abstract}
We prove a weighted summability criterion for the irrationality of
reciprocal Mersenne subseries.  For an integer base $b\ge2$ and a finite
nonempty prime set $\mathcal P$, let
$h(a)=\prod_{p\in\mathcal P}p^{v_p(a)}$ be the part of $a$ supported on
these primes.  Every infinite set $A\subseteq\Npos$ satisfying
\[
 \sum_{a\in A}\frac{h(a)}{a(b^{h(a)}-1)}<\infty
\]
has irrational $\sum_{a\in A}(b^a-1)^{-1}$.  The base-two condition gives
the conclusion at every integer base for every infinite subset of $A$,
and it permits an explicit example with divergent reciprocal sum.  The condition nevertheless excludes full
support and the set of all odd exponents.  The proof averages over multiples of a modulus divisible by a prescribed
finite part of the support.  A second average over dyadic observation
lengths controls the errors from incomplete periods.  Erd\H{o}s
already stated the weaker reciprocal-summable theorem without the
coprimality restriction; its complete ordinary proof is also printed here.
A finite estimate uniform in the moving modulus lets the weighted and
positive-cover arguments share one observation window.  This proves irrationality at every integer base for every infinite subset
of the union of two sets satisfying the respective criteria.

We also give finite-denominator and achievement-set calculations,
and formulate explicit conditional tests for the targets $1/2$
and $1/21$.  These tests do not decide either target or the universal
question.  The classical full-support, periodic-support and strict-tail
results used for comparison are attributed in the text.  Existing
work settles prime support at base~$2$ and squarefree support
at every power-of-two base; those cases are not claimed as new.
\end{abstract}

\section*{How to read this document}
Section~1 proves the support criteria and compares them with the
literature.  Sections~5--8 develop the finite arithmetic used in the
rational-membership tests.  Sections~9--12 distinguish the hypotheses
of conditional implications from the finite evidence for them.
Section~13 proves the arithmetic-sampling counterexample cited in the
short paper.  The global notation below is supplemented by definitions
where each specialised recurrence is used.

A sufficient condition need not be necessary; an equivalent condition
does not prove either side.  Counterexamples concern the specified
estimate, and a finite computation establishes only its stated range.

The short source tags distinguish the kind of evidence being cited.
\ev{Lean} identifies a cited formal-source declaration at the specified
revision and follows the status recorded in the supplied material.  The four
support endpoints cited in Part~I were replayed with their dependency closure
under Lean~4.29.1 at public commit \texttt{065e09523286}; their axiom audit
reports only \texttt{propext}, \texttt{Classical.choice}, and
\texttt{Quot.sound}.  Other supplied files retain their individually stated
status.  \ev{Cert} identifies a reported exact finite computation.
\ev{Math} identifies an ordinary argument, without asserting that no formal
counterpart exists.  \ev{Cited} refers to the literature, and \ev{Open}
marks an assumption or conclusion not proved here.  A statement may carry
more than one tag.  Parameter ranges are stated separately: a result at
specified values or below a bound is not a result at arbitrarily large
values.

The source links are for checking precise definitions and proof terms.
They are not additional hypotheses and do not replace an explanation of
the mathematics.  No independent verification of every other proof is
claimed.  Neither the universal problem nor either rational-target question
is resolved in this paper.
\subsection*{Notation used throughout}
For $A\subseteq\mathbb N_{>0}$ and an integer $b\ge2$, put
\[
 X_A(b)=\sum_{a\in A}\frac1{b^a-1},\qquad
 c_A(n)=\#\{a\in A:a\mid n\}\quad(n\ge1).
\]
At base two write $w_n=(2^n-1)^{-1}$, $R_n=\sum_{j>n}w_j$ and
$\mathcal A=\{X_A(2):A\subseteq\mathbb N_{>0}\}$.  For $x\ge0$ the
real greedy remainder is defined by $r_0(x)=x$ and
\[
 r_n(x)=r_{n-1}(x)-\begin{cases}
 w_n,&w_n\le r_{n-1}(x),\\
 0,&w_n>r_{n-1}(x).
 \end{cases}
\]
The greedy support $G_x$ consists of the selected indices; $G=G_{1/2}$.
For the half target index $1$ is skipped.  Thus a recurrence begun at index
$2$ with remainder $1/2$ is the same process, with its index shifted.

The finite quotient arguments use
\[
 q(M,d)=\left\lfloor\frac{2^M}{2^d-1}\right\rfloor,\qquad
 Q(D,M)=\sum_{d\in D}q(M,d).
\]
For $M\ge k$ their nonnegative binary remainder and signed next-step
expression are
\[
 \begin{aligned}
 S(D,k,M)&=\max\{2^{M-k}-Q(D,M)-1,0\},\\
 H(D,k,n)&=2S(D,k,n-1)+1-c_D(n)\quad(n>k).
 \end{aligned}
\]
The maximum matters: subtraction in the corresponding natural-number
source definition is truncated at zero.  Statements that remove the
maximum establish nonnegativity first.

The integer carry for a support omitting $1$ is
\[
 \operatorname{ihc}(A,N)=2^N-
 \sum_{j=2}^{N+1}2^{N+1-j}c_A(j),\qquad
 C_A(N)=\operatorname{ihc}(A,N)-1.
\]
Some finite-depth statements use $K_A(m)=\operatorname{ihc}(A,m-1)$.
The recurring bound $B(m)=2\lfloor\sqrt m\rfloor+4$ is integral; bounds
written with $\sqrt m$ rather than its floor are real-valued bounds.
Quantities specific to one argument are defined in that argument.
Some formal sources allow a support to be a subset of $\N$, including zero.
In all reciprocal sums in this paper only positive exponents contribute:
for such a support, $X_A(b)$ means $X_{A\cap\Npos}(b)$.  The expression
$1/(b^0-1)$ is never evaluated.

\clearpage
\tableofcontents
\clearpage
% ; - part a257_front ; -
% SPDX-FileCopyrightText: 2026 Will Cook
% SPDX-License-Identifier: CC-BY-4.0
%
% Front matter for the Erdos #257 reasoning-surface paper: the problem, the
% barrier analysis, the surviving routes, and the reading contract. Everything
% before the premise catalogue.

\clearpage
\part*{I. Support theorems and their proofs}
\addcontentsline{toc}{part}{I. Support theorems and their proofs}
\section{Support criteria and their proofs}
\label{sec:257-problem}

\subsection{Reciprocal-summable supports}

For every infinite set $A$ of positive integers with
$\sum_{a\in A}1/a<\infty$, the series
\[
 X_A(b)=\sum_{a\in A}\frac1{b^a-1}
\]
is irrational at every integer base $b\ge2$. The proof below is ordinary
mathematics and is the proof of the short paper's
Theorem~1.2. Erd\H{o}s's
\href{https://users.renyi.hu/~p_erdos/1969-09.pdf}{1968 paper, p.~222}
proves the pairwise-coprime case at every integer base and states the removal
of coprimality without printing the details. The proof below gives the stated extension by averaging; no comparison
with the omitted proof is possible.  The linked pairwise-coprime theorem
retains that extra hypothesis. The same
paper says that $\sum1/n_i<\infty$ could be replaced by a weaker but more
complicated condition (p.~222), and suggests that for pairwise coprime
supports Brun's method could probably replace it by
$\sum_{n_i<x}1/n_i=o(\log\log x)$ (p.~226).
Theorem~\ref{thm:257-weighted} below proves irrationality under one explicit
weaker condition without coprimality; no identification with his suggestions
is asserted.

The mechanism is to make positive displacements arbitrarily small. For $N>0$,
write $N=q_a a+r_a$, $0\le r_a<a$. The identity
\[
 \frac{b^N-1}{b^a-1}
 =b^{r_a}\sum_{j=0}^{q_a-1}b^{ja}
  +\frac{b^{r_a}-1}{b^a-1}
\]
shows that
\[
 \Delta_{b,A}(N):=\sum_{a\in A}\frac{b^{N\bmod a}-1}{b^a-1}
 =(b^N-1)X_A(b)-J_N,\qquad J_N\in\mathbb Z.
\]
The integer sum is finite, since $q_a=0$ for $a>N$. Every summand of
$\Delta_{b,A}(N)$ is nonnegative and an exponent $a>N$ makes it positive.
Thus $X_A(b)=p/q$ would imply $\Delta_{b,A}(N)\ge1/q$ for every $N>0$.

Put $w_{b,d}(N)=b^{N\bmod d}/(b^d-1)$. For a fixed positive modulus $Q$
and $g=(Q,d)$, the orbit of $Qm\bmod d$ has length $d/g$. Its geometric sum
gives
\[
 \lim_{X\to\infty}\frac1X\sum_{m=1}^Xw_{b,d}(Qm)
 =M_{Q,b}(d):=\frac{g}{d(b^g-1)}\le\frac1d.
\]
The important interchange is with the infinite support. Expanding the atom
as $\sum_{r\ge1}b^{-r}{\bf1}_{d\mid Qm+r}$ and counting positive multiples
of $d$ gives the bound
\[
 \frac1X\sum_{m=1}^Xw_{b,d}(Qm)
 \le\sum_{r\ge1}b^{-r}\frac{Q+r}{d}
 \le\frac{Q+2}{d}.
\]
Here $Q$ is fixed. Dominated convergence therefore identifies the limiting
mean of $\sum_{d\in A}w_{b,d}(Qm)$ with $\sum_{d\in A}M_{Q,b}(d)$.

Now choose $Q_t=\operatorname{lcm}(1,\ldots,t)$. For each fixed $d$,
eventually $d\mid Q_t$, so $M_{Q_t,b}(d)=(b^d-1)^{-1}$. The bound
$M_{Q_t,b}(d)\le1/d$ permits a second dominated-convergence argument:
\[
 \sum_{d\in A}M_{Q_t,b}(d)\longrightarrow X_A(b).
\]
Consequently the limiting mean of $\Delta_{b,A}(Q_tm)$ tends to zero.
Given $\varepsilon>0$, first choose $t$ and then a sufficiently long finite
average with mean below $\varepsilon$; one term of that average is below
$\varepsilon$. Its strict positivity contradicts the rational lattice.
This proves the assertion.  The exact all-base statement, with hypotheses
$b\ge2$, infinitude of $A$, and summability of $a\mapsto\mathbf1_A(a)/a$,
is kernel-checked as
\href{https://github.com/wcook04/plectis-erdos/blob/065e09523286894dfb57ba205e69666843817009/lean/Erdos249257/AllBaseReciprocalSupportIrrationality.lean#L395}{\texttt{irrational\_erdosSupportSeries\_of\_summable\_reciprocal}}. \ev{Lean}

The order of the two limits matters: the observation length increases with
$Q_t$ fixed before $t$ increases. No independence of different divisibility
events is assumed. The mechanism also has a definite limit. At base two and
full support, the contribution of exponents $a>N$ alone gives
$\Delta_{2,\mathbb N_{>0}}(N)>1-2^{-N}\ge1/2$. Full-support irrationality
therefore uses a different argument; its checked form appears below.

\subsection{A weighted condition on the support}
The reciprocal-summable case is a useful model, but powers of a fixed prime
can make an exponent cheaper than its reciprocal suggests.  For example,
with $a=2^km$ and $m$ odd, the weight below becomes
$1/[m(b^{2^k}-1)]$.  The exponential denominator permits a large collection
of odd factors at large $k$.


\begin{thm}[A weighted summability criterion]
\label{thm:257-weighted}
Let $b\ge2$ be an integer, let $\mathcal P$ be a finite nonempty set of
primes, and let $A\subseteq\mathbb N_{>0}$ be infinite.  Write
\[
 h(a)=\prod_{p\in\mathcal P}p^{v_p(a)},
\]
where $v_p(a)$ is the exponent of $p$ in $a$; thus $h(a)$ is the
$\mathcal P$-part of $a$.  If
\begin{equation}\label{eq:257-weighted-mass}
 W_{b,\mathcal P}(A):=
 \sum_{a\in A}\frac{h(a)}{a\bigl(b^{h(a)}-1\bigr)}<\infty,
\end{equation}
then
\[
 X_A(b)=\sum_{a\in A}\frac1{b^a-1}
\]
is irrational.  More precisely, for
\[
 \Delta_{b,A}(m)=
 \sum_{a\in A}\frac{b^{m\bmod a}-1}{b^a-1},
\]
for every $\varepsilon>0$ and every $N$ there is an $m\ge N$ with
$0<\Delta_{b,A}(m)<\varepsilon$.
\par\ev{Math}\scale{uniform}\coord{divisibility-weighted}
\end{thm}
\leannote{Lean: \lproof{ErdosProblems/Erdos257/PaperCompleteR8/WeightedReturn.lean}{120}{divisibilityWeightedClaim}, \lproof{ErdosProblems/Erdos257/PaperCompleteR8/WeightedReturn.lean}{100}{weighted_displacement_cofinal_close_return}.}

\paragraph{Examples and limits of the hypothesis.}
Squares, powers of two and factorials satisfy reciprocal summability and
therefore satisfy the base-two weighted condition.  Full support does not:
for any fixed finite prime set $P$, its subset coprime to $\prod_{p\in P}p$
has weight $1/[a(b-1)]$ and divergent reciprocal sum.  The set of all odd
exponents also fails, by restricting to integers coprime to
$2\prod_{p\in P}p$.  Inclusion--exclusion in the harmonic sum proves both
divergences.  This is a limitation of the sufficient condition, not an
argument for rationality; full support is already known to be irrational.

Writing $a=hm$ with $h$ its $\mathcal P$-part gives the exact decomposition
\[
 W_{b,\mathcal P}(A)=\sum_{h\in\mathcal H_{\mathcal P}}\frac1{b^h-1}
       \sum_{\substack{m\ge1,\ hm\in A\\(m,\prod_{p\in\mathcal P}p)=1}}\frac1m,
\]
where $\mathcal H_{\mathcal P}$ consists of $1$ and the positive
integers with no prime factor outside $\mathcal P$.  Each inner
reciprocal sum must be finite; the displayed outer sum imposes the
additional summability across different $\mathcal P$-parts.
For example, at base two with $\mathcal P=\{2\}$, choose in each layer
$2^k$, $k\ge1$, a finite set of odd factors whose reciprocal
sum is at least $2^{2^k}$.  Every inner sum is finite, but the
outer weighted sum diverges.  Such finite sets exist by divergence
of the harmonic sum over odd integers.  Prime support also fails
this criterion for every finite $\mathcal P$: its contribution
from primes outside $\mathcal P$ is a constant multiple of their
divergent reciprocal sum.  These failures concern the sufficient
hypothesis, not rationality of the resulting values.
The following construction shows, nonetheless, that this condition
is strictly weaker than reciprocal summability.
At base two take
\[
 A_\star=\{2^k m:k\ge1,\ m\text{ odd},\ m\le2^{2^k}\}.
\]
Write $H_N=\sum_{m=1}^N1/m$ and $N_k=2^{2^k}$.  The reciprocal mass
of the $k$th layer is
\[
 \rho_k=2^{-k}\!\sum_{\substack{m\le N_k\\m\text{ odd}}}\frac1m
 =2^{-k}\bigl(H_{N_k}-\tfrac12H_{N_k/2}\bigr)
 \longrightarrow\frac{\log2}{2}.
\]
Thus $\sum_{a\in A_\star}1/a$ diverges.  For $\mathcal P=\{2\}$,
the weighted contribution of the same layer is
\[
 \frac{2^k}{2^{2^k}-1}\rho_k.
\]
Since $H_N\le1+\log N$, these terms are at most
$(1+2^k\log2)2^{1-2^k}$, which is summable.  The theorem therefore proves
irrationality at every integer base for every infinite subset of this set,
although its reciprocal sum diverges.


\begin{proof}
We first record the rational obstruction and three finite estimates used to
defeat it.  If $m=qa+r$, $0\le r<a$, then
\[
 \frac{b^m-1}{b^a-1}-\frac{b^r-1}{b^a-1}
   =\sum_{j=1}^{q}b^{m-ja}.
\]
Consequently
\begin{equation}\label{eq:257-weighted-lattice}
 \Delta_{b,A}(m)=(b^m-1)X_A(b)-J_{b,A}(m),
 \qquad
 J_{b,A}(m)=
 \sum_{\substack{a\in A\\a\le m}}
 \sum_{j=1}^{\lfloor m/a\rfloor}b^{m-ja}\in\mathbb Z.
\end{equation}
The integer sum is finite.  For $m\ge1$, moreover,
$\Delta_{b,A}(m)>0$: since $A$ is infinite, some $a\in A$ exceeds $m$,
and its summand is positive.  At $m=0$ the displacement is zero.  Thus, if
$X_A(b)=u/v$ with $v>0$, equation~\eqref{eq:257-weighted-lattice} gives
$v\Delta_{b,A}(m)\in\mathbb Z$, whence
\begin{equation}\label{eq:257-weighted-gap}
 \Delta_{b,A}(m)\ge \frac1v\qquad(m\ge1).
\end{equation}
It is therefore enough to construct arbitrarily small positive displacements.

Put
\[
 \phi_a(m)=\frac{b^{m\bmod a}}{b^a-1},
 \qquad
 d_a(m)=\phi_a(m)-\frac1{b^a-1}.
\]
Then $d_a(m)\ge0$ and $\Delta_{b,A}(m)=\sum_{a\in A}d_a(m)$.
For $Q,a,T\ge1$, put $g=(Q,a)$.  The residues $tQ\bmod a$ have period
$a/g$ and visit $0,g,\ldots,a-g$ once in a period, so
\[
 \sum_{t=1}^{a/g}\phi_a(tQ)
 =\frac{1+b^g+\cdots+b^{a-g}}{b^a-1}
 =\frac1{b^g-1}.
\]
Complete periods and at most one incomplete period therefore give
\begin{equation}\label{eq:257-weighted-gcd-average}
 \frac1T\sum_{t=1}^{T}d_a(tQ)
 \le \frac{g}{a(b^g-1)}+\frac1{T(b^g-1)}.
\end{equation}
We shall also use that $n/(b^n-1)$ decreases for positive integers $n$;
indeed
\[
 n(b^{n+1}-1)-(n+1)(b^n-1)
 =b^n\bigl(n(b-1)-1\bigr)+1>0.
\]

The part of the support beyond the observation range has a uniform bound.  If
$Y=QT$ and $a>Y\ge tQ$, then $tQ\bmod a=tQ$ and
$d_a(tQ)\le2\,2^{tQ-a}$.  Hence
\[
 \sum_{a>Y}d_a(tQ)\le2\,2^{tQ-Y},
 \qquad
 \sum_{t=1}^{T}2^{tQ-Y}
 =\sum_{j=0}^{T-1}2^{-jQ}\le2,
\]
and therefore
\begin{equation}\label{eq:257-weighted-outer}
 \frac1T\sum_{t=1}^{T}
 \sum_{\substack{a\in A\\a>QT}}d_a(tQ)\le\frac4T.
\end{equation}

Finally, let $\alpha_a\ge0$ with $\sum_a\alpha_a/a<\infty$.  For $Q,M\ge1$,
finite rearrangement and a geometric sum give
\begin{equation}\label{eq:257-weighted-dyadic-average}
 \sum_{j=M}^{2M-1}\frac1{2^j}
       \sum_{a\le Q2^j}\alpha_a
 \le 2Q\sum_{a\ge1}\frac{\alpha_a}{a}.
\end{equation}
Indeed, for each fixed $a$ the admissible indices satisfy $2^j\ge a/Q$ and
\[
 \sum_{\substack{M\le j<2M\\a\le Q2^j}}2^{-j}\le\frac{2Q}{a}.
\]
Only finitely many terms occur on the left, so no infinite interchange is
hidden here.

We now choose the modulus and the observation lengths.  Fix $\varepsilon>0$ and a lower bound
$N$.  Select a finite $F\subseteq A$ such that
\begin{equation}\label{eq:257-weighted-tail-choice}
 \sum_{a\in A\mathbin{\backslash} F}
 \frac{h(a)}{a(b^{h(a)}-1)}<\varepsilon,
\end{equation}
and put $L=\operatorname{lcm}(F)$, with $L=1$ if $F$ is empty.  Write
$p_*=\max\mathcal P$ and $r=|\mathcal P|$.  For a large integer $H\ge2p_*$,
set
\[
 Q_0(H)=\prod_{p\in\mathcal P}p^{\lfloor\log_pH\rfloor},
 \qquad Q=LQ_0(H),
 \qquad G=\left\lfloor\frac H{p_*}\right\rfloor.
\]
Then $Q\le LH^r$, $G\ge1$, and every $a\in F$ divides $Q$, so its
$d_a(tQ)$ contribution vanishes.

Consider first $a\in A\mathbin{\backslash} F$ with $a\le QT$ and $h(a)\le H$.
Every prime-power factor of $h(a)$ occurs in $Q_0(H)$, hence
$h(a)\mid Q_0(H)$ and $(Q,a)\ge h(a)$.  Applying
\eqref{eq:257-weighted-gcd-average} and the monotonicity above, the total
contribution of these $a$ is at most
\begin{equation}\label{eq:257-weighted-small-part}
 \varepsilon+
 \frac1T\sum_{\substack{a\in A, a\le QT\\h(a)\le H}}
 \frac1{b^{h(a)}-1}.
\end{equation}
The first term is bounded using~\eqref{eq:257-weighted-tail-choice}.

For $h(a)>H$, one has $(Q,a)\ge G$.  If every $\mathcal P$-prime-power
component of $h(a)$ is at most $H$, then $h(a)\mid Q_0(H)$ and the gcd is
larger than $H$.  Otherwise some $p^{v_p(a)}>H$, and the gcd contains
$p^{\lfloor\log_pH\rfloor}>H/p\ge H/p_*$.  Thus
\eqref{eq:257-weighted-gcd-average}, enlarged from the support to all
$a\le QT$, bounds the large-$h(a)$ part by
\begin{equation}\label{eq:257-weighted-large-part}
 \frac{G(1+\log(QT))+Q}{b^G-1};
\end{equation}
we used $\sum_{a\le QT}a^{-1}\le1+\log(QT)$ and
$\#\{a\le QT\}/T=Q$.

Let
\[
 \mathcal M(Q,T)=\frac1T\sum_{t=1}^{T}\Delta_{b,A}(tQ).
\]
Combining \eqref{eq:257-weighted-outer},
\eqref{eq:257-weighted-small-part}, and
\eqref{eq:257-weighted-large-part} gives
\begin{equation}\label{eq:257-weighted-master-average}
 \mathcal M(Q,T)\le
 \varepsilon+
 \frac1T\sum_{\substack{a\in A, a\le QT\\h(a)\le H}}
       \frac1{b^{h(a)}-1}
 +\frac{G(1+\log(QT))+Q}{b^G-1}+\frac4T.
\end{equation}
After the second average, the competing errors will be of orders
$Q/M$ and $GM/b^G$.  With $L$ fixed, $Q$ grows polynomially in $H$
and $G$ grows linearly, so an intermediate exponential scale makes
both errors tend to zero.  Choose
\[
 M=\left\lfloor b^{G/2}\right\rfloor
\]
and average~\eqref{eq:257-weighted-master-average} over
$T=2^j$, $M\le j<2M$.  Apply
\eqref{eq:257-weighted-dyadic-average} with
$\alpha_a={\bf1}_A(a)/(b^{h(a)}-1)$.  Its hypothesis holds because
\[
 \sum_a\frac{\alpha_a}{a}
 =\sum_{a\in A}\frac1{a(b^{h(a)}-1)}
 \le W_{b,\mathcal P}(A).
\]
We obtain the explicit finite diagonal estimate
\begin{equation}\label{eq:257-weighted-final-diagonal}
 \frac1M\sum_{j=M}^{2M-1}\mathcal M(Q,2^j)
 \le \varepsilon+\frac{2QW_{b,\mathcal P}(A)}M
 +\frac{G(1+\log Q+2M\log2)+Q}{b^G-1}
 +4\,2^{-M}.
\end{equation}

As $H\to\infty$,
\[
 Q\le LH^r,\qquad G=H/p_*+O(1),\qquad M\asymp b^{G/2}.
\]
Consequently $Q/M\to0$, $GM/b^G\to0$,
$(G\log Q+Q)/b^G\to0$, and $2^{-M}\to0$.  For sufficiently large $H$,
we may also arrange $Q\ge N$ and make the right side of
\eqref{eq:257-weighted-final-diagonal} smaller than $2\varepsilon$.
The left side is a finite average of positive numbers
$\Delta_{b,A}(tQ)$.  One of them therefore satisfies
\[
 0<\Delta_{b,A}(tQ)<2\varepsilon,
 \qquad tQ\ge Q\ge N.
\]
The construction is valid for every positive $\varepsilon$; applying it with
$\varepsilon/2$ proves the stated cofinal bound with $\Delta_{b,A}(m)<\varepsilon$.
This contradicts~\eqref{eq:257-weighted-gap} if $X_A(b)$ is rational.
\end{proof}

Since $b^{h(a)}-1\ge2^{h(a)}-1$ for every $b\ge2$, one has
$W_{b,\mathcal P}(A)\le W_{2,\mathcal P}(A)$.  Thus the base-two weighted
hypothesis gives the conclusion simultaneously at every integer base.  The
hypothesis is also hereditary under passage to infinite subsets.

\paragraph{Averaging antecedents.}
Both averaging proofs in this section end by choosing a term no larger than a
finite mean.
\href{https://danielduverney.fr/documents/theorie-des-nombres/DuverneyTachiya190522.pdf}{Duverney
and Tachiya} choose their index in the same way in their refinement of the
Chowla--Erd\H{o}s method: they average a local coefficient mass over an
arithmetic progression built by the Chinese remainder theorem, and a term
attaining the minimum supplies the bounds that control the tail beyond it
(Section~2, (2.3)--(2.9), pp.~5--6 of the author preprint).
\href{https://arxiv.org/abs/2601.20743}{Kaneko, Suzuki and Tachiya} measure
the average decay of scaled tails by $R_c(q,x,z)$, their~(1.7) on p.~3 of
arXiv v1.  Under a counting condition on the support and a bound on $R_c$,
most scaled tails in a range are small, and arbitrarily small nonzero scaled
tails give irrationality (Lemmas~1 and~2, pp.~6--8).  The divisor-incidence
coefficients $\#\{a\in A:a\mid n\}$ are positive on every multiple of
$\min A$.  For any splitting into two sequences the two coefficient supports
therefore cover a set of positive lower density, which is incompatible with the
support-counting conditions
\cite[Theorem~1(iii), p.~3, and Theorem~3(iv), p.~5]{kanekosuzukitachiya}
requiring both counts to be $o(x_n/z_n)$ with $z_n\ge1$.  Their remote-tail
quantity $R_c(q,x,z)$ sums only offsets $j\ge z$, whereas the full scaled tail
\[
 \sum_{r\ge1}\frac{\#\{a\in A:a\mid m+r\}}{b^r}
 =\sum_{a\in A}\phi_a(m)=X_A(b)+\Delta_{b,A}(m)
\]
is at least $X_A(b)$, which gives no lower bound for $R_c$.  The proofs above
average the displacement instead,
along multiples of a modulus that freezes a finite part of $A$; the weighted
proof adds the dyadic average~\eqref{eq:257-weighted-dyadic-average}.
\ev{Cited}

\paragraph{Evidence boundary.}
The exact interface of Theorem~\ref{thm:257-weighted}, including its
fixed-base conclusion and hereditary all-base clause, is kernel-checked as
\href{https://github.com/wcook04/plectis-erdos/blob/065e09523286894dfb57ba205e69666843817009/lean/ErdosProblems/Erdos257/PaperCompleteR8/WeightedReturn.lean#L99}{\texttt{divisibilityWeightedClaim}}.  Its public dependency closure was
replayed under Lean~4.29.1.  This does not cover all infinite supports,
supply a quantitative irrationality measure, or assert theorem priority.

\subsection*{Formalisation coverage and remaining dependencies}
\label{sec:coverage}
\papersectiontarget{coverage}

Every theorem, lemma, proposition and corollary of this record has a Lean
statement of the same assertion, with the same hypotheses, checked by the Lean
kernel using only the axioms \texttt{propext}, \texttt{Classical.choice} and
\texttt{Quot.sound}, in the development at revision \texttt{181078b6b009}.
The snapshot link in the paragraph above identifies an earlier published
checkpoint; the coverage statement here is about the development revision just
named.  Statements the record takes from the literature are attributed where
they appear, and the \ev{Cited} tag marks a clause that invokes an external
result.  Theorem~\ref{thm:periodic-support} is attributed in this way to the
periodic-weight irrationality theorem of Luca and
Tachiya~\cite[Theorem~A, p.~139]{lucatachiya2017}.  No formal proof of their
theorem is given here; Theorem~\ref{thm:periodic-support} and the support
statements following from it each carry a separate formal proof of their own
assertion.

% BEGIN GENERATED CONCORDANCE (generated from the coverage ledger; do not edit by hand)
\paragraph{Concordance of statements and Lean declarations.}
Each result of this record that has a kernel-checked Lean statement of the same assertion is listed
below with the declarations that jointly state it.  Each name links to the source revision in its URL.  Where the Lean statement is stronger than the printed one and implies it by
an immediate specialisation, the entry says so.  An entry marked \emph{compared} was also checked
independently: a restatement of the same declarations against Mathlib alone, together with its
proof, was verified by Comparator (\texttt{leanprover/comparator}) in a clean continuous
integration environment, in the run named by its number.  Comparator trusts the restated
statement, so the correspondence between the printed result and that statement is the one this
concordance records.

\begingroup\small\raggedright
\par\noindent\hangindent=1.5em Theorem~\ref{thm:257-weighted} (the Lean statement is stronger): \href{https://github.com/wcook04/plectis-erdos/blob/a25cb360bef8dd818dde14b5fb752244304af354/lean/ErdosProblems/Erdos257/PaperCompleteR8/WeightedReturn.lean\#L120}{\nolinkurl{divisibilityWeightedClaim}}, \href{https://github.com/wcook04/plectis-erdos/blob/a25cb360bef8dd818dde14b5fb752244304af354/lean/ErdosProblems/Erdos257/PaperCompleteR8/WeightedReturn.lean\#L100}{\nolinkurl{weighted_displacement_cofinal_close_return}}.  \emph{Compared}, run \href{https://github.com/wcook04/plectis-erdos-lean/actions/runs/35544127144}{\texttt{35544127144}}.
\par\noindent\hangindent=1.5em Theorem~\ref{thm:257-variable-fractional-cover} (the Lean statement is stronger): \href{https://github.com/wcook04/plectis-erdos/blob/a25cb360bef8dd818dde14b5fb752244304af354/lean/ErdosProblems/Erdos257/PaperCompleteR8/PositiveCoverReturn.lean\#L241}{\nolinkurl{strengthenedPositiveCoverClaim}}.  \emph{Compared}, run \href{https://github.com/wcook04/plectis-erdos-lean/actions/runs/35544127144}{\texttt{35544127144}}.
\par\noindent\hangindent=1.5em Theorem~\ref{thm:257-mixed-supports} (the Lean statement is stronger): \href{https://github.com/wcook04/plectis-erdos/blob/a25cb360bef8dd818dde14b5fb752244304af354/lean/ErdosProblems/Erdos257/PaperCompleteR8/WeightedReturn.lean\#L126}{\nolinkurl{mixedSupportClaim}}, \href{https://github.com/wcook04/plectis-erdos/blob/a25cb360bef8dd818dde14b5fb752244304af354/lean/ErdosProblems/Erdos257/PaperCompleteR8/ArbitraryWeightMixedClaim.lean\#L101}{\nolinkurl{arbitraryWeightMixedSupport_allBase_hereditary}}.  \emph{Compared}, runs \href{https://github.com/wcook04/plectis-erdos-lean/actions/runs/35544127144}{\texttt{35544127144}}, \href{https://github.com/wcook04/plectis-erdos-lean/actions/runs/35674034595}{\texttt{35674034595}}.
\par\noindent\hangindent=1.5em Theorem~\ref{thm:geometry}: \href{https://github.com/wcook04/plectis-erdos/blob/a25cb360bef8dd818dde14b5fb752244304af354/lean/ErdosProblems/Erdos257/PaperCompleteR20/AchievementGeometry.lean\#L40}{\nolinkurl{paper_achievement_geometry}}.  \emph{Compared}, run \href{https://github.com/wcook04/plectis-erdos-lean/actions/runs/35643815458}{\texttt{35643815458}}.
\par\noindent\hangindent=1.5em Theorem~\ref{thm:supported-dichotomy} (the Lean statement is stronger): \href{https://github.com/wcook04/plectis-erdos/blob/a25cb360bef8dd818dde14b5fb752244304af354/lean/ErdosProblems/Erdos257/MersenneSubseriesRigidity.lean\#L397}{\nolinkurl{volume_supportedMersenneAchievementSet_dichotomy}}, \href{https://github.com/wcook04/plectis-erdos/blob/a25cb360bef8dd818dde14b5fb752244304af354/lean/ErdosProblems/Erdos257/MersenneSubseriesRigidity.lean\#L368}{\nolinkurl{volume_supportedMersenneAchievementSet_eq_zero_of_compl_infinite}}, \href{https://github.com/wcook04/plectis-erdos/blob/a25cb360bef8dd818dde14b5fb752244304af354/lean/ErdosProblems/Erdos257/MersenneSubseriesRigidity.lean\#L167}{\nolinkurl{perfect_supportedMersenneAchievementSet}}, \href{https://github.com/wcook04/plectis-erdos/blob/a25cb360bef8dd818dde14b5fb752244304af354/lean/ErdosProblems/Erdos257/MersenneSubseriesRigidity.lean\#L54}{\nolinkurl{supportedMersenneDigitValue_injective}}.  \emph{Compared}, run \href{https://github.com/wcook04/plectis-erdos-lean/actions/runs/35643815458}{\texttt{35643815458}}.
\par\noindent\hangindent=1.5em Theorem~\ref{thm:greedy-survival}: \href{https://github.com/wcook04/plectis-erdos/blob/a25cb360bef8dd818dde14b5fb752244304af354/lean/ErdosProblems/Erdos257/PaperCompleteR20/GeneralTargetGap.lean\#L183}{\nolinkurl{paper_greedy_survival}}.  \emph{Compared}, run \href{https://github.com/wcook04/plectis-erdos-lean/actions/runs/35674034595}{\texttt{35674034595}}.
\par\noindent\hangindent=1.5em Proposition~\ref{prop:canon} (the Lean statement is stronger): \href{https://github.com/wcook04/plectis-erdos/blob/a25cb360bef8dd818dde14b5fb752244304af354/lean/Erdos249257/HalfCutLocator.lean\#L442}{\nolinkurl{half_agrees_greedy}}, \href{https://github.com/wcook04/plectis-erdos/blob/a25cb360bef8dd818dde14b5fb752244304af354/lean/Erdos249257/BooleanMobiusCriticalCapacityCofinal.lean\#L50}{\nolinkurl{eq_halfGreedyPrefixSupport_of_critical_crossing}}, \href{https://github.com/wcook04/plectis-erdos/blob/a25cb360bef8dd818dde14b5fb752244304af354/lean/Erdos249257/BooleanMobiusGreedyReduction.lean\#L918}{\nolinkurl{remainder_lt_gap_iff_eq_integerGreedyBits}}.  \emph{Compared}, run \href{https://github.com/wcook04/plectis-erdos-lean/actions/runs/35935225572}{\texttt{35935225572}}.
\par\noindent\hangindent=1.5em Lemma~\ref{lem:collapse-mech} (the Lean statement is stronger): \href{https://github.com/wcook04/plectis-erdos/blob/a25cb360bef8dd818dde14b5fb752244304af354/lean/Erdos249257/HalfCarryReachability.lean\#L871}{\nolinkurl{integerHalfCarry_eq_scaled_residual_add_tail}}, \href{https://github.com/wcook04/plectis-erdos/blob/a25cb360bef8dd818dde14b5fb752244304af354/lean/Erdos249257/GenericTailOrbitRigidity.lean\#L78}{\nolinkurl{binaryCoeffTail_nonneg}}, \href{https://github.com/wcook04/plectis-erdos/blob/a25cb360bef8dd818dde14b5fb752244304af354/lean/Erdos249257/BooleanMobiusCarry.lean\#L290}{\nolinkurl{binaryCoeffTail_supportCoeff_le_two_sqrt_add_four}}.  \emph{Compared}, runs \href{https://github.com/wcook04/plectis-erdos-lean/actions/runs/35624228171}{\texttt{35624228171}}, \href{https://github.com/wcook04/plectis-erdos-lean/actions/runs/35643815458}{\texttt{35643815458}}, \href{https://github.com/wcook04/plectis-erdos-lean/actions/runs/35674034595}{\texttt{35674034595}}.
\par\noindent\hangindent=1.5em Proposition~\ref{prop:collapse}: \href{https://github.com/wcook04/plectis-erdos/blob/a25cb360bef8dd818dde14b5fb752244304af354/lean/ErdosProblems/Erdos257/PaperCompleteR20/CofinalCarryCollapse.lean\#L46}{\nolinkurl{half_of_cofinal_absolute_carry}}, \href{https://github.com/wcook04/plectis-erdos/blob/a25cb360bef8dd818dde14b5fb752244304af354/lean/ErdosProblems/Erdos257/PaperCompleteR20/CofinalCarryCollapse.lean\#L91}{\nolinkurl{greedy_half_of_cofinal_upper_carry}}.  \emph{Compared}, run \href{https://github.com/wcook04/plectis-erdos-lean/actions/runs/35674034595}{\texttt{35674034595}}.
\par\noindent\hangindent=1.5em Lemma~\ref{lem:sqrt-witness}: \href{https://github.com/wcook04/plectis-erdos/blob/a25cb360bef8dd818dde14b5fb752244304af354/lean/ErdosProblems/Erdos257/PaperCompleteR20/CarryCollapseCorrespondence.lean\#L8}{\nolinkurl{square_depth_witness}}.  \emph{Compared}, run \href{https://github.com/wcook04/plectis-erdos-lean/actions/runs/35674034595}{\texttt{35674034595}}.
\par\noindent\hangindent=1.5em Proposition~\ref{prop:collapsed-list}: \href{https://github.com/wcook04/plectis-erdos/blob/a25cb360bef8dd818dde14b5fb752244304af354/lean/ErdosProblems/Erdos257/PaperCompleteR20/SixMembershipConditions.lean\#L123}{\nolinkurl{six_membership_conditions}}.  \emph{Compared}, run \href{https://github.com/wcook04/plectis-erdos-lean/actions/runs/35674034595}{\texttt{35674034595}}.
\par\noindent\hangindent=1.5em Proposition~\ref{prop:local-void} (the Lean statement is stronger): \href{https://github.com/wcook04/plectis-erdos/blob/a25cb360bef8dd818dde14b5fb752244304af354/lean/Erdos249257/GenericTailOrbitRigidity.lean\#L307}{\nolinkurl{affineBinaryOrbit_mod_twoPow_eq}}, \href{https://github.com/wcook04/plectis-erdos/blob/a25cb360bef8dd818dde14b5fb752244304af354/lean/Erdos249257/GenericTailOrbitRigidity.lean\#L200}{\nolinkurl{balancedPulse_weighted_pair}}, \href{https://github.com/wcook04/plectis-erdos/blob/a25cb360bef8dd818dde14b5fb752244304af354/lean/Erdos249257/GenericTailOrbitRigidity.lean\#L209}{\nolinkurl{balancedPulse_endpoint_fanout}}, \href{https://github.com/wcook04/plectis-erdos/blob/a25cb360bef8dd818dde14b5fb752244304af354/lean/Erdos249257/GenericTailOrbitRigidity.lean\#L228}{\nolinkurl{balancedPulse_label_card_lower_bound}}, \href{https://github.com/wcook04/plectis-erdos/blob/a25cb360bef8dd818dde14b5fb752244304af354/lean/Erdos249257/GenericTailOrbitRigidity.lean\#L247}{\nolinkurl{balancedPulse_no_autonomous_decoder}}, \href{https://github.com/wcook04/plectis-erdos/blob/a25cb360bef8dd818dde14b5fb752244304af354/lean/ErdosProblems/Erdos257/PaperCompleteR21/CentredCompletionAndDecisionBoundary.lean\#L34}{\nolinkurl{paper_centred_completion_of_fixed_precision}}.  \emph{Compared}, runs \href{https://github.com/wcook04/plectis-erdos-lean/actions/runs/35643815458}{\texttt{35643815458}}, \href{https://github.com/wcook04/plectis-erdos-lean/actions/runs/35674034595}{\texttt{35674034595}}.
\par\noindent\hangindent=1.5em Proposition~\ref{prop:exponent-gap}: \href{https://github.com/wcook04/plectis-erdos/blob/a25cb360bef8dd818dde14b5fb752244304af354/lean/ErdosProblems/Erdos257/PaperCompleteR21/DenominatorBudget.lean\#L43}{\nolinkurl{skipSum_den_dvd_prod}}, \href{https://github.com/wcook04/plectis-erdos/blob/a25cb360bef8dd818dde14b5fb752244304af354/lean/ErdosProblems/Erdos257/PaperCompleteR21/DenominatorBudget.lean\#L78}{\nolinkurl{weighted_denominator_budget}}.  \emph{Compared}, run \href{https://github.com/wcook04/plectis-erdos-lean/actions/runs/35643815458}{\texttt{35643815458}}.
\par\noindent\hangindent=1.5em Theorem~\ref{thm:one-sided}: \href{https://github.com/wcook04/plectis-erdos/blob/a25cb360bef8dd818dde14b5fb752244304af354/lean/ErdosProblems/Erdos257/PaperCompleteR21/OneSidedCertificateHierarchy.lean\#L492}{\nolinkurl{paper_one_sidedness}}, \href{https://github.com/wcook04/plectis-erdos/blob/a25cb360bef8dd818dde14b5fb752244304af354/lean/ErdosProblems/Erdos257/PaperCompleteR21/OneSidedCertificateHierarchy.lean\#L291}{\nolinkurl{existsFatalHalfGap_iff_exists_certificate}}, \href{https://github.com/wcook04/plectis-erdos/blob/a25cb360bef8dd818dde14b5fb752244304af354/lean/ErdosProblems/Erdos257/PaperCompleteR21/OneSidedCertificateHierarchy.lean\#L174}{\nolinkurl{existsFatalHalfGap_of_certificate}}, \href{https://github.com/wcook04/plectis-erdos/blob/a25cb360bef8dd818dde14b5fb752244304af354/lean/ErdosProblems/Erdos257/PaperCompleteR21/OneSidedCertificateHierarchy.lean\#L243}{\nolinkurl{certificate_of_existsFatalHalfGap}}, \href{https://github.com/wcook04/plectis-erdos/blob/a25cb360bef8dd818dde14b5fb752244304af354/lean/ErdosProblems/Erdos257/PaperCompleteR21/OneSidedCertificateHierarchy.lean\#L103}{\nolinkurl{scaledMersenneWeight_cast}}, \href{https://github.com/wcook04/plectis-erdos/blob/a25cb360bef8dd818dde14b5fb752244304af354/lean/ErdosProblems/Erdos257/PaperCompleteR21/OneSidedCertificateHierarchy.lean\#L133}{\nolinkurl{certifiedWordValue_cast}}, \href{https://github.com/wcook04/plectis-erdos/blob/a25cb360bef8dd818dde14b5fb752244304af354/lean/ErdosProblems/Erdos257/PaperCompleteR21/OneSidedCertificateHierarchy.lean\#L144}{\nolinkurl{certifiedTailBound_cast}}, \href{https://github.com/wcook04/plectis-erdos/blob/a25cb360bef8dd818dde14b5fb752244304af354/lean/ErdosProblems/Erdos257/PaperCompleteR21/OneSidedCertificateHierarchy.lean\#L160}{\nolinkurl{mersenneTail_eq_sum_add}}.  \emph{Compared}, runs \href{https://github.com/wcook04/plectis-erdos-lean/actions/runs/35643815458}{\texttt{35643815458}}, \href{https://github.com/wcook04/plectis-erdos-lean/actions/runs/35682858662}{\texttt{35682858662}}.
\par\noindent\hangindent=1.5em Theorem~\ref{thm:full-support}: \href{https://github.com/wcook04/plectis-erdos/blob/a25cb360bef8dd818dde14b5fb752244304af354/lean/Erdos249257/CertificateKernel.lean\#L8328}{\nolinkurl{irrational_erdosSum_full_support}}.  \emph{Compared}, run \href{https://github.com/wcook04/plectis-erdos-lean/actions/runs/35624228171}{\texttt{35624228171}}.
\par\noindent\hangindent=1.5em Theorem (Purely periodic support): \href{https://github.com/wcook04/plectis-erdos/blob/a25cb360bef8dd818dde14b5fb752244304af354/lean/Erdos249257/CertificateKernel.lean\#L11590}{\nolinkurl{irrational_erdosSupportSeries_periodic}}.  \emph{Compared}, run \href{https://github.com/wcook04/plectis-erdos-lean/actions/runs/35624228171}{\texttt{35624228171}}.
\par\noindent\hangindent=1.5em Theorem (Eventually periodic support): \href{https://github.com/wcook04/plectis-erdos/blob/a25cb360bef8dd818dde14b5fb752244304af354/lean/Erdos249257/CertificateKernel.lean\#L11604}{\nolinkurl{irrational_erdosSupportSeries_eventuallyPeriodic}}.  \emph{Compared}, run \href{https://github.com/wcook04/plectis-erdos-lean/actions/runs/35624228171}{\texttt{35624228171}}.
\par\noindent\hangindent=1.5em Theorem (Residue-class support): \href{https://github.com/wcook04/plectis-erdos/blob/a25cb360bef8dd818dde14b5fb752244304af354/lean/ErdosProblems/Erdos257/PaperCompleteR21/ResidueClassSupport.lean\#L21}{\nolinkurl{irrational_residueClass_positive_support}}.  \emph{Compared}, run \href{https://github.com/wcook04/plectis-erdos-lean/actions/runs/35624228171}{\texttt{35624228171}}.
\par\noindent\hangindent=1.5em Theorem (Odd support): \href{https://github.com/wcook04/plectis-erdos/blob/a25cb360bef8dd818dde14b5fb752244304af354/lean/Erdos249257/CertificateKernel.lean\#L11686}{\nolinkurl{irrational_erdosSupportSeries_odd}}.  \emph{Compared}, run \href{https://github.com/wcook04/plectis-erdos-lean/actions/runs/35624228171}{\texttt{35624228171}}.
\par\noindent\hangindent=1.5em Theorem~\ref{thm:topology}: \href{https://github.com/wcook04/plectis-erdos/blob/a25cb360bef8dd818dde14b5fb752244304af354/lean/ErdosProblems/Erdos257/PaperCompleteR20/AchievementGeometry.lean\#L40}{\nolinkurl{paper_achievement_geometry}}, \href{https://github.com/wcook04/plectis-erdos/blob/a25cb360bef8dd818dde14b5fb752244304af354/lean/ErdosProblems/Erdos257/PaperCompleteR20/MersenneConstantDecimal.lean\#L21}{\nolinkurl{mersenne_topology_quantitative}}.  \emph{Compared}, run \href{https://github.com/wcook04/plectis-erdos-lean/actions/runs/35643815458}{\texttt{35643815458}}.
\par\noindent\hangindent=1.5em Theorem~\ref{thm:greedy-survival-record}: \href{https://github.com/wcook04/plectis-erdos/blob/a25cb360bef8dd818dde14b5fb752244304af354/lean/Erdos249257/GreedyAchievementSet.lean\#L1458}{\nolinkurl{mem_mersenneAchievementSet_iff_greedy_survival}}, \href{https://github.com/wcook04/plectis-erdos/blob/a25cb360bef8dd818dde14b5fb752244304af354/lean/Erdos249257/GreedyAchievementSet.lean\#L1375}{\nolinkurl{greedy_survives_of_mem_mersenneAchievementSet}}, \href{https://github.com/wcook04/plectis-erdos/blob/a25cb360bef8dd818dde14b5fb752244304af354/lean/Erdos249257/GreedyAchievementSet.lean\#L1394}{\nolinkurl{mem_mersenneAchievementSet_of_greedy_survival}}.  \emph{Compared}, run \href{https://github.com/wcook04/plectis-erdos-lean/actions/runs/35674034595}{\texttt{35674034595}}.
\par\noindent\hangindent=1.5em Theorem~\ref{thm:master-identity}: \href{https://github.com/wcook04/plectis-erdos/blob/a25cb360bef8dd818dde14b5fb752244304af354/lean/ErdosProblems/Erdos257/PaperCompleteR20/QuotientRowIdentity.lean\#L89}{\nolinkurl{paper_master_identity_floors}}.  \emph{Compared}, run \href{https://github.com/wcook04/plectis-erdos-lean/actions/runs/35643815458}{\texttt{35643815458}}.
\par\noindent\hangindent=1.5em Theorem~\ref{thm:real-form}: \href{https://github.com/wcook04/plectis-erdos/blob/a25cb360bef8dd818dde14b5fb752244304af354/lean/ErdosProblems/Erdos257/PaperCompleteR20/QuotientRowReal.lean\#L84}{\nolinkurl{paper_real_quotient_core}}, \href{https://github.com/wcook04/plectis-erdos/blob/a25cb360bef8dd818dde14b5fb752244304af354/lean/ErdosProblems/Erdos257/PaperCompleteR20/QuotientRowReal.lean\#L103}{\nolinkurl{paper_real_quotient_margins}}, \href{https://github.com/wcook04/plectis-erdos/blob/a25cb360bef8dd818dde14b5fb752244304af354/lean/ErdosProblems/Erdos257/PaperCompleteR20/QuotientRowReal.lean\#L94}{\nolinkurl{row_constant_eq_tail}}, \href{https://github.com/wcook04/plectis-erdos/blob/a25cb360bef8dd818dde14b5fb752244304af354/lean/ErdosProblems/Erdos257/PaperCompleteR20/MersenneConstantDecimal.lean\#L8}{\nolinkurl{mersenne_constant_decimal}}.  \emph{Compared}, run \href{https://github.com/wcook04/plectis-erdos-lean/actions/runs/35643815458}{\texttt{35643815458}}.
\par\noindent\hangindent=1.5em Theorem~\ref{thm:dynamics}: \href{https://github.com/wcook04/plectis-erdos/blob/a25cb360bef8dd818dde14b5fb752244304af354/lean/ErdosProblems/Erdos257/PaperCompleteR21/ThreeBranchRowDynamics.lean\#L424}{\nolinkurl{paper_dynamics}}, \href{https://github.com/wcook04/plectis-erdos/blob/a25cb360bef8dd818dde14b5fb752244304af354/lean/ErdosProblems/Erdos257/PaperCompleteR21/ThreeBranchRowDynamics.lean\#L398}{\nolinkurl{paper_rowLower_existsUnique}}, \href{https://github.com/wcook04/plectis-erdos/blob/a25cb360bef8dd818dde14b5fb752244304af354/lean/ErdosProblems/Erdos257/PaperCompleteR21/ThreeBranchRowDynamics.lean\#L404}{\nolinkurl{paper_rowUpper_existsUnique}}, \href{https://github.com/wcook04/plectis-erdos/blob/a25cb360bef8dd818dde14b5fb752244304af354/lean/ErdosProblems/Erdos257/PaperCompleteR21/ThreeBranchRowDynamics.lean\#L279}{\nolinkurl{paper_greedySupport_isRowLower}}, \href{https://github.com/wcook04/plectis-erdos/blob/a25cb360bef8dd818dde14b5fb752244304af354/lean/ErdosProblems/Erdos257/PaperCompleteR21/ThreeBranchRowDynamics.lean\#L364}{\nolinkurl{paper_upperSupport_isRowUpper}}, \href{https://github.com/wcook04/plectis-erdos/blob/a25cb360bef8dd818dde14b5fb752244304af354/lean/ErdosProblems/Erdos257/PaperCompleteR21/ThreeBranchRowDynamics.lean\#L379}{\nolinkurl{paper_isRowLower_unique}}, \href{https://github.com/wcook04/plectis-erdos/blob/a25cb360bef8dd818dde14b5fb752244304af354/lean/ErdosProblems/Erdos257/PaperCompleteR21/ThreeBranchRowDynamics.lean\#L388}{\nolinkurl{paper_isRowUpper_unique}}, \href{https://github.com/wcook04/plectis-erdos/blob/a25cb360bef8dd818dde14b5fb752244304af354/lean/ErdosProblems/Erdos257/PaperCompleteR21/ThreeBranchRowDynamics.lean\#L301}{\nolinkurl{paper_greedySupport_greedy_rule}}, \href{https://github.com/wcook04/plectis-erdos/blob/a25cb360bef8dd818dde14b5fb752244304af354/lean/ErdosProblems/Erdos257/PaperCompleteR21/ThreeBranchRowDynamics.lean\#L254}{\nolinkurl{paper_greedySupport_mem}}, \href{https://github.com/wcook04/plectis-erdos/blob/a25cb360bef8dd818dde14b5fb752244304af354/lean/ErdosProblems/Erdos257/PaperCompleteR21/ThreeBranchRowDynamics.lean\#L244}{\nolinkurl{paper_greedy_step}}, \href{https://github.com/wcook04/plectis-erdos/blob/a25cb360bef8dd818dde14b5fb752244304af354/lean/ErdosProblems/Erdos257/PaperCompleteR21/ThreeBranchRowDynamics.lean\#L69}{\nolinkurl{paper_consecutive_not_both_divisible}}, \href{https://github.com/wcook04/plectis-erdos/blob/a25cb360bef8dd818dde14b5fb752244304af354/lean/ErdosProblems/Erdos257/PaperCompleteR21/ThreeBranchRowDynamics.lean\#L50}{\nolinkurl{paper_rowPulse_eq_indicators}}, \href{https://github.com/wcook04/plectis-erdos/blob/a25cb360bef8dd818dde14b5fb752244304af354/lean/ErdosProblems/Erdos257/PaperCompleteR21/ThreeBranchRowDynamics.lean\#L58}{\nolinkurl{paper_rowQuotient_eq_weightSum}}, \href{https://github.com/wcook04/plectis-erdos/blob/a25cb360bef8dd818dde14b5fb752244304af354/lean/ErdosProblems/Erdos257/PaperCompleteR21/ThreeBranchRowDynamics.lean\#L41}{\nolinkurl{paper_rowWeight_eq_floor}}, \href{https://github.com/wcook04/plectis-erdos/blob/a25cb360bef8dd818dde14b5fb752244304af354/lean/ErdosProblems/Erdos257/PaperCompleteR21/ThreeBranchRowDynamics.lean\#L45}{\nolinkurl{paper_rowTarget_eq}}.  \emph{Compared}, run \href{https://github.com/wcook04/plectis-erdos-lean/actions/runs/35935225572}{\texttt{35935225572}}.
\par\noindent\hangindent=1.5em Theorem~\ref{thm:fatal-absorbing}: \href{https://github.com/wcook04/plectis-erdos/blob/a25cb360bef8dd818dde14b5fb752244304af354/lean/ErdosProblems/Erdos257/PaperCompleteR21/GreedyGapCriteria.lean\#L27}{\nolinkurl{fatal_absorbing}}.  \emph{Compared}, run \href{https://github.com/wcook04/plectis-erdos-lean/actions/runs/35674034595}{\texttt{35674034595}}.
\par\noindent\hangindent=1.5em Theorem~\ref{thm:seam-limit}: \href{https://github.com/wcook04/plectis-erdos/blob/a25cb360bef8dd818dde14b5fb752244304af354/lean/ErdosProblems/Erdos257/PaperCompleteR21/SeamPrefixStabilityLimit.lean\#L560}{\nolinkurl{paper_seam_limit_unconditional}}, \href{https://github.com/wcook04/plectis-erdos/blob/a25cb360bef8dd818dde14b5fb752244304af354/lean/ErdosProblems/Erdos257/PaperCompleteR21/SeamPrefixStabilityLimit.lean\#L438}{\nolinkurl{tendsto_seamGreedyFiniteValue_greedyHalfTargetValue}}, \href{https://github.com/wcook04/plectis-erdos/blob/a25cb360bef8dd818dde14b5fb752244304af354/lean/ErdosProblems/Erdos257/PaperCompleteR21/SeamPrefixStabilityLimit.lean\#L537}{\nolinkurl{tendsto_seamGreedyNormalizedRemainder}}, \href{https://github.com/wcook04/plectis-erdos/blob/a25cb360bef8dd818dde14b5fb752244304af354/lean/ErdosProblems/Erdos257/PaperCompleteR21/SeamPrefixStabilityLimit.lean\#L370}{\nolinkurl{eventually_seamSupport_agrees}}, \href{https://github.com/wcook04/plectis-erdos/blob/a25cb360bef8dd818dde14b5fb752244304af354/lean/ErdosProblems/Erdos257/PaperCompleteR21/SeamPrefixStabilityLimit.lean\#L259}{\nolinkurl{seamScaledRem_eq_tailGreedyRemainder}}, \href{https://github.com/wcook04/plectis-erdos/blob/a25cb360bef8dd818dde14b5fb752244304af354/lean/ErdosProblems/Erdos257/PaperCompleteR21/SeamPrefixStabilityLimit.lean\#L332}{\nolinkurl{mem_seamGreedySupport_iff_scaled}}, \href{https://github.com/wcook04/plectis-erdos/blob/a25cb360bef8dd818dde14b5fb752244304af354/lean/ErdosProblems/Erdos257/PaperCompleteR21/SeamPrefixStabilityLimit.lean\#L398}{\nolinkurl{abs_supportValue_sub_le_mersenneTail}}.  \emph{Compared}, runs \href{https://github.com/wcook04/plectis-erdos-lean/actions/runs/35643815458}{\texttt{35643815458}}, \href{https://github.com/wcook04/plectis-erdos-lean/actions/runs/35840809568}{\texttt{35840809568}}, \href{https://github.com/wcook04/plectis-erdos-lean/actions/runs/35871539470}{\texttt{35871539470}}.
\par\noindent\hangindent=1.5em Theorem~\ref{thm:two-channel-cap}: \href{https://github.com/wcook04/plectis-erdos/blob/a25cb360bef8dd818dde14b5fb752244304af354/lean/Erdos249257/GreedyAchievementSet.lean\#L1425}{\nolinkurl{half_mem_mersenneAchievementSet_of_skipped_twoChannelCap}}, \href{https://github.com/wcook04/plectis-erdos/blob/a25cb360bef8dd818dde14b5fb752244304af354/lean/Erdos249257/GreedyAchievementSet.lean\#L1444}{\nolinkurl{half_mem_mersenneAchievementSet_of_skipped_dyadicCap}}.  \emph{Compared}, run \href{https://github.com/wcook04/plectis-erdos-lean/actions/runs/35674034595}{\texttt{35674034595}}.
\par\noindent\hangindent=1.5em Theorem~\ref{thm:second-channel}: \href{https://github.com/wcook04/plectis-erdos/blob/a25cb360bef8dd818dde14b5fb752244304af354/lean/Erdos249257/GreedyAchievementSet.lean\#L3118}{\nolinkurl{half_mem_mersenneAchievementSet_of_secondChannelSeparation}}, \href{https://github.com/wcook04/plectis-erdos/blob/a25cb360bef8dd818dde14b5fb752244304af354/lean/Erdos249257/GreedyAchievementSet.lean\#L3151}{\nolinkurl{half_mem_mersenneAchievementSet_of_secondChannelSeparationRat_from_seven}}.  \emph{Compared}, run \href{https://github.com/wcook04/plectis-erdos-lean/actions/runs/35674034595}{\texttt{35674034595}}.
\par\noindent\hangindent=1.5em Theorem~\ref{thm:straddle-closed-set}: \href{https://github.com/wcook04/plectis-erdos/blob/a25cb360bef8dd818dde14b5fb752244304af354/lean/ErdosProblems/Erdos257/PaperCompleteR21/GreedyGapCriteria.lean\#L42}{\nolinkurl{straddle_all_depths_iff_mem}}, \href{https://github.com/wcook04/plectis-erdos/blob/a25cb360bef8dd818dde14b5fb752244304af354/lean/ErdosProblems/Erdos257/PaperCompleteR21/GreedyGapCriteria.lean\#L103}{\nolinkurl{straddle_limiting_support_inputs}}.  \emph{Compared}, run \href{https://github.com/wcook04/plectis-erdos-lean/actions/runs/35643815458}{\texttt{35643815458}}.
\par\noindent\hangindent=1.5em Theorem~\ref{thm:largest-skip-late}: \href{https://github.com/wcook04/plectis-erdos/blob/a25cb360bef8dd818dde14b5fb752244304af354/lean/Erdos249257/HalfCylinderLargestSkipInduction.lean\#L167}{\nolinkurl{half_mem_mersenneAchievementSet_of_largestSkipLateStepSocket}}, \href{https://github.com/wcook04/plectis-erdos/blob/a25cb360bef8dd818dde14b5fb752244304af354/lean/Erdos249257/HalfCylinderLargestSkipInduction.lean\#L73}{\nolinkurl{largestSkipLateAt_fourteen}}, \href{https://github.com/wcook04/plectis-erdos/blob/a25cb360bef8dd818dde14b5fb752244304af354/lean/Erdos249257/HalfCylinderHalfMembershipClassification.lean\#L57}{\nolinkurl{seamGreedy_terminal_false_iff_upperOrMiddle}}.  \emph{Compared}, run \href{https://github.com/wcook04/plectis-erdos-lean/actions/runs/35935225572}{\texttt{35935225572}}.
\par\noindent\hangindent=1.5em Theorem~\ref{thm:middle-producer-escape}: \href{https://github.com/wcook04/plectis-erdos/blob/a25cb360bef8dd818dde14b5fb752244304af354/lean/Erdos249257/HalfCylinderMiddleCarryLowerBound.lean\#L2364}{\nolinkurl{half_mem_mersenneAchievementSet_of_middleProducerCardEscape}}, \href{https://github.com/wcook04/plectis-erdos/blob/a25cb360bef8dd818dde14b5fb752244304af354/lean/Erdos249257/HalfCylinderMiddleCarryLowerBound.lean\#L2371}{\nolinkurl{half_mem_mersenneAchievementSet_of_middleProducerRowEscape}}.  \emph{Compared}, run \href{https://github.com/wcook04/plectis-erdos-lean/actions/runs/35935225572}{\texttt{35935225572}}.
\par\noindent\hangindent=1.5em Theorem~\ref{thm:middle-allright-defect} (the Lean statement is stronger): \href{https://github.com/wcook04/plectis-erdos/blob/a25cb360bef8dd818dde14b5fb752244304af354/lean/Erdos249257/HalfCylinderMiddleCarryLowerBound.lean\#L1814}{\nolinkurl{middleProducer_allRight_forces_carry_lt_tail}}, \href{https://github.com/wcook04/plectis-erdos/blob/a25cb360bef8dd818dde14b5fb752244304af354/lean/Erdos249257/HalfCylinderMiddleCarryLowerBound.lean\#L1888}{\nolinkurl{middleProducer_allRight_forces_rational_skip}}, \href{https://github.com/wcook04/plectis-erdos/blob/a25cb360bef8dd818dde14b5fb752244304af354/lean/Erdos249257/HalfCylinderProducerLowerBound.lean\#L150}{\nolinkurl{producerCarry_insert_seamBelowSupport_eq_middleCoordinate}}, \href{https://github.com/wcook04/plectis-erdos/blob/a25cb360bef8dd818dde14b5fb752244304af354/lean/Erdos249257/GenericTailOrbitRigidity.lean\#L78}{\nolinkurl{binaryCoeffTail_nonneg}}, \href{https://github.com/wcook04/plectis-erdos/blob/a25cb360bef8dd818dde14b5fb752244304af354/lean/Erdos249257/HalfCylinderMiddleCarryLowerBound.lean\#L222}{\nolinkurl{binaryCoeffTail_supportCoeff_coe_finset_le_card}}.  \emph{Compared}, run \href{https://github.com/wcook04/plectis-erdos-lean/actions/runs/35935225572}{\texttt{35935225572}}.
\par\noindent\hangindent=1.5em Theorem~\ref{thm:two-sided-dyadic} (the Lean statement is stronger): \href{https://github.com/wcook04/plectis-erdos/blob/a25cb360bef8dd818dde14b5fb752244304af354/lean/Erdos249257/HalfCylinderMiddleCarryLowerBound.lean\#L4448}{\nolinkurl{twoSided}}.  \emph{Compared}, run \href{https://github.com/wcook04/plectis-erdos-lean/actions/runs/35840809568}{\texttt{35840809568}}.
\par\noindent\hangindent=1.5em Theorem~\ref{thm:upper-reset-band} (the Lean statement is stronger): \href{https://github.com/wcook04/plectis-erdos/blob/a25cb360bef8dd818dde14b5fb752244304af354/lean/Erdos249257/HalfCylinderMiddleCarryLowerBound.lean\#L4790}{\nolinkurl{half_mem_mersenneAchievementSet_of_upperResetDyadicBandEscape}}, \href{https://github.com/wcook04/plectis-erdos/blob/a25cb360bef8dd818dde14b5fb752244304af354/lean/Erdos249257/HalfCylinderUpperResetBandCertificates.lean\#L78}{\nolinkurl{seamUpperResetDyadicBandEscape_through_thirty}}.  \emph{Compared}, run \href{https://github.com/wcook04/plectis-erdos-lean/actions/runs/35935225572}{\texttt{35935225572}}.
\par\noindent\hangindent=1.5em Theorem~\ref{thm:mobius-centred-nonneg}: \href{https://github.com/wcook04/plectis-erdos/blob/a25cb360bef8dd818dde14b5fb752244304af354/lean/Erdos249257/HalfCylinderFinalMiddleCellEscape.lean\#L94}{\nolinkurl{mobiusCenteredHalfCarry_nonneg_of_supportSeries_lt_half}}, \href{https://github.com/wcook04/plectis-erdos/blob/a25cb360bef8dd818dde14b5fb752244304af354/lean/Erdos249257/HalfCarryReachability.lean\#L871}{\nolinkurl{integerHalfCarry_eq_scaled_residual_add_tail}}.  \emph{Compared}, run \href{https://github.com/wcook04/plectis-erdos-lean/actions/runs/35674034595}{\texttt{35674034595}}.
\par\noindent\hangindent=1.5em Theorem~\ref{thm:sqrt-bound-route}: \href{https://github.com/wcook04/plectis-erdos/blob/a25cb360bef8dd818dde14b5fb752244304af354/lean/Erdos249257/HalfCarryReachability.lean\#L919}{\nolinkurl{greedy_mobiusCenteredHalfCarry_nonneg}}, \href{https://github.com/wcook04/plectis-erdos/blob/a25cb360bef8dd818dde14b5fb752244304af354/lean/Erdos249257/HalfCarryReachability.lean\#L834}{\nolinkurl{greedy_half_infinite_of_mobiusCenteredHalfCarry_sqrtBound}}, \href{https://github.com/wcook04/plectis-erdos/blob/a25cb360bef8dd818dde14b5fb752244304af354/lean/Erdos249257/HalfCarryReachability.lean\#L817}{\nolinkurl{infinite_support_half_of_mobiusCenteredHalfCarry_sqrtBound}}.  \emph{Compared}, run \href{https://github.com/wcook04/plectis-erdos-lean/actions/runs/35674034595}{\texttt{35674034595}}.
\par\noindent\hangindent=1.5em Theorem~\ref{thm:frozen-margin}: \href{https://github.com/wcook04/plectis-erdos/blob/a25cb360bef8dd818dde14b5fb752244304af354/lean/Erdos249257/HalfCylinderFullShellSeamBridge.lean\#L633}{\nolinkurl{half_mem_mersenneAchievementSet_of_skippedFullShellNonnegative}}, \href{https://github.com/wcook04/plectis-erdos/blob/a25cb360bef8dd818dde14b5fb752244304af354/lean/Erdos249257/HalfCylinderFullShellSeamBridge.lean\#L671}{\nolinkurl{skippedSeamAlignmentZero_iff_skippedFullShellNonnegative}}, \href{https://github.com/wcook04/plectis-erdos/blob/a25cb360bef8dd818dde14b5fb752244304af354/lean/Erdos249257/HalfCylinderFullShellSeamBridge.lean\#L722}{\nolinkurl{half_mem_mersenneAchievementSet_of_skippedSeamEscape}}.  \emph{Compared}, runs \href{https://github.com/wcook04/plectis-erdos-lean/actions/runs/35674034595}{\texttt{35674034595}}, \href{https://github.com/wcook04/plectis-erdos-lean/actions/runs/35840809568}{\texttt{35840809568}}.
\par\noindent\hangindent=1.5em Theorem~\ref{thm:full-support-catalogue}: \href{https://github.com/wcook04/plectis-erdos/blob/a25cb360bef8dd818dde14b5fb752244304af354/lean/Erdos249257/CertificateKernel.lean\#L8328}{\nolinkurl{irrational_erdosSum_full_support}}.  \emph{Compared}, run \href{https://github.com/wcook04/plectis-erdos-lean/actions/runs/35624228171}{\texttt{35624228171}}.
\par\noindent\hangindent=1.5em Theorem~\ref{thm:pairwise-coprime} (the Lean statement is stronger): \href{https://github.com/wcook04/plectis-erdos/blob/a25cb360bef8dd818dde14b5fb752244304af354/lean/Erdos249257/CertificateKernel.lean\#L10776}{\nolinkurl{irrational_erdosSupportSeries_pairwise_coprime}}.  \emph{Compared}, run \href{https://github.com/wcook04/plectis-erdos-lean/actions/runs/35624228171}{\texttt{35624228171}}.
\par\noindent\hangindent=1.5em Theorem~\ref{thm:weighted-coeff-engine}: \href{https://github.com/wcook04/plectis-erdos/blob/a25cb360bef8dd818dde14b5fb752244304af354/lean/Erdos249257/CertificateKernel.lean\#L8665}{\nolinkurl{irrational_coeff_series_of_weighted_coeff_block_certificates}}.  \emph{Compared}, run \href{https://github.com/wcook04/plectis-erdos-lean/actions/runs/35624228171}{\texttt{35624228171}}.
\par\noindent\hangindent=1.5em Theorem~\ref{thm:lcm-gap-engine}: \href{https://github.com/wcook04/plectis-erdos/blob/a25cb360bef8dd818dde14b5fb752244304af354/lean/Erdos249257/CertificateKernel.lean\#L5883}{\nolinkurl{irrational_erdosSum_of_lcm_gap}}.  \emph{Compared}, run \href{https://github.com/wcook04/plectis-erdos-lean/actions/runs/35624228171}{\texttt{35624228171}}.
\par\noindent\hangindent=1.5em Theorem~\ref{thm:factorial-twopow-support} (the Lean statement is stronger): \href{https://github.com/wcook04/plectis-erdos/blob/a25cb360bef8dd818dde14b5fb752244304af354/lean/Erdos249257/CertificateKernel.lean\#L6035}{\nolinkurl{irrational_erdosSum_factorial_support}}, \href{https://github.com/wcook04/plectis-erdos/blob/a25cb360bef8dd818dde14b5fb752244304af354/lean/Erdos249257/CertificateKernel.lean\#L6059}{\nolinkurl{irrational_erdosSum_two_pow_support}}.  \emph{Compared}, run \href{https://github.com/wcook04/plectis-erdos-lean/actions/runs/35624228171}{\texttt{35624228171}}.
\par\noindent\hangindent=1.5em Theorem~\ref{thm:multiples-support} (the Lean statement is stronger): \href{https://github.com/wcook04/plectis-erdos/blob/a25cb360bef8dd818dde14b5fb752244304af354/lean/Erdos249257/CertificateKernel.lean\#L9054}{\nolinkurl{erdosSupportSeries_multiples_eq_pow_base_full_support}}, \href{https://github.com/wcook04/plectis-erdos/blob/a25cb360bef8dd818dde14b5fb752244304af354/lean/Erdos249257/CertificateKernel.lean\#L9103}{\nolinkurl{irrational_erdosSupportSeries_multiples}}.  \emph{Compared}, run \href{https://github.com/wcook04/plectis-erdos-lean/actions/runs/35624228171}{\texttt{35624228171}}.
\par\noindent\hangindent=1.5em The statement at~\ref{thm:periodic-support} (the Lean statement is stronger): \href{https://github.com/wcook04/plectis-erdos/blob/7f79e63d0b36b5b4f0b47b6368342b4a50824f4e/lean/ErdosProblems/Erdos257/PaperCompleteR20/PositivePeriodicSupport.lean\#L49}{\nolinkurl{irrational_erdosSupportSeries_positivePeriodic}}. \emph{Comparator replay pending; Palomar entry pending.}
\par\noindent\hangindent=1.5em Theorem~\ref{thm:eventually-periodic} (the Lean statement is stronger): \href{https://github.com/wcook04/plectis-erdos/blob/a25cb360bef8dd818dde14b5fb752244304af354/lean/Erdos249257/CertificateKernel.lean\#L11604}{\nolinkurl{irrational_erdosSupportSeries_eventuallyPeriodic}}.  \emph{Compared}, run \href{https://github.com/wcook04/plectis-erdos-lean/actions/runs/35624228171}{\texttt{35624228171}}.
\par\noindent\hangindent=1.5em Theorem~\ref{thm:residue-odd} (the Lean statement is stronger): \href{https://github.com/wcook04/plectis-erdos/blob/a25cb360bef8dd818dde14b5fb752244304af354/lean/Erdos249257/CertificateKernel.lean\#L11672}{\nolinkurl{irrational_erdosSupportSeries_residueClass}}, \href{https://github.com/wcook04/plectis-erdos/blob/a25cb360bef8dd818dde14b5fb752244304af354/lean/Erdos249257/CertificateKernel.lean\#L11686}{\nolinkurl{irrational_erdosSupportSeries_odd}}.  \emph{Compared}, run \href{https://github.com/wcook04/plectis-erdos-lean/actions/runs/35624228171}{\texttt{35624228171}}.
\par\noindent\hangindent=1.5em Theorem~\ref{thm:signed-periodic} (the Lean statement is stronger): \href{https://github.com/wcook04/plectis-erdos/blob/a25cb360bef8dd818dde14b5fb752244304af354/lean/Erdos249257/CertificateKernel.lean\#L14175}{\nolinkurl{irrational_or_bpow_mul_eq_intCast_intWeightedErdosSeries_periodic}}, \href{https://github.com/wcook04/plectis-erdos/blob/a25cb360bef8dd818dde14b5fb752244304af354/lean/Erdos249257/CertificateKernel.lean\#L14583}{\nolinkurl{irrational_intWeightedErdosSeries_periodic_of_coeff_nonneg_of_frequently_ne_zero}}, \href{https://github.com/wcook04/plectis-erdos/blob/a25cb360bef8dd818dde14b5fb752244304af354/lean/Erdos249257/CertificateKernel.lean\#L14643}{\nolinkurl{irrational_intWeightedErdosSeries_periodic_of_coeff_nonpos_of_frequently_ne_zero}}.  \emph{Compared}, run \href{https://github.com/wcook04/plectis-erdos-lean/actions/runs/35624228171}{\texttt{35624228171}}.
\par\noindent\hangindent=1.5em Theorem~\ref{thm:mersenne-channel-survival}: \href{https://github.com/wcook04/plectis-erdos/blob/a25cb360bef8dd818dde14b5fb752244304af354/lean/ErdosProblems/Erdos257/PaperCompleteR21/MersenneChannelSurvivalAllHeights.lean\#L210}{\nolinkurl{paper_mersenne_channel_survival}}, \href{https://github.com/wcook04/plectis-erdos/blob/a25cb360bef8dd818dde14b5fb752244304af354/lean/ErdosProblems/Erdos257/PaperCompleteR21/MersenneChannelSurvivalAllHeights.lean\#L234}{\nolinkurl{paper_mersenne_channel_survival_of_coprime_scale}}, \href{https://github.com/wcook04/plectis-erdos/blob/a25cb360bef8dd818dde14b5fb752244304af354/lean/ErdosProblems/Erdos257/PaperCompleteR21/MersenneChannelSurvivalAllHeights.lean\#L248}{\nolinkurl{paper_channel_factor_gcd_eq_one}}, \href{https://github.com/wcook04/plectis-erdos/blob/a25cb360bef8dd818dde14b5fb752244304af354/lean/ErdosProblems/Erdos257/PaperCompleteR21/MersenneChannelSurvivalAllHeights.lean\#L260}{\nolinkurl{paper_channel_factors_pairwise_coprime}}, \href{https://github.com/wcook04/plectis-erdos/blob/a25cb360bef8dd818dde14b5fb752244304af354/lean/ErdosProblems/Erdos257/PaperCompleteR21/MersenneChannelSurvivalAllHeights.lean\#L118}{\nolinkurl{paperB_eq_divInt_paperA}}, \href{https://github.com/wcook04/plectis-erdos/blob/a25cb360bef8dd818dde14b5fb752244304af354/lean/ErdosProblems/Erdos257/PaperCompleteR21/MersenneChannelSurvivalAllHeights.lean\#L100}{\nolinkurl{paperB_eq_baseMobiusShadow}}, \href{https://github.com/wcook04/plectis-erdos/blob/a25cb360bef8dd818dde14b5fb752244304af354/lean/ErdosProblems/Erdos257/PaperCompleteR21/MersenneChannelSurvivalAllHeights.lean\#L83}{\nolinkurl{paperA_eq_mobiusNumerator}}, \href{https://github.com/wcook04/plectis-erdos/blob/a25cb360bef8dd818dde14b5fb752244304af354/lean/ErdosProblems/Erdos257/PaperCompleteR21/MersenneChannelSurvivalAllHeights.lean\#L163}{\nolinkurl{channelProduct_coprime_mobiusNumerator_of_one_le}}, \href{https://github.com/wcook04/plectis-erdos/blob/a25cb360bef8dd818dde14b5fb752244304af354/lean/ErdosProblems/Erdos257/PaperCompleteR21/MersenneChannelSurvivalAllHeights.lean\#L187}{\nolinkurl{upperHalfChannel_survivorProduct_dvd_den_of_one_le}}.  \emph{Compared}, run \href{https://github.com/wcook04/plectis-erdos-lean/actions/runs/35643815458}{\texttt{35643815458}}.
\par\noindent\hangindent=1.5em Theorem~\ref{thm:mersenne-channel-growth}: \href{https://github.com/wcook04/plectis-erdos/blob/a25cb360bef8dd818dde14b5fb752244304af354/lean/Erdos249257/MersenneShadowDenominatorGrowth.lean\#L85}{\nolinkurl{lcmHeight_scaledMobiusShadow_den_lower_bound}}, \href{https://github.com/wcook04/plectis-erdos/blob/a25cb360bef8dd818dde14b5fb752244304af354/lean/Erdos249257/MersenneShadowDenominatorGrowth.lean\#L60}{\nolinkurl{upperHalfMersenneProduct_lower_bound}}, \href{https://github.com/wcook04/plectis-erdos/blob/a25cb360bef8dd818dde14b5fb752244304af354/lean/Erdos249257/MersenneShadowDenominatorGrowth.lean\#L147}{\nolinkurl{lcmHeight_scaledMobiusShadow_den_exact}}.  \emph{Compared}, run \href{https://github.com/wcook04/plectis-erdos-lean/actions/runs/35643815458}{\texttt{35643815458}}.
\par\noindent\hangindent=1.5em Theorem~\ref{thm:mobius-lambert-identity} (the Lean statement is stronger): \href{https://github.com/wcook04/plectis-erdos/blob/a25cb360bef8dd818dde14b5fb752244304af354/lean/Erdos249257/MersenneLambertLadder.lean\#L587}{\nolinkurl{tsum_moebius_div_two_pow_sub_one_eq_half}}.  \emph{Compared}, run \href{https://github.com/wcook04/plectis-erdos-lean/actions/runs/35715779405}{\texttt{35715779405}}.
\par\noindent\hangindent=1.5em Corollary~\ref{cor:negative-mobius-overshoot} (the Lean statement is stronger): \href{https://github.com/wcook04/plectis-erdos/blob/a25cb360bef8dd818dde14b5fb752244304af354/lean/Erdos249257/MobiusSignSupportNoGo.lean\#L111}{\nolinkurl{tsum_negativeMobius_eq_half_add_positiveMobiusTail}}, \href{https://github.com/wcook04/plectis-erdos/blob/a25cb360bef8dd818dde14b5fb752244304af354/lean/Erdos249257/MobiusSignSupportNoGo.lean\#L164}{\nolinkurl{half_lt_tsum_negativeMobius}}.  \emph{Compared}, run \href{https://github.com/wcook04/plectis-erdos-lean/actions/runs/35643815458}{\texttt{35643815458}}.
\par\noindent\hangindent=1.5em Theorem~\ref{thm:half-skip-dichotomy} (the Lean statement is stronger): \href{https://github.com/wcook04/plectis-erdos/blob/a25cb360bef8dd818dde14b5fb752244304af354/lean/Erdos249257/GreedyAchievementSet.lean\#L2583}{\nolinkurl{half_mem_mersenneAchievementSet_iff_greedySkippedSupport_infinite}}, \href{https://github.com/wcook04/plectis-erdos/blob/a25cb360bef8dd818dde14b5fb752244304af354/lean/Erdos249257/GreedyAchievementSet.lean\#L2469}{\nolinkurl{irrational_erdosBorweinMersenneConstant}}.  \emph{Compared}, runs \href{https://github.com/wcook04/plectis-erdos-lean/actions/runs/35643815458}{\texttt{35643815458}}, \href{https://github.com/wcook04/plectis-erdos-lean/actions/runs/35674034595}{\texttt{35674034595}}.
\par\noindent\hangindent=1.5em Theorem~\ref{thm:nine-way-hub} (the Lean statement is stronger): \href{https://github.com/wcook04/plectis-erdos/blob/a25cb360bef8dd818dde14b5fb752244304af354/lean/Erdos249257/HalfCylinderHalfMembershipClassification.lean\#L112}{\nolinkurl{half_mem_mersenneAchievementSet_iff_not_seamGreedyEventuallyRight}}, \href{https://github.com/wcook04/plectis-erdos/blob/a25cb360bef8dd818dde14b5fb752244304af354/lean/Erdos249257/HalfCylinderHalfMembershipClassification.lean\#L126}{\nolinkurl{half_mem_mersenneAchievementSet_iff_unboundedTerminalFalse}}, \href{https://github.com/wcook04/plectis-erdos/blob/a25cb360bef8dd818dde14b5fb752244304af354/lean/Erdos249257/HalfCylinderHalfMembershipClassification.lean\#L156}{\nolinkurl{half_mem_mersenneAchievementSet_iff_unboundedUpperOrMiddle}}, \href{https://github.com/wcook04/plectis-erdos/blob/a25cb360bef8dd818dde14b5fb752244304af354/lean/Erdos249257/HalfCylinderHalfMembershipClassification.lean\#L203}{\nolinkurl{half_mem_mersenneAchievementSet_iff_cofinalTerminalFalse}}, \href{https://github.com/wcook04/plectis-erdos/blob/a25cb360bef8dd818dde14b5fb752244304af354/lean/Erdos249257/HalfCylinderHalfMembershipClassification.lean\#L213}{\nolinkurl{half_mem_mersenneAchievementSet_iff_exists_unboundedSkippedRanksAlong}}, \href{https://github.com/wcook04/plectis-erdos/blob/a25cb360bef8dd818dde14b5fb752244304af354/lean/Erdos249257/HalfCylinderHalfMembershipClassification.lean\#L228}{\nolinkurl{unboundedTerminalFalse_iff_greedyMersenneSkippedSupport_infinite}}, \href{https://github.com/wcook04/plectis-erdos/blob/a25cb360bef8dd818dde14b5fb752244304af354/lean/Erdos249257/HalfCylinderHalfMembershipClassification.lean\#L235}{\nolinkurl{half_mem_mersenneAchievementSet_iff_no_lastHalfGreedySkip}}.  \emph{Compared}, run \href{https://github.com/wcook04/plectis-erdos-lean/actions/runs/35935225572}{\texttt{35935225572}}.
\par\noindent\hangindent=1.5em Lemma~\ref{lem:eventually-right-impossible} (the Lean statement is stronger): \href{https://github.com/wcook04/plectis-erdos/blob/a25cb360bef8dd818dde14b5fb752244304af354/lean/Erdos249257/HalfCylinderFatalGapRightTail.lean\#L402}{\nolinkurl{prefix_add_mersenneTail_lt_half_of_eventually_right}}, \href{https://github.com/wcook04/plectis-erdos/blob/a25cb360bef8dd818dde14b5fb752244304af354/lean/Erdos249257/HalfCylinderFatalGapRightTail.lean\#L627}{\nolinkurl{half_lt_upper_competitor_of_eventually_right}}.  \emph{Compared}, run \href{https://github.com/wcook04/plectis-erdos-lean/actions/runs/35840809568}{\texttt{35840809568}}.
\par\noindent\hangindent=1.5em Lemma~\ref{lem:mersenne-tail-weight} (the Lean statement is stronger): \href{https://github.com/wcook04/plectis-erdos/blob/a25cb360bef8dd818dde14b5fb752244304af354/lean/Erdos249257/GreedyAchievementSet.lean\#L114}{\nolinkurl{mersenneTail_eq_weight_add}}, \href{https://github.com/wcook04/plectis-erdos/blob/a25cb360bef8dd818dde14b5fb752244304af354/lean/Erdos249257/GreedyAchievementSet.lean\#L1226}{\nolinkurl{halfTwoChannelCap_lt_mersenneTail}}, \href{https://github.com/wcook04/plectis-erdos/blob/a25cb360bef8dd818dde14b5fb752244304af354/lean/Erdos249257/GreedyAchievementSet.lean\#L155}{\nolinkurl{mersenneTail_le_two_mul_weight}}, \href{https://github.com/wcook04/plectis-erdos/blob/a25cb360bef8dd818dde14b5fb752244304af354/lean/Erdos249257/GreedyAchievementSet.lean\#L125}{\nolinkurl{two_mul_mersenneWeight_succ_lt}}, \href{https://github.com/wcook04/plectis-erdos/blob/a25cb360bef8dd818dde14b5fb752244304af354/lean/Erdos249257/GreedyAchievementSet.lean\#L180}{\nolinkurl{mersenneTail_lt_weight}}.  \emph{Compared}, run \href{https://github.com/wcook04/plectis-erdos-lean/actions/runs/35643815458}{\texttt{35643815458}}.
\par\noindent\hangindent=1.5em Theorem~\ref{thm:greedy-survival-catalogue} (the Lean statement is stronger): \href{https://github.com/wcook04/plectis-erdos/blob/a25cb360bef8dd818dde14b5fb752244304af354/lean/Erdos249257/GreedyAchievementSet.lean\#L1458}{\nolinkurl{mem_mersenneAchievementSet_iff_greedy_survival}}.  \emph{Compared}, run \href{https://github.com/wcook04/plectis-erdos-lean/actions/runs/35674034595}{\texttt{35674034595}}.
\par\noindent\hangindent=1.5em Lemma~\ref{lem:rank-step-trichotomy} (the Lean statement is stronger): \href{https://github.com/wcook04/plectis-erdos/blob/a25cb360bef8dd818dde14b5fb752244304af354/lean/Erdos249257/HalfCutLocator.lean\#L205}{\nolinkurl{isStraddlePrefix_step_trichotomy}}, \href{https://github.com/wcook04/plectis-erdos/blob/a25cb360bef8dd818dde14b5fb752244304af354/lean/Erdos249257/HalfCutLocator.lean\#L301}{\nolinkurl{half_step_forced}}.  \emph{Compared}, run \href{https://github.com/wcook04/plectis-erdos-lean/actions/runs/35715779405}{\texttt{35715779405}}.
\par\noindent\hangindent=1.5em Lemma~\ref{lem:fatal-gap-exclusion}: \href{https://github.com/wcook04/plectis-erdos/blob/a25cb360bef8dd818dde14b5fb752244304af354/lean/ErdosProblems/Erdos257/PaperCompleteR21/GreedyGapCriteria.lean\#L174}{\nolinkurl{fatal_gap_excludes_every_representation}}, \href{https://github.com/wcook04/plectis-erdos/blob/a25cb360bef8dd818dde14b5fb752244304af354/lean/ErdosProblems/Erdos257/PaperCompleteR21/GreedyGapCriteria.lean\#L114}{\nolinkurl{fatal_gap_endpoint_bounds}}, \href{https://github.com/wcook04/plectis-erdos/blob/a25cb360bef8dd818dde14b5fb752244304af354/lean/ErdosProblems/Erdos257/PaperCompleteR21/GreedyGapCriteria.lean\#L144}{\nolinkurl{depth_prefix_interval_disjoint}}, \href{https://github.com/wcook04/plectis-erdos/blob/a25cb360bef8dd818dde14b5fb752244304af354/lean/ErdosProblems/Erdos257/PaperCompleteR21/GreedyGapCriteria.lean\#L162}{\nolinkurl{fatal_gap_within_prefix_interval}}.  \emph{Compared}, run \href{https://github.com/wcook04/plectis-erdos-lean/actions/runs/35643815458}{\texttt{35643815458}}.
\par\noindent\hangindent=1.5em Lemma~\ref{lem:half-endpoint-kills} (the Lean statement is stronger): \href{https://github.com/wcook04/plectis-erdos/blob/a25cb360bef8dd818dde14b5fb752244304af354/lean/Erdos249257/HalfCutLocator.lean\#L243}{\nolinkurl{positiveMersenneSupportValue_coe_finset_ne_half}}, \href{https://github.com/wcook04/plectis-erdos/blob/a25cb360bef8dd818dde14b5fb752244304af354/lean/Erdos249257/HalfCutLocator.lean\#L263}{\nolinkurl{half_ne_coe_finset_add_mersenneTail}}.  \emph{Compared}, run \href{https://github.com/wcook04/plectis-erdos-lean/actions/runs/35643815458}{\texttt{35643815458}}.
\par\noindent\hangindent=1.5em Lemma~\ref{lem:straddle-agrees-greedy} (the Lean statement is stronger): \href{https://github.com/wcook04/plectis-erdos/blob/a25cb360bef8dd818dde14b5fb752244304af354/lean/Erdos249257/HalfCutLocator.lean\#L442}{\nolinkurl{half_agrees_greedy}}.  \emph{Compared}, run \href{https://github.com/wcook04/plectis-erdos-lean/actions/runs/35935225572}{\texttt{35935225572}}.
\par\noindent\hangindent=1.5em Theorem~\ref{thm:last-skip-iff-fatal} (the Lean statement is stronger): \href{https://github.com/wcook04/plectis-erdos/blob/a25cb360bef8dd818dde14b5fb752244304af354/lean/Erdos249257/HalfCylinderFixedTailSocket.lean\#L22}{\nolinkurl{isLastHalfGreedySkip_iff_skip_and_fatal}}, \href{https://github.com/wcook04/plectis-erdos/blob/a25cb360bef8dd818dde14b5fb752244304af354/lean/Erdos249257/HalfCylinderFixedTailSocket.lean\#L73}{\nolinkurl{half_mem_iff_every_actual_skip_survives}}.  \emph{Compared}, run \href{https://github.com/wcook04/plectis-erdos-lean/actions/runs/35674034595}{\texttt{35674034595}}.
\par\noindent\hangindent=1.5em Lemma~\ref{lem:seam-upper-or-middle} (the Lean statement is stronger): \href{https://github.com/wcook04/plectis-erdos/blob/a25cb360bef8dd818dde14b5fb752244304af354/lean/Erdos249257/HalfCylinderHalfMembershipClassification.lean\#L57}{\nolinkurl{seamGreedy_terminal_false_iff_upperOrMiddle}}.  \emph{Compared}, run \href{https://github.com/wcook04/plectis-erdos-lean/actions/runs/35840809568}{\texttt{35840809568}}.
\par\noindent\hangindent=1.5em Lemma~\ref{lem:largest-false-rank-algebra}: \href{https://github.com/wcook04/plectis-erdos/blob/a25cb360bef8dd818dde14b5fb752244304af354/lean/ErdosProblems/Erdos257/PaperCompleteR21/SeamRowGapAndCarry.lean\#L39}{\nolinkurl{largest_false_rank_algebra}}, \href{https://github.com/wcook04/plectis-erdos/blob/a25cb360bef8dd818dde14b5fb752244304af354/lean/ErdosProblems/Erdos257/PaperCompleteR21/SeamRowGapAndCarry.lean\#L98}{\nolinkurl{paper_largest_false_rank_algebra}}.  \emph{Compared}, runs \href{https://github.com/wcook04/plectis-erdos-lean/actions/runs/35643815458}{\texttt{35643815458}}, \href{https://github.com/wcook04/plectis-erdos-lean/actions/runs/35840809568}{\texttt{35840809568}}.
\par\noindent\hangindent=1.5em Theorem~\ref{thm:critical-dyadic-band}: \href{https://github.com/wcook04/plectis-erdos/blob/a25cb360bef8dd818dde14b5fb752244304af354/lean/Erdos249257/HalfUpperResetCriticalBand.lean\#L46}{\nolinkurl{exists_criticalDyadicBandIndex}}, \href{https://github.com/wcook04/plectis-erdos/blob/a25cb360bef8dd818dde14b5fb752244304af354/lean/Erdos249257/HalfUpperResetCriticalBand.lean\#L108}{\nolinkurl{dyadicBandEscape_iff_exists_critical}}, \href{https://github.com/wcook04/plectis-erdos/blob/a25cb360bef8dd818dde14b5fb752244304af354/lean/Erdos249257/HalfUpperResetCriticalBand.lean\#L883}{\nolinkurl{seamUpperResetCriticalBandEscape_iff}}.  \emph{Compared}, run \href{https://github.com/wcook04/plectis-erdos-lean/actions/runs/35935225572}{\texttt{35935225572}}.
\par\noindent\hangindent=1.5em Theorem~\ref{thm:final-middle-cell}: \href{https://github.com/wcook04/plectis-erdos/blob/a25cb360bef8dd818dde14b5fb752244304af354/lean/Erdos249257/HalfCylinderFinalMiddleCellEscape.lean\#L587}{\nolinkurl{finalMiddleCell_neg_three_not_last}}, \href{https://github.com/wcook04/plectis-erdos/blob/a25cb360bef8dd818dde14b5fb752244304af354/lean/Erdos249257/HalfCylinderFinalMiddleCellEscape.lean\#L39}{\nolinkurl{mobiusCenteredHalfCarry_add_two}}.  \emph{Compared}, run \href{https://github.com/wcook04/plectis-erdos-lean/actions/runs/35935225572}{\texttt{35935225572}}.
\par\noindent\hangindent=1.5em Lemma~\ref{lem:skipped-endpoint-trichotomy}: \href{https://github.com/wcook04/plectis-erdos/blob/a25cb360bef8dd818dde14b5fb752244304af354/lean/Erdos249257/HalfCylinderSkippedEndpointClassifier.lean\#L246}{\nolinkurl{halfGreedy_skipped_endpoint_trichotomy}}.  \emph{Compared}, run \href{https://github.com/wcook04/plectis-erdos-lean/actions/runs/35840809568}{\texttt{35840809568}}.
\par\noindent\hangindent=1.5em Lemma~\ref{lem:reverse-carry-word}: \href{https://github.com/wcook04/plectis-erdos/blob/a25cb360bef8dd818dde14b5fb752244304af354/lean/ErdosProblems/Erdos257/PaperCompleteR21/SeamRowGapAndCarry.lean\#L240}{\nolinkurl{paper_reverse_carry_word}}, \href{https://github.com/wcook04/plectis-erdos/blob/a25cb360bef8dd818dde14b5fb752244304af354/lean/ErdosProblems/Erdos257/PaperCompleteR21/SeamRowGapAndCarry.lean\#L209}{\nolinkurl{reverse_carry_word_common_bound_sharp}}.  \emph{Compared}, run \href{https://github.com/wcook04/plectis-erdos-lean/actions/runs/35643815458}{\texttt{35643815458}}.
\par\noindent\hangindent=1.5em Lemma~\ref{lem:linear-channel-nogo} (the Lean statement is stronger): \href{https://github.com/wcook04/plectis-erdos/blob/a25cb360bef8dd818dde14b5fb752244304af354/lean/Erdos249257/AdelicHeightObstruction.lean\#L120}{\nolinkurl{linearDescender_eq_smul_eval}}, \href{https://github.com/wcook04/plectis-erdos/blob/a25cb360bef8dd818dde14b5fb752244304af354/lean/Erdos249257/HalfTrappingReturnCarry.lean\#L42}{\nolinkurl{relationInvariantLinearChannels_det_eq_zero}}.  \emph{Compared}, runs \href{https://github.com/wcook04/plectis-erdos-lean/actions/runs/35624228171}{\texttt{35624228171}}, \href{https://github.com/wcook04/plectis-erdos-lean/actions/runs/35643815458}{\texttt{35643815458}}.
\par\noindent\hangindent=1.5em Theorem~\ref{thm:two-thirds-band}: \href{https://github.com/wcook04/plectis-erdos/blob/a25cb360bef8dd818dde14b5fb752244304af354/lean/ErdosProblems/Erdos257/PaperCompleteR21/PostTakeBandLocalisation.lean\#L96}{\nolinkurl{paper_two_thirds_band}}.  \emph{Compared}, run \href{https://github.com/wcook04/plectis-erdos-lean/actions/runs/35643815458}{\texttt{35643815458}}.
\par\noindent\hangindent=1.5em Theorem~\ref{thm:sharp-fatal-gap}: \href{https://github.com/wcook04/plectis-erdos/blob/a25cb360bef8dd818dde14b5fb752244304af354/lean/ErdosProblems/Erdos257/PaperCompleteR21/GreedyGapCriteria.lean\#L299}{\nolinkurl{paper_sharp_fatal_gap}}.  \emph{Compared}, run \href{https://github.com/wcook04/plectis-erdos-lean/actions/runs/35643815458}{\texttt{35643815458}}.
\par\noindent\hangindent=1.5em Lemma~\ref{lem:gap-mass-summability} (the Lean statement is stronger): \href{https://github.com/wcook04/plectis-erdos/blob/a25cb360bef8dd818dde14b5fb752244304af354/lean/Erdos249257/GreedyAchievementSet.lean\#L2346}{\nolinkurl{mersenneGap_pos}}, \href{https://github.com/wcook04/plectis-erdos/blob/a25cb360bef8dd818dde14b5fb752244304af354/lean/Erdos249257/HalfGapMass.lean\#L77}{\nolinkurl{summable_mersenneGap_succ}}, \href{https://github.com/wcook04/plectis-erdos/blob/a25cb360bef8dd818dde14b5fb752244304af354/lean/Erdos249257/HalfGapMass.lean\#L83}{\nolinkurl{mersenneGap_tail_le}}, \href{https://github.com/wcook04/plectis-erdos/blob/a25cb360bef8dd818dde14b5fb752244304af354/lean/Erdos249257/HalfGapMass.lean\#L104}{\nolinkurl{tendsto_mersenneGap_tail_zero}}.  \emph{Compared}, run \href{https://github.com/wcook04/plectis-erdos-lean/actions/runs/35643815458}{\texttt{35643815458}}.
\par\noindent\hangindent=1.5em Lemma~\ref{lem:half-divisor-unit-drop} (the Lean statement is stronger): \href{https://github.com/wcook04/plectis-erdos/blob/a25cb360bef8dd818dde14b5fb752244304af354/lean/Erdos249257/HalfCylinderIntegerGreedy.lean\#L882}{\nolinkurl{supportCoeff_insert_eq_add_indicator}}, \href{https://github.com/wcook04/plectis-erdos/blob/a25cb360bef8dd818dde14b5fb752244304af354/lean/Erdos249257/HalfDivisorUnitDrop.lean\#L20}{\nolinkurl{supportCoeff_extend_true_eq_false_add_one_at_double}}.  \emph{Compared}, runs \href{https://github.com/wcook04/plectis-erdos-lean/actions/runs/35624228171}{\texttt{35624228171}}, \href{https://github.com/wcook04/plectis-erdos-lean/actions/runs/35643815458}{\texttt{35643815458}}.
\par\noindent\hangindent=1.5em Theorem~\ref{thm:tempered-orbit-rigidity} (the Lean statement is stronger): \href{https://github.com/wcook04/plectis-erdos/blob/a25cb360bef8dd818dde14b5fb752244304af354/lean/Erdos249257/GenericTailOrbitRigidity.lean\#L426}{\nolinkurl{binaryCoeffSeries_rational_iff_exists_temperedBinaryOrbit}}, \href{https://github.com/wcook04/plectis-erdos/blob/a25cb360bef8dd818dde14b5fb752244304af354/lean/Erdos249257/GenericTailOrbitRigidity.lean\#L339}{\nolinkurl{temperedBinaryOrbit_eq_scaledTail}}.  \emph{Compared}, run \href{https://github.com/wcook04/plectis-erdos-lean/actions/runs/35643815458}{\texttt{35643815458}}.
\par\noindent\hangindent=1.5em Lemma~\ref{lem:tail-transfer} (the Lean statement is stronger): \href{https://github.com/wcook04/plectis-erdos/blob/a25cb360bef8dd818dde14b5fb752244304af354/lean/Erdos249257/CertificateKernel.lean\#L9467}{\nolinkurl{irrational_erdosSupportSeries_of_tail}}, \href{https://github.com/wcook04/plectis-erdos/blob/a25cb360bef8dd818dde14b5fb752244304af354/lean/Erdos249257/CertificateKernel.lean\#L9476}{\nolinkurl{irrational_erdosSupportSeries_tail_of_irrational}}.  \emph{Compared}, run \href{https://github.com/wcook04/plectis-erdos-lean/actions/runs/35624228171}{\texttt{35624228171}}.
\par\noindent\hangindent=1.5em Lemma~\ref{lem:dyadic-excess-reformulation} (the Lean statement is stronger): \href{https://github.com/wcook04/plectis-erdos/blob/a25cb360bef8dd818dde14b5fb752244304af354/lean/Erdos249257/DyadicPrefixCompression.lean\#L198}{\nolinkurl{divInt_le_nextDyadic_iff_excess_nonpos}}, \href{https://github.com/wcook04/plectis-erdos/blob/a25cb360bef8dd818dde14b5fb752244304af354/lean/Erdos249257/DyadicPrefixCompression.lean\#L1044}{\nolinkurl{greedyHalf_mem_nextMersenneDyadicSliver_iff_excess}}.  \emph{Compared}, runs \href{https://github.com/wcook04/plectis-erdos-lean/actions/runs/35624228171}{\texttt{35624228171}}, \href{https://github.com/wcook04/plectis-erdos-lean/actions/runs/35674034595}{\texttt{35674034595}}.
\par\noindent\hangindent=1.5em Lemma~\ref{lem:denominator-sandwich} (the Lean statement is stronger): \href{https://github.com/wcook04/plectis-erdos/blob/a25cb360bef8dd818dde14b5fb752244304af354/lean/Erdos249257/DyadicPrefixCompression.lean\#L118}{\nolinkurl{dyadicResidual_denominator_sandwich}}.  \emph{Compared}, run \href{https://github.com/wcook04/plectis-erdos-lean/actions/runs/35624228171}{\texttt{35624228171}}.
\par\noindent\hangindent=1.5em Lemma~\ref{lem:denominator-survival} (the Lean statement is stronger): \href{https://github.com/wcook04/plectis-erdos/blob/a25cb360bef8dd818dde14b5fb752244304af354/lean/Erdos249257/RationalDenominatorSurvival.lean\#L17}{\nolinkurl{divisor_dvd_divInt_den}}, \href{https://github.com/wcook04/plectis-erdos/blob/a25cb360bef8dd818dde14b5fb752244304af354/lean/Erdos249257/RationalDenominatorSurvival.lean\#L38}{\nolinkurl{survivingDivisor_dvd_scaled_divInt_den}}.  \emph{Compared}, run \href{https://github.com/wcook04/plectis-erdos-lean/actions/runs/35643815458}{\texttt{35643815458}}.
\par\noindent\hangindent=1.5em Lemma~\ref{lem:mixed-prime-power-layer} (the Lean statement is stronger): \href{https://github.com/wcook04/plectis-erdos/blob/a25cb360bef8dd818dde14b5fb752244304af354/lean/Erdos249257/MaximalOmegaLayer.lean\#L29}{\nolinkurl{primePowerLayer_comm}}, \href{https://github.com/wcook04/plectis-erdos/blob/a25cb360bef8dd818dde14b5fb752244304af354/lean/Erdos249257/MaximalOmegaLayer.lean\#L39}{\nolinkurl{mixedPrimePowerLayerTwo_supportCoeffInt}}.  \emph{Compared}, runs \href{https://github.com/wcook04/plectis-erdos-lean/actions/runs/35624228171}{\texttt{35624228171}}, \href{https://github.com/wcook04/plectis-erdos-lean/actions/runs/35643815458}{\texttt{35643815458}}.
\par\noindent\hangindent=1.5em Proposition~\ref{prop:achievement-set-topology}: \href{https://github.com/wcook04/plectis-erdos/blob/a25cb360bef8dd818dde14b5fb752244304af354/lean/ErdosProblems/Erdos257/PaperCompleteR21/AchievementSetTopologyAndFiniteHalf.lean\#L27}{\nolinkurl{paper_achievement_set_topology}}.  \emph{Compared}, run \href{https://github.com/wcook04/plectis-erdos-lean/actions/runs/35643815458}{\texttt{35643815458}}.
\par\noindent\hangindent=1.5em Theorem~\ref{thm:master-dichotomy} (the Lean statement is stronger): \href{https://github.com/wcook04/plectis-erdos/blob/a25cb360bef8dd818dde14b5fb752244304af354/lean/Erdos249257/HalfCutLocator.lean\#L654}{\nolinkurl{half_mem_mersenneAchievementSet_iff_no_existsFatalHalfGap}}, \href{https://github.com/wcook04/plectis-erdos/blob/a25cb360bef8dd818dde14b5fb752244304af354/lean/Erdos249257/HalfCutLocator.lean\#L623}{\nolinkurl{half_mem_mersenneAchievementSet_or_exists_fatal_gap}}.  \emph{Compared}, run \href{https://github.com/wcook04/plectis-erdos-lean/actions/runs/35643815458}{\texttt{35643815458}}.
\par\noindent\hangindent=1.5em Theorem~\ref{thm:perturbed-family-maximality} (the Lean statement is stronger): \href{https://github.com/wcook04/plectis-erdos/blob/a25cb360bef8dd818dde14b5fb752244304af354/lean/Erdos249257/HalfCylinderIntegerGreedy.lean\#L1390}{\nolinkurl{prefixChoice_maximal}}, \href{https://github.com/wcook04/plectis-erdos/blob/a25cb360bef8dd818dde14b5fb752244304af354/lean/Erdos249257/HalfCylinderIntegerGreedy.lean\#L1463}{\nolinkurl{nextRemainder_trichotomy}}.
\par\noindent\hangindent=1.5em Theorem~\ref{record:257bm-c1}: \href{https://github.com/wcook04/plectis-erdos/blob/a25cb360bef8dd818dde14b5fb752244304af354/lean/Erdos249257/BooleanMobiusCofinalExactRows.lean\#L58}{\nolinkurl{abs_exactLocalMersenneRowValue_sub_half_le}}, \href{https://github.com/wcook04/plectis-erdos/blob/a25cb360bef8dd818dde14b5fb752244304af354/lean/Erdos249257/BooleanMobiusCofinalExactRows.lean\#L71}{\nolinkurl{half_mem_mersenneAchievementSet_of_cofinalExactLocalRows}}, \href{https://github.com/wcook04/plectis-erdos/blob/a25cb360bef8dd818dde14b5fb752244304af354/lean/Erdos249257/HalfCarryReachability.lean\#L589}{\nolinkurl{finite_boolSupport_ne_half}}.  \emph{Compared}, run \href{https://github.com/wcook04/plectis-erdos-lean/actions/runs/35643815458}{\texttt{35643815458}}.
\par\noindent\hangindent=1.5em Theorem~\ref{record:257bm-c2}: \href{https://github.com/wcook04/plectis-erdos/blob/a25cb360bef8dd818dde14b5fb752244304af354/lean/Erdos249257/BooleanMobiusSkipRowCofinal.lean\#L55}{\nolinkurl{exactLocalMersenneHalfRow_of_positiveHalfGreedySkip}}, \href{https://github.com/wcook04/plectis-erdos/blob/a25cb360bef8dd818dde14b5fb752244304af354/lean/Erdos249257/BooleanMobiusSkipRowCofinal.lean\#L84}{\nolinkurl{cofinalExactLocalMersenneHalfRows_of_positiveHalfGreedySkips}}, \href{https://github.com/wcook04/plectis-erdos/blob/a25cb360bef8dd818dde14b5fb752244304af354/lean/Erdos249257/BooleanMobiusSkipRowCofinal.lean\#L97}{\nolinkurl{half_mem_mersenneAchievementSet_of_positiveHalfGreedySkips}}, \href{https://github.com/wcook04/plectis-erdos/blob/a25cb360bef8dd818dde14b5fb752244304af354/lean/Erdos249257/BooleanMobiusSkipRowCofinal.lean\#L110}{\nolinkurl{cofinalPositiveHalfGreedySkips_iff_half_mem}}.  \emph{Compared}, runs \href{https://github.com/wcook04/plectis-erdos-lean/actions/runs/35544127144}{\texttt{35544127144}}, \href{https://github.com/wcook04/plectis-erdos-lean/actions/runs/35674034595}{\texttt{35674034595}}.
\par\noindent\hangindent=1.5em Theorem~\ref{record:257bm-c3}: \href{https://github.com/wcook04/plectis-erdos/blob/a25cb360bef8dd818dde14b5fb752244304af354/lean/ErdosProblems/Erdos257/PaperCompleteR21/CompatibleFiniteRowFamily.lean\#L95}{\nolinkurl{paper_compatible_bit_stable}}, \href{https://github.com/wcook04/plectis-erdos/blob/a25cb360bef8dd818dde14b5fb752244304af354/lean/ErdosProblems/Erdos257/PaperCompleteR21/CompatibleFiniteRowFamily.lean\#L111}{\nolinkurl{paper_compatible_finite_row_conditions}}, \href{https://github.com/wcook04/plectis-erdos/blob/a25cb360bef8dd818dde14b5fb752244304af354/lean/ErdosProblems/Erdos257/PaperCompleteR21/CompatibleFiniteRowFamily.lean\#L146}{\nolinkurl{paper_compatible_first_condition_gives_nonneg}}, \href{https://github.com/wcook04/plectis-erdos/blob/a25cb360bef8dd818dde14b5fb752244304af354/lean/ErdosProblems/Erdos257/PaperCompleteR21/CompatibleFiniteRowFamily.lean\#L158}{\nolinkurl{paper_compatible_rows_agree_with_limit}}.  \emph{Compared}, runs \href{https://github.com/wcook04/plectis-erdos-lean/actions/runs/35674034595}{\texttt{35674034595}}, \href{https://github.com/wcook04/plectis-erdos-lean/actions/runs/35682858662}{\texttt{35682858662}}, \href{https://github.com/wcook04/plectis-erdos-lean/actions/runs/35840809568}{\texttt{35840809568}}.
\par\noindent\hangindent=1.5em Theorem~\ref{record:257bm-c4}: \href{https://github.com/wcook04/plectis-erdos/blob/a25cb360bef8dd818dde14b5fb752244304af354/lean/Erdos249257/BooleanMobiusCriticalCapacityCofinal.lean\#L1309}{\nolinkurl{cofinalExactLocalMersenneHalfRows_of_criticalQuotientSupply}}, \href{https://github.com/wcook04/plectis-erdos/blob/a25cb360bef8dd818dde14b5fb752244304af354/lean/Erdos249257/BooleanMobiusCriticalCapacityCofinal.lean\#L1321}{\nolinkurl{half_mem_mersenneAchievementSet_of_criticalQuotientSupply}}, \href{https://github.com/wcook04/plectis-erdos/blob/a25cb360bef8dd818dde14b5fb752244304af354/lean/Erdos249257/HalfCarryReachability.lean\#L589}{\nolinkurl{finite_boolSupport_ne_half}}.  \emph{Compared}, run \href{https://github.com/wcook04/plectis-erdos-lean/actions/runs/35643815458}{\texttt{35643815458}}.
\par\noindent\hangindent=1.5em Theorem~\ref{record:257bm-c5}: \href{https://github.com/wcook04/plectis-erdos/blob/a25cb360bef8dd818dde14b5fb752244304af354/lean/Erdos249257/BooleanMobiusCriticalCapacityCofinal.lean\#L1087}{\nolinkurl{ProtectedExactLocalMersenneRow}}, \href{https://github.com/wcook04/plectis-erdos/blob/a25cb360bef8dd818dde14b5fb752244304af354/lean/Erdos249257/BooleanMobiusCriticalCapacityCofinal.lean\#L1139}{\nolinkurl{exists_laterProtectedExactLocalMersenneRow}}, \href{https://github.com/wcook04/plectis-erdos/blob/a25cb360bef8dd818dde14b5fb752244304af354/lean/Erdos249257/BooleanMobiusCriticalCapacityCofinal.lean\#L1309}{\nolinkurl{cofinalExactLocalMersenneHalfRows_of_criticalQuotientSupply}}.  \emph{Compared}, run \href{https://github.com/wcook04/plectis-erdos-lean/actions/runs/35935225572}{\texttt{35935225572}}.
\par\noindent\hangindent=1.5em Theorem~\ref{record:257bm-c6}: \href{https://github.com/wcook04/plectis-erdos/blob/a25cb360bef8dd818dde14b5fb752244304af354/lean/Erdos249257/BooleanMobiusCriticalCapacityCofinal.lean\#L1068}{\nolinkurl{skippedCoreCriticalQuotientSupply_iff_halfGreedySkipped}}.  \emph{Compared}, run \href{https://github.com/wcook04/plectis-erdos-lean/actions/runs/35674034595}{\texttt{35674034595}}.
\par\noindent\hangindent=1.5em Theorem~\ref{record:257bm-c6a}: \href{https://github.com/wcook04/plectis-erdos/blob/a25cb360bef8dd818dde14b5fb752244304af354/lean/Erdos249257/BooleanMobiusCriticalCapacityCofinal.lean\#L682}{\nolinkurl{halfGreedy_precriticalSuffix_lt_of_next_skip}}, \href{https://github.com/wcook04/plectis-erdos/blob/a25cb360bef8dd818dde14b5fb752244304af354/lean/Erdos249257/BooleanMobiusCriticalCapacityCofinal.lean\#L1001}{\nolinkurl{halfGreedySkippedCriticalQuotientSupply_of_precriticalSuffix}}, \href{https://github.com/wcook04/plectis-erdos/blob/a25cb360bef8dd818dde14b5fb752244304af354/lean/Erdos249257/BooleanMobiusCriticalCapacityCofinal.lean\#L838}{\nolinkurl{halfGreedySkippedPrecriticalSuffixSupply_iff_preTake}}.  \emph{Compared}, run \href{https://github.com/wcook04/plectis-erdos-lean/actions/runs/35674034595}{\texttt{35674034595}}.
\par\noindent\hangindent=1.5em Theorem~\ref{record:257bm-c6b}: \href{https://github.com/wcook04/plectis-erdos/blob/a25cb360bef8dd818dde14b5fb752244304af354/lean/Erdos249257/BooleanMobiusCriticalCapacityCofinal.lean\#L603}{\nolinkurl{halfGreedy_precriticalSuffix_lt_of_future_skip_after_takenBlock}}, \href{https://github.com/wcook04/plectis-erdos/blob/a25cb360bef8dd818dde14b5fb752244304af354/lean/Erdos249257/BooleanMobiusCriticalCapacityCofinal.lean\#L472}{\nolinkurl{precriticalCrossingTax_of_futureThreshold}}, \href{https://github.com/wcook04/plectis-erdos/blob/a25cb360bef8dd818dde14b5fb752244304af354/lean/Erdos249257/BooleanMobiusCriticalCapacityCofinal.lean\#L432}{\nolinkurl{sub_two_le_two_pow_sub_four}}.  \emph{Compared}, runs \href{https://github.com/wcook04/plectis-erdos-lean/actions/runs/35643815458}{\texttt{35643815458}}, \href{https://github.com/wcook04/plectis-erdos-lean/actions/runs/35674034595}{\texttt{35674034595}}.
\par\noindent\hangindent=1.5em Theorem~\ref{record:257bm-c7}: \href{https://github.com/wcook04/plectis-erdos/blob/a25cb360bef8dd818dde14b5fb752244304af354/lean/Erdos249257/BooleanMobiusSkipRow.lean\#L238}{\nolinkurl{exists_exactRowStrictUpperFill_of_skippedCoreSharpCapacity}}.  \emph{Compared}, run \href{https://github.com/wcook04/plectis-erdos-lean/actions/runs/35643815458}{\texttt{35643815458}}.
\par\noindent\hangindent=1.5em Theorem~\ref{record:257bm-c8}: \href{https://github.com/wcook04/plectis-erdos/blob/a25cb360bef8dd818dde14b5fb752244304af354/lean/Erdos249257/BooleanMobiusSkipRow.lean\#L332}{\nolinkurl{exactLocalMersenneHalfRow_two_mul_sub_two_of_skippedCoreSharpCapacity}}.  \emph{Compared}, run \href{https://github.com/wcook04/plectis-erdos-lean/actions/runs/35643815458}{\texttt{35643815458}}.
\par\noindent\hangindent=1.5em Proposition~\ref{record:257bm-c9}: \href{https://github.com/wcook04/plectis-erdos/blob/a25cb360bef8dd818dde14b5fb752244304af354/lean/Erdos249257/BooleanMobiusLocalRepair.lean\#L26}{\nolinkurl{localMersenneQuotient_eq_two_pow_sub_of_half_lt}}, \href{https://github.com/wcook04/plectis-erdos/blob/a25cb360bef8dd818dde14b5fb752244304af354/lean/Erdos249257/BooleanMobiusLocalRepair.lean\#L525}{\nolinkurl{exists_boolean_word_of_lt_two_pow}}.  \emph{Compared}, run \href{https://github.com/wcook04/plectis-erdos-lean/actions/runs/35643815458}{\texttt{35643815458}}.
\par\noindent\hangindent=1.5em Theorem~\ref{record:257bm-c10}: \href{https://github.com/wcook04/plectis-erdos/blob/a25cb360bef8dd818dde14b5fb752244304af354/lean/ErdosProblems/Erdos257/PaperCompleteR21/MersenneQuotientRowRecurrences.lean\#L264}{\nolinkurl{paper_exact_row_from_skipped_prefix}}.  \emph{Compared}, run \href{https://github.com/wcook04/plectis-erdos-lean/actions/runs/35643815458}{\texttt{35643815458}}.
\par\noindent\hangindent=1.5em Theorem~\ref{record:257bm-c11}: \href{https://github.com/wcook04/plectis-erdos/blob/a25cb360bef8dd818dde14b5fb752244304af354/lean/ErdosProblems/Erdos257/PaperCompleteR21/DyadicBandAndTwoSidedBounds.lean\#L26}{\nolinkurl{paper_critical_dyadic_band_index_unique}}, \href{https://github.com/wcook04/plectis-erdos/blob/a25cb360bef8dd818dde14b5fb752244304af354/lean/ErdosProblems/Erdos257/PaperCompleteR21/DyadicBandAndTwoSidedBounds.lean\#L52}{\nolinkurl{paper_critical_dyadic_boundary_is_smallest}}, \href{https://github.com/wcook04/plectis-erdos/blob/a25cb360bef8dd818dde14b5fb752244304af354/lean/ErdosProblems/Erdos257/PaperCompleteR21/DyadicBandAndTwoSidedBounds.lean\#L71}{\nolinkurl{paper_critical_dyadic_band_index_eq_top_of_le_two}}, \href{https://github.com/wcook04/plectis-erdos/blob/a25cb360bef8dd818dde14b5fb752244304af354/lean/ErdosProblems/Erdos257/PaperCompleteR21/DyadicBandAndTwoSidedBounds.lean\#L89}{\nolinkurl{paper_dyadic_band_escape_iff_single_test}}.  \emph{Compared}, run \href{https://github.com/wcook04/plectis-erdos-lean/actions/runs/35643815458}{\texttt{35643815458}}.
\par\noindent\hangindent=1.5em Proposition~\ref{record:257bm-c12}: \href{https://github.com/wcook04/plectis-erdos/blob/a25cb360bef8dd818dde14b5fb752244304af354/lean/Erdos249257/HalfCylinderMiddleCarryLowerBound.lean\#L3542}{\nolinkurl{seamUpperBranch_remainder_add_resetCharge_eq}}, \href{https://github.com/wcook04/plectis-erdos/blob/a25cb360bef8dd818dde14b5fb752244304af354/lean/Erdos249257/HalfUpperResetCriticalBand.lean\#L131}{\nolinkurl{seamUpperResetCharge_le}}.  \emph{Compared}, run \href{https://github.com/wcook04/plectis-erdos-lean/actions/runs/35935225572}{\texttt{35935225572}}.
\par\noindent\hangindent=1.5em Theorem~\ref{record:257bm-c14}: \href{https://github.com/wcook04/plectis-erdos/blob/a25cb360bef8dd818dde14b5fb752244304af354/lean/Erdos249257/HalfCylinderFullShellSeamBridge.lean\#L569}{\nolinkurl{skipped_fullShell_neg_iff_alignment_and_seamRemainder_pos}}, \href{https://github.com/wcook04/plectis-erdos/blob/a25cb360bef8dd818dde14b5fb752244304af354/lean/Erdos249257/HalfCylinderFullShellSeamBridge.lean\#L531}{\nolinkurl{greedyHalfFrozenMargin_fullShell_eq_neg_seamRemainder_of_alignment}}, \href{https://github.com/wcook04/plectis-erdos/blob/a25cb360bef8dd818dde14b5fb752244304af354/lean/Erdos249257/HalfCylinderFullShellSeamBridge.lean\#L671}{\nolinkurl{skippedSeamAlignmentZero_iff_skippedFullShellNonnegative}}.  \emph{Compared}, run \href{https://github.com/wcook04/plectis-erdos-lean/actions/runs/35840809568}{\texttt{35840809568}}.
\par\noindent\hangindent=1.5em Theorem~\ref{record:257bm-c15}: \href{https://github.com/wcook04/plectis-erdos/blob/a25cb360bef8dd818dde14b5fb752244304af354/lean/ErdosProblems/Erdos257/PaperCompleteR21/SeamEscapeAndTerminalStrip.lean\#L36}{\nolinkurl{paper_seam_escape_forces_remainder_band}}, \href{https://github.com/wcook04/plectis-erdos/blob/a25cb360bef8dd818dde14b5fb752244304af354/lean/ErdosProblems/Erdos257/PaperCompleteR21/SeamEscapeAndTerminalStrip.lean\#L47}{\nolinkurl{paper_seam_escape_implies_full_shell_nonnegative}}, \href{https://github.com/wcook04/plectis-erdos/blob/a25cb360bef8dd818dde14b5fb752244304af354/lean/ErdosProblems/Erdos257/PaperCompleteR21/SeamEscapeAndTerminalStrip.lean\#L61}{\nolinkurl{paper_seam_escape_implies_half_membership}}.  \emph{Compared}, run \href{https://github.com/wcook04/plectis-erdos-lean/actions/runs/35840809568}{\texttt{35840809568}}.
\par\noindent\hangindent=1.5em Theorem~\ref{record:257rig-c16}: \href{https://github.com/wcook04/plectis-erdos/blob/a25cb360bef8dd818dde14b5fb752244304af354/lean/Erdos249257/HalfCarryReachability.lean\#L919}{\nolinkurl{greedy_mobiusCenteredHalfCarry_nonneg}}, \href{https://github.com/wcook04/plectis-erdos/blob/a25cb360bef8dd818dde14b5fb752244304af354/lean/Erdos249257/HalfCarryReachability.lean\#L953}{\nolinkurl{greedy_half_infinite_of_mobiusCenteredHalfCarry_upperBound}}.  \emph{Compared}, run \href{https://github.com/wcook04/plectis-erdos-lean/actions/runs/35674034595}{\texttt{35674034595}}.
\par\noindent\hangindent=1.5em Theorem~\ref{record:257rig-c17} (the Lean statement is stronger): \href{https://github.com/wcook04/plectis-erdos/blob/a25cb360bef8dd818dde14b5fb752244304af354/lean/ErdosProblems/Erdos257/PaperCompleteR21/TerminalStripExactRowGap.lean\#L44}{\nolinkurl{paper_terminal_strip_forces_half_membership}}, \href{https://github.com/wcook04/plectis-erdos/blob/a25cb360bef8dd818dde14b5fb752244304af354/lean/ErdosProblems/Erdos257/PaperCompleteR21/TerminalStripExactRowGap.lean\#L88}{\nolinkurl{paper_integerHalfCarry_eq_two_pow_sub_localPrefixQuotient}}, \href{https://github.com/wcook04/plectis-erdos/blob/a25cb360bef8dd818dde14b5fb752244304af354/lean/ErdosProblems/Erdos257/PaperCompleteR21/TerminalStripExactRowGap.lean\#L106}{\nolinkurl{paper_exact_row_integerHalfCarry_eq_one}}, \href{https://github.com/wcook04/plectis-erdos/blob/a25cb360bef8dd818dde14b5fb752244304af354/lean/ErdosProblems/Erdos257/PaperCompleteR21/TerminalStripExactRowGap.lean\#L122}{\nolinkurl{paper_terminal_strip_witness_six}}, \href{https://github.com/wcook04/plectis-erdos/blob/a25cb360bef8dd818dde14b5fb752244304af354/lean/ErdosProblems/Erdos257/PaperCompleteR21/TerminalStripExactRowGap.lean\#L160}{\nolinkurl{paper_both_cofinal_statements_iff_half_membership}}.  \emph{Compared}, runs \href{https://github.com/wcook04/plectis-erdos-lean/actions/runs/35643815458}{\texttt{35643815458}}, \href{https://github.com/wcook04/plectis-erdos-lean/actions/runs/35674034595}{\texttt{35674034595}}.
\par\noindent\hangindent=1.5em Theorem~\ref{record:257rig-c18}: \href{https://github.com/wcook04/plectis-erdos/blob/a25cb360bef8dd818dde14b5fb752244304af354/lean/Erdos249257/CofinalStripReturn.lean\#L130}{\nolinkurl{greedy_half_infinite_of_cofinalStripReturn}}, \href{https://github.com/wcook04/plectis-erdos/blob/a25cb360bef8dd818dde14b5fb752244304af354/lean/Erdos249257/BooleanMobiusCriticalCapacityCofinal.lean\#L949}{\nolinkurl{halfGreedy_precriticalSuffix_lt_iff_futureSkipCoverage}}.  \emph{Compared}, run \href{https://github.com/wcook04/plectis-erdos-lean/actions/runs/35674034595}{\texttt{35674034595}}.
\par\noindent\hangindent=1.5em Theorem~\ref{record:257bm-c19} (the Lean statement is stronger): \href{https://github.com/wcook04/plectis-erdos/blob/a25cb360bef8dd818dde14b5fb752244304af354/lean/ErdosProblems/Erdos257/PaperCompleteR21/EventualNonnegativeMargin.lean\#L56}{\nolinkurl{paper_halfGreedyPrefixSupport_eq_greedy_inter_Icc}}, \href{https://github.com/wcook04/plectis-erdos/blob/a25cb360bef8dd818dde14b5fb752244304af354/lean/ErdosProblems/Erdos257/PaperCompleteR21/EventualNonnegativeMargin.lean\#L131}{\nolinkurl{paper_eventual_nonnegative_margin_equivalence}}, \href{https://github.com/wcook04/plectis-erdos/blob/a25cb360bef8dd818dde14b5fb752244304af354/lean/ErdosProblems/Erdos257/PaperCompleteR21/EventualNonnegativeMargin.lean\#L165}{\nolinkurl{paper_frozen_margin_normalised_value}}, \href{https://github.com/wcook04/plectis-erdos/blob/a25cb360bef8dd818dde14b5fb752244304af354/lean/ErdosProblems/Erdos257/PaperCompleteR21/EventualNonnegativeMargin.lean\#L180}{\nolinkurl{paper_frozen_margin_normalised_monotone}}, \href{https://github.com/wcook04/plectis-erdos/blob/a25cb360bef8dd818dde14b5fb752244304af354/lean/ErdosProblems/Erdos257/PaperCompleteR21/EventualNonnegativeMargin.lean\#L200}{\nolinkurl{paper_eta_eq_coeffTail_sub_carry}}, \href{https://github.com/wcook04/plectis-erdos/blob/a25cb360bef8dd818dde14b5fb752244304af354/lean/ErdosProblems/Erdos257/PaperCompleteR21/EventualNonnegativeMargin.lean\#L221}{\nolinkurl{paper_frozen_margin_normalised_tendsto}}, \href{https://github.com/wcook04/plectis-erdos/blob/a25cb360bef8dd818dde14b5fb752244304af354/lean/ErdosProblems/Erdos257/PaperCompleteR21/EventualNonnegativeMargin.lean\#L239}{\nolinkurl{paper_frozen_margin_limit_pos_iff}}, \href{https://github.com/wcook04/plectis-erdos/blob/a25cb360bef8dd818dde14b5fb752244304af354/lean/ErdosProblems/Erdos257/PaperCompleteR21/EventualNonnegativeMargin.lean\#L92}{\nolinkurl{paper_greedyHalfRemainder_ne_dyadicCap}}, \href{https://github.com/wcook04/plectis-erdos/blob/a25cb360bef8dd818dde14b5fb752244304af354/lean/ErdosProblems/Erdos257/PaperCompleteR21/EventualNonnegativeMargin.lean\#L258}{\nolinkurl{paper_coeffTail_le_index_add_two}}, \href{https://github.com/wcook04/plectis-erdos/blob/a25cb360bef8dd818dde14b5fb752244304af354/lean/ErdosProblems/Erdos257/PaperCompleteR21/EventualNonnegativeMargin.lean\#L266}{\nolinkurl{paper_effective_horizon_test}}, \href{https://github.com/wcook04/plectis-erdos/blob/a25cb360bef8dd818dde14b5fb752244304af354/lean/ErdosProblems/Erdos257/PaperCompleteR21/EventualNonnegativeMargin.lean\#L311}{\nolinkurl{paper_eta_hasRationalValue}}.  \emph{Compared}, runs \href{https://github.com/wcook04/plectis-erdos-lean/actions/runs/35643815458}{\texttt{35643815458}}, \href{https://github.com/wcook04/plectis-erdos-lean/actions/runs/35674034595}{\texttt{35674034595}}.
\par\noindent\hangindent=1.5em Theorem~\ref{record:257bm-c20}: \href{https://github.com/wcook04/plectis-erdos/blob/a25cb360bef8dd818dde14b5fb752244304af354/lean/Erdos249257/GreedyAchievementSet.lean\#L2583}{\nolinkurl{half_mem_mersenneAchievementSet_iff_greedySkippedSupport_infinite}}, \href{https://github.com/wcook04/plectis-erdos/blob/a25cb360bef8dd818dde14b5fb752244304af354/lean/Erdos249257/GreedyAchievementSet.lean\#L1528}{\nolinkurl{mem_mersenneAchievementSet_of_greedySkippedSupport_infinite}}.  \emph{Compared}, run \href{https://github.com/wcook04/plectis-erdos-lean/actions/runs/35674034595}{\texttt{35674034595}}.
\par\noindent\hangindent=1.5em Theorem~\ref{record:257bm-i1a}: \href{https://github.com/wcook04/plectis-erdos/blob/a25cb360bef8dd818dde14b5fb752244304af354/lean/ErdosProblems/Erdos257/PaperCompleteR21/MersenneQuotientRowRecurrences.lean\#L37}{\nolinkurl{paper_next_floor_quotient}}, \href{https://github.com/wcook04/plectis-erdos/blob/a25cb360bef8dd818dde14b5fb752244304af354/lean/ErdosProblems/Erdos257/PaperCompleteR21/MersenneQuotientRowRecurrences.lean\#L74}{\nolinkurl{paper_next_floor_quotient_no_fixed_point}}, \href{https://github.com/wcook04/plectis-erdos/blob/a25cb360bef8dd818dde14b5fb752244304af354/lean/ErdosProblems/Erdos257/PaperCompleteR21/MersenneQuotientRowRecurrences.lean\#L82}{\nolinkurl{paper_next_quotient_sum}}, \href{https://github.com/wcook04/plectis-erdos/blob/a25cb360bef8dd818dde14b5fb752244304af354/lean/ErdosProblems/Erdos257/PaperCompleteR21/MersenneQuotientRowRecurrences.lean\#L44}{\nolinkurl{paper_floor_quotient_geometric_sum}}.  \emph{Compared}, run \href{https://github.com/wcook04/plectis-erdos-lean/actions/runs/35643815458}{\texttt{35643815458}}.
\par\noindent\hangindent=1.5em Theorem~\ref{record:257bm-i1b}: \href{https://github.com/wcook04/plectis-erdos/blob/a25cb360bef8dd818dde14b5fb752244304af354/lean/Erdos249257/BooleanMobiusExactTransition.lean\#L131}{\nolinkurl{localPrefixQuotient_succ}}.  \emph{Compared}, run \href{https://github.com/wcook04/plectis-erdos-lean/actions/runs/35643815458}{\texttt{35643815458}}.
\par\noindent\hangindent=1.5em Theorem~\ref{record:257bm-i1c}: \href{https://github.com/wcook04/plectis-erdos/blob/a25cb360bef8dd818dde14b5fb752244304af354/lean/ErdosProblems/Erdos257/PaperCompleteR21/MersenneQuotientRowRecurrences.lean\#L92}{\nolinkurl{paper_signed_endpoint_recurrence}}, \href{https://github.com/wcook04/plectis-erdos/blob/a25cb360bef8dd818dde14b5fb752244304af354/lean/ErdosProblems/Erdos257/PaperCompleteR21/MersenneQuotientRowRecurrences.lean\#L99}{\nolinkurl{paper_endpoint_term_counts_divisors}}, \href{https://github.com/wcook04/plectis-erdos/blob/a25cb360bef8dd818dde14b5fb752244304af354/lean/ErdosProblems/Erdos257/PaperCompleteR21/MersenneQuotientRowRecurrences.lean\#L106}{\nolinkurl{paper_signed_endpoint_defect_succ}}, \href{https://github.com/wcook04/plectis-erdos/blob/a25cb360bef8dd818dde14b5fb752244304af354/lean/ErdosProblems/Erdos257/PaperCompleteR21/MersenneQuotientRowRecurrences.lean\#L114}{\nolinkurl{paper_repair_integer_eq_endpoint_defect}}.  \emph{Compared}, runs \href{https://github.com/wcook04/plectis-erdos-lean/actions/runs/35624228171}{\texttt{35624228171}}, \href{https://github.com/wcook04/plectis-erdos-lean/actions/runs/35643815458}{\texttt{35643815458}}.
\par\noindent\hangindent=1.5em Theorem~\ref{record:257bm-i5}: \href{https://github.com/wcook04/plectis-erdos/blob/a25cb360bef8dd818dde14b5fb752244304af354/lean/Erdos249257/BooleanMobiusSkippedCoreCriticalCapacity.lean\#L22}{\nolinkurl{localBinarySuffix_two_mul_sub_two_lt_criticalCapacity_iff}}.  \emph{Compared}, run \href{https://github.com/wcook04/plectis-erdos-lean/actions/runs/35643815458}{\texttt{35643815458}}.
\par\noindent\hangindent=1.5em Theorem~\ref{record:257bm-i6}: \href{https://github.com/wcook04/plectis-erdos/blob/a25cb360bef8dd818dde14b5fb752244304af354/lean/ErdosProblems/Erdos257/PaperCompleteR21/MersenneQuotientRowRecurrences.lean\#L126}{\nolinkurl{paper_unconditional_bound_one_extra_bit}}, \href{https://github.com/wcook04/plectis-erdos/blob/a25cb360bef8dd818dde14b5fb752244304af354/lean/ErdosProblems/Erdos257/PaperCompleteR21/MersenneQuotientRowRecurrences.lean\#L141}{\nolinkurl{paper_sharper_additive_estimate}}, \href{https://github.com/wcook04/plectis-erdos/blob/a25cb360bef8dd818dde14b5fb752244304af354/lean/ErdosProblems/Erdos257/PaperCompleteR21/MersenneQuotientRowRecurrences.lean\#L236}{\nolinkurl{paper_capacity_band_exclusion}}.  \emph{Compared}, run \href{https://github.com/wcook04/plectis-erdos-lean/actions/runs/35643815458}{\texttt{35643815458}}.
\par\noindent\hangindent=1.5em Theorem~\ref{record:257bm-i7}: \href{https://github.com/wcook04/plectis-erdos/blob/a25cb360bef8dd818dde14b5fb752244304af354/lean/Erdos249257/BooleanMobiusExactRowDoubling.lean\#L44}{\nolinkurl{localBinarySuffix_two_mul_sub_one_lt_upperWindow_of_exact_below}}, \href{https://github.com/wcook04/plectis-erdos/blob/a25cb360bef8dd818dde14b5fb752244304af354/lean/Erdos249257/BooleanMobiusExactRowDoubling.lean\#L185}{\nolinkurl{exists_exactRowStrictUpperExtension_two_mul_sub_one_of_exact_below}}.  \emph{Compared}, run \href{https://github.com/wcook04/plectis-erdos-lean/actions/runs/35643815458}{\texttt{35643815458}}.
\par\noindent\hangindent=1.5em Proposition~\ref{record:257bm-i-rank2}: \href{https://github.com/wcook04/plectis-erdos/blob/a25cb360bef8dd818dde14b5fb752244304af354/lean/Erdos249257/BooleanMobiusExactRowRankTwo.lean\#L23}{\nolinkurl{two_mem_of_exact_localMersenneQuotient}}.  \emph{Compared}, run \href{https://github.com/wcook04/plectis-erdos-lean/actions/runs/35643815458}{\texttt{35643815458}}.
\par\noindent\hangindent=1.5em Proposition~\ref{record:257bm-i9}: \href{https://github.com/wcook04/plectis-erdos/blob/a25cb360bef8dd818dde14b5fb752244304af354/lean/Erdos249257/BooleanMobiusGlobalRepair.lean\#L243}{\nolinkurl{abs_localMersennePrefixValue_sub_half_le}}.  \emph{Compared}, run \href{https://github.com/wcook04/plectis-erdos-lean/actions/runs/35643815458}{\texttt{35643815458}}.
\par\noindent\hangindent=1.5em Proposition~\ref{record:257bm-i10}: \href{https://github.com/wcook04/plectis-erdos/blob/a25cb360bef8dd818dde14b5fb752244304af354/lean/ErdosProblems/Erdos257/PaperCompleteR21/OddReciprocalDenominators.lean\#L54}{\nolinkurl{paper_finite_sum_inv_odd_den_odd}}, \href{https://github.com/wcook04/plectis-erdos/blob/a25cb360bef8dd818dde14b5fb752244304af354/lean/ErdosProblems/Erdos257/PaperCompleteR21/OddReciprocalDenominators.lean\#L73}{\nolinkurl{paper_finite_sum_inv_odd_ne_half}}, \href{https://github.com/wcook04/plectis-erdos/blob/a25cb360bef8dd818dde14b5fb752244304af354/lean/ErdosProblems/Erdos257/PaperCompleteR21/OddReciprocalDenominators.lean\#L94}{\nolinkurl{paper_finiteErdosSum_den_odd}}, \href{https://github.com/wcook04/plectis-erdos/blob/a25cb360bef8dd818dde14b5fb752244304af354/lean/ErdosProblems/Erdos257/PaperCompleteR21/OddReciprocalDenominators.lean\#L114}{\nolinkurl{paper_finite_support_series_ne_half}}, \href{https://github.com/wcook04/plectis-erdos/blob/a25cb360bef8dd818dde14b5fb752244304af354/lean/ErdosProblems/Erdos257/PaperCompleteR21/OddReciprocalDenominators.lean\#L142}{\nolinkurl{paper_half_representing_support_is_infinite}}.  \emph{Compared}, runs \href{https://github.com/wcook04/plectis-erdos-lean/actions/runs/35624228171}{\texttt{35624228171}}, \href{https://github.com/wcook04/plectis-erdos-lean/actions/runs/35643815458}{\texttt{35643815458}}.
\par\noindent\hangindent=1.5em Theorem~\ref{record:257bm-i11a}: \href{https://github.com/wcook04/plectis-erdos/blob/a25cb360bef8dd818dde14b5fb752244304af354/lean/Erdos249257/BooleanMobiusGreedyReduction.lean\#L918}{\nolinkurl{remainder_lt_gap_iff_eq_integerGreedyBits}}.  \emph{Compared}, run \href{https://github.com/wcook04/plectis-erdos-lean/actions/runs/35674034595}{\texttt{35674034595}}.
\par\noindent\hangindent=1.5em Proposition~\ref{record:257bm-i11b}: \href{https://github.com/wcook04/plectis-erdos/blob/a25cb360bef8dd818dde14b5fb752244304af354/lean/Erdos249257/BooleanMobiusGreedyReduction.lean\#L684}{\nolinkurl{localMersenneWeightsFrom_gapDominates}}, \href{https://github.com/wcook04/plectis-erdos/blob/a25cb360bef8dd818dde14b5fb752244304af354/lean/Erdos249257/BooleanMobiusGreedyReduction.lean\#L846}{\nolinkurl{localMersenneWeights_gapDominates_even}}, \href{https://github.com/wcook04/plectis-erdos/blob/a25cb360bef8dd818dde14b5fb752244304af354/lean/Erdos249257/BooleanMobiusGreedyReduction.lean\#L855}{\nolinkurl{localMersenneWeights_gapDominates_odd}}.  \emph{Compared}, run \href{https://github.com/wcook04/plectis-erdos-lean/actions/runs/35682858662}{\texttt{35682858662}}.
\par\noindent\hangindent=1.5em Proposition~\ref{record:257bm-i11c}: \href{https://github.com/wcook04/plectis-erdos/blob/a25cb360bef8dd818dde14b5fb752244304af354/lean/Erdos249257/BooleanMobiusGreedyReduction.lean\#L997}{\nolinkurl{localMersenneHalfTarget_lower_word_eq_greedy_and_remainder_eq}}.  \emph{Compared}, run \href{https://github.com/wcook04/plectis-erdos-lean/actions/runs/35682858662}{\texttt{35682858662}}.
\par\noindent\hangindent=1.5em Theorem~\ref{record:257bm-i12}: \href{https://github.com/wcook04/plectis-erdos/blob/a25cb360bef8dd818dde14b5fb752244304af354/lean/Erdos249257/BooleanMobiusCriticalCapacityGeometric.lean\#L27}{\nolinkurl{localMersenneQuotient_eq_geometric}}, \href{https://github.com/wcook04/plectis-erdos/blob/a25cb360bef8dd818dde14b5fb752244304af354/lean/Erdos249257/BooleanMobiusCriticalCapacityGeometric.lean\#L196}{\nolinkurl{localBinarySuffix_two_mul_sub_two_lt_criticalCapacity_iff_geometric}}, \href{https://github.com/wcook04/plectis-erdos/blob/a25cb360bef8dd818dde14b5fb752244304af354/lean/Erdos249257/BooleanMobiusCriticalCapacityGeometric.lean\#L208}{\nolinkurl{localBinarySuffix_two_mul_sub_two_lt_criticalCapacity_iff_geometricCore}}.  \emph{Compared}, run \href{https://github.com/wcook04/plectis-erdos-lean/actions/runs/35643815458}{\texttt{35643815458}}.
\par\noindent\hangindent=1.5em Theorem~\ref{record:257bm-i2} (the Lean statement is stronger): \href{https://github.com/wcook04/plectis-erdos/blob/a25cb360bef8dd818dde14b5fb752244304af354/lean/Erdos249257/GenericTailOrbitRigidity.lean\#L85}{\nolinkurl{binaryCoeffTail_le}}, \href{https://github.com/wcook04/plectis-erdos/blob/a25cb360bef8dd818dde14b5fb752244304af354/lean/Erdos249257/GenericTailOrbitRigidity.lean\#L91}{\nolinkurl{binaryCoeffTail_div_pow_tendsto_zero}}.  \emph{Compared}, run \href{https://github.com/wcook04/plectis-erdos-lean/actions/runs/35643815458}{\texttt{35643815458}}.
\par\noindent\hangindent=1.5em Theorem~\ref{record:257bm-i-t7} (the Lean statement is stronger): \href{https://github.com/wcook04/plectis-erdos/blob/a25cb360bef8dd818dde14b5fb752244304af354/lean/Erdos249257/GenericTailOrbitRigidity.lean\#L435}{\nolinkurl{not_irrational_binaryCoeffSeries_iff_exists_temperedBinaryOrbit}}.  \emph{Compared}, run \href{https://github.com/wcook04/plectis-erdos-lean/actions/runs/35544127144}{\texttt{35544127144}}.
\par\noindent\hangindent=1.5em Theorem~\ref{record:257bm-i-mob} (the Lean statement is stronger): \href{https://github.com/wcook04/plectis-erdos/blob/a25cb360bef8dd818dde14b5fb752244304af354/lean/Erdos249257/BooleanMobiusCarry.lean\#L95}{\nolinkurl{moebius_mul_supportCoeffAF}}, \href{https://github.com/wcook04/plectis-erdos/blob/a25cb360bef8dd818dde14b5fb752244304af354/lean/Erdos249257/BooleanMobiusCarry.lean\#L112}{\nolinkurl{mobius_supportCoeff_eq_one_iff}}, \href{https://github.com/wcook04/plectis-erdos/blob/a25cb360bef8dd818dde14b5fb752244304af354/lean/Erdos249257/BooleanMobiusCarry.lean\#L118}{\nolinkurl{mobius_supportCoeff_boolean}}.  \emph{Compared}, run \href{https://github.com/wcook04/plectis-erdos-lean/actions/runs/35624228171}{\texttt{35624228171}}.
\par\noindent\hangindent=1.5em Theorem~\ref{record:257bm-i-bridge} (the Lean statement is stronger): \href{https://github.com/wcook04/plectis-erdos/blob/a25cb360bef8dd818dde14b5fb752244304af354/lean/Erdos249257/BooleanMobiusCarry.lean\#L377}{\nolinkurl{erdosSupportSeries_two_eq_binaryCoeffSeries}}, \href{https://github.com/wcook04/plectis-erdos/blob/a25cb360bef8dd818dde14b5fb752244304af354/lean/Erdos249257/CertificateKernel.lean\#L8860}{\nolinkurl{supportCoeff_le_card_divisors}}, \href{https://github.com/wcook04/plectis-erdos/blob/a25cb360bef8dd818dde14b5fb752244304af354/lean/Erdos249257/CertificateKernel.lean\#L8868}{\nolinkurl{supportCoeff_le_self}}, \href{https://github.com/wcook04/plectis-erdos/blob/a25cb360bef8dd818dde14b5fb752244304af354/lean/Erdos249257/BooleanMobiusCarry.lean\#L384}{\nolinkurl{erdosSupportSeries_rational_iff_exists_temperedCarry}}.  \emph{Compared}, run \href{https://github.com/wcook04/plectis-erdos-lean/actions/runs/35624228171}{\texttt{35624228171}}.
\par\noindent\hangindent=1.5em Theorem~\ref{record:257rig-i2} (the Lean statement is stronger): \href{https://github.com/wcook04/plectis-erdos/blob/a25cb360bef8dd818dde14b5fb752244304af354/lean/Erdos249257/RationalSupportCarrySkeleton.lean\#L2124}{\nolinkurl{one_lt_reciprocalMass_of_dyadic_support_fraction_of_two_pos_mem}}, \href{https://github.com/wcook04/plectis-erdos/blob/a25cb360bef8dd818dde14b5fb752244304af354/lean/Erdos249257/RationalSupportCarrySkeleton.lean\#L2210}{\nolinkurl{dyadic_support_fraction_reciprocalMass_diverges_or_gt_one}}.  \emph{Compared}, runs \href{https://github.com/wcook04/plectis-erdos-lean/actions/runs/35544127144}{\texttt{35544127144}}, \href{https://github.com/wcook04/plectis-erdos-lean/actions/runs/35624228171}{\texttt{35624228171}}.
\par\noindent\hangindent=1.5em Theorem~\ref{record:257rig-i3} (the Lean statement is stronger): \href{https://github.com/wcook04/plectis-erdos/blob/a25cb360bef8dd818dde14b5fb752244304af354/lean/Erdos249257/RationalSupportCarrySkeleton.lean\#L2237}{\nolinkurl{one_add_mul_card_le_two_mul_shifted_state}}, \href{https://github.com/wcook04/plectis-erdos/blob/a25cb360bef8dd818dde14b5fb752244304af354/lean/Erdos249257/RationalSupportCarrySkeleton.lean\#L2327}{\nolinkurl{shifted_state_unbounded_of_infinite_support}}, \href{https://github.com/wcook04/plectis-erdos/blob/a25cb360bef8dd818dde14b5fb752244304af354/lean/Erdos249257/RationalSupportCarrySkeleton.lean\#L2383}{\nolinkurl{exists_unbounded_shifted_odd_tail_nat_state_of_support_fraction}}.  \emph{Compared}, runs \href{https://github.com/wcook04/plectis-erdos-lean/actions/runs/35544127144}{\texttt{35544127144}}, \href{https://github.com/wcook04/plectis-erdos-lean/actions/runs/35624228171}{\texttt{35624228171}}.
\par\noindent\hangindent=1.5em Theorem~\ref{record:257rig-i4a}: \href{https://github.com/wcook04/plectis-erdos/blob/a25cb360bef8dd818dde14b5fb752244304af354/lean/ErdosProblems/Erdos257/PaperCompleteR21/SkipSafetyAndDivisorZeroRuns.lean\#L219}{\nolinkurl{paper_zero_run_le_eps_logb}}, \href{https://github.com/wcook04/plectis-erdos/blob/a25cb360bef8dd818dde14b5fb752244304af354/lean/ErdosProblems/Erdos257/PaperCompleteR21/SkipSafetyAndDivisorZeroRuns.lean\#L237}{\nolinkurl{paper_zero_run_le_of_mem}}.  \emph{Compared}, run \href{https://github.com/wcook04/plectis-erdos-lean/actions/runs/35643815458}{\texttt{35643815458}}.
\par\noindent\hangindent=1.5em Proposition~\ref{record:257rig-i4b} (the Lean statement is stronger): \href{https://github.com/wcook04/plectis-erdos/blob/a25cb360bef8dd818dde14b5fb752244304af354/lean/Erdos249257/SublogDivisorCoverage.lean\#L107}{\nolinkurl{card_divisors_pow_le_divisorSubpowerConst_pow_mul}}, \href{https://github.com/wcook04/plectis-erdos/blob/a25cb360bef8dd818dde14b5fb752244304af354/lean/Erdos249257/SublogDivisorCoverage.lean\#L142}{\nolinkurl{card_divisors_le_divisorSubpowerConst_mul_rpow}}.  \emph{Compared}, run \href{https://github.com/wcook04/plectis-erdos-lean/actions/runs/35624228171}{\texttt{35624228171}}.
\par\noindent\hangindent=1.5em Proposition~\ref{record:257rig-i5}: \href{https://github.com/wcook04/plectis-erdos/blob/a25cb360bef8dd818dde14b5fb752244304af354/lean/Erdos249257/MaximalOmegaLayer.lean\#L29}{\nolinkurl{primePowerLayer_comm}}, \href{https://github.com/wcook04/plectis-erdos/blob/a25cb360bef8dd818dde14b5fb752244304af354/lean/Erdos249257/MaximalOmegaLayer.lean\#L39}{\nolinkurl{mixedPrimePowerLayerTwo_supportCoeffInt}}, \href{https://github.com/wcook04/plectis-erdos/blob/a25cb360bef8dd818dde14b5fb752244304af354/lean/Erdos249257/MaximalOmegaLayer.lean\#L65}{\nolinkurl{mixedPrimePowerLayerTwo_twelve_fixture}}.  \emph{Compared}, runs \href{https://github.com/wcook04/plectis-erdos-lean/actions/runs/35624228171}{\texttt{35624228171}}, \href{https://github.com/wcook04/plectis-erdos-lean/actions/runs/35643815458}{\texttt{35643815458}}.
\par\noindent\hangindent=1.5em Theorem~\ref{record:257bm-i-cross2}: \href{https://github.com/wcook04/plectis-erdos/blob/a25cb360bef8dd818dde14b5fb752244304af354/lean/Erdos249257/BooleanMobiusExactRowCrossing.lean\#L31}{\nolinkurl{exists_first_localMersenne_crossing}}.  \emph{Compared}, run \href{https://github.com/wcook04/plectis-erdos-lean/actions/runs/35643815458}{\texttt{35643815458}}.
\par\noindent\hangindent=1.5em Theorem~\ref{record:257hg-i6}: \href{https://github.com/wcook04/plectis-erdos/blob/a25cb360bef8dd818dde14b5fb752244304af354/lean/Erdos249257/HalfCarryReachability.lean\#L871}{\nolinkurl{integerHalfCarry_eq_scaled_residual_add_tail}}, \href{https://github.com/wcook04/plectis-erdos/blob/a25cb360bef8dd818dde14b5fb752244304af354/lean/Erdos249257/HalfCylinderFinalMiddleCellEscape.lean\#L94}{\nolinkurl{mobiusCenteredHalfCarry_nonneg_of_supportSeries_lt_half}}.  \emph{Compared}, run \href{https://github.com/wcook04/plectis-erdos-lean/actions/runs/35674034595}{\texttt{35674034595}}.
\par\noindent\hangindent=1.5em Theorem~\ref{record:257hg-i7}: \href{https://github.com/wcook04/plectis-erdos/blob/a25cb360bef8dd818dde14b5fb752244304af354/lean/Erdos249257/HalfCylinderSkippedEndpointClassifier.lean\#L246}{\nolinkurl{halfGreedy_skipped_endpoint_trichotomy}}.  \emph{Compared}, run \href{https://github.com/wcook04/plectis-erdos-lean/actions/runs/35840809568}{\texttt{35840809568}}.
\par\noindent\hangindent=1.5em Theorem~\ref{record:257hg-i2}: \href{https://github.com/wcook04/plectis-erdos/blob/a25cb360bef8dd818dde14b5fb752244304af354/lean/Erdos249257/HalfGapMass.lean\#L40}{\nolinkurl{mersenneGap_le}}, \href{https://github.com/wcook04/plectis-erdos/blob/a25cb360bef8dd818dde14b5fb752244304af354/lean/Erdos249257/HalfGapMass.lean\#L69}{\nolinkurl{summable_mersenneGap_shift}}, \href{https://github.com/wcook04/plectis-erdos/blob/a25cb360bef8dd818dde14b5fb752244304af354/lean/Erdos249257/HalfGapMass.lean\#L83}{\nolinkurl{mersenneGap_tail_le}}, \href{https://github.com/wcook04/plectis-erdos/blob/a25cb360bef8dd818dde14b5fb752244304af354/lean/Erdos249257/HalfGapMass.lean\#L104}{\nolinkurl{tendsto_mersenneGap_tail_zero}}.  \emph{Compared}, run \href{https://github.com/wcook04/plectis-erdos-lean/actions/runs/35643815458}{\texttt{35643815458}}.
\par\noindent\hangindent=1.5em Theorem~\ref{record:257hg-i4}: \href{https://github.com/wcook04/plectis-erdos/blob/a25cb360bef8dd818dde14b5fb752244304af354/lean/ErdosProblems/Erdos257/PaperCompleteR21/SkipSafetyAndDivisorZeroRuns.lean\#L30}{\nolinkurl{paper_dyadic_skip_test_iff}}, \href{https://github.com/wcook04/plectis-erdos/blob/a25cb360bef8dd818dde14b5fb752244304af354/lean/ErdosProblems/Erdos257/PaperCompleteR21/SkipSafetyAndDivisorZeroRuns.lean\#L54}{\nolinkurl{paper_sharp_skip_safe_lb3}}, \href{https://github.com/wcook04/plectis-erdos/blob/a25cb360bef8dd818dde14b5fb752244304af354/lean/ErdosProblems/Erdos257/PaperCompleteR21/SkipSafetyAndDivisorZeroRuns.lean\#L62}{\nolinkurl{paper_sharp_skip_safe_actual_tail}}, \href{https://github.com/wcook04/plectis-erdos/blob/a25cb360bef8dd818dde14b5fb752244304af354/lean/ErdosProblems/Erdos257/PaperCompleteR21/SkipSafetyAndDivisorZeroRuns.lean\#L71}{\nolinkurl{paper_three_channel_margin_identity}}, \href{https://github.com/wcook04/plectis-erdos/blob/a25cb360bef8dd818dde14b5fb752244304af354/lean/ErdosProblems/Erdos257/PaperCompleteR21/SkipSafetyAndDivisorZeroRuns.lean\#L92}{\nolinkurl{paper_sharp_weaker_than_dyadic}}, \href{https://github.com/wcook04/plectis-erdos/blob/a25cb360bef8dd818dde14b5fb752244304af354/lean/ErdosProblems/Erdos257/PaperCompleteR21/SkipSafetyAndDivisorZeroRuns.lean\#L97}{\nolinkurl{paper_sharp_strictly_weaker_realizable}}, \href{https://github.com/wcook04/plectis-erdos/blob/a25cb360bef8dd818dde14b5fb752244304af354/lean/ErdosProblems/Erdos257/PaperCompleteR21/SkipSafetyAndDivisorZeroRuns.lean\#L104}{\nolinkurl{paper_sharp_gives_three_over_three_t_sub_one}}, \href{https://github.com/wcook04/plectis-erdos/blob/a25cb360bef8dd818dde14b5fb752244304af354/lean/ErdosProblems/Erdos257/PaperCompleteR21/SkipSafetyAndDivisorZeroRuns.lean\#L131}{\nolinkurl{paper_two_channels_insufficient}}, \href{https://github.com/wcook04/plectis-erdos/blob/a25cb360bef8dd818dde14b5fb752244304af354/lean/ErdosProblems/Erdos257/PaperCompleteR21/SkipSafetyAndDivisorZeroRuns.lean\#L138}{\nolinkurl{paper_exact_mass_threshold}}.  \emph{Compared}, run \href{https://github.com/wcook04/plectis-erdos-lean/actions/runs/35643815458}{\texttt{35643815458}}.
\par\noindent\hangindent=1.5em Proposition~\ref{record:257hg-i5}: \href{https://github.com/wcook04/plectis-erdos/blob/a25cb360bef8dd818dde14b5fb752244304af354/lean/Erdos249257/HalfGreedyFatalGap.lean\#L135}{\nolinkurl{unitNumerator_skipSafe}}, \href{https://github.com/wcook04/plectis-erdos/blob/a25cb360bef8dd818dde14b5fb752244304af354/lean/Erdos249257/HalfGreedyFatalGap.lean\#L161}{\nolinkurl{two_le_of_fatal}}, \href{https://github.com/wcook04/plectis-erdos/blob/a25cb360bef8dd818dde14b5fb752244304af354/lean/Erdos249257/HalfGreedyFatalGap.lean\#L173}{\nolinkurl{three_le_of_fatal_of_odd}}, \href{https://github.com/wcook04/plectis-erdos/blob/a25cb360bef8dd818dde14b5fb752244304af354/lean/Erdos249257/HalfGreedyFatalGap.lean\#L246}{\nolinkurl{unitNumerator_skipSafe_actualTail}}.  \emph{Compared}, run \href{https://github.com/wcook04/plectis-erdos-lean/actions/runs/35682858662}{\texttt{35682858662}}.
\par\noindent\hangindent=1.5em Theorem~\ref{record:257bm-i13}: \href{https://github.com/wcook04/plectis-erdos/blob/a25cb360bef8dd818dde14b5fb752244304af354/lean/ErdosProblems/Erdos257/PaperCompleteR21/DyadicBandAndTwoSidedBounds.lean\#L105}{\nolinkurl{paper_two_sided_dyadic_bound}}.  \emph{Compared}, run \href{https://github.com/wcook04/plectis-erdos-lean/actions/runs/35935225572}{\texttt{35935225572}}.
\par\noindent\hangindent=1.5em Proposition~\ref{record:257bm-i14}: \href{https://github.com/wcook04/plectis-erdos/blob/a25cb360bef8dd818dde14b5fb752244304af354/lean/ErdosProblems/Erdos257/PaperCompleteR21/UpperResetBandCertificate.lean\#L29}{\nolinkurl{paper_finite_band_check_thirteen_to_thirty}}, \href{https://github.com/wcook04/plectis-erdos/blob/a25cb360bef8dd818dde14b5fb752244304af354/lean/ErdosProblems/Erdos257/PaperCompleteR21/UpperResetBandCertificate.lean\#L44}{\nolinkurl{paper_successor_remainders_fourteen_through_thirtyone}}, \href{https://github.com/wcook04/plectis-erdos/blob/a25cb360bef8dd818dde14b5fb752244304af354/lean/ErdosProblems/Erdos257/PaperCompleteR21/UpperResetBandCertificate.lean\#L77}{\nolinkurl{paper_universal_band_condition_unfolded}}, \href{https://github.com/wcook04/plectis-erdos/blob/a25cb360bef8dd818dde14b5fb752244304af354/lean/ErdosProblems/Erdos257/PaperCompleteR21/UpperResetBandCertificate.lean\#L70}{\nolinkurl{paper_universal_band_condition_would_close_half}}.  \emph{Compared}, run \href{https://github.com/wcook04/plectis-erdos-lean/actions/runs/35935225572}{\texttt{35935225572}}.
\par\noindent\hangindent=1.5em Theorem~\ref{record:257bm-i15}: \href{https://github.com/wcook04/plectis-erdos/blob/a25cb360bef8dd818dde14b5fb752244304af354/lean/ErdosProblems/Erdos257/PaperCompleteR21/UpperResetBandCertificate.lean\#L97}{\nolinkurl{paper_perturbed_order_preservation}}, \href{https://github.com/wcook04/plectis-erdos/blob/a25cb360bef8dd818dde14b5fb752244304af354/lean/ErdosProblems/Erdos257/PaperCompleteR21/UpperResetBandCertificate.lean\#L106}{\nolinkurl{paper_perturbed_prefixChoice_maximal}}, \href{https://github.com/wcook04/plectis-erdos/blob/a25cb360bef8dd818dde14b5fb752244304af354/lean/ErdosProblems/Erdos257/PaperCompleteR21/UpperResetBandCertificate.lean\#L124}{\nolinkurl{paper_perturbed_nextRemainder_three_branches}}, \href{https://github.com/wcook04/plectis-erdos/blob/a25cb360bef8dd818dde14b5fb752244304af354/lean/ErdosProblems/Erdos257/PaperCompleteR21/UpperResetBandCertificate.lean\#L140}{\nolinkurl{paper_perturbed_separation_global_maximality}}, \href{https://github.com/wcook04/plectis-erdos/blob/a25cb360bef8dd818dde14b5fb752244304af354/lean/ErdosProblems/Erdos257/PaperCompleteR21/UpperResetBandCertificate.lean\#L158}{\nolinkurl{paper_perturbed_two_stage_is_global_maximum}}, \href{https://github.com/wcook04/plectis-erdos/blob/a25cb360bef8dd818dde14b5fb752244304af354/lean/ErdosProblems/Erdos257/PaperCompleteR21/UpperResetBandCertificate.lean\#L245}{\nolinkurl{paper_perturbed_weak_cap_counterexample}}, \href{https://github.com/wcook04/plectis-erdos/blob/a25cb360bef8dd818dde14b5fb752244304af354/lean/ErdosProblems/Erdos257/PaperCompleteR21/UpperResetBandCertificate.lean\#L259}{\nolinkurl{prefixChoice_eq_below}}.  \emph{Compared}, run \href{https://github.com/wcook04/plectis-erdos-lean/actions/runs/35840809568}{\texttt{35840809568}}.
\par\noindent\hangindent=1.5em Theorem~\ref{record:257bm-i16}: \href{https://github.com/wcook04/plectis-erdos/blob/a25cb360bef8dd818dde14b5fb752244304af354/lean/Erdos249257/HalfTrappingReturnCarry.lean\#L124}{\nolinkurl{overlappingReverseCarryWords_carryDifference_eq_twoPow_mul_odd}}, \href{https://github.com/wcook04/plectis-erdos/blob/a25cb360bef8dd818dde14b5fb752244304af354/lean/Erdos249257/HalfTrappingReturnCarry.lean\#L191}{\nolinkurl{overlappingReverseCarryWords_twoPow_le_realBound}}, \href{https://github.com/wcook04/plectis-erdos/blob/a25cb360bef8dd818dde14b5fb752244304af354/lean/Erdos249257/HalfTrappingReturnCarry.lean\#L244}{\nolinkurl{overlappingMidpointReturns_twoPow_le_realBound}}, \href{https://github.com/wcook04/plectis-erdos/blob/a25cb360bef8dd818dde14b5fb752244304af354/lean/ErdosProblems/Erdos257/PaperCompleteR21/SeamRowGapAndCarry.lean\#L240}{\nolinkurl{paper_reverse_carry_word}}.  \emph{Compared}, runs \href{https://github.com/wcook04/plectis-erdos-lean/actions/runs/35643815458}{\texttt{35643815458}}, \href{https://github.com/wcook04/plectis-erdos-lean/actions/runs/35674034595}{\texttt{35674034595}}.
\par\noindent\hangindent=1.5em Proposition~\ref{record:257bm-i17} (the Lean statement is stronger): \href{https://github.com/wcook04/plectis-erdos/blob/a25cb360bef8dd818dde14b5fb752244304af354/lean/Erdos249257/HalfCylinderFiniteShadow.lean\#L642}{\nolinkurl{supportCoeff_insert_divisor}}, \href{https://github.com/wcook04/plectis-erdos/blob/a25cb360bef8dd818dde14b5fb752244304af354/lean/Erdos249257/HalfDivisorUnitDrop.lean\#L20}{\nolinkurl{supportCoeff_extend_true_eq_false_add_one_at_double}}, \href{https://github.com/wcook04/plectis-erdos/blob/a25cb360bef8dd818dde14b5fb752244304af354/lean/Erdos249257/HalfDivisorUnitDrop.lean\#L35}{\nolinkurl{supportCoeff_boundaryPair_unitDrop_at_double}}.  \emph{Compared}, run \href{https://github.com/wcook04/plectis-erdos-lean/actions/runs/35643815458}{\texttt{35643815458}}.
\par\noindent\hangindent=1.5em Theorem~\ref{record:257bm-k1}: \href{https://github.com/wcook04/plectis-erdos/blob/a25cb360bef8dd818dde14b5fb752244304af354/lean/ErdosProblems/Erdos257/PaperCompleteR21/ExactRowDichotomyCountermodels.lean\#L84}{\nolinkurl{paper_bounded_double_or_recycle_countermodel}}, \href{https://github.com/wcook04/plectis-erdos/blob/a25cb360bef8dd818dde14b5fb752244304af354/lean/ErdosProblems/Erdos257/PaperCompleteR21/ExactRowDichotomyCountermodels.lean\#L102}{\nolinkurl{paper_exists_seeded_bounded_double_or_recycle_model}}.  \emph{Compared}, run \href{https://github.com/wcook04/plectis-erdos-lean/actions/runs/35643815458}{\texttt{35643815458}}.
\par\noindent\hangindent=1.5em Theorem~\ref{record:257bm-k-dich}: \href{https://github.com/wcook04/plectis-erdos/blob/a25cb360bef8dd818dde14b5fb752244304af354/lean/ErdosProblems/Erdos257/PaperCompleteR21/ExactRowDichotomyCountermodels.lean\#L33}{\nolinkurl{paper_exact_row_double_or_recycle}}, \href{https://github.com/wcook04/plectis-erdos/blob/a25cb360bef8dd818dde14b5fb752244304af354/lean/ErdosProblems/Erdos257/PaperCompleteR21/ExactRowDichotomyCountermodels.lean\#L41}{\nolinkurl{paper_finite_row_value_ne_half}}, \href{https://github.com/wcook04/plectis-erdos/blob/a25cb360bef8dd818dde14b5fb752244304af354/lean/ErdosProblems/Erdos257/PaperCompleteR21/ExactRowDichotomyCountermodels.lean\#L55}{\nolinkurl{paper_exact_row_example_six_and_eleven}}.  \emph{Compared}, run \href{https://github.com/wcook04/plectis-erdos-lean/actions/runs/35643815458}{\texttt{35643815458}}.
\par\noindent\hangindent=1.5em Proposition~\ref{record:257bm-k2}: \href{https://github.com/wcook04/plectis-erdos/blob/a25cb360bef8dd818dde14b5fb752244304af354/lean/ErdosProblems/Erdos257/PaperCompleteR21/ExactRowDichotomyCountermodels.lean\#L118}{\nolinkurl{paper_fractional_mass_bound_not_necessary}}, \href{https://github.com/wcook04/plectis-erdos/blob/a25cb360bef8dd818dde14b5fb752244304af354/lean/ErdosProblems/Erdos257/PaperCompleteR21/ExactRowDichotomyCountermodels.lean\#L139}{\nolinkurl{paper_fractional_mass_bound_suffices_for_sharp_capacity}}.  \emph{Compared}, run \href{https://github.com/wcook04/plectis-erdos-lean/actions/runs/35643815458}{\texttt{35643815458}}.
\par\noindent\hangindent=1.5em Theorem~\ref{record:257bm-k4}: \href{https://github.com/wcook04/plectis-erdos/blob/a25cb360bef8dd818dde14b5fb752244304af354/lean/ErdosProblems/Erdos257/PaperCompleteR21/ExactRowDichotomyCountermodels.lean\#L163}{\nolinkurl{paper_skipped_core_recycling_witness_bounded}}, \href{https://github.com/wcook04/plectis-erdos/blob/a25cb360bef8dd818dde14b5fb752244304af354/lean/ErdosProblems/Erdos257/PaperCompleteR21/ExactRowDichotomyCountermodels.lean\#L172}{\nolinkurl{paper_returning_endpoint_may_fail_to_grow}}, \href{https://github.com/wcook04/plectis-erdos/blob/a25cb360bef8dd818dde14b5fb752244304af354/lean/ErdosProblems/Erdos257/PaperCompleteR21/ExactRowDichotomyCountermodels.lean\#L102}{\nolinkurl{paper_exists_seeded_bounded_double_or_recycle_model}}, \href{https://github.com/wcook04/plectis-erdos/blob/a25cb360bef8dd818dde14b5fb752244304af354/lean/ErdosProblems/Erdos257/PaperCompleteR21/ExactRowDichotomyCountermodels.lean\#L184}{\nolinkurl{paper_protected_row_crossing_beyond_cutoff}}, \href{https://github.com/wcook04/plectis-erdos/blob/a25cb360bef8dd818dde14b5fb752244304af354/lean/ErdosProblems/Erdos257/PaperCompleteR21/ExactRowDichotomyCountermodels.lean\#L218}{\nolinkurl{paper_protected_row_endpoint_growth}}.  \emph{Compared}, runs \href{https://github.com/wcook04/plectis-erdos-lean/actions/runs/35643815458}{\texttt{35643815458}}, \href{https://github.com/wcook04/plectis-erdos-lean/actions/runs/35674034595}{\texttt{35674034595}}.
\par\noindent\hangindent=1.5em Theorem~\ref{record:257rig-k6}: \href{https://github.com/wcook04/plectis-erdos/blob/a25cb360bef8dd818dde14b5fb752244304af354/lean/ErdosProblems/Erdos257/PaperCompleteR21/ExactRowDichotomyCountermodels.lean\#L231}{\nolinkurl{paper_critical_crossing_support_is_greedy_prefix}}.  \emph{Compared}, run \href{https://github.com/wcook04/plectis-erdos-lean/actions/runs/35674034595}{\texttt{35674034595}}.
\par\noindent\hangindent=1.5em Theorem~\ref{record:257bm-k9}: \href{https://github.com/wcook04/plectis-erdos/blob/a25cb360bef8dd818dde14b5fb752244304af354/lean/ErdosProblems/Erdos257/PaperCompleteR21/LinearChannelAndMiddleCellExclusion.lean\#L37}{\nolinkurl{paper_relationInvariant_channels_rank_le_one}}, \href{https://github.com/wcook04/plectis-erdos/blob/a25cb360bef8dd818dde14b5fb752244304af354/lean/ErdosProblems/Erdos257/PaperCompleteR21/LinearChannelAndMiddleCellExclusion.lean\#L63}{\nolinkurl{paper_relationInvariant_channels_det_eq_zero}}.  \emph{Compared}, run \href{https://github.com/wcook04/plectis-erdos-lean/actions/runs/35643815458}{\texttt{35643815458}}.
\par\noindent\hangindent=1.5em Theorem~\ref{record:257hg-k12}: \href{https://github.com/wcook04/plectis-erdos/blob/a25cb360bef8dd818dde14b5fb752244304af354/lean/ErdosProblems/Erdos257/PaperCompleteR21/LinearChannelAndMiddleCellExclusion.lean\#L100}{\nolinkurl{paper_final_middle_cell_ne_neg_three}}, \href{https://github.com/wcook04/plectis-erdos/blob/a25cb360bef8dd818dde14b5fb752244304af354/lean/ErdosProblems/Erdos257/PaperCompleteR21/LinearChannelAndMiddleCellExclusion.lean\#L118}{\nolinkurl{paper_final_middle_cell_remaining_negative_values}}, \href{https://github.com/wcook04/plectis-erdos/blob/a25cb360bef8dd818dde14b5fb752244304af354/lean/ErdosProblems/Erdos257/PaperCompleteR21/LinearChannelAndMiddleCellExclusion.lean\#L155}{\nolinkurl{paper_mobiusCenteredHalfCarry_add_two}}, \href{https://github.com/wcook04/plectis-erdos/blob/a25cb360bef8dd818dde14b5fb752244304af354/lean/ErdosProblems/Erdos257/PaperCompleteR21/LinearChannelAndMiddleCellExclusion.lean\#L163}{\nolinkurl{paper_cofiniteRightTail_ne_zero_centeredEndpoint}}.  \emph{Compared}, run \href{https://github.com/wcook04/plectis-erdos-lean/actions/runs/35935225572}{\texttt{35935225572}}.
\par\noindent\hangindent=1.5em Proposition (The unsafe middle range is exactly three integers): \href{https://github.com/wcook04/plectis-erdos/blob/a25cb360bef8dd818dde14b5fb752244304af354/lean/ErdosProblems/Erdos257/PaperCompleteR21/LinearChannelAndMiddleCellExclusion.lean\#L83}{\nolinkurl{paper_unsafe_middle_range_is_three_integers}}.  \emph{Compared}, run \href{https://github.com/wcook04/plectis-erdos-lean/actions/runs/35840809568}{\texttt{35840809568}}.
\par\noindent\hangindent=1.5em Proposition (Exact Lebesgue-measure dichotomy): \href{https://github.com/wcook04/plectis-erdos/blob/a25cb360bef8dd818dde14b5fb752244304af354/lean/ErdosProblems/Erdos257/PaperCompleteR21/SharedPrefixFamiliesAndMeasureDichotomy.lean\#L68}{\nolinkurl{paper_volume_supportedMersenneAchievementSet_dichotomy}}.  \emph{Compared}, run \href{https://github.com/wcook04/plectis-erdos-lean/actions/runs/35643815458}{\texttt{35643815458}}.
\par\noindent\hangindent=1.5em Proposition~\ref{prop:one-orbit}: \href{https://github.com/wcook04/plectis-erdos/blob/a25cb360bef8dd818dde14b5fb752244304af354/lean/ErdosProblems/Erdos257/PaperCompleteR21/GreedyOrbitNoTies.lean\#L272}{\nolinkurl{paper_one_orbit_stability}}, \href{https://github.com/wcook04/plectis-erdos/blob/a25cb360bef8dd818dde14b5fb752244304af354/lean/ErdosProblems/Erdos257/PaperCompleteR21/GreedyOrbitNoTies.lean\#L182}{\nolinkurl{approx_orbit_induction}}, \href{https://github.com/wcook04/plectis-erdos/blob/a25cb360bef8dd818dde14b5fb752244304af354/lean/ErdosProblems/Erdos257/PaperCompleteR21/GreedyOrbitNoTies.lean\#L166}{\nolinkurl{tailGreedyRemainder_mersenne}}.  \emph{Compared}, run \href{https://github.com/wcook04/plectis-erdos-lean/actions/runs/35674034595}{\texttt{35674034595}}.
\par\noindent\hangindent=1.5em Lemma~\ref{lem:no-ties}: \href{https://github.com/wcook04/plectis-erdos/blob/a25cb360bef8dd818dde14b5fb752244304af354/lean/ErdosProblems/Erdos257/PaperCompleteR21/GreedyOrbitNoTies.lean\#L126}{\nolinkurl{paper_no_ties}}, \href{https://github.com/wcook04/plectis-erdos/blob/a25cb360bef8dd818dde14b5fb752244304af354/lean/ErdosProblems/Erdos257/PaperCompleteR21/GreedyOrbitNoTies.lean\#L71}{\nolinkurl{paper_no_ties_take}}, \href{https://github.com/wcook04/plectis-erdos/blob/a25cb360bef8dd818dde14b5fb752244304af354/lean/ErdosProblems/Erdos257/PaperCompleteR21/GreedyOrbitNoTies.lean\#L114}{\nolinkurl{paper_no_ties_skip}}, \href{https://github.com/wcook04/plectis-erdos/blob/a25cb360bef8dd818dde14b5fb752244304af354/lean/ErdosProblems/Erdos257/PaperCompleteR21/GreedyOrbitNoTies.lean\#L44}{\nolinkurl{irrational_mersenneTail}}.  \emph{Compared}, runs \href{https://github.com/wcook04/plectis-erdos-lean/actions/runs/35643815458}{\texttt{35643815458}}, \href{https://github.com/wcook04/plectis-erdos-lean/actions/runs/35674034595}{\texttt{35674034595}}.
\par\noindent\hangindent=1.5em Lemma~\ref{lem:tr-forced-greedy} (the Lean statement is stronger): \href{https://github.com/wcook04/plectis-erdos/blob/a25cb360bef8dd818dde14b5fb752244304af354/lean/ErdosProblems/Erdos257/PaperCompleteR21/TruncatedRungGreedyDecision.lean\#L486}{\nolinkurl{paper_forced_greedy_unique_support_and_criterion}}, \href{https://github.com/wcook04/plectis-erdos/blob/a25cb360bef8dd818dde14b5fb752244304af354/lean/ErdosProblems/Erdos257/PaperCompleteR21/TruncatedRungGreedyDecision.lean\#L567}{\nolinkurl{paper_forced_greedy_low_ranks}}.  \emph{Compared}, runs \href{https://github.com/wcook04/plectis-erdos-lean/actions/runs/35643815458}{\texttt{35643815458}}, \href{https://github.com/wcook04/plectis-erdos-lean/actions/runs/35674034595}{\texttt{35674034595}}.
\par\noindent\hangindent=1.5em Lemma~\ref{lem:tr-parity}: \href{https://github.com/wcook04/plectis-erdos/blob/a25cb360bef8dd818dde14b5fb752244304af354/lean/ErdosProblems/Erdos257/PaperCompleteR21/TruncatedRungWitnessHorizon.lean\#L176}{\nolinkurl{paper_parity_excludes_finite_support}}.  \emph{Compared}, run \href{https://github.com/wcook04/plectis-erdos-lean/actions/runs/35643815458}{\texttt{35643815458}}.
\par\noindent\hangindent=1.5em Theorem~\ref{thm:tr-witness-exclusion}: \href{https://github.com/wcook04/plectis-erdos/blob/a25cb360bef8dd818dde14b5fb752244304af354/lean/ErdosProblems/Erdos257/PaperCompleteR21/TruncatedRungWitnessHorizon.lean\#L250}{\nolinkurl{paper_witness_exclusion}}.  \emph{Compared}, run \href{https://github.com/wcook04/plectis-erdos-lean/actions/runs/35643815458}{\texttt{35643815458}}.
\par\noindent\hangindent=1.5em Corollary~\ref{cor:tr-half-lcm}: \href{https://github.com/wcook04/plectis-erdos/blob/a25cb360bef8dd818dde14b5fb752244304af354/lean/ErdosProblems/Erdos257/PaperCompleteR21/TruncatedRungWitnessHorizon.lean\#L467}{\nolinkurl{paper_half_lcm_horizon}}.  \emph{Compared}, run \href{https://github.com/wcook04/plectis-erdos-lean/actions/runs/35643815458}{\texttt{35643815458}}.
\par\noindent\hangindent=1.5em Lemma~\ref{lem:tr-mod12}: \href{https://github.com/wcook04/plectis-erdos/blob/a25cb360bef8dd818dde14b5fb752244304af354/lean/ErdosProblems/Erdos257/PaperCompleteR21/TruncatedRungWitnessHorizon.lean\#L544}{\nolinkurl{paper_mod_twelve_filter}}.  \emph{Compared}, run \href{https://github.com/wcook04/plectis-erdos-lean/actions/runs/35643815458}{\texttt{35643815458}}.
\par\noindent\hangindent=1.5em Theorem~\ref{thm:tr-finite-decision}: \href{https://github.com/wcook04/plectis-erdos/blob/a25cb360bef8dd818dde14b5fb752244304af354/lean/ErdosProblems/Erdos257/PaperCompleteR21/TruncatedRungGreedyDecision.lean\#L896}{\nolinkurl{paper_rung_finite_decision}}.  \emph{Compared}, run \href{https://github.com/wcook04/plectis-erdos-lean/actions/runs/35674034595}{\texttt{35674034595}}.
\par\noindent\hangindent=1.5em Theorem (A: a sufficient lower bound at every reset): \href{https://github.com/wcook04/plectis-erdos/blob/a25cb360bef8dd818dde14b5fb752244304af354/lean/ErdosProblems/Erdos257/PaperCompleteR21/ResetSqrtEscapeHalfMembership.lean\#L595}{\nolinkurl{paper_theoremA_half_membership}}, \href{https://github.com/wcook04/plectis-erdos/blob/a25cb360bef8dd818dde14b5fb752244304af354/lean/ErdosProblems/Erdos257/PaperCompleteR21/ResetSqrtEscapeHalfMembership.lean\#L545}{\nolinkurl{paperResetSqrtEscape_iff_square}}, \href{https://github.com/wcook04/plectis-erdos/blob/a25cb360bef8dd818dde14b5fb752244304af354/lean/ErdosProblems/Erdos257/PaperCompleteR21/ResetSqrtEscapeHalfMembership.lean\#L121}{\nolinkurl{paper_theoremA_crossing_bound_square_le}}, \href{https://github.com/wcook04/plectis-erdos/blob/a25cb360bef8dd818dde14b5fb752244304af354/lean/ErdosProblems/Erdos257/PaperCompleteR21/ResetSqrtEscapeHalfMembership.lean\#L261}{\nolinkurl{paper_theoremA_right_branch_forces_small_deviation}}, \href{https://github.com/wcook04/plectis-erdos/blob/a25cb360bef8dd818dde14b5fb752244304af354/lean/ErdosProblems/Erdos257/PaperCompleteR21/ResetSqrtEscapeHalfMembership.lean\#L526}{\nolinkurl{half_mem_mersenneAchievementSet_of_resetSqrtEscape}}.  \emph{Compared}, run \href{https://github.com/wcook04/plectis-erdos-lean/actions/runs/35935225572}{\texttt{35935225572}}.
\par\noindent\hangindent=1.5em Theorem (C ;  the one-sided finite decision boundary) (the Lean statement is stronger): \href{https://github.com/wcook04/plectis-erdos/blob/a25cb360bef8dd818dde14b5fb752244304af354/lean/ErdosProblems/Erdos257/PaperCompleteR21/CentredCompletionAndDecisionBoundary.lean\#L77}{\nolinkurl{paper_one_sided_finite_decision_boundary}}.  \emph{Compared}, run \href{https://github.com/wcook04/plectis-erdos-lean/actions/runs/35674034595}{\texttt{35674034595}}.
\par\noindent\hangindent=1.5em Proposition~\ref{prop:2adic-nogo}: \href{https://github.com/wcook04/plectis-erdos/blob/a25cb360bef8dd818dde14b5fb752244304af354/lean/ErdosProblems/Erdos257/PaperCompleteR21/CentredCompletionAndDecisionBoundary.lean\#L34}{\nolinkurl{paper_centred_completion_of_fixed_precision}}.  \emph{Compared}, run \href{https://github.com/wcook04/plectis-erdos-lean/actions/runs/35643815458}{\texttt{35643815458}}.
\par\noindent\hangindent=1.5em Proposition~\ref{prop:upper-unconditional}: \href{https://github.com/wcook04/plectis-erdos/blob/a25cb360bef8dd818dde14b5fb752244304af354/lean/ErdosProblems/Erdos257/PaperCompleteR21/BranchCellHorizonExclusions.lean\#L31}{\nolinkurl{paper_upper_branch_needs_no_exceptional_cell}}.  \emph{Compared}, run \href{https://github.com/wcook04/plectis-erdos-lean/actions/runs/35840809568}{\texttt{35840809568}}.
\par\noindent\hangindent=1.5em Theorem~\ref{thm:cd-neg3-impossible}: \href{https://github.com/wcook04/plectis-erdos/blob/a25cb360bef8dd818dde14b5fb752244304af354/lean/ErdosProblems/Erdos257/PaperCompleteR21/BranchCellHorizonExclusions.lean\#L50}{\nolinkurl{paper_final_middle_cell_at_least_neg_two}}.  \emph{Compared}, run \href{https://github.com/wcook04/plectis-erdos-lean/actions/runs/35840809568}{\texttt{35840809568}}.
\par\noindent\hangindent=1.5em Corollary~\ref{cor:cd-remaining}: \href{https://github.com/wcook04/plectis-erdos/blob/a25cb360bef8dd818dde14b5fb752244304af354/lean/ErdosProblems/Erdos257/PaperCompleteR21/BranchCellHorizonExclusions.lean\#L81}{\nolinkurl{paper_final_middle_cell_remaining_cells}}.  \emph{Compared}, run \href{https://github.com/wcook04/plectis-erdos-lean/actions/runs/35840809568}{\texttt{35840809568}}.
\par\noindent\hangindent=1.5em Proposition~\ref{prop:finite-state-nogo} (the Lean statement is stronger): \href{https://github.com/wcook04/plectis-erdos/blob/a25cb360bef8dd818dde14b5fb752244304af354/lean/ErdosProblems/Erdos257/PaperCompleteR21/BalancedPulseFanOutCount.lean\#L113}{\nolinkurl{paper_pulse_family_no_autonomous_decoder}}, \href{https://github.com/wcook04/plectis-erdos/blob/a25cb360bef8dd818dde14b5fb752244304af354/lean/ErdosProblems/Erdos257/PaperCompleteR21/BalancedPulseFanOutCount.lean\#L61}{\nolinkurl{paper_balanced_pulse_fanout_is_radius_succ}}, \href{https://github.com/wcook04/plectis-erdos/blob/a25cb360bef8dd818dde14b5fb752244304af354/lean/ErdosProblems/Erdos257/PaperCompleteR21/BalancedPulseFanOutCount.lean\#L95}{\nolinkurl{paper_balanced_pulse_fanout_unbounded_corrected}}, \href{https://github.com/wcook04/plectis-erdos/blob/a25cb360bef8dd818dde14b5fb752244304af354/lean/ErdosProblems/Erdos257/PaperCompleteR21/BalancedPulseFanOutCount.lean\#L123}{\nolinkurl{paper_pulse_family_finite_state_card}}, \href{https://github.com/wcook04/plectis-erdos/blob/a25cb360bef8dd818dde14b5fb752244304af354/lean/ErdosProblems/Erdos257/PaperCompleteR21/BalancedPulseFanOutCount.lean\#L52}{\nolinkurl{balancedPulseCoeff_injective}}, \href{https://github.com/wcook04/plectis-erdos/blob/a25cb360bef8dd818dde14b5fb752244304af354/lean/Erdos249257/GenericTailOrbitRigidity.lean\#L200}{\nolinkurl{balancedPulse_weighted_pair}}.  \emph{Compared}, runs \href{https://github.com/wcook04/plectis-erdos-lean/actions/runs/35624228171}{\texttt{35624228171}}, \href{https://github.com/wcook04/plectis-erdos-lean/actions/runs/35643815458}{\texttt{35643815458}}.
\par\noindent\hangindent=1.5em Proposition~\ref{prop:mobius-nogo}: \href{https://github.com/wcook04/plectis-erdos/blob/a25cb360bef8dd818dde14b5fb752244304af354/lean/ErdosProblems/Erdos257/PaperCompleteR21/MobiusSignAndFiniteCertificates.lean\#L70}{\nolinkurl{paper_mobius_support_overshoots_half}}, \href{https://github.com/wcook04/plectis-erdos/blob/a25cb360bef8dd818dde14b5fb752244304af354/lean/ErdosProblems/Erdos257/PaperCompleteR21/MobiusSignAndFiniteCertificates.lean\#L48}{\nolinkurl{paper_first_positiveMobius_tail_term}}.  \emph{Compared}, run \href{https://github.com/wcook04/plectis-erdos-lean/actions/runs/35643815458}{\texttt{35643815458}}.
\par\noindent\hangindent=1.5em Proposition~\ref{prop:finite-boolSupport-and-onesided}: \href{https://github.com/wcook04/plectis-erdos/blob/a25cb360bef8dd818dde14b5fb752244304af354/lean/ErdosProblems/Erdos257/PaperCompleteR21/MobiusSignAndFiniteCertificates.lean\#L101}{\nolinkurl{paper_finite_support_and_onesided_certificate}}.  \emph{Compared}, run \href{https://github.com/wcook04/plectis-erdos-lean/actions/runs/35674034595}{\texttt{35674034595}}.
\par\noindent\hangindent=1.5em Proposition~\ref{prop:carry-survivor-extinction}: \href{https://github.com/wcook04/plectis-erdos/blob/a25cb360bef8dd818dde14b5fb752244304af354/lean/ErdosProblems/Erdos257/PaperCompleteR21/TotientPeriodCertificateSupply.lean\#L40}{\nolinkurl{paper_carry_survivor_extinction}}, \href{https://github.com/wcook04/plectis-erdos/blob/a25cb360bef8dd818dde14b5fb752244304af354/lean/ErdosProblems/Erdos257/PaperCompleteR21/TotientPeriodCertificateSupply.lean\#L23}{\nolinkurl{paper_periodLcm_is_prefix_lcm}}.  \emph{Compared}, run \href{https://github.com/wcook04/plectis-erdos-lean/actions/runs/35674034595}{\texttt{35674034595}}.
\par\noindent\hangindent=1.5em Lemma~\ref{lem:scalar-localization}: \href{https://github.com/wcook04/plectis-erdos/blob/a25cb360bef8dd818dde14b5fb752244304af354/lean/ErdosProblems/Erdos257/PaperCompleteR21/ScalarLocalisationHeightObstruction.lean\#L29}{\nolinkurl{paper_scalar_localization}}, \href{https://github.com/wcook04/plectis-erdos/blob/a25cb360bef8dd818dde14b5fb752244304af354/lean/ErdosProblems/Erdos257/PaperCompleteR21/ScalarLocalisationHeightObstruction.lean\#L53}{\nolinkurl{paper_scalar_localization_size_bound}}, \href{https://github.com/wcook04/plectis-erdos/blob/a25cb360bef8dd818dde14b5fb752244304af354/lean/ErdosProblems/Erdos257/PaperCompleteR21/ScalarLocalisationHeightObstruction.lean\#L62}{\nolinkurl{paper_scalar_localization_zero_degenerate}}.  \emph{Compared}, run \href{https://github.com/wcook04/plectis-erdos-lean/actions/runs/35643815458}{\texttt{35643815458}}.
\par\noindent\hangindent=1.5em Corollary~\ref{cor:mersenne-height}: \href{https://github.com/wcook04/plectis-erdos/blob/a25cb360bef8dd818dde14b5fb752244304af354/lean/ErdosProblems/Erdos257/PaperCompleteR21/ScalarLocalisationHeightObstruction.lean\#L77}{\nolinkurl{paper_mersenne_height}}.  \emph{Compared}, run \href{https://github.com/wcook04/plectis-erdos-lean/actions/runs/35643815458}{\texttt{35643815458}}.
\par\noindent\hangindent=1.5em Proposition~\ref{prop:critical-band-index}: \href{https://github.com/wcook04/plectis-erdos/blob/a25cb360bef8dd818dde14b5fb752244304af354/lean/ErdosProblems/Erdos257/PaperCompleteR21/CriticalDyadicBandCollapse.lean\#L20}{\nolinkurl{paper_critical_band_index_collapse}}.  \emph{Compared}, run \href{https://github.com/wcook04/plectis-erdos-lean/actions/runs/35935225572}{\texttt{35935225572}}.
\par\noindent\hangindent=1.5em Lemma~\ref{lem:odometer} (the Lean statement is stronger): \href{https://github.com/wcook04/plectis-erdos/blob/a25cb360bef8dd818dde14b5fb752244304af354/lean/ErdosProblems/Erdos257/PaperCompleteR21/ShortWindowDivisorPhase.lean\#L417}{\nolinkurl{theta_eq_tsum_divisorResidue}}, \href{https://github.com/wcook04/plectis-erdos/blob/a25cb360bef8dd818dde14b5fb752244304af354/lean/ErdosProblems/Erdos257/PaperCompleteR21/ShortWindowDivisorPhase.lean\#L442}{\nolinkurl{theta_eq_tsum_geometricForm}}, \href{https://github.com/wcook04/plectis-erdos/blob/a25cb360bef8dd818dde14b5fb752244304af354/lean/ErdosProblems/Erdos257/PaperCompleteR21/ShortWindowDivisorPhase.lean\#L368}{\nolinkurl{theta_eq_divisorResidueSum}}, \href{https://github.com/wcook04/plectis-erdos/blob/a25cb360bef8dd818dde14b5fb752244304af354/lean/ErdosProblems/Erdos257/PaperCompleteR21/ShortWindowDivisorPhase.lean\#L378}{\nolinkurl{theta_eq_geometricForm}}, \href{https://github.com/wcook04/plectis-erdos/blob/a25cb360bef8dd818dde14b5fb752244304af354/lean/ErdosProblems/Erdos257/PaperCompleteR21/ShortWindowDivisorPhase.lean\#L116}{\nolinkurl{residue_condition_iff}}, \href{https://github.com/wcook04/plectis-erdos/blob/a25cb360bef8dd818dde14b5fb752244304af354/lean/ErdosProblems/Erdos257/PaperCompleteR21/ShortWindowDivisorPhase.lean\#L303}{\nolinkurl{card_divisors_sub_one}}, \href{https://github.com/wcook04/plectis-erdos/blob/a25cb360bef8dd818dde14b5fb752244304af354/lean/ErdosProblems/Erdos257/PaperCompleteR21/ShortWindowDivisorPhase.lean\#L132}{\nolinkurl{iLeast_mem_Icc}}, \href{https://github.com/wcook04/plectis-erdos/blob/a25cb360bef8dd818dde14b5fb752244304af354/lean/ErdosProblems/Erdos257/PaperCompleteR21/ShortWindowDivisorPhase.lean\#L138}{\nolinkurl{dvd_add_iLeast}}, \href{https://github.com/wcook04/plectis-erdos/blob/a25cb360bef8dd818dde14b5fb752244304af354/lean/ErdosProblems/Erdos257/PaperCompleteR21/ShortWindowDivisorPhase.lean\#L147}{\nolinkurl{not_dvd_of_lt_iLeast}}, \href{https://github.com/wcook04/plectis-erdos/blob/a25cb360bef8dd818dde14b5fb752244304af354/lean/ErdosProblems/Erdos257/PaperCompleteR21/ShortWindowDivisorPhase.lean\#L160}{\nolinkurl{iLeast_congr}}, \href{https://github.com/wcook04/plectis-erdos/blob/a25cb360bef8dd818dde14b5fb752244304af354/lean/ErdosProblems/Erdos257/PaperCompleteR21/ShortWindowDivisorPhase.lean\#L278}{\nolinkurl{mCount_eq_zero_of_lt_iLeast}}, \href{https://github.com/wcook04/plectis-erdos/blob/a25cb360bef8dd818dde14b5fb752244304af354/lean/ErdosProblems/Erdos257/PaperCompleteR21/ShortWindowDivisorPhase.lean\#L293}{\nolinkurl{geometric_term_eq_zero_of_lt_iLeast}}, \href{https://github.com/wcook04/plectis-erdos/blob/a25cb360bef8dd818dde14b5fb752244304af354/lean/ErdosProblems/Erdos257/PaperCompleteR21/ShortWindowDivisorPhase.lean\#L339}{\nolinkurl{theta_eq_Psi}}, \href{https://github.com/wcook04/plectis-erdos/blob/a25cb360bef8dd818dde14b5fb752244304af354/lean/ErdosProblems/Erdos257/PaperCompleteR21/ShortWindowDivisorPhase.lean\#L467}{\nolinkurl{Psi_eq_of_residues_eq}}, \href{https://github.com/wcook04/plectis-erdos/blob/a25cb360bef8dd818dde14b5fb752244304af354/lean/ErdosProblems/Erdos257/PaperCompleteR21/ShortWindowDivisorPhase.lean\#L523}{\nolinkurl{theta_one_one}}, \href{https://github.com/wcook04/plectis-erdos/blob/a25cb360bef8dd818dde14b5fb752244304af354/lean/ErdosProblems/Erdos257/PaperCompleteR21/ShortWindowDivisorPhase.lean\#L528}{\nolinkurl{theta_one_three}}, \href{https://github.com/wcook04/plectis-erdos/blob/a25cb360bef8dd818dde14b5fb752244304af354/lean/ErdosProblems/Erdos257/PaperCompleteR21/ShortWindowDivisorPhase.lean\#L538}{\nolinkurl{residue_cutoff_reading_fails}}.  \emph{Compared}, run \href{https://github.com/wcook04/plectis-erdos-lean/actions/runs/35643815458}{\texttt{35643815458}}.
\par\noindent\hangindent=1.5em Lemma~\ref{lem:sqwitness}: \href{https://github.com/wcook04/plectis-erdos/blob/a25cb360bef8dd818dde14b5fb752244304af354/lean/ErdosProblems/Erdos257/PaperCompleteR21/SquareDepthAndHalfMembershipEquivalences.lean\#L27}{\nolinkurl{paper_square_depth_terminal_bound}}, \href{https://github.com/wcook04/plectis-erdos/blob/a25cb360bef8dd818dde14b5fb752244304af354/lean/ErdosProblems/Erdos257/PaperCompleteR20/CarryCollapseCorrespondence.lean\#L8}{\nolinkurl{square_depth_witness}}.  \emph{Compared}, run \href{https://github.com/wcook04/plectis-erdos-lean/actions/runs/35674034595}{\texttt{35674034595}}.
\par\noindent\hangindent=1.5em Proposition~\ref{prop:squarefree}: \href{https://github.com/wcook04/plectis-erdos/blob/a25cb360bef8dd818dde14b5fb752244304af354/lean/ErdosProblems/Erdos257/PaperCompleteR21/SquarefreeSupportEngineCeiling.lean\#L20}{\nolinkurl{paper_squarefree_support_engine_ceiling}}.  \emph{Compared}, run \href{https://github.com/wcook04/plectis-erdos-lean/actions/runs/35624228171}{\texttt{35624228171}}.
\par\noindent\hangindent=1.5em Proposition~\ref{prop:cpgs-equiv}: \href{https://github.com/wcook04/plectis-erdos/blob/a25cb360bef8dd818dde14b5fb752244304af354/lean/ErdosProblems/Erdos257/PaperCompleteR21/SquareDepthAndHalfMembershipEquivalences.lean\#L44}{\nolinkurl{paper_cpgs_equiv}}.  \emph{Compared}, run \href{https://github.com/wcook04/plectis-erdos-lean/actions/runs/35674034595}{\texttt{35674034595}}.
\par\noindent\hangindent=1.5em Proposition~\ref{prop:strip-equiv}: \href{https://github.com/wcook04/plectis-erdos/blob/a25cb360bef8dd818dde14b5fb752244304af354/lean/ErdosProblems/Erdos257/PaperCompleteR21/SquareDepthAndHalfMembershipEquivalences.lean\#L130}{\nolinkurl{paper_terminal_strip_equiv}}, \href{https://github.com/wcook04/plectis-erdos/blob/a25cb360bef8dd818dde14b5fb752244304af354/lean/ErdosProblems/Erdos257/PaperCompleteR21/SquareDepthAndHalfMembershipEquivalences.lean\#L112}{\nolinkurl{paper_relaxed_constant_six_every_depth}}.  \emph{Compared}, run \href{https://github.com/wcook04/plectis-erdos-lean/actions/runs/35674034595}{\texttt{35674034595}}.
\par\noindent\hangindent=1.5em Theorem~\ref{thm:257-logarithmic-counterexample}: \href{https://github.com/wcook04/plectis-erdos/blob/a25cb360bef8dd818dde14b5fb752244304af354/lean/ErdosProblems/Erdos257/PaperCompleteR21/ArithmeticCounterexampleAssembly.lean\#L1645}{\nolinkurl{arithmetic_logarithmic_counterexample}}, \href{https://github.com/wcook04/plectis-erdos/blob/a25cb360bef8dd818dde14b5fb752244304af354/lean/ErdosProblems/Erdos257/PaperCompleteR21/ArithmeticCounterexampleAssembly.lean\#L1859}{\nolinkurl{no_absolute_dyadic_kappaOne_constant}}.  \emph{Compared}, run \href{https://github.com/wcook04/plectis-erdos-lean/actions/runs/35624228171}{\texttt{35624228171}}.
\par\noindent\hangindent=1.5em Corollary~\ref{cor:257-logarithmic-separation}: \href{https://github.com/wcook04/plectis-erdos/blob/a25cb360bef8dd818dde14b5fb752244304af354/lean/ErdosProblems/Erdos257/PaperCompleteR21/ArithmeticCoverLowerBoundPaperForm.lean\#L40}{\nolinkurl{sup_condExceedProb_le_paperCoverCost}}, \href{https://github.com/wcook04/plectis-erdos/blob/a25cb360bef8dd818dde14b5fb752244304af354/lean/ErdosProblems/Erdos257/PaperCompleteR21/ArithmeticCounterexampleAssembly.lean\#L1900}{\nolinkurl{exists_support_paperCoverCost_ge}}, \href{https://github.com/wcook04/plectis-erdos/blob/a25cb360bef8dd818dde14b5fb752244304af354/lean/ErdosProblems/Erdos257/PaperCompleteR21/ArithmeticCounterexampleAssembly.lean\#L1912}{\nolinkurl{no_absolute_paperCoverCost_kappaOne_constant}}.  \emph{Compared}, run \href{https://github.com/wcook04/plectis-erdos-lean/actions/runs/35840809568}{\texttt{35840809568}}.
\par\noindent\hangindent=1.5em Proposition~\ref{prop:257-logarithmic-initial-interval} (the Lean statement is stronger): \href{https://github.com/wcook04/plectis-erdos/blob/a25cb360bef8dd818dde14b5fb752244304af354/lean/ErdosProblems/Erdos257/PaperCompleteR21/LogarithmicInitialInterval.lean\#L491}{\nolinkurl{logarithmic_initial_interval}}.  \emph{Compared}, run \href{https://github.com/wcook04/plectis-erdos-lean/actions/runs/35624228171}{\texttt{35624228171}}.
\par\endgroup
% END GENERATED CONCORDANCE

\subsection{A finite estimate on common observation scales}\label{sec:257-common-kernel}
For $1<B\le2$, positive integers $L,d,M$, and an integer $R\ge0$, put
\[
 w_{B,d}(n)=\frac{B^{n\bmod d}}{B^d-1},\qquad
 \mathscr D_{L;R,M}F
 =\frac1M\sum_{j=R}^{R+M-1}\frac1{2^j}
      \sum_{m=1}^{2^j}F(Lm).
\]
The finite estimate
\begin{equation}
 \mathscr D_{L;R,M}w_{B,d}
 \le\frac{1+4L/M}{d(B-1)}
 \label{eq:257-mixed-finite-kernel}
\end{equation}
equivalently bounds $(B-1)\mathscr D_{L;R,M}w_{B,d}$ by
$(1+4L/M)/d$.  This normalized bound is uniform as $B\downarrow1$;
the unnormalized right side grows like $(B-1)^{-1}$.
The ratio $L/M$ measures the cost of incomplete modular periods.

To prove it, fix $T=2^j$ and average over $1\le m\le T$.
When $d\le LT$, the complete orbit has length $d/g$, with $g=(L,d)$,
and total weight $1/(B^g-1)$.  Complete cycles and one remaining piece give
\[
 \frac1T\sum_{m=1}^T w_{B,d}(Lm)
 \le\frac{g}{d(B^g-1)}+\frac1{T(B^g-1)}
 \le\frac1{d(B-1)}+\frac1{T(B-1)}.
\]
Across dyadic lengths satisfying $d\le L2^j$, the reciprocal-length
errors sum to at most $2L/[d(B-1)]$.
When $d>2LT$, there is no wrap and $2Lm\le d-1$; hence
$\sum_{i=0}^{d-1}B^i\ge dB^{(d-1)/2}\ge dB^{Lm}$.
Each atom is then at most $1/[d(B-1)]$.
Finally, at most one dyadic length satisfies $LT<d\le2LT$.
For that length the geometric sum and convexity give
\[
 \frac1T\sum_{m=1}^T w_{B,d}(Lm)
 =\frac{B^L}{T(B^L-1)}\frac{B^{LT}-1}{B^d-1}
 \le\frac{LB^L}{d(B^L-1)}
 \le\frac{2L}{d(B-1)}.
\]
Summing the main terms and these two error budgets proves~\eqref{eq:257-mixed-finite-kernel}.
The release snapshot contains a proof body for the full estimate in
\href{https://github.com/wcook04/plectis-erdos-lean/blob/52f29ad173b04e3bac941b3663f2b9aebe5de0bb/ErdosProblems/Erdos257/PaperCompleteR8/DyadicKernel.lean#L207}{the finite dyadic averaging estimate}.
Its complete-cycle and no-wrap scalar ingredients are in
\href{https://github.com/wcook04/plectis-erdos-lean/blob/52f29ad173b04e3bac941b3663f2b9aebe5de0bb/ErdosProblems/Erdos257/PaperCompleteR7/CoverKernel.lean#L130}{the complete-period and no-wrap estimates}.
These are source locators, not a fresh compilation or axiom-audit receipt;
the argument displayed here is an ordinary proof.
Nonnegative interchange permits summation against any coefficients $c_d$
with $\sum_dc_d/d<\infty$.

\subsection{Positive divisor majorants}
For a finite set $F\subseteq\Npos$, put
$f_F(n)=\#\{a\in F:a\mid n\}$.  The next theorem bounds a fractional power
of this count by a nonnegative sum over divisors.

\begin{thm}[A summable family of divisor majorants]\label{thm:257-variable-fractional-cover}
For each $j\ge1$, let $F_j\subseteq\Npos$ be finite, let
$0<\alpha_j\le1$, and let $c_{j,d}\ge0$ satisfy
\[
 f_{F_j}(n)^{\alpha_j}\le\sum_{d\mid n}c_{j,d}\quad(n\ge1).
\]
Set $C_j=\sum_{d\ge1}c_{j,d}/d$.  If
\begin{equation}
 \sum_{j\ge1}\frac{C_j2^{j\alpha_j}}{2^{\alpha_j}-1}<\infty,
 \label{eq:257-strengthened-cover}
\end{equation}
then $X_A(b)$ is irrational for every infinite
$A\subseteq\bigcup_jF_j$ and every integer $b\ge2$.
\end{thm}
\leannote{Lean: \lproof{ErdosProblems/Erdos257/PaperCompleteR8/PositiveCoverReturn.lean}{241}{strengthenedPositiveCoverClaim}.}

Every finite set admits such a majorant: take $\alpha_j=1$ and
$c_{j,d}=\mathbf1_{F_j}(d)$.  The restrictive part is summability over the
whole family, not the existence of a majorant for an individual set.
For example, $F_j=\{4^j\}$ with these choices contributes $2^{-j}$
to~\eqref{eq:257-strengthened-cover}, so the theorem applies.  Full support
cannot admit a cover satisfying the theorem, because the proof gives
arbitrarily small positive displacements and the full-support displacement
is at least $1/2$.

\begin{proof}
Write $B_j=2^{\alpha_j}$ and
\[
 U_j(N)=\sum_{r\ge1}2^{-r}f_{F_j}(N+r),\qquad
 V_j(N)=\sum_{d\ge1}c_{j,d}w_{B_j,d}(N).
\]
The choice $B_j=2^{\alpha_j}$ matches the decay after taking a
fractional power.  Subadditivity and the divisor majorant give
\[
 U_j(N)^{\alpha_j}
 \le\sum_{r\ge1}B_j^{-r}f_{F_j}(N+r)^{\alpha_j}
 \le\sum_{d\ge1}c_{j,d}
       \sum_{\substack{r\ge1\\d\mid N+r}}B_j^{-r}
 =V_j(N).
\]
The last identity sums a geometric progression in each residue class.
The geometric-series identity also gives
\[
 \Delta_{2,F_j}(N)=U_j(N)-X_{F_j}(2)\le U_j(N).
\]
Consequently, once the first $J$ covering sets have zero displacement, nonnegativity
bounds the displacement of their union by $\sum_{j>J}U_j(N)$.

Fix $\varepsilon>0$, set $t_j=\varepsilon2^{-j}$, and choose $J$ with
\[
 K_J:=\sum_{j>J}\frac{C_jt_j^{-\alpha_j}}{B_j-1}<\frac14.
\]
This is possible since $\varepsilon^{-\alpha_j}\le\max(1,\varepsilon^{-1})$.
Choose $L$ divisible by every member of the first $J$ covering sets.  Equation~\eqref{eq:257-mixed-finite-kernel} bounds the finite mean of
$S_J(N):=\sum_{j>J}t_j^{-\alpha_j}V_j(N)$ by
$(1+4L/M)K_J<1/2$ whenever $M\ge4L$.
One sample therefore has $S_J(N)<1$, forcing $U_j(N)<t_j$ for every $j>J$.
Every exponent in the first $J$ finite sets divides $L$, so those terms
have zero displacement.  Consequently,
\[
 0<\Delta_{2,A}(N)\le\sum_{j>J}U_j(N)\le\varepsilon.
\]
The fixed rational lattice excludes rationality at base two.
For $0\le r<d$, the atom $(b^r-1)/(b^d-1)$ is nonincreasing
in $b>1$.  For $r>0$, cancellation of $b-1$ writes it as
$A(b)/(A(b)+C(b))$, with $A(b)=\sum_{i<r}b^i$ and
$C(b)=\sum_{r\le j<d}b^j$.  Its derivative has the sign of
$A'C-AC'=\sum_{i<r\le j<d}(i-j)b^{i+j-1}<0$.
Thus $\Delta_{b,A}(N)\le\Delta_{2,A}(N)$ for real $b\ge2$;
for integer $b$ the rational-lattice argument gives irrationality.
This ordinary argument is sharper than the factor-two bound in
\href{https://github.com/wcook04/plectis-erdos-lean/blob/52f29ad173b04e3bac941b3663f2b9aebe5de0bb/ErdosProblems/Erdos257/PaperCompleteR7/CoverKernel.lean#L232}{the scalar comparison between bases}.
This locator is not a fresh Lean verification.
Any prescribed positive divisor can be included in $L$, so the witnesses
can also be required to be arbitrarily large.
\end{proof}

The same proof permits any positive weights $\eta_j$ with $\sum_j\eta_j=1$:
replace $2^{j\alpha_j}$ in~\eqref{eq:257-strengthened-cover} by
$\eta_j^{-\alpha_j}$ and take $t_j=\varepsilon\eta_j$.
\subsection{Combining the two support criteria}
Separate small-displacement witnesses need not occur at the same index.
For example, a sequence small only at even indices and one small only at odd
indices need never have a small sum.  The useful feature of~\eqref{eq:257-mixed-finite-kernel} is its
uniformity in the moving modulus: the positive-cover argument can use the
exact observation window selected by the weighted proof.
\begin{thm}[mixed weighted and cover supports]\label{thm:257-mixed-supports}
Let $E,V\subseteq\Npos$.  Suppose $E$ has finite weighted mass
\eqref{eq:257-weighted-mass} at $b=2$ for a finite nonempty prime set
$\mathcal P$.  Suppose also that $V\subseteq\bigcup_jF_j$ for sets and
majorants satisfying Theorem~\ref{thm:257-variable-fractional-cover},
with~\eqref{eq:257-strengthened-cover} or its positive-weight variant.
Then $X_A(b)$ is irrational for every
infinite $A\subseteq E\cup V$ and every integer $b\ge2$.
\end{thm}
\leannote{Lean: \lproof{ErdosProblems/Erdos257/PaperCompleteR8/WeightedReturn.lean}{126}{mixedSupportClaim}, \lproof{ErdosProblems/Erdos257/PaperCompleteR8/ArbitraryWeightMixedClaim.lean}{101}{arbitraryWeightMixedSupport_allBase_hereditary}.}
\begin{proof}
Fix $\varepsilon>0$ and an integer $N_0\ge1$, and put $\rho=\varepsilon/3$.
Use the cover notation $B_j,U_j,V_j$ above, with $\eta_j=2^{-j}$ in
the case~\eqref{eq:257-strengthened-cover}.  Set $t_j=\rho\eta_j$ and choose
$J$ so that
\[
 K_J=\sum_{j>J}\frac{C_jt_j^{-\alpha_j}}{B_j-1}<\frac1{16}.
\]
Choose a finite $F\subseteq E$ whose complementary weighted mass is
$\kappa<\rho/16$.  Let $L$ be a positive common multiple of $N_0$, all
members of $F$, and all members of the first $J$ covering sets.
For large $H$, choose the modulus and averaging length from the base-two weighted proof:
\[
 Q=L\prod_{p\in\mathcal P}p^{\lfloor\log_pH\rfloor},\qquad
 G=\lfloor H/\max\mathcal P\rfloor,\qquad M=\lfloor2^{G/2}\rfloor.
\]
The finite estimate~\eqref{eq:257-weighted-final-diagonal}, with complementary
weighted mass $\kappa$, remains valid for this $L$: its proof requires only
that every member of $F$ divide $L$.  It also applies to finite or empty $E$, since positivity was
used only after the averaging estimate.  Thus
\[
 \mathscr D_{Q;M,M}(\Delta_{2,E}/\rho)<\frac18
\]
for sufficiently large $H$.  For the same finite distribution, \eqref{eq:257-mixed-finite-kernel} and
nonnegative interchange give
\[
 \mathscr D_{Q;M,M}S_J\le(1+4Q/M)K_J<\frac18,
 \qquad S_J=\sum_{j>J}t_j^{-\alpha_j}V_j,
\]
because $Q/M\to0$.  Hence some sample $N=Qm$ satisfies
$\Delta_{2,E}(N)/\rho+S_J(N)<1$.  At this index the weighted displacement is below $\rho$ and
$U_j(N)<t_j$ for every $j>J$.  The earlier terms vanish because every
exponent in their finite sets divides $L$, and hence divides $N$.  Therefore
\[
 \Delta_{2,A}(N)\le\Delta_{2,E}(N)+\Delta_{2,V}(N)<2\rho<\varepsilon,
 \qquad N\ge Q\ge L\ge N_0.
\]
Overlaps between the supports only improve the inequality.  Infinitude of
$A$ makes its displacement positive.  The atom comparison used in the cover
proof transfers arbitrarily small displacements to every integer base;
the fixed rational lattice then excludes rationality.
\end{proof}

The strengthened cover conclusion is kernel-checked as
\href{https://github.com/wcook04/plectis-erdos/blob/065e09523286894dfb57ba205e69666843817009/lean/ErdosProblems/Erdos257/PaperCompleteR8/PositiveCoverReturn.lean#L241}{\texttt{strengthenedPositiveCoverClaim}}, and the common-witness conclusion is
kernel-checked as
\href{https://github.com/wcook04/plectis-erdos/blob/065e09523286894dfb57ba205e69666843817009/lean/ErdosProblems/Erdos257/PaperCompleteR8/WeightedReturn.lean#L105}{\texttt{mixedSupportClaim}}.  Their public dependency closure was replayed
under Lean~4.29.1.  The arbitrary-positive-weight cover and mixed-support
conclusions are also kernel-checked as
\href{https://github.com/wcook04/plectis-erdos/blob/f3f54798508b29ded6e854e0fa5784a560bd6cc7/lean/ErdosProblems/Erdos257/PaperCompleteR8/ArbitraryWeightMixedClaim.lean#L19}{\texttt{arbitraryWeightPositiveCover\_allBase\_hereditary}} and
\href{https://github.com/wcook04/plectis-erdos/blob/f3f54798508b29ded6e854e0fa5784a560bd6cc7/lean/ErdosProblems/Erdos257/PaperCompleteR8/ArbitraryWeightMixedClaim.lean#L101}{\texttt{arbitraryWeightMixedSupport\_allBase\_hereditary}}; the latter gives the theorem's hereditary all-base conclusion for the arbitrary-weight branch.
This theorem does not establish strict enlargement over both individual
classes.  The weighted non-cover host has a separate source body in
\href{https://github.com/wcook04/plectis-erdos-lean/blob/52f29ad173b04e3bac941b3663f2b9aebe5de0bb/ErdosProblems/Erdos257/PaperCompleteR8/AnalyticSeparationReturn.lean#L16}{the weighted support without a strengthened cover};
the proposed cover-only host is not established here.
It also leaves universal irrationality and the quantitative and prime-power
thinning claims open here. \ev{Math}

With arbitrary positive cover weights, the class of subsets of such mixed
hosts is closed under finite unions and finite changes.  For weighted
supports use the union of the two finite prime sets: the prime part grows
and $h/(2^h-1)$ decreases.  For two covers, interleave their finite sets with
weights $\eta_j/2$ and $\theta_j/2$; the total cost grows by at most two,
since $2^{\alpha_j}\le2$.  Subsets inherit the same hosts, and finite sets
have finite weighted mass.  This argument uses the positive-weight variant;
it does not silently reindex the dyadic weights in~\eqref{eq:257-strengthened-cover}.
Countable unions require a tail budget.  Every prime singleton is admitted,
but for the full prime support $\mathcal P$ one has
$\Delta_{2,\mathcal P}(N)>1/3$ for every $N\ge1$: indeed
$\sum_{r\ge1}2^{-r}\omega(N+r)\ge1$, whereas
$X_{\mathcal P}(2)\le\sum_{a\ge2}(2^a-1)^{-1}<2/3$.
Here $\omega(n)$ is the number of distinct prime divisors of $n$.

\paragraph{The elementary independence boundary.}
For the divisor counts $c_A$ and a fixed $a\in A$,
the identity $c_A(an)-c_A(n)=1-\mathbf1_{a\mid n}$ holds for every
$n\ge1$ exactly when $a$ is coprime to every other member of $A$.
Every summand $\mathbf1_{t\mid an}-\mathbf1_{t\mid n}$ is nonnegative;
$t=a$ supplies the right side, and coprimality makes all other terms zero.
If $g=(a,t)>1$ for some $t\ne a$, take $n=t/g$ to obtain an additional
positive term.  For distinct odd primes $p,q$, the covariance of
$\mathbf1_{2p\mid n}$ and $\mathbf1_{2q\mid n}$ over a common period is
$1/(2pq)-1/(4pq)=1/(4pq)$.  Thus dilating a prime support already destroys
the independence used in the prime-incidence argument.  This brief ordinary
calculation marks a boundary of that argument, not an irrationality criterion.

\subsection{The universal problem and the half-value question}

For an integer $n \ge 1$ set $x_n := 1/(2^n-1)$, and for $A \subseteq \{1,2,3,\dots\}$ write
$x_A := \sum_{n \in A} x_n$; the sum always converges, since $x_n = O(2^{-n})$. Erd\H{o}s's
Problem \#257 asks a single universally quantified question about this family.

\begin{defn}[Universal \#257, written (U)]
\label{defn:U}
Is $x_A = \sum_{n \in A} 1/(2^n-1)$ irrational for \emph{every} infinite $A \subseteq \mathbb N_{>0}$?
\textbf{OPEN.} \ev{Open}
\end{defn}

\begin{defn}[The half-value question, written (H)]
\label{defn:H}
Let the \emph{Mersenne achievement set} be
$\mathcal{A} := \{\, y \in \R : \exists A \subseteq \N,\ 0 \notin A,\ y = \sum_{n \in A} x_n \,\}$
(\lean{mersenneAchievementSet}{GreedyAchievementSet.lean:571}). Is $1/2 \in \mathcal{A}$?
\textbf{OPEN.} \ev{Open}
\end{defn}

The two are joined by one elementary and completely one-sided implication. Every $2^n-1$ is odd,
so every \emph{finite} subset sum of $\{x_n\}$ has odd reduced denominator and cannot equal $1/2$
(\lean{positiveMersenneSupportValue\_coe\_finset\_ne\_half}{HalfCutLocator.lean:243}). Hence a
witness to $1/2 \in \mathcal{A}$ is necessarily an infinite $A$ with $x_A = 1/2$ rational, which
refutes (U) outright. Conversely, (U) implies $1/2 \notin \mathcal{A}$ as one instance among
uncountably many, and a proof of $1/2 \notin \mathcal{A}$ leaves (U) entirely open. Thus (H) is one candidate counterexample to (U), not a restatement of it.
Section~\ref{bar:asym} explains the corresponding difference between a
finite certificate of nonmembership and an unbounded survival condition. \ev{Math}

The reciprocal-summable and finite-prime weighted theorems, together with the
checked support theorems below, leave (U) and (H) open. The later countermodels
identify the particular arithmetic information lost by individual approaches;
they do not prove that all such failures have one common cause.

\subsection{Known support theorems}

Full-support irrationality is Erd\H{o}s's theorem~\cite{erdos1948}.  The
pairwise-coprime case is Erd\H{o}s's theorem of 1968, where he also stated that
the coprimality condition can be removed~\cite[p.~222]{erdos1968}.  Luca and
Tachiya proved irrationality for every purely periodic integer weight that is
not identically zero~\cite[Theorem~A, p.~139]{lucatachiya2017}; the purely
periodic, eventually periodic, residue-class, odd-support and periodic-weight
rows below follow from their theorem, and odd support is their
Example~2~\cite[p.~140]{lucatachiya2017}.  The lcm-gap supports, including
factorials and powers of two, have summable reciprocals, so they fall under the
reciprocal-summable theorem stated by Erd\H{o}s.  The Lean declarations in the
table are formal proofs of these statements.

The following table includes two proof mechanisms: block certificates
for the divisor counts $c_A(n)$, and rational-approximation estimates
for a growing least-common-multiple gap.  The latter need not be
presented as instances of the former.  The notation
$\operatorname{sc}_A(n)$ in later source excerpts means the same
function as $c_A(n)$.  All bases in the table are integers $b\ge2$.
For the lcm-gap row write
$L_k=\operatorname{lcm}(a_0,\ldots,a_{k-1})$, with $L_0=1$.

\begin{center}
\begin{longtable}{@{}p{0.30\textwidth}p{0.32\textwidth}p{0.30\textwidth}@{}}
\hline
\textbf{Support class} & \textbf{Exact hypotheses} & \textbf{Site} \ev{Lean} \\
\hline
\endhead
Full support $A = \Npos$ & $b \ge 2$ & \lean{irrational\_erdosSum\_full\_support}{CertificateKernel.lean:8328} \\
Multiples $A = d\Npos$ & $b \ge 2$, $d \ge 1$ & \lean{irrational\_erdosSupportSeries\_multiples}{CertificateKernel.lean:9103} \\
Purely periodic & $m\ge1$, $m$-periodic, contains a positive element & \lean{irrational\_erdosSupportSeries\_periodic}{CertificateKernel.lean:11590} \\
Eventually periodic & $m\ge1$, $m$-periodic from $N_0$, $A$ infinite & \lean{irrational\_erdosSupportSeries\_eventuallyPeriodic}{CertificateKernel.lean:11604} \\
Residue class & any $m \ge 1$, any $c$ & \lean{irrational\_erdosSupportSeries\_residueClass}{CertificateKernel.lean:11672} \\
Odd support (density $1/2$) & $b\ge2$ & \lean{irrational\_erdosSupportSeries\_odd}{CertificateKernel.lean:11686} \\
Growing lcm gap & $a_k$ strictly increasing and positive; $a_k-L_k\to\infty$ & \lean{irrational\_erdosSum\_of\_lcm\_gap}{CertificateKernel.lean:5883} \\
Factorials; powers of two & instances of the lcm-gap theorem & \lean{irrational\_erdosSum\_factorial\_support}{CertificateKernel.lean:6035}, \lean{irrational\_erdosSum\_two\_pow\_support}{CertificateKernel.lean:6059} \\
Pairwise coprime (Erd\H{o}s 1968) & $A$ infinite, pairwise coprime, \emph{and} $\sum_{a \in A} 1/a < \infty$ & \lean{irrational\_erdosSupportSeries\_pairwise\_coprime}{CertificateKernel.lean:10776} \\
Signed periodic weights & periodic integer weights; their divisor-transformed coefficients are nonnegative and frequently nonzero & \lean{irrational\_intWeightedErdosSeries\_periodic\_of\_coeff\_nonneg\_of\_frequently\_ne\_zero}{CertificateKernel.lean:14583} \\
\hline
\end{longtable}
\end{center}

The lcm-gap hypothesis does imply reciprocal summability.  Eventually
$a_{k+2}>\operatorname{lcm}(a_k,a_{k+1})$.  If
$a_{k+1}\ge2a_k$, monotonicity gives $a_{k+2}>2a_k$.  Otherwise
$a_k\nmid a_{k+1}$, so
$\operatorname{lcm}(a_k,a_{k+1})\ge2a_{k+1}>2a_k$.
Thus each of the even- and odd-indexed subsequences eventually grows
geometrically, proving $\sum_k1/a_k<\infty$.  This explains the
inclusion claimed above, rather than using ``lacunary'' as a substitute
for a growth hypothesis.  For eventually periodic weights, the
periodic tail must not be identically zero: an eventually zero weight
sequence gives only a finite rational sum.

Full-support irrationality for every integer base $b\ge2$ is classical (Erd\H{o}s, 1948). The
contribution recorded here is its Lean formalisation and the explicitly delimited
structured-support developments, not a new extension from base two to arbitrary integer bases.
Each support family retains its stated hypotheses; none establishes irrationality for every
infinite support. The $b=2$ case of the first row recovers the Erd\H{o}s--Borwein constant
$E = \sum_{n \ge 1} 1/(2^n-1)$ as a one-line corollary of that classical theorem. Separately,
the independence and tail estimates used by a particular certificate
construction must be checked for that support. Barrier~\ref{bar:box}
discusses those construction requirements; the averaging proof above does
not assume pairwise coprimality or independence of divisor events.

A separate conditional result allows common factors confined to divisors
of one fixed integer $Q$, while the remaining factors must be pairwise
coprime and have summable reciprocals
(\lean{irrational\_erdosSupportSeries\_of\_orthogonalPetalBouquet}{SupportSunflowerDichotomy.lean:540}).
It also assumes that for every $K\ge1$ there is an $N\ge0$ such that
\[
 2^K\mid\sum_{r=1}^K c_A(N+r)2^{K-r},\qquad
 \sum_{r\ge1}c_A(N+K+r)2^{-r}\le16.
\]
Thus relaxing coprimality does not remove the simultaneous divisibility
and tail-bound requirement.  This additional hypothesis is not proved
here, and the source file does not establish a dichotomy for arbitrary
supports. \ev{Lean}\ev{Open}

The record also gives a logarithmic bound on runs of zero divisor counts
under rationality: for each $\varepsilon>0$, a suitable constant $B$ bounds
the length by $\varepsilon\log_2N+B$ beyond the stated threshold
(\lean{supportCoeffZeroWindow\_length\_le\_eps\_logb}{SublogDivisorCoverage.lean:435}).
For these supports there is already a simpler unconditional bound.  If
$a_0=\min A$, every $a_0$ consecutive positive integers contain a multiple
of $a_0$, where $c_A$ is positive.  Hence every zero run has length at most
$a_0-1$, whether or not $X_A(2)$ is rational.  In particular the logarithmic
bound alone does not distinguish rational-valued supports.  The separate
unbounded-tail result
\lean{shifted\_state\_unbounded\_of\_infinite\_support}{RationalSupportCarrySkeleton.lean:2327}
also requires interpretation: an unbounded sequence can still satisfy
$u(N)=o(2^N)$. \ev{Lean}\ev{Math}

\begin{rem}[What is and is not covered in the universal direction]
The pinned Lean corpus contains no theorem for $A=$ the primes, but this is a
formalisation boundary, not a mathematical open case: Tao and
Ter\"av\"ainen prove the base-$2$ prime-support value irrational by first
identifying it with $\sum_{n\ge1}\omega(n)2^{-n}$.  Duverney and Tachiya
likewise settle squarefree support at every power-of-two base, and the
perfect $i$th powers at every integer base (Corollary~1.2 and Example~1.1,
p.~4 of the author preprint).  Neither
result is formalised here.  The statement of Problem~\#257 has been formalised
before, as a conjecture with an unfilled proof, in the \emph{Formal
Conjectures}
\href{https://github.com/google-deepmind/formal-conjectures/blob/f776d2f2039351b00737ffcafb9d7d7666e1d9af/FormalConjectures/ErdosProblems/257.lean}{collection};
that file also proves the Lambert identity for the full-support value and
states its irrationality, the 1948 theorem of Erd\H{o}s, with proof
\texttt{sorry}.  It is statement-level prior art, and no theorem recorded here
is derived from it.  The corpus still has no general theorem for an
arbitrary positive-density nonperiodic or divisor-dense support, and the
universal statement remains open.
\end{rem}

\paragraph{Periodic weights: the two hypothesis sets.}
The periodic sequence theorem belongs to Luca and Tachiya
\cite{lucatachiya2014periodic}; their later author-written account states
it as Theorem~A~\cite[p.~139]{lucatachiya2017}.  In the displayed Lean
row above, nonnegativity and frequent nonvanishing are assumptions on
$c_w(n)=\sum_{d\mid n}w(d)$, not on $w(n)$ itself.  The formal row must
therefore be read as a separately specified instance, not as an exact
transcription of every clause of the published theorem.

\paragraph{A separate denominator-growth criterion.}
Erd\H{o}s's growth theorem~\cite{erdos1975}, as restated and recovered
by Barreto, Kang, Kim, Kova\v{c} and Zhang
\cite[Remark~4(4), p.~5]{bkkkz2026}, gives irrationality of
$\sum_n1/a_n$ for an increasing integer sequence with
$a_n\ge n^{1+\tau}$ for some $\tau>0$ and
$\limsup_n a_n^{1/2^n}=\infty$.  For $a_n=b^{c_n}-1$, where
$c_n$ are increasing positive exponents and $b\ge2$ is an integer,
the growth hypothesis becomes
\[
 \limsup_{n\to\infty}\frac{c_n}{2^n}=\infty.
\]
Indeed $c_n\ge n$, so the polynomial lower bound holds eventually,
and $\log_b(b^{c_n}-1)=c_n+O(1)$.  An eventual lower bound suffices:
deleting finitely many terms subtracts a rational number and changes
$2^n$ only by a fixed factor after reindexing.

This is not the reciprocal-summability condition of
the reciprocal-support theorem in Section~1.1.  The squares $c_n=n^2$ have
convergent reciprocal sum but $c_n/2^n\to0$.  In the other direction,
set $n_1=1$ and $n_{k+1}=n_k+4^{n_k}$, and let block $k$ consist of
all integers from $4^{n_k}$ to $2\cdot4^{n_k}-1$.  The blocks are
successive disjoint increasing blocks, and their union has increasing
enumeration satisfying $c_{n_k}=4^{n_k}$.  Thus
$c_{n_k}/2^{n_k}=2^{n_k}\to\infty$, whereas each block has reciprocal
sum at least $1/2$.  Its reciprocal sum therefore diverges.
These examples separate these two hypotheses, not every support
criterion in this paper.  \ev{Cited}\ev{Math}

\subsection{Geometry and rational-target membership}

\begin{thm}[Achievement-set geometry]
\label{thm:geometry}
$\mathcal{A}$ is compact, closed, perfect, totally disconnected and nowhere dense, and
$\operatorname{volume}(\mathcal{A}) = 1$.  Thus its measure is positive
although it contains no interval.  Its convex hull is $[0,E]$, where
$E=\sum_{n\ge1}w_n$.  The positive-index digit coding onto
$\mathcal{A}$ is injective: each achievable real has \emph{exactly one} support.
\coord{achievement-set} \scale{uniform} \ev{Lean}\\
\textit{Sources:} \lean{isCompact\_mersenneAchievementSet}{GreedyAchievementSet.lean:656},
\lean{perfect\_mersenneAchievementSet}{GreedyAchievementSet.lean:1656},
\lean{isTotallyDisconnected\_mersenneAchievementSet}{GreedyAchievementSet.lean:1672},
\lean{isNowhereDense\_mersenneAchievementSet}{GreedyAchievementSet.lean:1681},
\lean{volume\_mersenneAchievementSet}{GreedyAchievementSet.lean:996},
\lean{positiveMersenneDigitValue\_injective}{GreedyAchievementSet.lean:1582}.
\end{thm}
\leannote{Lean: \lproof{ErdosProblems/Erdos257/PaperCompleteR20/AchievementGeometry.lean}{40}{paper_achievement_geometry}.}

\begin{thm}[Support-restricted refinement]
\label{thm:supported-dichotomy}
Use zero-based indices in this statement: coordinate $j\in\N$ carries
weight $w_{j+1}$.  For $J\subseteq\N$, consider the sums that use only
coordinates in $J$.  If $\N\smallsetminus J$ is finite, this achievement set
has measure $2^{-|\N\smallsetminus J|}$; if infinitely many coordinates are
omitted, its measure is zero.  Injectivity survives every restriction;
perfectness is proved when $J$ is infinite.  No perfectness claim is made for
finite $J$, whose coding range is finite.
\coord{achievement-set} \scale{uniform} \ev{Lean}\\
\textit{Sources:} \lean{volume\_supportedMersenneAchievementSet\_dichotomy}{ErdosProblems/Erdos257/MersenneSubseriesRigidity.lean:397},
\lean{volume\_supportedMersenneAchievementSet\_eq\_zero\_of\_compl\_infinite}{ErdosProblems/Erdos257/MersenneSubseriesRigidity.lean:368},
\lean{perfect\_supportedMersenneAchievementSet}{ErdosProblems/Erdos257/MersenneSubseriesRigidity.lean:167},
\lean{supportedMersenneDigitValue\_injective}{ErdosProblems/Erdos257/MersenneSubseriesRigidity.lean:54}.
\end{thm}

\paragraph{Hausdorff dimension of a restricted set.}
For example, retaining the positive exponents $2,4,6,\ldots$ gives
a set of measure zero but Hausdorff dimension $1/2$.  More generally,
let $J\subseteq\{1,2,\ldots\}$ and put
\[
 S_J=\left\{\sum_{n\in J}\varepsilon_nw_n:
                  \varepsilon_n\in\{0,1\}\right\},\qquad
 s_N=\#(J\cap\{1,\ldots,N\}),\qquad
 d=\liminf_{N\to\infty}s_N/N.
\]
Then $\dim_H S_J=d$.  Here is a proof, separate from the supplied
formal measure statements.  If $J$ is finite, both sides are zero,
so suppose $J$ is infinite.

Write $w_n=2^{-n}+\delta_n$, where
$\delta_n=\sum_{r\ge2}2^{-rn}$.  A geometric sum gives
\[
 \sum_{k>n}\delta_k
 =\sum_{r\ge2}\frac{2^{-rn}}{2^r-1}\le\delta_n/3<\delta_n.
\]
The unique Mersenne digit coding therefore defines a map from $S_J$
onto the binary digit-restricted set
$B_J=\{\sum_{n\in J}\varepsilon_n2^{-n}\}$ by keeping the same digits.
This map is $1$-Lipschitz.  In fact, at the first differing digit $n$,
the binary difference is nonnegative after orienting the larger digit
first, and the Mersenne difference exceeds it by at least
$\delta_n-\sum_{k>n}\delta_k>0$.

The $2^{s_N}$ Mersenne prefix intervals have length at most
$R_N\le2^{1-N}$.  For any $t>d$, a subsequence with
$s_N/N<(d+t)/2$ makes their total $t$-power length tend to zero.
Hence $\dim_H S_J\le d$.
For the reverse inequality, give the allowed binary digits independent
probabilities $1/2$.  Each dyadic interval of length $2^{-N}$ has
measure at most $2^{-s_N}$; endpoints have measure zero because $J$
is infinite.  An interval $I$ with
$2^{-N-1}<|I|\le2^{-N}$ meets at most three such dyadic intervals.
For $0<s<d$ and all sufficiently large $N$, its measure is therefore
at most $3\cdot2^{-s_N}\le3\cdot2^s|I|^s$.
Summing this inequality over any sufficiently fine cover of $B_J$
gives a positive lower bound for the sum of its $s$-power lengths.
Thus $\dim_H B_J\ge s$; letting $s\uparrow d$ proves
$\dim_H B_J\ge d$.  A Lipschitz map cannot increase Hausdorff
dimension, since it sends a cover to one of no larger diameters.
The map above now yields $\dim_H S_J\ge d$, as required.
For $d=0$ the lower bound is automatic.  \ev{Math}

\begin{rem}[Achievement-set terminology and source boundary]
We use the modern term \emph{achievement set} for the subsum set that
\href{https://doi.org/10.11429/ptmps1907.7.14_250}{Kakeya's 1914 note} describes
as a set of partial sums.  For an infinite summable sequence of positive
terms, Kakeya proves perfectness, and nowhere density under an all-index strict-tail
condition; \href{https://arxiv.org/abs/2406.17593v4}{Kova\v{c}--Tao,
Remark~4.1} verifies that strict-tail condition for the fixed-base Mersenne
weights.  Nitecki proves that a subsum set whose terms all exceed their tails
is a Cantor set of measure $\lim_n2^nR_n$, and credits this to
Hornich~\cite[Theorem~4(1), p.~9]{nitecki2013}; this is the classical source for
the topology and measure used here.  The local
volume dichotomy, injectivity, and greedy-membership statements above and
below are Lean-checked claims of this corpus.  The sparse power-series criteria
of \href{https://arxiv.org/abs/2601.20743}{Kaneko--Suzuki--Tachiya}
(Theorems~1--3, pp.~3--5 of arXiv v1) assume sparse coefficient supports;
their relation to the divisor-incidence coefficients of these subseries is
recorded after Theorem~\ref{thm:257-weighted}. \ev{Cited}
\end{rem}

\begin{thm}[Membership equals greedy survival; the fatal-gap dichotomy]
\label{thm:greedy-survival}
For a real target $x\ge0$, let $r_n(x)$ be the remainder after the
greedy rule has processed weights $w_1,\ldots,w_n$, and let
$R_n=\sum_{j>n}w_j$, with $r_0(x)=x$ and $R_0=E$.  Then
\[
 x\in\mathcal A\quad\Longleftrightarrow\quad
 x\ge0\ \text{ and }\ r_n(x)\le R_n\ \text{for every }n\ge0.
\]
For every $x\in[0,E]$, nonmembership is equivalent to a finite
strict gap between the two next-prefix intervals; see
Observation~\ref{obs:general-target-gap}.
\coord{greedy recurrence} \scale{uniform} \ev{Lean}\\
\textit{Sources:} \lean{mem\_mersenneAchievementSet\_iff\_greedy\_survival}{GreedyAchievementSet.lean:1458};
\lean{half\_mem\_mersenneAchievementSet\_or\_exists\_fatal\_gap}{HalfCutLocator.lean:623};
\lean{existsFatalHalfGap\_iff\_half\_not\_mem\_mersenneAchievementSet}{HalfCutLocator.lean:643};
\lean{half\_mem\_mersenneAchievementSet\_iff\_no\_existsFatalHalfGap}{HalfCutLocator.lean:654};
witness type \lean{ExistsFatalHalfGap}{HalfCutLocator.lean:526}.
The general-target statement is checked by
\lean{paper\_greedy\_survival}{lean/ErdosProblems/Erdos257/PaperCompleteR20/GeneralTargetGap.lean:183}.
\end{thm}

Two boundary degeneracies are excluded at every rank, unconditionally: $r_{k-1}(1/2)\ne w_k$ by the
odd-denominator argument above, and $r_{k-1}(1/2)\ne R_k$ because the tail differs from $E$ by a
rational while $E$ is irrational. Both greedy decisions are therefore strict at every rank.
\ev{Math}

The supplied computation record reports the following finite checks.
The underlying tests use exact integer arithmetic; decimal margins are
summaries of those data, not proofs of an asymptotic estimate.  The scan to row $200{,}000$ and the long grouped computation were
not rerun for this revision.  The quotient calculation through row
$2500$ was recomputed with exact integers, as specified below.
Reset margins are measured by
$\log_2|\mathrm{rem}(r+1)-2^{r+1}|-(r+5)/2$.
The separate real-greedy scan uses
$\log_2(4^k(r_{k-1}(1/2)-w_k))$ at a take and
$\log_2(4^k(R_k-r_{k-1}(1/2)))$ at a safe skip.
These are different normalisations, not one common margin statistic.
Truncation rungs $J = 3,\dots,22$ all \emph{proved} to survive (not merely computed). Seam orbit
certified to row $200{,}000$: branch counts $\mathrm{R{:}M{:}U} = 100197{:}49899{:}49898$, the
target inequality failing at exactly one reset row, $r = 7$; maximal pure-R run $19$, attained at
row $158{,}096$.

Independently, rows $6 \le n \le 2500$: $1209$ resets, zero classification
anomalies, minimum in-scope margin $+1.1119$ bits at row $14$ growing to $\approx 1246$ bits by row
$2500$ (mean $627$), never re-approaching the threshold.

Sharp tail margin, ranks $2$ to $3000$: no
fatal skip, take fraction $0.4992$, tightest take-margin $+0.0340$ bits at rank $7$ (the near-tie
$1/126>1/127$).  In the stated units this is
$\log_2(4^7(1/126-1/127))=0.034035\ldots$.
It is the smallest take-margin reported in this finite scan. The recorded finite greedy prefix through $m=200{,}000$ satisfies
$0\le 1/2-\sum_{d\in G\cap[1,m]}w_d<2^{-199999}$.
The sum here is finite; this statement does not evaluate the infinite greedy sum.

A half-trapping macro receipt reaches depth
$13{,}548{,}057$ (labelled in its own source as a private exact finite computation receipt, not
publication authority). Kernel-checked corroboration of the low rows:
\lean{seamUpperResetDyadicBandEscape\_through\_thirty}{HalfCylinderUpperResetBandCertificates.lean:78}
verifies rows $13$--$30$ by \texttt{decide}. \ev{Cert}

\begin{rem}[Literature, exactly]
$E$ is irrational~\cite{erdos1948}. Zudilin proves the irrationality measure bound
$\mu(E)\le2.46497868\ldots$~\cite[Theorem~1, p.~154]{zudilin2004}, correcting an earlier
computation (see the remark at the end of his Section~3, p.~159). Transcendence of $E$ is
unknown~\cite[p.~2]{duverneytachiya}, and the Erd\H{o}s--Graham conjecture that every such subset
sum is irrational, i.e.\ (U), is itself open. For the binary digits of $E$, Erd\H{o}s's proof
produces arbitrarily long runs of zeros. Crandall asked whether the block \texttt{11} occurs
infinitely often in the binary expansion of $E$, and Campbell proves that it
does~\cite[Theorem~1, p.~12]{campbell2026}. \ev{Cited}
\end{rem}

The cited digit results establish the occurrence of particular patterns.
An irrationality-measure estimate gives a different kind of information:
a run of $L$ equal digits after position $N$ approximates $E$ by a dyadic
rational with denominator at most $2^N$, giving
$L<(\mu(E)-1+\varepsilon)N$ for all sufficiently large $N$.
The reset argument considered here needs a bound for runs in its own
inhomogeneous recurrence, not simply in the binary digits of $E$.
No implication from the cited digit results to that bound is proved here.
Such a bound is an input to this particular approach to (H), not a
necessary ingredient of every possible proof of membership.

\clearpage
\part*{II. Limitations and extensions of recorded methods}
\addcontentsline{toc}{part}{II. Limitations and extensions of recorded methods}
\section{Limitations of the recorded methods}
\label{sec:257-wall}

\subsection{What the obstructions do and do not show}

This section studies the scope of particular reformulations and
certificate mechanisms.  Several share the greedy representation, but no
proof identifies a single necessary obstruction to all methods.  The
statements below are useful insofar as they specify which hypotheses fail,
which implications are equivalences, and which parameters must be uniform.
The first structural fact is uniqueness at a fixed target and depth.

\begin{obs}[Uniqueness within the fixed-target coding]
\label{obs:no-ensemble}
Each Mersenne weight exceeds the sum of all subsequent weights.
Consequently the digit coding is injective (Theorem~\ref{thm:geometry}):
a represented target has only one support, recovered by the greedy rule.
For the integer weights, both the scale and the cutoff matter.  If
$2\le k\le M$ and
$q_{M,d}=\lfloor2^M/(2^d-1)\rfloor$ for $2\le d\le k$, then
\[
 q_{M,d}-\sum_{e=d+1}^{k}q_{M,e}\ge2^{M-k}.
\]
Indeed, with $w_d=(2^d-1)^{-1}$ and $R_k=\sum_{e>k}w_e$,
\[
 w_d-\sum_{e=d+1}^{k}w_e=w_d-R_d+R_k>R_k>2^{-k}.
\]
Rounding the head down loses less than one, while rounding the later
terms down can only increase the difference.  The resulting integer
is greater than $2^{M-k}-1$, proving the bound.  Comparing two words at
their first differing rank gives the same separation for their sums.
The rows with $M=2s$, $k=s-1$ and $s\ge3$ therefore have gap at least
$2^{s+1}$; the full rows with $k=M$ have gap at least $1$.
The first bound must not be applied after enlarging the cutoff:
at $M=6$, adjoining rank $3$ introduces the subset sum
$q_{6,3}=9<16=2^{3+1}$.

When a finite construction is proved to encode a fixed greedy target,
its choices are constrained by that support.  Injectivity alone does
not identify an arbitrary finite approximation with a real greedy
prefix.  Nor does it exclude averaging over indices, moduli, nearby
targets or auxiliary constructions.
\coord{meta} \scale{uniform} \ev{Lean}\ev{Math}\\
\textit{Sources:} \lean{weightedBoolSum\_separated}{HalfCylinderIntegerGreedy.lean:447},
\lean{three\_mul\_tailWeight\_add\_exactLateGap\_eq\_three\_mul\_headWeight}{HalfCylinderLargestSkipGap.lean:230},
\lean{positiveMersenneDigitValue\_injective}{GreedyAchievementSet.lean:1582}.
\end{obs}

The six limitations below require separate arguments.  Fixed-target
uniqueness explains canonicalisation (\ref{bar:canon}); the recorded
identities determine which sufficient conditions are equivalent to their
goal (\ref{bar:collapse}).  Measure and category (\ref{bar:measure}),
observation scales (\ref{bar:scale}), finite versus infinite membership
(\ref{bar:asym}), and the hypotheses of the support certificate
(\ref{bar:box}) do not follow from injectivity alone.

\subsection{Uniqueness of the specified greedy words}
\label{bar:canon}

\begin{prop}[Uniqueness under the stated finite conditions]
\label{prop:canon}
Put $G=\mathrm{greedyMersenneSupport}(1/2)$,
$w_n=(2^n-1)^{-1}$ and $R_d=\sum_{n>d}w_n$.
The following are three separate uniqueness statements.
\begin{enumerate}[nosep,label=(\roman*)]
\item If $D\subseteq\{1,\ldots,d\}$ and
$X_D(2)\le1/2\le X_D(2)+R_d$, then $D=G\cap\{1,\ldots,d\}$.
\item If $c\ge4$, $D\subseteq\{2,\ldots,c-1\}$ and
$0<1/2-X_D(2)<w_c$, then $D=G\cap\{1,\ldots,c-1\}$.
\item Let positive integer weights $v_1,\ldots,v_k$ satisfy
$v_i\ge g+\sum_{j>i}v_j$ for a positive integer $g$.
For a nonnegative integer target $T$, let $y^*$ be the word obtained by
visiting these weights in order and taking each when it fits.
An admissible Boolean word $y$ has
$0\le T-\sum_i y_iv_i<g$ if and only if $y=y^*$ and
$0\le T-\sum_i y_i^*v_i<g$.
\end{enumerate}
Part (iii) gives uniqueness under the small-remainder condition,
not existence.  Theorem~\ref{record:257bm-i11a} gives a
one-weight counterexample to dropping that condition.
\coord{straddle-prefix ;  integer-greedy} \scale{uniform} \ev{Lean}\\
\textit{Sources:} \lean{IsStraddlePrefix.half\_agrees\_greedy}{HalfCutLocator.lean:442};
\lean{eq\_halfGreedyPrefixSupport\_of\_critical\_crossing}{BooleanMobiusCriticalCapacityCofinal.lean:50};
\lean{remainder\_lt\_gap\_iff\_eq\_integerGreedyBits}{BooleanMobiusGreedyReduction.lean:918}.
\end{prop}
\begin{proof}
For (i), truncate $D$ at each successive rank.  The target remains
in the corresponding prefix interval because all discarded terms
are bounded by the complete tail.  The two child intervals are
disjoint: the lower one forces a greedy skip and the upper one
forces a greedy take.  Induction identifies every digit.
Part (ii) implies (i) at depth $c-1$, since $w_c<R_{c-1}$.
For (iii), consider the first disagreement with the integer greedy
word.  Taking a rejected weight would exceed $T$; rejecting an accepted
weight leaves a deficit at least $g$, even after taking every later
weight.  The converse is immediate.
\end{proof}

\textbf{Excludes at fixed depth.} Under the three canonicalisation
hypotheses just stated, a proof cannot gain freedom by choosing a different
admissible word for the same target and depth.  This says nothing about
families satisfying different approximation conditions, or about averaging
over depths or auxiliary data.

\textbf{What the finite tests show.} Three different calculations are
recorded here.  For unrestricted prefix residues, the proposed
$2^{t/2}$ exclusion radius is reported to fail in $74\%$ of the tested
ranks $t\in[8,26]$.  Here the finite test uses the truncated,
scaled fractional parts
\[
 b_d(t)=\left\lfloor 2^t\left\{\frac{2^t}{2^d-1}\right\}\right\rfloor,
 \qquad \{y\}=y-\lfloor y\rfloor.
\]
At $t=15$, the values for $d=2,3,5,6,10$ are respectively
$21845,4681,1057,4161,1025$.  Their sum is $32769=2^{15}+1$,
so the sum has circular distance one from $0$ modulo $2^{15}$.
This is a statement about these integer truncations, not the exact
fractional-part sum.  A heuristic count of near-hits has
size $2^{t/6}$; that estimate is not a theorem about the greedy support.
\ev{Cert}

Long prescribed digit strings also occur for unrestricted finite words:
the recorded examples at $t=20$ and $t=22$ have respectively eighteen
leading ones and twenty leading zeros.  A separate construction supplies
about $t/12$ independently adjustable bits, with the reported finite check
extending to $t=601$.  These examples disprove the proposed static
exclusions, not an assertion about the actual greedy word.  \ev{Cert}

The integer-window search reported no counterexamples for $t\le22$ or at
$t=40,42,45$ over the tested cut positions.  That absence is a finite
observation.  Proposition~\ref{prop:canon} explains why certain
\emph{specified} approximation conditions leave only the greedy word; it
does not identify all three searches with one another or rule out averaging
over other data.  \ev{Cert}

\begin{rem}[A recorded hazard]
The natural sufficient condition via fractional split-mass is \emph{not} necessary on real crossing
cores: the fixture $D = \{2,3\}$, $c = 5$ is an explicit counterexample
(\lean{splitFractionMass\_one\_bound\_not\_necessary\_fixture}{BooleanMobiusSkippedCoreCriticalCapacity.lean:183}).
The example refutes necessity of that sufficient condition.  It does not
refute using the condition on a class of cores for which it can be proved.
\ev{Lean}
\end{rem}

\subsection{Equivalent formulations of half-membership}
\label{bar:collapse}

The six conditions listed below are equivalent to half-membership.
Their quantifier patterns can look less restrictive than a single
compatible infinite construction, but the equivalences must be checked
mathematically.  They follow here from a carry identity and an unconditional
tail estimate.  No claim is made that every conceivable sufficient
condition is equivalent to the goal.

\begin{lem}[Carry and residual value]
\label{lem:collapse-mech}
Let $A \subseteq \N$ with $1 \notin A$, put $\delta := 1/2 - x_A$, and let
$\operatorname{ihc}(A,N)$ denote the integer half-carry at level $N$. Then
\[
  \operatorname{ihc}(A,N) \;=\; 2^{N+1}\,\delta \;+\; \mathrm{T}(N+1),
  \qquad 0 \;\le\; \mathrm{T}(N+1) \;\le\; 2\sqrt{N+1} + 4 ,
\]
where $\mathrm T(m)=\sum_{r\ge1}c_A(m+r)2^{-r}$.
\coord{half-carry recurrence} \scale{uniform} \ev{Lean}\\
\textit{Sources:} \lean{integerHalfCarry\_eq\_scaled\_residual\_add\_tail}{HalfCarryReachability.lean:871};
\lean{binaryCoeffTail\_nonneg}{GenericTailOrbitRigidity.lean:78};
\lean{binaryCoeffTail\_supportCoeff\_le\_two\_sqrt\_add\_four}{BooleanMobiusCarry.lean:290}.
\end{lem}

\begin{prop}[A cofinal carry bound]
\label{prop:collapse}
If $|\operatorname{ihc}(A,N)|\le C\sqrt{N+1}+C'$ for arbitrarily
large $N$, with fixed constants $C,C'$, then $\delta=0$.
For the particular greedy support $A=G$ of $1/2$, a one-sided upper bound
$\operatorname{ihc}(G,N)\le C\sqrt{N+1}+C'$ already suffices.
\end{prop}
\leannote{Lean: \lproof{ErdosProblems/Erdos257/PaperCompleteR20/CofinalCarryCollapse.lean}{46}{half_of_cofinal_absolute_carry}, \lproof{ErdosProblems/Erdos257/PaperCompleteR20/CofinalCarryCollapse.lean}{91}{greedy_half_of_cofinal_upper_carry}.}

\begin{proof}
Lemma~\ref{lem:collapse-mech} gives
\[
 2^{N+1}|\delta|
 \le |\operatorname{ihc}(A,N)|+\mathrm T(N+1).
\]
The right side is $O(\sqrt{N+1})$ along the stated sequence, so division by
$2^{N+1}$ yields $\delta=0$.  For $A=G$, the greedy partial sums never
exceed $1/2$, hence $\delta\ge0$
(\lean{erdosSupportSeries\_greedyMersenneSupport\_le}{HalfCarryReachability.lean:888}).
Since $\mathrm T(N+1)\ge0$, the one-sided upper bound gives the same
conclusion.  Conversely, $\delta=0$ makes the carry equal to
$\mathrm T(N+1)$, so the bound holds with $C=2$ and $C'=4$.
This converse does not assert the bound for arbitrary prescribed constants.
For a prescribed $A$, the conclusion is $x_A=1/2$; it is not merely
membership witnessed by some other support.  For $A=G$, this equality is
equivalent to $1/2\in\mathcal A$
by \lean{mem\_mersenneAchievementSet\_iff\_greedy\_survival}{GreedyAchievementSet.lean:1458}.
\end{proof}

\begin{lem}[The terminal bound at square depths]
\label{lem:sqrt-witness}
Suppose $x_A = 1/2$ and $1 \notin A$. Then for every $k \ge 1$,
\[
  \operatorname{ihc}(A, k^2 - 1) \;=\; \mathrm{T}(k^2) \;\le\; 2k + 4 \;=\; \ensuremath{B}(k^2).
\]
\coord{half-carry recurrence} \scale{cofinal} \ev{Math}
\end{lem}
\leannote{Lean: \lproof{ErdosProblems/Erdos257/PaperCompleteR20/CarryCollapseCorrespondence.lean}{8}{square_depth_witness}.}

\begin{proof}
By Lemma~\ref{lem:collapse-mech} with $\delta = 0$, $\operatorname{ihc}(A,k^2-1) = \mathrm{T}(k^2)
\le 2\sqrt{k^2} + 4$. At a perfect square the real and truncated square roots agree exactly,
$\sqrt{k^2} = k = \lfloor \sqrt{k^2} \rfloor$, so the real envelope coincides with the discrete
strip $\ensuremath{B}(n) = 2\lfloor\sqrt n\rfloor + 4$
(\lean{halfStripBound}{HalfCarryReachability.lean:32}) with no slack required.
\end{proof}

At the depths $M=k^2$, the real square root and the integer square root
agree.  Thus the reverse implication for the terminal bound uses the
original constant $4$, not a relaxed constant $6$.  Truncating an achieving
support at $M=k^2$ leaves the carry at that depth unchanged, by
\lean{integerHalfCarry\_inter\_Iic\_eq\_of\_succ\_le}{HalfCarryReachability.lean:49}.
An achieving half-support excludes rank one because $w_1=1>1/2$.

\begin{prop}[Six equivalent membership conditions]
\label{prop:collapsed-list}
Each of the following is logically \emph{equivalent} to $1/2 \in \mathcal{A}$, not strictly weaker:
\begin{enumerate}[nosep,label=(\alph*)]
\item $C_G(N)\le2\sqrt N+4$ for every $N\ge0$
  (\lean{greedy\_half\_infinite\_of\_mobiusCenteredHalfCarry\_sqrtBound}{HalfCarryReachability.lean:834};
  \lean{greedy\_half\_infinite\_of\_mobiusCenteredHalfCarry\_upperBound}{HalfCarryReachability.lean:953});
\item there are arbitrarily large positive exponents outside $G$
  (\lean{CofinalPositiveHalfGreedySkips}{BooleanMobiusSkipRowCofinal.lean:22});
\item at arbitrarily large depths $n$, some $D\subseteq\{2,\ldots,n\}$
  satisfies $Q(D,n)=2^{n-1}-1$
  (\lean{CofinalExactLocalMersenneHalfRows}{BooleanMobiusCofinalExactRows.lean:38});
  no agreement between the sets $D$ at different depths is required;
\item at arbitrarily large $M$,
  $\operatorname{ihc}(G,M)\le B(M+1)$
  (\lean{GreedyHalfCarryCofinalStripReturn}{CofinalStripReturn.lean:76});
\item at arbitrarily large positive depths $M$, some
  $D\subseteq\{2,\ldots,M\}$ satisfies
  $|\operatorname{ihc}(D,M-1)|\le B(M)$
  (\lean{HalfCarryCofinalTerminalOnlyStrip}{TerminalOnlyCofinal.lean:34});
\item the half-target greedy remainder never exceeds the full remaining
  tail, so no finite fatal gap exists
  (\lean{ExistsFatalHalfGap}{HalfCutLocator.lean:526}); the seven equivalent
  forms in the later classification express this same condition.
\end{enumerate}
Forward directions are supplied formal (e.g.\ \lean{half\_mem\_mersenneAchievementSet\_iff\_greedySkippedSupport\_infinite}{GreedyAchievementSet.lean:2583},
\lean{half\_mem\_mersenneAchievementSet\_of\_cofinalExactLocalRows}{BooleanMobiusCofinalExactRows.lean:71},
\lean{existsFatalHalfGap\_iff\_half\_not\_mem\_mersenneAchievementSet}{HalfCutLocator.lean:643}).
For (d) and (e), the reverse directions use the square depths in
Lemma~\ref{lem:sqrt-witness}.  All six equivalences, including these reverse
directions, are now Lean-checked in
\lean{six\_membership\_conditions}{lean/ErdosProblems/Erdos257/PaperCompleteR20/SixMembershipConditions.lean:123}.
For (a), membership gives
$C_G(N)\le2\sqrt{N+1}+3\le2\sqrt N+4$ when $N\ge1$, and $C_G(0)=0$.
The forward implication still needs the greedy inequality $x_G\le1/2$.
\coord{half-carry recurrence ;  exact-row} \scale{cofinal} \ev{Lean}\ev{Math}
\end{prop}

\textbf{Scope of the equivalences.} These conditions are equivalent to membership, not strictly
weaker substitutes. Equivalence does not make a formulation useless as an invariant, but it does
mean that proving any of them is the same problem as (H). The equivalent conditions remain
unproved. In particular, mutually \emph{incompatible} witnesses at different endpoints do not buy
a weaker sufficient condition, and the $\Pi^0_2$ shape of these hypotheses is not evidence of extra
room.

\textbf{Scope limit, preserved exactly.} Item (a) does \emph{not} extend to every $A$ with
$1 \notin A$. Dropping the greedy-specific fact $x_G \le 1/2$ yields only the one-sided statement
$\texttt{hbound} \iff (1/2 - x_A) \le 0$. Thus the conclusion for a general support is a one-sided inequality,
not equality with $1/2$.

\subsection{The size of the required error bounds}
\label{bar:scale}

One sufficient half-membership argument requires the reset inequality
\[
 |\mathrm{rem}(r+1)-2^{r+1}|>2^{(r+5)/2}
 \qquad(r\ge10\text{ an upper or middle reset}).
\]
For the normalised deviation
$\lambda_{r+1}=(\mathrm{rem}(r+1)-2^{r+1})/2^{r+1}$, this is
$|\lambda_{r+1}|>2^{(3-r)/2}$.  A near-integer estimate for a
weighted divisor sum is a proposed way to obtain information about
this inequality, not an equivalent condition established here.
Section~\ref{sec:invent-R1} displays the exact recurrence and the
additional, changing skip-set term.  Each obstruction below has its
own hypotheses; none excludes every method at the stated scale.

\textbf{Class (a): fixed precision or bounded state.} The following
three statements concern abstract recurrence models.  Their
conclusions apply to the Mersenne problem only within the specified
models; they do not use its full divisor-count constraints.

\begin{prop}[Limits of the specified local summaries]
\label{prop:local-void}
(i) After $L$ common steps the endpoint residue mod $2^L$ of an affine binary orbit is
\emph{independent} of the initial carry: $u(L) - v(L) = 2^L(u_0 - v_0)$
(\lean{affineBinaryOrbit\_mod\_twoPow\_eq}{GenericTailOrbitRigidity.lean:307}).
(ii) Fix $m\ge2$ and put $h=\lfloor(m+1)/2\rfloor$.  Consider the
coefficient sequences
\[
 c_r(m)=h-r,\qquad c_r(m+1)=2r,\qquad
 c_r(j)=0\ (j\ne m,m+1),\qquad 0\le r\le h.
\]
Their binary sums are all $h2^{-m}$, and their scaled tails before
position $m$ agree, but their scaled tails immediately after $m$ are
$r$.  Thus any exact label-and-decoder system recovering those latter
tails must have at least $h+1=\lfloor(m+1)/2\rfloor+1$ labels.  A state
determined only by the common preceding history, with no new input,
cannot distinguish them
(\lean{balancedPulse\_no\_autonomous\_decoder}{GenericTailOrbitRigidity.lean:247}). (iii) For any
starting carry and any finite list of prescribed valuations and
odd unit residues modulo $2^u$, with $u\ge1$, there are integer input
coefficients with those data whose successor carries lie in their
prescribed centred intervals.  The unrestricted higher bits of the
input coefficients may be chosen separately at each step, as in
Proposition~\ref{prop:2adic-nogo}
(\lean{fixedPrecisionTropicalNoGo}{TropicalCurvatureCarry.lean:137}), so these fixed-precision symbols alone do not exclude a completion in
that specified family.  This says nothing about extra constraints imposed
by divisor counts from one common support. \coord{carry} \scale{uniform} \ev{Lean}
\end{prop}
\leannote{Lean: \lproof{Erdos249257/GenericTailOrbitRigidity.lean}{307}{affineBinaryOrbit_mod_twoPow_eq}, \lproof{Erdos249257/GenericTailOrbitRigidity.lean}{200}{balancedPulse_weighted_pair}, \lproof{Erdos249257/GenericTailOrbitRigidity.lean}{209}{balancedPulse_endpoint_fanout}, \lproof{Erdos249257/GenericTailOrbitRigidity.lean}{228}{balancedPulse_label_card_lower_bound}, and 2 further declarations in the \hyperref[sec:coverage]{coverage section}.}

For an integer recurrence $J(M+1)=2J(M)-\delta(M+1)$,
reduction modulo $2$ gives $J(M+1)\equiv\delta(M+1)\pmod2$.
For the particular recurrence $K(M+1)=2K(M)-(\tau(M+1)-1)$,
with $K(2)=1$, this gives, for $M\ge2$, that $K(M)$ is even
exactly when $M$ is a square: divisors pair off except at a square.
This is a genuine restriction at a specified index, but it is already
encoded in the known forcing term.  It supplies no additional
run-length estimate by itself. \ev{Math}

\textbf{Class (b): global averages and Diophantine data.}

\begin{prop}[Weighted denominator budget]
\label{prop:exponent-gap}
Let $n\ge2$, let $\mathrm{Skip}_n\subseteq\{2,\ldots,n-1\}$,
and let $D_n$ be the reduced denominator of the associated finite sum.
Since $D_n$ divides the product of its Mersenne denominators,
\[
 \log_2D_n\le\sum_{d\in\mathrm{Skip}_n}\log_2(2^d-1)
 \le\sum_{d\in\mathrm{Skip}_n}d\le\frac{n(n-1)}2-1.
\]
The middle inequality is strict when the skip set is nonempty; for an
empty skip set both sums are zero.  This is an upper bound, not a
quadratic asymptotic for $D_n$.  Combined with the cited fixed
irrationality-measure exponent, it does not supply the required
$2^{-3n/2}$-scale estimate.  Even an asymptotic for the exponent sum would
need additional information about the skip distribution.
\coord{diophantine} \scale{cofinal} \ev{Math}
\end{prop}
\leannote{Lean: \lproof{ErdosProblems/Erdos257/PaperCompleteR21/DenominatorBudget.lean}{43}{skipSum_den_dvd_prod}, \lproof{ErdosProblems/Erdos257/PaperCompleteR21/DenominatorBudget.lean}{78}{weighted_denominator_budget}.}

\textbf{Scope of these obstructions.} The common pre-index history in
Proposition~\ref{prop:local-void}(ii) does not determine the post-index
tail in its countermodel family.  This does not exclude a state using
additional arithmetic input from the actual support.  The displayed
denominator budget combined with the cited fixed irrationality-measure
exponent does not yield the required run bound. These calculations do not
exclude averaging over shifts, residues, or observation scales, nor a sharper analysis of the
actual skip distribution.

\textbf{Prime steps and bounded windows.} In the universal recurrence,
a prime step subtracts $1$, so it sends $K(M)$ to $2K(M)-1$; it is
not exact doubling.  Observation~\ref{obs:prime-doubling} gives an
explicit frozen-row drift with values between $2$ and $6$ at the eleven
indices $607,\ldots,617$, including the three prime indices $607$, $613$,
and $617$.  This finite example shows that prime steps can occur within
a small-value window.  It neither proves statistical independence nor
rules out every argument using prime-counting information. \ev{Cert}

\begin{rem}[Status of the exhaustion claim]
The preceding counterexamples concern the specified finite summaries and
estimates.  They do not classify all available tools, even within the
broader categories of local information and global averages.  No exhaustion
claim is used in a subsequent deduction.
\end{rem}

\subsection{What finite certificates decide}
\label{bar:asym}

\begin{thm}[One-sidedness]
\label{thm:one-sided}
Non-membership of $1/2$ in $\mathcal A$ has an effectively checkable
finite-certificate formulation, hence a $\Sigma^0_1$ formulation.
Membership has the complementary $\Pi^0_1$ formulation. Survival through a tested finite
depth alone does not establish membership; a uniform theorem or inductive invariant could.
The arithmetical-hierarchy form by itself proves neither undecidability nor the absence of finite
proofs.
\coord{meta} \scale{uniform} \ev{Lean}\ev{Math}
The linked declarations establish the fatal-gap equivalences.  The
arithmetical-hierarchy classification also uses the effective tail
estimate explained below; it is an ordinary computability deduction,
not a separately replayed Lean theorem.
\end{thm}
\leannote{Lean: \lproof{ErdosProblems/Erdos257/PaperCompleteR21/OneSidedCertificateHierarchy.lean}{492}{paper_one_sidedness}, \lproof{ErdosProblems/Erdos257/PaperCompleteR21/OneSidedCertificateHierarchy.lean}{291}{existsFatalHalfGap_iff_exists_certificate}, \lproof{ErdosProblems/Erdos257/PaperCompleteR21/OneSidedCertificateHierarchy.lean}{174}{existsFatalHalfGap_of_certificate}, \lproof{ErdosProblems/Erdos257/PaperCompleteR21/OneSidedCertificateHierarchy.lean}{243}{certificate_of_existsFatalHalfGap}, and 4 further declarations in the \hyperref[sec:coverage]{coverage section}.}

\begin{proof}
The dichotomy $1/2 \in \mathcal{A}$ or $\exists$ a fatal gap is unconditional
(\lean{half\_mem\_mersenneAchievementSet\_or\_exists\_fatal\_gap}{HalfCutLocator.lean:623}), and the
two alternatives are exactly complementary
(\lean{existsFatalHalfGap\_iff\_half\_not\_mem\_mersenneAchievementSet}{HalfCutLocator.lean:643}),
with \lean{ExistsFatalHalfGap}{HalfCutLocator.lean:526} a finite witness type. In the other
direction, certified death implies non-membership
(\lean{certifiedGreedyMersenneDeath\_not\_mem}{GreedyAchievementSet.lean:1762}, over
\lean{CertifiedGreedyMersenneDeath}{GreedyAchievementSet.lean:1756}; $3/4$ is certified dead at
level $1$, lookahead $0$, by
\lean{three\_fourths\_certifiedGreedyMersenneDeath}{GreedyAchievementSet.lean:1778}), whereas the
absence of a certificate through depth $D$ alone gives no conclusion
about later depths.  To make the search effective, note that a greedy
remainder at a finite rank is rational.  The tail $R_n$ has rational
upper bounds obtained by summing through a cutoff $M$ and bounding the
remaining sum by $2^{-M}+(2/3)4^{-M}$.  These upper bounds converge to
$R_n$.  Thus the strict inequality $r_n(1/2)>R_n$, when true, is witnessed
by some finite cutoff and a comparison of rational numbers.  Enumerating
ranks and cutoffs gives the asserted finite-certificate search.
Membership requires the survival inequality at every rank
(Theorem~\ref{thm:greedy-survival}).  Checking a finite initial segment
alone does not supply the later inequalities; a separate uniform argument
or induction would be needed.
\end{proof}

\textbf{Finite search versus infinite survival.}
A discovered fatal-gap certificate proves non-membership.  Further
surviving prefixes give larger verified ranges, but do not establish
membership without an argument covering all ranks.  This distinction
is unchanged whether one increases the geometric truncation depth,
the number of quotient rows, or the depth of a grouped computation.
The different finite parameters should not be reported as a single
verification depth.

\subsection{Why measure and category do not decide a rational target}
\label{bar:measure}

\begin{rem}[Achievement-set geometry does not decide rational-target membership]
\label{prop:soft}
The achievement set has positive measure and is nowhere dense
(Theorem~\ref{thm:geometry}). Neither fact alone decides membership of a specified rational.
Injectivity fixes the support of an achievable value, but it does not exclude averaging over
indices, shifts, residues, observation scales or auxiliary constructions. The recorded
calculations obstruct the particular models and estimates analysed here; no impossibility theorem
for all probabilistic, measure-theoretic or averaging methods is claimed.
\coord{achievement-set} \scale{uniform} \ev{Math}
\end{rem}

The geometry itself is Lean-checked in Theorem~\ref{thm:geometry}: $\mathcal{A}$ has Lebesgue
measure exactly $1$ inside an ambient span of $E \approx 1.6067$, a relative density of about
$62\%$, and is simultaneously nowhere dense, perfect and totally disconnected. Those facts may hold
together. They do not yield competing theorems about specified rationals. Showing that
$\mathcal{A}$ is ``small'' cannot decide (U) by measure: it is $62\%$ of its ambient interval.

\textbf{Recorded probabilistic heuristic.}
The supplied record reports an expected count of about
$2\times10^{-4}$ beyond $n=31$ in a Bernoulli$(1/2)$ digit model.
This number is not a probability assigned to membership of the fixed
point $1/2$.  Such an interpretation would require a specified
probability space, a definition of the failure events, and a link
between those events and the deterministic greedy orbit.  None is
provided here.  Section~\ref{sec:invent-R1} separates a reproducible
finite phase experiment from its unproved probabilistic application.
\ev{Open}

\subsection{Hypotheses of the CRT block-certificate argument}
\label{bar:box}

\begin{obs}[Two inputs to the CRT certificate argument]
\label{obs:box}
The CRT block-certificate arguments considered here require both
compatible divisibility prescriptions and a tail estimate strong enough
for the height inequality $q(C+N+L+2)<b^L$.  Pairwise coprimality is one
way to obtain the first input; reciprocal summability is one way to obtain
the second.  This observation concerns those certificate arguments only.
It does not describe all theorems in Section~\ref{sec:257-problem}:
the reciprocal-summable and weighted averaging proofs do not require
pairwise coprimality.
\coord{block-certificate hypotheses} \scale{uniform} \ev{Lean}\ev{Math}
\end{obs}

\textbf{What remains to check.} One cannot apply a conditional CRT
certificate theorem to arbitrary infinite $A$ without establishing its
hypotheses.  Shared factors impose compatibility conditions on the
prescribed residues; they do not make every CRT construction impossible.
The tail budget in the cited proof is used at
\texttt{CertificateKernel.lean:9586--9598}.  Deleting that assumption does
not complete the proof: one needs a replacement estimate for that step.
This leaves open different constructions and different averaging arguments.

\textbf{Explains only a method boundary.} The primes are infinite and pairwise coprime, so they satisfy the
first input.  The identity
\lean{compositeDilationDefect\_eq\_zero\_of\_prime\_support}{CompositeDilationDefect.lean:103} makes
the dilation algebra exact there; yet this certificate theorem does not
apply because $\sum_p1/p$ diverges, so the stated reciprocal-sum
hypothesis fails.  This explains
the absence of a \emph{Lean theorem from this engine}; it does not make the
base-$2$ value open, since Tao--Ter\"av\"ainen settle it by a different
analytic method.  The sunflower generalisation inherits both axes and
therefore does not recover that external theorem.

\subsection{The distinct estimates still required}

The half-target analysis involves a geometrically weighted divisor sum
$\sum_{i\le L}(\tau(M+i)-1)2^{-i}$ and possible near-integer estimates.
One universal-support certificate instead asks for an opening divisible
block together with a controlled middle window.  These inputs are related
through incidence identities, but the displayed arguments do not prove
them equivalent for all supports.  They should therefore be retained as
separate obligations.  Digit-pattern occurrence results do not by
themselves supply either required uniform estimate.

\section{Unproved inputs for further arguments}
\label{sec:257-survivors}

The following proposals isolate inputs not supplied by the preceding
results.  Some are sufficient conditions, some are equivalent formulations,
and one concerns reproducing an already known prime-support theorem.
They are not an exhaustive list of possible approaches.  The relevant
question in each case is what must be proved, and whether the cited
obstruction applies to that particular statement.

\subsection*{A claimed quantitative support condition}
The absent analytic supplement claims a second criterion admitting
reciprocal-divergent supports.
Put $H_A(x)=\sum_{a\in A,\ a\le x}1/a$, $T_0=2$, $T_{j+1}=2^{T_j}$,
and $\ell(x)=\min\{j:x\le T_j\}$.  The condition
$H_A(x)=o(\ell(x))$ is claimed to imply all-base irrationality, with no
positive density assumption on $A$.  Neither a proof of this implication
nor the claimed sharper rational-phase lower bound is present in this
checkout.


\subsection*{Two further analytic tests}
Small displacements at a real base do not by themselves give the fixed
rational lattice used in the integer-base proof.  For
$\beta=(3+\sqrt5)/2$, let $u=\beta^{-1}$ and
$H=\{2^km:k\ge2,\ m\text{ odd},\ m\le2^{2^k}\}$.
The finite prefix identity and an algebraic norm show that the condition
\[
 \Delta_{\beta,H}(N)
 \left(1+\sum_{\substack{a\in H\\a\le N}}
                 \sum_{j=1}^{\lfloor N/a\rfloor}u^{N-ja}\right)
 \longrightarrow0
\]
along some sequence of positive integers $N\to\infty$ would exclude $X_A(\beta)$ from $\Q(\sqrt5)$ for every
infinite $A\subseteq H$.  The second factor controls the algebraic
conjugate of the finite prefix; smallness of the displacement alone does
not control this product.  The product estimate remains unproved.

To justify the conditional implication, suppose instead that
$X_A(\beta)\in\Q(\sqrt5)$, and choose an integer $q\ge1$ with
$qX_A(\beta)\in\Z[\beta]$.  The finite-prefix identity gives, for
$N\ge1$,
\[
 Y_N:=q\Delta_{\beta,A}(N)
 =q(\beta^N-1)X_A(\beta)
  -q\sum_{\substack{a\in A\\a\le N}}
           \sum_{j=1}^{\lfloor N/a\rfloor}\beta^{N-ja}
 \in\Z[\beta].
\]
The infinite support makes $Y_N>0$.  Since
$\beta^2-3\beta+1=0$, conjugation $\sigma$ sends $\beta$ to
$u=3-\beta\in(0,1)$.  For $m,n\in\Z$ the norm of $m+n\beta$
is $m^2+3mn+n^2$, so the nonzero $Y_N$ has nonzero integer norm.
Moreover,
\[
 |\sigma(Y_N)|\le q\left(|\sigma(X_A(\beta))|
       +\sum_{\substack{a\in H\\a\le N}}
             \sum_{j=1}^{\lfloor N/a\rfloor}u^{N-ja}\right).
\]
Using $0<\Delta_{\beta,A}(N)\le\Delta_{\beta,H}(N)$, we obtain
\[
 1\le |Y_N\sigma(Y_N)|
 \le q^2\max\{1,|\sigma(X_A(\beta))|\}\,
       \Delta_{\beta,H}(N)
       \left(1+\sum_{\substack{a\in H\\a\le N}}
                  \sum_{j=1}^{\lfloor N/a\rfloor}u^{N-ja}\right).
\]
The asserted limit would contradict this inequality.  The multiplier
$q$ and the conjugate of the putative value are fixed, independently
of $N$.  This proves only the implication from the displayed product
estimate, not that estimate.  \ev{Math}

A different extension requires squarefree remote-tail control and a
weighted displacement to be small at the same indices.  CRT establishes
compatibility of the finite opening divisibility conditions, but not this
simultaneous smallness.  Neither argument is claimed to establish its
required analytic hypothesis.


\subsection{A lower bound at reset ranks}
\label{ssec:route-1}

\textbf{What it needs.} Either of the first two sufficient conditions
below. Their different scales do not establish an implication between
them. The sign condition (iii) is a separate question; no implication
from it to half-membership is proved here.
\begin{enumerate}[nosep,label=(\roman*)]
\item \emph{Reset-deviation bound.} For every upper or middle reset
  row $r\ge10$, require
  $|\mathrm{rem}(r+1)-2^{r+1}|>2^{(r+5)/2}$.
  Section~\ref{sec:o4} states the exact conditional implication to
  membership.  A run-length estimate is a possible way to approach
  this bound, not an established exact reformulation of it.
\item \emph{Lean-facing form, already proved as a conditional implication.} \texttt{SeamUpperResetDyadicBandEscape}: for every
  $d \ge 13$ with \texttt{successorCarries} and every $j \le d$, the reset charge
  $4\cdot\text{overshoot} + \text{abovePulse}$ avoids the linear-width band
  $(2^{d-j+1} - 2(d+j),\, 2^{d-j+1}]$.
\item \emph{Sign law alone.} At every reset, M-resets have $\mathrm{dev} > 0$ and U-resets have
  $\mathrm{dev} < 0$. This specifies the sign but not a quantitative distance from zero.
  Neither an implication to half-membership nor a strict comparison with
  the first two conditions is established here.
\end{enumerate}
\textit{Sources.} \lean{SeamUpperResetDyadicBandEscape}{HalfCylinderMiddleCarryLowerBound.lean:4566}
with conditional implication
\lean{half\_mem\_mersenneAchievementSet\_of\_upperResetDyadicBandEscape}{HalfCylinderMiddleCarryLowerBound.lean:4790};
quantifier collapse \lean{dyadicBandEscape\_iff\_exists\_critical}{HalfUpperResetCriticalBand.lean:108},
which reduces the $\forall j \in [0,d]$ band check to one nearest-boundary index per row; finite
verification \lean{seamUpperResetDyadicBandEscape\_through\_thirty}{HalfCylinderUpperResetBandCertificates.lean:78};
alternative implication \lean{LargestSkipLateStepSocket}{HalfCylinderLargestSkipInduction.lean:34} with
conditional implication at \texttt{:167} and base case
\lean{largestSkipLateAt\_fourteen}{HalfCylinderLargestSkipInduction.lean:73}.

\textbf{Logical strength.} The reset bound is a sufficient condition for
half-membership in the cited argument.  No converse is established here.
The fact that it is quantitative, while membership is qualitative, does
not prove that the condition is strictly stronger: that would require a
separate implication or counterexample.  The carry identity of
Proposition~\ref{prop:collapse} does not itself furnish this lower bound
on the reset deviation.

\textbf{What the finite checks establish.} The reported computation
checks the reset inequality over its stated range, with margins of
$1.1119$ bits at row $14$ and about $1246$ bits at row $2500$; the reported
mean is $627$ bits.  These are finite observations, not a proof of an
asymptotic lower bound or of the inequality at every later reset.
\ev{Cert}

Uniqueness in Section~\ref{bar:canon} does not prevent this argument:
the proposed bound concerns the actual greedy sequence.  It does not
require alternative supports for the same target.

\textbf{Run lengths.} The record reports a maximum right-branch run of
length $19$ at row $158{,}096$, far below the proposed threshold
$(r-3)/2=79{,}046.5$ there.  Observing larger runs at larger ranks does not
refute the existence of some uniform finite bound.  Neither a uniform
bound nor the weaker asymptotic assertion $L_r=o(r)$ is proved by these
data.  The stated implication needs the required inequality beyond the
verified initial ranks.  \ev{Cert}

The missing step is the quantitative lower bound itself.  The finite
observations do not supply the growing-precision estimate discussed in
Section~\ref{bar:scale}. \ev{Open}

\subsection{Increasing the depths of exact finite sums}
\label{ssec:route-2}

\textbf{What it needs.} Rule out the recycle branch returning to a bounded endpoint forever. Any
one of:
\begin{enumerate}[nosep,label=(\alph*)]
\item \codeid{SkippedCoreCriticalQuotientSupply}: for every below-half core $D$ bounded by $[2,c)$
  with genuine crossing deficit, $2^{(2c-2)-1} \le \ensuremath{Q}\,(\mathrm{insert}\
  c\ D)\,(2c-2)$;
\item a direct proof that the crossing rank $c$ produced by the recycle branch cannot repeat the
  value $4$, or any bounded value, infinitely often;
\item a growth law for $c$ as a function of the incoming endpoint $n$.
\end{enumerate}
By Proposition~\ref{prop:canon}(ii), form (a) need only be checked along the \emph{one} canonical
orbit \texttt{halfGreedyPrefixSupport}, not searched over $D$; and by the recorded sharp-capacity
biconditional it is exactly the $c-2$ bit capacity test
$\ensuremath{S}\,D\,1\,(2c-2) < 2^{c-2}$. The fixture $D=\{2,3\}$, $c=5$ shows that the fractional-mass
sufficient condition is not necessary.  It does not rule out proving that
condition for a suitably restricted family.

\textit{Sources.} \lean{exactLocalMersenneHalfRow\_double\_or\_recycle}{BooleanMobiusExactRowDichotomy.lean:26};
the formal falsifier
\lean{exists\_seeded\_bounded\_double\_or\_recycle\_model}{BooleanMobiusExactRowDichotomy.lean:91}
(the model $n \mapsto (n = 6)$ satisfies the seed and the exact transition shape yet is not
cofinal, so the schema plus an endpoint-six seed provably does \emph{not} give cofinality);
conditional strict progress
\lean{skippedCoreSharpCapacity\_below\_or\_strictly\_laterExactRow}{BooleanMobiusExactRowCrossing.lean:233};
unconditional recycling required input
\lean{exists\_skippedCoreExactRow\_of\_value\_above}{BooleanMobiusExactRowCrossing.lean:155};
hypothesis \lean{SkippedCoreCriticalQuotientSupply}{BooleanMobiusCriticalCapacityCofinal.lean:30};
capacity iff
\lean{localBinarySuffix\_two\_mul\_sub\_two\_lt\_criticalCapacity\_iff}{BooleanMobiusSkippedCoreCriticalCapacity.lean:22};
seed fixture $\{2,3,6\}$ at
\lean{exactLocalMersenneHalfRow\_six}{BooleanMobiusExactRowSeed.lean:34}.

\textbf{The finite-row test and its unique candidate.}
The real-prefix argument in Proposition~\ref{prop:canon}(ii) does not
by itself identify an exact-row witness.  The integer weights have a
separate, decisive property: for fixed $n\ge2$,
\[
 q_{n,d}=\left\lfloor\frac{2^n}{2^d-1}\right\rfloor
 \quad(2\le d\le n)
\]
form a decreasing sequence in which each weight exceeds the sum of all
later weights by at least $1$.  Apply the integer greedy rule with
capacity $2^{n-1}-1$: visit $d=2,\ldots,n$ and subtract $q_{n,d}$ when
it fits.  By Theorem~\ref{record:257bm-i11a} and Propositions~\ref{record:257bm-i11b}--\ref{record:257bm-i11c},
any exact-row witness has zero remainder and must be this greedy word.
Thus there is at most one witness at a fixed depth, and its existence is
tested by checking whether the final greedy remainder is zero.
An exhaustive search over $2^{n-1}$ subsets is unnecessary.

This integer greedy word should not be silently identified with the
real greedy prefix $G\cap[2,n]$: the two rules use different weights.
The separation theorem proves uniqueness within the integer problem,
not that identification.

\textbf{Finite decidability and the cofinal requirement.}
The test decides the exact-row predicate at any specified depth.  The
claim that successful depths are arbitrarily large is still equivalent
to half-membership, by Proposition~\ref{prop:collapsed-list}(c).
Neither uniqueness nor finite decidability supplies that cofinal
assertion.  The double-or-recycle argument must prove genuine progress
in the depths, not rely on a purported freedom to choose many witnesses
at one depth.

\subsection{The known prime-support theorem}
\label{ssec:route-3}

The base-$2$ prime-support value is not an open target:
\[
  \sum_{p\ {\rm prime}}\frac1{2^p-1}
  =\sum_{n\ge1}\frac{\omega(n)}{2^n}
\]
by expanding each reciprocal Mersenne term geometrically and collecting
divisors.  Tao and Ter\"av\"ainen prove the right-hand series irrational in
Theorem~1.3 of
\href{https://arxiv.org/abs/2512.01739}{\emph{Quantitative correlations and
some problems on prime factors of consecutive integers}}.  Their proof is
analytic and is cited here; it is not formalised in the pinned Lean corpus.

The local pairwise-coprime certificate theorem still does not cover the
primes because it assumes summability of $\sum_p1/p$.  That is a boundary of
this engine, not a boundary of the mathematics.  The previously proposed
windowed-tail repair is therefore withdrawn as a research contribution.  It
may remain useful as an audit question about whether the existing Lean engine
can reproduce the known theorem, but reproducing a cited result is not new
progress on \#257.  After Theorem~1.3 the authors state the
prime-support extension to every integer base $b\ge2$ and, separately,
the full prime-power result at base $2$, leaving those modifications to
the reader.  These statements should not be combined into an attributed
all-base theorem for arbitrary prime-power subsets.

\subsection{Equivalent membership conditions}
\label{ssec:route-4}

Proposition~\ref{prop:collapsed-list} lists the equivalent conditions:
a pointwise upper bound on $C_G$, infinitely many skipped ranks, exact
quotient sums at unbounded depths, and the two cofinal carry bounds.
For the carry conditions, the residual identity and its tail estimates are
the relevant inputs:
\lean{integerHalfCarry\_eq\_scaled\_residual\_add\_tail}{HalfCarryReachability.lean:871},
\lean{binaryCoeffTail\_nonneg}{GenericTailOrbitRigidity.lean:78},
\lean{binaryCoeffTail\_supportCoeff\_le\_two\_sqrt\_add\_four}{BooleanMobiusCarry.lean:290}, and, for the terminal-strip conditions, evaluation of the square root at
a perfect square.  The one-sided scope limitation in
Proposition~\ref{prop:collapsed-list} still applies to general supports.

\textbf{Purpose of the equivalences.} These results identify the exact
logical content of the proposed hypotheses.  They do not settle
half-membership, nor do they make an equivalent formulation useless: a
different expression of the same condition may expose an invariant or a
more convenient induction.  The formalisation task is separate from the
unproved arithmetic assertion on either side.

\subsection{Why zero windows do not split the support class}
\label{ssec:route-5}

One proposed split asks whether every infinite support satisfies
\[
 \begin{gathered}
 \text{$c_A$ has arbitrarily late super-logarithmic zero windows}\\
 \text{or a weighted block certificate exists at every precision.}
 \end{gathered}
\]
For the present divisor counts, the first alternative never occurs.
Writing $a_0=\min A$, every run of zeros has length at most $a_0-1$,
since each $a_0$ consecutive integers contain a multiple of $a_0$.
The logarithmic estimate
\lean{supportCoeffZeroWindow\_length\_le\_eps\_logb}{SublogDivisorCoverage.lean:435}
therefore does not split the support class into two useful cases.
The lcm-gap theorem
\lean{irrational\_erdosSum\_of\_lcm\_gap}{CertificateKernel.lean:5883}
remains an independent irrationality result, not an instance of this
zero-window alternative.  The proposed assertion consequently requires
block certificates for every infinite support: for each required precision,
one must construct an opening block with
$2^K\mid\sum_{r=1}^Kc_A(N+r)2^{K-r}$, a bounded middle contribution $C$,
and $q(C+N+L+2)<2^L$.

This reformulation does not remove the CRT compatibility and tail-estimate
requirements of Section~\ref{bar:box}.  It asks for a new construction that
meets them for general supports.  Nor does it identify that construction
with the reset estimate used for the half-target: the two arguments use
related divisor-count identities, but no equivalence of their missing
hypotheses is proved here. \ev{Open}

\subsection{Comparison of the remaining inputs}

Route~4 packages recorded equivalences; it does not establish their
unproved existence hypotheses.  Route~3 must be compared with the external
prime-support theorem of Tao and Ter\"av\"ainen
\cite[Theorem~1.3]{taoteravainen2025}; reproducing it in this framework
requires its own analytic input.  Route~2 offers a different finite
search space under the conditions stated there.  Routes~1 and~5 share
some incidence identities, but this is not a proof that every missing
hypothesis is equivalent.  The appropriate comparison is the input each
route needs, its quantifiers, and what a successful estimate would imply.

\section{Reading the detailed record}
\label{sec:257-howto}

The opening note has already defined the evidence bands, scale tags,
coordinates, and catalogue terminology.  The longer record has two different
reading modes, and they should not be confused.

\paragraph{Read linearly.} Sections~\ref{sec:257-wall} and
\ref{sec:257-survivors} form one argument: the first identifies the recurring
obstructions, and the second isolates the routes those obstructions do not
exclude.  Read them in order before selecting a new attack.

\paragraph{Use as reference.} The mathematical-ingredient catalogue may be
entered by coordinate: the certificate method and universal-direction record;
the greedy and achievement-set geometry; the integer-quotient model, quotient identity, and
dynamics; the half-carry and Boolean--M\"obius coordinates; and the support
rigidity results.  The later implication tables record what chains with what,
while the near-miss index gives the exact mismatch at each open obligation.

\paragraph{Scale and coordinate discipline.} The direct certified instances
remain fixed or bounded; rungs $J \le 22$, seam rows $\le 200{,}000$, ranks
$\le 3000$, and macro depth $\le 13{,}548{,}057$; whereas the open
obligations are cofinal.  Theorem~\ref{thm:one-sided} shows that extending a
bounded range does not establish a cofinal statement.  Coordinate tags are
equally literal: an obstruction in \coord{half-carry recurrence}, \coord{quotient-row}, or
\coord{diophantine} does not transfer to another representation without a
proved map.  Coordinate-free barriers are marked explicitly.
% ; - part a257_p0 ; -
\clearpage
\part*{III. Detailed statement and arithmetic catalogue}
\addcontentsline{toc}{part}{III. Detailed statement and arithmetic catalogue}
\section{The series, the greedy rule and finite identities}

\subsection{The object}

In this part, write $x_n=w_n=(2^n-1)^{-1}$ for every positive
integer $n$, and $x_A=X_A(2)=\sum_{n\in A}x_n$ for
$A\subseteq\mathbb N_{>0}$.  These series converge by comparison
with $\sum_n2^{-n}$.  Exponent $1$ contributes the rational number
$x_1=1$, so deleting or adjoining it does not affect irrationality.
It cannot occur in a representation of $1/2$; those calculations
therefore start at exponent $2$.

\begin{defn}[Universal Problem \#257]
\textbf{Is $x_A = \sum_{n \in A} 1/(2^n-1)$ irrational for \emph{every} infinite
$A \subseteq \mathbb{N}_{>0}$?}
The quantifier ranges over all infinite sets of positive exponents. It is
\textbf{OPEN}. \ev{Open}
\end{defn}

The support theorems prove irrationality for specified infinite
families.  The following question instead asks whether one particular
rational number is represented.  A positive answer would disprove
the universal assertion; a negative answer would settle only that
target.  The later analysis also treats $1/21$.

\begin{defn}[The distinguished half-value question]
\label{defn:half-question}
Define the \emph{Mersenne achievement set}
$\mathcal{A} := \{\,y \in \mathbb{R} : \exists A \subseteq \mathbb{N},\ 0 \notin A,\ y =
\sum_{n \in A} x_n\,\}$, the set of reals realizable as a subset sum of the positive-index Mersenne
weights $\{x_n\}_{n \ge 1}$. \textbf{Is $1/2 \in \mathcal{A}$?} This is \textbf{OPEN}. \ev{Open}
\end{defn}

\begin{obs}[Why the half-value question controls universal \#257]
\label{obs:half-controls-universal}
Any finite subset sum of $\{x_n\}$ has odd reduced denominator, since every $2^n - 1$ is odd; hence
no finite $A$ can sum to $1/2$ (\lean{positiveMersenneSupportValue_coe_finset_ne_half}{HalfCutLocator.lean:243}).
Consequently, \emph{if} $1/2 \in \mathcal{A}$, the witnessing support $A$ is necessarily infinite, and
$x_A = 1/2$ is rational. Such an $A$ is then a single infinite set for which the universal-\#257
sum is rational -- an outright refutation of universal \#257. Conversely, universal \#257 being true
forces $1/2 \notin \mathcal{A}$ as one instance among uncountably many. The half-value question is therefore one concrete test of the negation,
not a restatement of the universal problem.  Excluding $1/2$ would
not exclude every other rational value, even though the support
theorems already settle many infinite families.  Representing $1/2$
would give an infinite-support counterexample to the universal
assertion. \scale{n/a} \ev{Math}
\end{obs}

\subsection{What is proved}

\subsubsection{Full-support irrationality, for every base}

\begin{thm}[Full-support irrationality, unconditional, every base $b \ge 2$]
\label{thm:full-support}
For every integer $b \ge 2$,
\[
\sum_{k=0}^{\infty} \frac{1}{b^{k+1}-1} \quad \text{is irrational.}
\]
\textbf{Hypotheses:} $b \in \mathbb{N}$, $b \ge 2$. \textbf{Conclusion:} unconditional irrationality
of the full-support Erd\H{o}s--Borwein-type series at every base, not merely base $2$.
\coord{full-support} \scale{uniform} \ev{Lean}\\
\textit{Source:} \lean{irrational\_erdosSum\_full\_support}{CertificateKernel.lean:8328}
(`$b : \mathbb{N}$, hb : $2 \le b$'), recorded via a certificate machine (bounded Bertrand/CRT
first-block frame, middle-window divisor-pair average with pigeonhole selection, explicit parameter
closure). The base-$2$ instance recovering the classical Erd\H{o}s--Borwein constant's irrationality
is \lean{irrational\_erdosBorwein\_series}{CertificateKernel.lean:8335}, a one-line corollary.
\end{thm}

This formalizes the positive-integer-base scope of
\href{https://www.renyi.hu/~p_erdos/1948-04.pdf}{Erd\H{o}s's classical 1948 full-support theorem},
which already covered integer bases of absolute value greater than one. The base-two
Erd\H{o}s--Borwein constant is a particular case; the Lean theorem is recorded with no open
hypotheses.

\subsubsection{Structured infinite supports}

Beyond full support ($A = \{1,2,3,\dots\}$, i.e.\ $m=1$ below), the following named infinite
families are proved irrational for every integer base $b\ge2$.  The purely
periodic, eventually periodic, residue-class and odd supports follow from the periodic theorem of
Luca and Tachiya~\cite[Theorem~A, p.~139]{lucatachiya2017}, and the declarations below are formal
proofs of these cases:

\begin{thm}[Purely periodic support]
For integers $b\ge2$ and $m\ge1$, and every $m$-periodic $A\subseteq\mathbb N$ (i.e.\
$n+m \in A \Leftrightarrow n \in A$ for all $n$) containing a positive element, the support series
$\sum_{a\in A,\ a\ge1}(b^a-1)^{-1}$ is irrational. \textbf{Hypotheses:} $b \ge 2$, $m \ge 1$, $A$
$m$-periodic, $\exists a > 0,\ a \in A$. \coord{periodic-support} \scale{uniform} \ev{Lean}\\
\textit{Source:} \lean{irrational\_erdosSupportSeries\_periodic}{CertificateKernel.lean:11590}.
\end{thm}

\begin{thm}[Eventually periodic support]
For integers $b\ge2$ and $m\ge1$, if $A\subseteq\mathbb N$ is infinite and its membership is $m$-periodic from
some threshold $N_0$ on (i.e.\ $n+m \in A \Leftrightarrow n \in A$ for all $n \ge N_0$), then
$\sum_{a\in A,\ a\ge1}(b^a-1)^{-1}$ is irrational, by transferring irrationality across the finite
symmetric difference from the shifted purely periodic set $A_{\mathrm{pure}} := \{n : n + m N_0 \in
A\}$. \textbf{Hypotheses:} $b \ge 2$, $m \ge 1$, $N_0 \in \mathbb{N}$, eventual $m$-periodicity from
$N_0$, $A$ infinite. \coord{eventually-periodic-support} \scale{uniform} \ev{Lean}\\
\textit{Source:} \lean{irrational\_erdosSupportSeries\_eventuallyPeriodic}{CertificateKernel.lean:11604}.
\end{thm}

\begin{thm}[Residue-class support]
For integers $b\ge2$, $m\ge1$ and $c$, the series
$\sum_{n\ge1,\ n\equiv c\pmod m}(b^n-1)^{-1}$ is irrational. \coord{residue-class-support} \scale{uniform} \ev{Lean}\\
\textit{Source:} \lean{irrational\_erdosSupportSeries\_residueClass}{CertificateKernel.lean:11672}
(specializes purely periodic support).
\end{thm}
\leannote{Lean: \lproof{ErdosProblems/Erdos257/PaperCompleteR21/ResidueClassSupport.lean}{21}{irrational_residueClass_positive_support}.}

\begin{thm}[Odd support]
For every integer $b\ge2$, the series
$\sum_{n\ge1,\ n\text{ odd}}(b^n-1)^{-1}$ is irrational. This density-$1/2$ support is
the case treated explicitly by Luca and Tachiya~\cite[Example~2, p.~140]{lucatachiya2017}.
\coord{odd-support} \scale{uniform} \ev{Lean}\\
\textit{Source:} \lean{irrational\_erdosSupportSeries\_odd}{CertificateKernel.lean:11686}
(specializes residue-class support at $m=2,c=1$).
\end{thm}

A support representing $1/2$ cannot be eventually periodic, by these
irrationality theorems.  This is a restriction on a hypothetical
representation, not a proof that the actual greedy support is nonperiodic:
if the greedy construction fails, it eventually selects every exponent.
The periodic-support results therefore decide neither alternative for
$1/2\in\mathcal A$.

\subsubsection{Topology and measure of the achievement set}

\begin{thm}[Achievement-set topology and measure]
\label{thm:topology}
$\mathcal{A}$ (Definition~\ref{defn:half-question}) is compact, closed, perfect, totally
disconnected, and nowhere dense; its Lebesgue measure is exactly $1$:
$\operatorname{volume}(\mathcal{A}) = 1$. \textbf{Hypotheses:} none. \textbf{Conclusion:}
$\mathcal A$ is a Cantor set of positive measure, often called a
fat Cantor set.  Its measure is $1$, not the length of its ambient
interval $[0,E]$, where $E\approx1.6067$.
\coord{achievement-set-topology}
\scale{n/a} \ev{Lean}\\
\textit{Source:} \lean{isCompact\_mersenneAchievementSet}{GreedyAchievementSet.lean:656},
\lean{isClosed\_mersenneAchievementSet}{GreedyAchievementSet.lean:660},
\lean{perfect\_mersenneAchievementSet}{GreedyAchievementSet.lean:1656},
\lean{isTotallyDisconnected\_mersenneAchievementSet}{GreedyAchievementSet.lean:1672},
\lean{isNowhereDense\_mersenneAchievementSet}{GreedyAchievementSet.lean:1681},
\lean{volume\_mersenneAchievementSet}{GreedyAchievementSet.lean:996}. Strict separation and summability give the compactness, unique coding, and Cantor topology, and
Hornich's strict-tail theorem, as proved by Nitecki~\cite[Theorem~4(1), p.~9]{nitecki2013}, gives
the measure as $\lim_N2^NR_N$ with $R_N=\sum_{n>N}x_n$: each level-$N$ cylinder has length $R_N$,
and the $2^N$ disjoint cylinders have total length $2^NR_N$. The Mersenne-specific input is
$2^NR_N\to1$. Other weight sequences require their own tail asymptotic.\\
\textit{Use:} Theorem~\ref{thm:greedy-survival-record} below (compactness is exactly what powers
every ``limit of achieved points is achieved'' argument used downstream, including the seam-limit
route of Part~2).
\end{thm}
\leannote{Lean: \lproof{ErdosProblems/Erdos257/PaperCompleteR20/AchievementGeometry.lean}{40}{paper_achievement_geometry}, \lproof{ErdosProblems/Erdos257/PaperCompleteR20/MersenneConstantDecimal.lean}{21}{mersenne_topology_quantitative}.}

\begin{thm}[Membership equals greedy survival at every level]
\label{thm:greedy-survival-record}
For a real $x$, let $r_n(x)$ be its greedy remainder after rank $n$,
and let $R_n=\sum_{j>n}w_j$, as in the initial notation.  Then
\[
 x\in\mathcal A\quad\Longleftrightarrow\quad
 x\ge0\ \text{and}\ r_n(x)\le R_n\text{ for every }n\ge0.
\]
The inequality says that the unselected tail has enough total mass
at every rank.  A failure at any one rank excludes membership;
the inequalities at all ranks give a representation. \coord{greedy-survival} \scale{uniform} \ev{Lean}\\
\textit{Source:} \lean{mem\_mersenneAchievementSet\_iff\_greedy\_survival}{GreedyAchievementSet.lean:1458}
(forward direction \lean{greedy\_survives\_of\_mem\_mersenneAchievementSet}{GreedyAchievementSet.lean:1375},
reverse \lean{mem\_mersenneAchievementSet\_of\_greedy\_survival}{GreedyAchievementSet.lean:1394}).
\end{thm}

\subsection{The greedy rule for one half}
\label{ssec:greedy-orbit}

Fix the target $1/2$ and the weights $x_n = 1/(2^n-1)$. Here $x_1=1$, so exponent one cannot occur
in a representation of $1/2$ by positive summands; the half-target discussion starts at exponent
two. (The reciprocal formula is undefined only at exponent zero.) The greedy
rule is deterministic.  Strict tail domination makes it recover the
support of any represented target
(Theorem~\ref{thm:greedy-survival-record}):

\begin{defn}[The half-value greedy orbit]
\label{defn:greedy-orbit}
Let $\rho_k=r_{k-1}(1/2)$ be the residual entering rank $k$, so
$\rho_2=1/2$.  In this subsection $x_k=w_k$ and
$T_{k+1}=R_k=\sum_{j>k}w_j$.  At rank $k\ge2$:
\begin{itemize}[nosep]
\item \textbf{Take} iff $\rho_k \ge x_k$; if taken, $\rho_{k+1} := \rho_k - x_k$.
\item \textbf{Skip} otherwise; $\rho_{k+1} := \rho_k$.
\end{itemize}
A skip at rank $k$ is:
\begin{itemize}[nosep]
\item \textbf{safe} iff $\rho_k \le T_{k+1} := \sum_{j>k} x_j$
(the total remaining mass does not yet rule out a representation; this
inequality alone does not guarantee one);
\item \textbf{fatal} iff $\rho_k \in (T_{k+1}, x_k)$ (the residual is too large for the remaining
tail to ever reach, yet too small to have been taken at $k$ -- an impossible gap).
\end{itemize}
\end{defn}

For example, ranks $2$ and $3$ are selected, leaving
$1/2-1/3-1/7=1/42$.  Ranks $4$ and $5$ are skipped, while rank $6$
is selected and leaves $1/126$.  These decisions illustrate the rule;
they do not establish its infinite survival.

By Theorem~\ref{thm:greedy-survival-record}, $1/2 \in \mathcal{A}$ iff no rank of this orbit is ever fatal.
Two boundary degeneracies are excluded unconditionally and independently of any open hypothesis:

\begin{rem}[No-ties: both greedy-orbit boundaries are strict]
\label{rem:no-ties}
(a) $\rho_k \ne x_k$ at any rank: a finite Mersenne sum has odd reduced denominator (each $2^n-1$ is
odd), so the residual $1/2$ minus a finite prefix cannot equal the odd-denominator
number $x_k$, sharpening the general
denominator-parity fact used in Observation~\ref{obs:half-controls-universal}.
(b) $\rho_k \ne T_{k+1}$ at any rank: the tail $T_{k+1} = E - 1 - \sum_{2 \le j \le k} x_j$, where
$E := \sum_{n\ge1} x_n$ is the Erd\H{o}s--Borwein constant, is irrational by Erd\H{o}s's own 1948
theorem (the base-$2$ case of Theorem~\ref{thm:full-support}), while $\rho_k$ is always rational;
equality is impossible. \textbf{Conclusion:} both greedy-orbit boundaries are provably strict at
every rank, so the orbit's decisions stabilize after finitely many strict inequalities at any fixed
truncation depth. \coord{no-ties} \scale{uniform} \ev{Math}
\end{rem}

\subsection{Integer quotients and their remainder identity}
\label{ssec:seam-model}

The quotient rows are finite integer approximations to the real
greedy rule, not identical finite stages of that rule.  Their
normalised weights tend to the Mersenne weights and their targets
tend to $1/2$.  Proposition~\ref{prop:one-orbit} proves agreement
at each fixed prefix once the row index is sufficiently large;
it does not assert simultaneous agreement at all ranks.
Fix a row $n\ge6$. The integer quotient weights at row $n$ are
\[
w(n,d) := \left\lfloor \frac{4^n}{2^d-1} \right\rfloor, \qquad d = 2,\dots,n-1,
\]
and the integer target is $T_n := 2^{2n-1} - 2^n$. The integer greedy process takes ranks $d = 2,\dots,n-1$
in ascending order against $T_n$ exactly as in Definition~\ref{defn:greedy-orbit}, producing a take
set $D_n \subseteq \{2,\dots,n-1\}$ and skip set
$\mathrm{Skip}_n := \{2,\dots,n-1\} \mathbin{\backslash} D_n$.
Define the \emph{integer remainder} $\mathrm{rem}(n) := T_n - \sum_{d \in D_n} w(n,d)$ and the
\emph{scaled remainder deviation} $\Delta_n := \mathrm{rem}(n) - 2^n$.

\begin{defn}[The full-support quotient recurrence]
For $M \ge 2$, define
\[
K(M) := 2^{M-1} - \sum_{d=2}^{M} \left\lfloor \frac{2^M}{2^d-1} \right\rfloor.
\]
This depends only on $M$, not on the selected support in a quotient
row.  The floor-quotient recurrence gives
\[
 K(2)=1,\qquad K(M+1)=2K(M)-(\tau(M+1)-1).
\]
Here $\tau(n)$ counts the positive divisors of $n$.  To see the correction term, double the quotients for $2\le d\le M$
and add one for each such $d$ dividing $M+1$.  The new endpoint
$d=M+1$ contributes one more quotient.  Together these are
exactly the $\tau(M+1)-1$ divisors at least two.
This is an identity for every $M\ge2$, not merely a checked
finite recurrence. \coord{universal-recursion}
\scale{uniform} \ev{Math}\ev{Cert}
\end{defn}

\begin{thm}[Decomposition of the integer remainder]
\label{thm:master-identity}
For every quotient row $n \ge 6$,
\begin{equation*}
\Delta_n \;=\; K(2n) \;+\; \sum_{d \in \mathrm{Skip}_n} \left\lfloor \frac{4^n}{2^d-1} \right\rfloor.
\tag{I}
\end{equation*}
Here the take and skip sets, capacity and remainder are those defined
at the start of this subsection.  The identity separates the contribution
of the full exponent range from that of the omitted ranks.  It makes no
claim that there are few omitted ranks; controlling their sum is a
separate part of any application.
\coord{quotient-row ;  integer-greedy (Mersenne weights, binary digits of $K$)} \scale{fixed} \ev{Math}\ev{Cert}\\
\textbf{Proof.} For $d=n$ the quotient is $2^n+1$, and for
$n<d\le2n$ it is $2^{2n-d}$.  Hence
\[
 \sum_{d=n}^{2n}\left\lfloor\frac{4^n}{2^d-1}\right\rfloor
 =2^n+1+\sum_{j=0}^{n-1}2^j=2^{n+1}.
\]
Substitute this in the definition of $K(2n)$, split
$\{2,\ldots,n-1\}=D_n\sqcup\mathrm{Skip}_n$, and use
$\Delta_n=2^{2n-1}-2^{n+1}-\sum_{d\in D_n}w(n,d)$.
This gives (I).  \textbf{Certification:} verified rows
$6\le n\le200$, zero mismatches, reported as reproduced in two separate computations (\S\ref{ssec:anchors}).\\
\textit{Use:} (I) permits substitution of information about $K(2n)$.
A resulting estimate for $\Delta_n$ must also retain or bound the
skip-set sum; information about $K$ alone need not control their
cancellation.
\end{thm}
\leannote{Lean: \lproof{ErdosProblems/Erdos257/PaperCompleteR20/QuotientRowIdentity.lean}{89}{paper_master_identity_floors}.}

\subsection{A real-valued form of the quotient identity}

\begin{thm}[Real (non-integer) form of the quotient identity]
\label{thm:real-form}
Let $x_d := 1/(2^d-1)$ and define the constant
\[
C \;:=\; \sum_{d \ge 2} x_d \;-\; \tfrac12 \;=\; E - \tfrac32 \;=\; 0.1066951524152917\ldots,
\]
where $E$ is the Erd\H{o}s--Borwein constant. Then, for every quotient row $n \ge 6$,
\begin{equation*}
\Delta_n \;=\; 4^n \Big( \sum_{d \in \mathrm{Skip}_n} x_d \;-\; C \Big) \;+\; \eta_n, \qquad
|\eta_n| < 2n+2.
\tag{II}
\end{equation*}
\textbf{Proof.} Splitting the full series at $n$ gives
\[
 \eta_n=4^nR_{n-1}-2^{n+1}
       +\sum_{d\in D_n}\left\{\frac{4^n}{2^d-1}\right\}.
\]
The tail bounds $2^{-(n-1)}<R_{n-1}\le
2^{-(n-1)}+\tfrac23 4^{-(n-1)}$ and $|D_n|\le n-2$ show
$0<\eta_n<n+2/3<2n+2$.

\textbf{Use of the error bound.} For $H>2n+2$, the condition
$|\sum_{d\in\mathrm{Skip}_n}x_d-C|>(H+2n+2)/4^n$ is sufficient
for $|\Delta_n|>H$.  Conversely, $|\Delta_n|>H$ implies
$|\sum_{d\in\mathrm{Skip}_n}x_d-C|>(H-2n-2)/4^n$.
Thus a square-root-exponential deviation corresponds to a separation
of order $2^{-3n/2}$, with the displayed additive error retained.
This is not an exact equivalence after simply discarding that error.
\coord{real approximation and integer remainder} \scale{uniform}
\ev{Math}\\
\textbf{Note:} $C = E - 3/2$ is Mersenne-specific, but the shape ``scaled deviation $=$
$4^n \cdot$(finite skip-sum $-$ target constant) $+ O(n)$'' is a template, not yet matched to any
analogous constant on the \#249 side.
\end{thm}
\leannote{Lean: \lproof{ErdosProblems/Erdos257/PaperCompleteR20/QuotientRowReal.lean}{84}{paper_real_quotient_core}, \lproof{ErdosProblems/Erdos257/PaperCompleteR20/QuotientRowReal.lean}{103}{paper_real_quotient_margins}, \lproof{ErdosProblems/Erdos257/PaperCompleteR20/QuotientRowReal.lean}{94}{row_constant_eq_tail}, \lproof{ErdosProblems/Erdos257/PaperCompleteR20/MersenneConstantDecimal.lean}{8}{mersenne_constant_decimal}.}

\subsection{Dynamics}

\begin{thm}[Exact recurrences for the three branches]
\label{thm:dynamics}
Fix $n\ge5$.  Among the quotient sums over subsets of
$\{2,\ldots,n-1\}$, let $D_n$ give the largest sum at most $T_n$
and let $B_n$ give the smallest sum strictly greater than $T_n$.
Thus $D_n$ is the greedy support already defined.  Put
\[
 r=\mathrm{rem}(n),\qquad o=Q(B_n,2n)-T_n,
\]
and define the two correction terms by
\[
 p^- =\sum_{d\in D_n}\bigl(2\mathbf1_{d\mid2n+1}
                              +\mathbf1_{d\mid2n+2}\bigr),\qquad
 p^+ =\sum_{d\in B_n}\bigl(2\mathbf1_{d\mid2n+1}
                              +\mathbf1_{d\mid2n+2}\bigr).
\]
Both lie in $[0,2(n-2)]$, since $d\ge2$ cannot divide both
consecutive integers.  The next remainder is
\[
 \mathrm{rem}(n+1)=
 \begin{cases}
 2^{n+1}-4o-p^+,&\text{if }4o+p^+\le2^{n+1}\quad(\mathrm U),\\
 4r+2^{n+1}-p^-,&\text{if not }\mathrm U\text{ and }
                   4r+2^{n+1}-p^-<2^{n+2}+4\quad(\mathrm M),\\
 4r-2^{n+1}-p^--4,&\text{otherwise}\quad(\mathrm R).
 \end{cases}
\]
Consequently, for $\lambda_n=(\mathrm{rem}(n)-2^n)/2^n$,
\begin{equation*}
 \lambda_{n+1}=
 \begin{cases}
 -(4o+p^+)/2^{n+1},&\mathrm U,\\
 2\lambda_n+2-p^-/2^{n+1},&\mathrm M,\\
 2\lambda_n-(p^-+4)/2^{n+1},&\mathrm R.
 \end{cases}
 \tag{III}
\end{equation*}
\coord{integer quotient recurrences}\scale{uniform}\ev{Math}
\end{thm}
\leannote{Lean: \lproof{ErdosProblems/Erdos257/PaperCompleteR21/ThreeBranchRowDynamics.lean}{424}{paper_dynamics}, \lproof{ErdosProblems/Erdos257/PaperCompleteR21/ThreeBranchRowDynamics.lean}{398}{paper_rowLower_existsUnique}, \lproof{ErdosProblems/Erdos257/PaperCompleteR21/ThreeBranchRowDynamics.lean}{404}{paper_rowUpper_existsUnique}, \lproof{ErdosProblems/Erdos257/PaperCompleteR21/ThreeBranchRowDynamics.lean}{279}{paper_greedySupport_isRowLower}, and 11 further declarations in the \hyperref[sec:coverage]{coverage section}.}

\begin{proof}
Applying the one-step quotient identity twice gives
\[
 q(2n+2,d)=4q(2n,d)+2\mathbf1_{d\mid2n+1}
                         +\mathbf1_{d\mid2n+2}.
\]
The new capacity is $T_{n+1}=4T_n+2^{n+1}$, and the new terminal
weight, at exponent $n$, is $q(2n+2,n)=2^{n+2}+4$.
Write $g=2^{n+1}$.  The old quotient sums are separated by at
least $g$, so the corrected sums remain separated by at least
$4g-2(n-2)>2g+4$.  This exceeds the new terminal weight.
The upper word exists: already the terms at $d=2,3,4$ have
sum greater than $T_n$ for $n\ge5$, since
$1/3+1/7+1/15=19/35>1/2$ and their total floor error is less than $3$.
Consequently, adding the terminal weight to a lower old word cannot
pass the next corrected old word.  Moreover, the increment
$T_{n+1}-4T_n=g$ is smaller than this separation, so no old word
beyond the adjacent upper one can fit.  It is therefore enough
to compare the two adjacent old words and the terminal weight.  The upper word fits precisely under
the condition $\mathrm U$; its remainder is then too small to take
the terminal weight.  Otherwise the lower word is used, and
comparison with the terminal weight gives $\mathrm M$ or
$\mathrm R$.  Substitution yields the displayed formulas.
The supplied release gives the corresponding
\href{https://github.com/wcook04/plectis-erdos-lean/blob/52f29ad173b04e3bac941b3663f2b9aebe5de0bb/Erdos257PeriodNoncollapse/HalfCylinderIntegerGreedy.lean\#L1463}{abstract adjacent-cut identity}
and
\href{https://github.com/wcook04/plectis-erdos-lean/blob/52f29ad173b04e3bac941b3663f2b9aebe5de0bb/Erdos257PeriodNoncollapse/HalfCylinderConcreteSeamAdapter.lean\#L699}{its identification with the concrete greedy row}.
Both are listed as \texttt{release\_only} in the supplied index;
these links identify the inspected source, not a new kernel replay.
\end{proof}

Only the $\mathrm M$ and $\mathrm R$ formulas have the asserted
small correction to a fixed affine map.  To rewrite $\mathrm U$
in that form, let
$\gamma_n=r+o-2^{n+1}\ge0$ be the excess of the actual adjacent
spacing over its lower bound.  Then
\[
 \lambda_{n+1}=2\lambda_n-2-
                   \frac{4\gamma_n+p^+}{2^{n+1}}\qquad(\mathrm U).
\]
There is no $O(n)$ bound on $\gamma_n$ supplied by the separation
inequality.  At $n=13$, exact calculation gives
\[
 r=15147,\quad o=3997,\quad p^+=4,\quad
 \gamma_{13}=2760,\quad \mathrm{rem}(14)=392.
\]
Thus $\lambda_{14}-2\lambda_{13}+2=-2761/4096$.
In particular, the proposed uniform error bound
$|\delta_n|\le(2n+6)2^{-n}$ in a three-branch signed-digit formula
is false already at this row.  The adjacent-spacing term cannot
be absorbed into the divisor correction.

Iterating the $\mathrm R$ formula is still legitimate along a run
of right branches.  Its accumulated correction and endpoint indices
must be retained in a reset estimate.  These exact recurrences do
not identify the proposed ceiling $L_r<(r-3)/2$ with a converse to
the conditional reset square-root estimate.

\subsection{Reported finite checks}
\label{ssec:anchors}

The following values and finite ranges are reported in the supplied
record for the quotient-greedy sequence of Section~\ref{ssec:seam-model}.
They provide test cases for the formulas above. The low-row remainders
at $n=14,15,32$ were reproduced by exact integer arithmetic in this
revision's accompanying checks. The larger computation ranges remain
reports from the supplied record, not new computations in this revision.

\begin{obs}[Exact low-row values]
\label{obs:anchors-low}
At row $n=14$: $\mathrm{rem}(14) = 392$, with take set $D_{14} = \{2,3,6,7\}$. At row $n=15$:
$\mathrm{rem}(15) = 34333$. At row $n=32$: $\mathrm{rem}(32) = 3865046005$ and $\Delta_{32} =
-429921291$. \coord{quotient row, exact integer} \scale{fixed} \ev{Cert}\\
\textit{Source:} the linked proof
\lean{seamUpperResetDyadicBandEscape\_through\_thirty}{HalfCylinderUpperResetBandCertificates.lean:78}
uses exact successor remainders at rows $14$--$31$ to verify the
band-escape condition for reset indices $13\le d\le30$.
The two ranges differ by one: the remainder used for reset $d$ is at
row $d+1$. The source uses \texttt{decide +kernel}; its Lean build was
not replayed for this revision.
\end{obs}

\begin{obs}[Certified computation depth, all four coordinates]
\label{obs:anchors-depth}
The quotient identity (Theorem~\ref{thm:master-identity}) is verified rows $6 \le n \le 200$ with zero
mismatches, reported as reproduced in two separate computations, and the quotient-greedy orbit itself is extended and
certified to row $2500$ (and separately, in a longer single run, to row $200{,}000$ -- \S SE-7 of the
underlying record, $207.7\mathrm{s}$ runtime, branch counts $\mathrm{R{:}M{:}U} =
100197{:}49899{:}49898$, TARGET failing at exactly one reset row $\le 200{,}000$, namely $r=7$).
\coord{multiple: quotient row, binary digit stream of $C$, greedy set $G$} \scale{fixed} \ev{Cert}
\end{obs}

\begin{obs}[The four coordinate depths, side by side]
\label{obs:anchors-coords}
Four related computations are reported, with different targets or
integer representations, at the following ranges: truncation rung $J = 3,\dots,22$ (all $20$ rungs proved to
survive, not merely computed -- Part~3 of the truncation-rung record); quotient row $6 \le n \le 2500$
(Theorem~\ref{thm:master-identity} and the branch data underlying
Theorem~\ref{thm:dynamics}); integer margin scan to $m_c = 5\times10^5$
(advisory Type~B scan, not yet promoted to a certified theorem); sharp tail margin, ranks $2$ through
$3000$ (no fatal skip at any rank; take fraction $0.4992$; tightest certified take-margin
$0.0340$ bits at rank $7$, the near-tie $1/126 \ge 1/127$). \coord{all four coordinates}
\scale{fixed} \ev{Cert}\\
\textit{Use:} truncated weights define a different greedy problem,
and the quotient-row construction uses integer weights rather than the
original real weights.  These are not automatically four readings of
the same decision stream.  Their relevance to full-weight membership
comes from the approximation, limit, or identification theorems cited
for each construction.  A certificate at one range does not transfer
to another without such a theorem, and these reports do not decide the
universal problem.
\end{obs}

The large computation ranges in this section are inherited from their
identified sources.  The source tags retain the distinction between
formal proofs and reported finite checks; neither should be read as a
new replay of those computations.
% ; - part a257_p1a ; -
\section{Detailed results and their hypotheses}

This section catalogues machinery for Erd\H{o}s \#257 (is $\sum_{n\in A} 1/(2^n-1)$
irrational for every infinite $A\subseteq\mathbb{N}_{\ge1}$?) drawn from the half-membership
route (is $1/2$ in the Mersenne achievement set?) and the support-rigidity /
full-support route (irrationality of $\sum_{a\in A} 1/(2^a-1)$ for named families $A$).
Every entry states exact hypotheses, exact conclusion, coordinate, scale, evidence band,
and Lean source location.  None of these results decides \#257.  Entries
are grouped into \textbf{conditional implications} (with any unproved
hypothesis identified) and \textbf{equivalences and exact identities}
(including structural properties, definitions and exact finite values).
The scale labels distinguish a single instance, a bounded range,
unbounded depths and a uniform assertion.  They are descriptions of
quantifiers, not a total ordering of the strength of the results.

\subsection{Conditional membership tests}

The finite tests below use $g_n=w_n-R_n$.  When written as functions,
$R(n)=R_n$, $g(n)=g_n$, $r(x,n)=r_n(x)$ and $C(A,N)=C_A(N)$.
For a finite set $D$, $V(D)=X_D(2)$.  For $n\ge2$, the integer remainder
$s(n)=\mathrm{rem}(n)$ uses weights $q(2n,d)$, $2\le d<n$, and
target $T_n=2^{2n-1}-2^n$; it is not the real remainder $r_n(1/2)$.
In the branch conditions below, $D_s,B_s,o_s,p_s^-,p_s^+$ denote the
lower and upper supports, overshoot and correction terms of
Theorem~\ref{thm:dynamics}, with the row added as a subscript.
The labels $\mathrm U,\mathrm M,\mathrm R$ mean precisely its three
arithmetic cases.  Finally, $F(k,J)=F_k(J)$ denotes the expression in
Theorem~\ref{thm:frozen-margin}.



\begin{thm}[Infinitely many greedy skips force membership]
\label{thm:fatal-absorbing}
Fix $x\ge0$.  If $r_n(x)>R_n$ at some rank $n$, every later
rank is selected and $r_{n+k}(x)>R_{n+k}$ for every $k\ge0$.
Consequently, infinitely many omitted positive ranks imply
$x\in\mathcal A$.
\par\ev{Lean}\scale{cofinal}\coord{greedy recurrence}
\lean{mem\_mersenneAchievementSet\_of\_greedySkippedSupport\_infinite}{GreedyAchievementSet.lean:1528}
\end{thm}
\leannote{Lean: \lproof{ErdosProblems/Erdos257/PaperCompleteR21/GreedyGapCriteria.lean}{27}{fatal_absorbing}.}
\begin{proof}
The identity $R_n=w_{n+1}+R_{n+1}$ gives
$r_n(x)>w_{n+1}$, so the rule takes rank $n+1$ and leaves
$r_{n+1}(x)=r_n(x)-w_{n+1}>R_{n+1}$.  Induction proves persistence.
Infinitely many skips therefore force $0\le r_n(x)\le R_n$ at
every rank; since $R_n\to0$, the greedy sum equals $x$.
\end{proof}

\begin{thm}[The quotient remainders converge to the greedy deficit]
\label{thm:seam-limit}
Let $G$ be the real greedy support for $1/2$, and let $D_s$ be the
integer-greedy support at row $s$.  Then
\[
 X_{D_s}(2)\longrightarrow X_G(2),\qquad
 \frac{\mathrm{rem}(s)}{4^s}\longrightarrow
 \delta:=\frac12-X_G(2)\ge0.
\]
In particular, the following conditions are equivalent: $1/2\in\mathcal A$;
the full sequence $\mathrm{rem}(s)/4^s$ tends to zero; and there
exist $s_j\to\infty$ with $\mathrm{rem}(s_j)/4^{s_j}\to0$.
\par\ev{Math}\scale{cofinal}\coord{integer quotients}
\end{thm}
\leannote{Lean: \lproof{ErdosProblems/Erdos257/PaperCompleteR21/SeamPrefixStabilityLimit.lean}{560}{paper_seam_limit_unconditional}, \lproof{ErdosProblems/Erdos257/PaperCompleteR21/SeamPrefixStabilityLimit.lean}{438}{tendsto_seamGreedyFiniteValue_greedyHalfTargetValue}, \lproof{ErdosProblems/Erdos257/PaperCompleteR21/SeamPrefixStabilityLimit.lean}{537}{tendsto_seamGreedyNormalizedRemainder}, \lproof{ErdosProblems/Erdos257/PaperCompleteR21/SeamPrefixStabilityLimit.lean}{370}{eventually_seamSupport_agrees}, and 3 further declarations in the \hyperref[sec:coverage]{coverage section}.}
\begin{proof}
For each fixed $d$, the scaled weight
$4^{-s}q(2s,d)$ tends to $w_d$, and the scaled target tends to
$1/2$.  The finite-prefix stability argument of
Proposition~\ref{prop:one-orbit} therefore gives
$D_s\cap\{2,\ldots,K\}=G\cap\{2,\ldots,K\}$ for all sufficiently
large $s$, for each fixed $K$.  At such rows,
\[
 \bigl|X_{D_s}(2)-X_G(2)\bigr|\le R_K.
\]
First letting $s\to\infty$ and then $K\to\infty$ proves convergence
of the values.  This uses summability, not an unproved uniform
rate of prefix agreement.

For $s\ge5$, put
$\Phi_s=\sum_{d\in D_s}\{4^s/(2^d-1)\}$, so
$0\le\Phi_s<s-2$.  The quotient identity gives exactly
\[
 X_{D_s}(2)=\frac12-2^{-s}
            -\frac{\mathrm{rem}(s)}{4^s}+\frac{\Phi_s}{4^s}.
\]
The two error terms tend to zero, proving the asserted limit.
The real greedy sums never exceed $1/2$, so $\delta\ge0$.
By strict tail separation, a support representing $1/2$ must agree
with the greedy rule, as in Theorem~\ref{thm:greedy-survival-catalogue}.
Thus membership is equivalent to $\delta=0$, which is also equivalent
to a zero limit along any one sequence $s_j\to\infty$.
\end{proof}
The source theorem \lean{}{HalfCylinderSeamLimit.lean:169} proves the original forward implication from
cofinal vanishing to membership, using
$0\le1/2-X_{D_s}(2)\le2^{-s}+\mathrm{rem}(s)/4^s$ and closedness:
\lean{half\_mem\_mersenneAchievementSet\_of\_subquadraticAlong}{HalfCylinderSeamLimit.lean:179}.
The full-limit identification and converse above are ordinary arguments,
not a claim about that declaration's formal statement.  Convergence
itself is automatic; proving that its limit is zero is exactly the
unresolved membership question.

\begin{thm}[Two-channel and dyadic cap sufficiency]
\label{thm:two-channel-cap}
Suppose that for every $n\ge0$ at which the next weight is skipped,
\[
 w_{n+1}>r_n(1/2)\quad\Longrightarrow\quad
 r_n(1/2)\le 2^{-(n+1)}+\frac13\,4^{-(n+1)}.
\]
Then $1/2\in\mathcal A$.  The stronger bound
$r_n(1/2)\le2^{-(n+1)}$ at the same skipped ranks also suffices.
The first bound retains two positive geometric terms of the full tail;
it is larger, and hence less restrictive, than the second.  Both are
sufficient tests on the specified greedy orbit.  The result does not
assert that either bound holds at all its skipped ranks.
\par\ev{Lean}\scale{uniform}\coord{two-channel-dyadic-cap}
\lean{half\_mem\_mersenneAchievementSet\_of\_skipped\_twoChannelCap}{GreedyAchievementSet.lean:1425}
\lean{half\_mem\_mersenneAchievementSet\_of\_skipped\_dyadicCap}{GreedyAchievementSet.lean:1444}
\end{thm}

\begin{thm}[Second-channel phase separation sufficiency]
\label{thm:second-channel}
Define $P_n=4^n\bigl(2r_n(1/2)-2^{-n}\bigr)$.  If
\[
 \frac16+\frac{37}{56}\,2^{-n}\le
 \left|P_n-\frac13\right|\qquad\text{for every }n\ge1,
\]
then $1/2\in\mathcal A$.  These are exact rational inequalities because
the target and every finite greedy remainder are rational.  The cited
finite calculation establishes them for $1\le n\le6$; it therefore
suffices to prove them for every $n\ge7$.

Unlike Theorem~\ref{thm:two-channel-cap}, this hypothesis concerns every
positive rank, not only skipped ranks.  It excludes an explicit interval
around $1/3$.  For example, $P_n\notin(0,1)$ implies the displayed
inequality when $n\ge2$, but no such avoidance theorem for the whole
orbit is proved here.
\par\ev{Lean}\scale{uniform}\coord{second-channel-phase}
\lean{half\_mem\_mersenneAchievementSet\_of\_secondChannelSeparation}{GreedyAchievementSet.lean:3118}
\lean{half\_mem\_mersenneAchievementSet\_of\_secondChannelSeparationRat\_from\_seven}{GreedyAchievementSet.lean:3151}
\end{thm}

\begin{thm}[Straddle-prefix closed-set criterion]
\label{thm:straddle-closed-set}
Suppose that for every $d\ge0$ there is a finite support
$D_d\subseteq\{1,\ldots,d\}$ such that
\[
 X_{D_d}(2)\le t\le X_{D_d}(2)+R_d.
\]
Then $t\in\mathcal A$.  The supports need not agree at different depths:
compactness, together with $R_d\to0$, supplies a limiting support.
Conversely, truncating a support representing $t$ supplies such finite
supports.  Thus this condition is an exact finite-approximation
formulation of membership, not a criterion already known to hold for an
unresolved target.
\par\ev{Lean}\scale{uniform}\coord{dedekind-cut}
\lean{mem\_mersenneAchievementSet\_of\_straddle\_all\_depths}{HalfCutLocator.lean:173}
\end{thm}
\leannote{Lean: \lproof{ErdosProblems/Erdos257/PaperCompleteR21/GreedyGapCriteria.lean}{42}{straddle_all_depths_iff_mem}, \lproof{ErdosProblems/Erdos257/PaperCompleteR21/GreedyGapCriteria.lean}{103}{straddle_limiting_support_inputs}.}

\begin{thm}[Keeping the largest skipped rank beyond two thirds]
\label{thm:largest-skip-late}
Let $D_s\subseteq\{2,\ldots,s-1\}$ be the support of the integer-greedy
row defined in Section~\ref{ssec:seam-model}.  Assume the following implication
for every $s\ge14$ and every largest omitted rank
$d=\max(\{2,\ldots,s-1\}\smallsetminus D_s)$:
\[
 2s<3d\quad\Longrightarrow\quad
 2(s+1)<3d\quad\text{or}\quad s\notin D_{s+1}.
\]
Then $1/2\in\mathcal A$.  The hypothesis says that either the same
omitted rank remains beyond two thirds of the next row, or the next
row omits its terminal rank.  The latter alternative is exactly an
upper or middle transition in the preceding classification.

The cited finite calculation establishes the initial inequality at
row $14$.  The displayed implication then preserves the inequality
by induction and supplies omitted ranks tending to infinity.
It is an unproved condition on every later row; the verified base
case alone does not establish it.
\par\ev{Lean}\scale{uniform}\coord{integer quotients}
\lean{half\_mem\_mersenneAchievementSet\_of\_largestSkipLateStepSocket}{HalfCylinderLargestSkipInduction.lean:167}
\lean{largestSkipLateAt\_fourteen}{HalfCylinderLargestSkipInduction.lean:73}
\end{thm}
\leannote{Lean: \lproof{Erdos249257/HalfCylinderLargestSkipInduction.lean}{167}{half_mem_mersenneAchievementSet_of_largestSkipLateStepSocket}, \lproof{Erdos249257/HalfCylinderLargestSkipInduction.lean}{73}{largestSkipLateAt_fourteen}, \lproof{Erdos249257/HalfCylinderHalfMembershipClassification.lean}{57}{seamGreedy_terminal_false_iff_upperOrMiddle}.}

\begin{thm}[Two sufficient bounds at middle transitions]
\label{thm:middle-producer-escape}
Suppose that at every middle transition with row $s\ge13$,
\[
 |D_s|+p_s^-+5<4\,\mathrm{rem}(s).
\]
Then $1/2\in\mathcal A$.  The stronger condition
$\mathrm{rem}(s)\ge s$ at all such rows also suffices.
\par\ev{Lean}\scale{uniform}\coord{integer quotients}
\lean{half\_mem\_mersenneAchievementSet\_of\_middleProducerCardEscape}{HalfCylinderMiddleCarryLowerBound.lean:2364}
\lean{half\_mem\_mersenneAchievementSet\_of\_middleProducerRowEscape}{HalfCylinderMiddleCarryLowerBound.lean:2371}
\end{thm}
Indeed, $|D_s|\le s-2$ and $p_s^-\le2(s-2)$ give
$|D_s|+p_s^-+5\le3s-1<4s$.  This explains the implication
between the two hypotheses.  The implication from the first bound to
membership is supplied by the cited source.  Neither all-row hypothesis
is established by the finite computations.

\begin{thm}[The tail inequality at a final middle transition]
\label{thm:middle-allright-defect}
Suppose row $D\ge13$ is middle and all rows after it are right.
Let $F_D=D_D\cup\{D\}$ and put
\[
 C_D=4\,\mathrm{rem}(D)-p_D^--4,\qquad
 \Theta_D=\sum_{j\ge1}c_{F_D}(2D+2+j)2^{-j}.
\]
Then $C_D<\Theta_D$.  Since $0\le\Theta_D\le|F_D|$,
an estimate in the opposite direction would exclude this scenario.
No such reverse estimate is assumed or proved here.
Remark~\ref{rem:tail-dominance-open} states the corresponding sufficient
hypothesis with the value $-3$ excepted.
\par\ev{Lean}\scale{uniform}\coord{integer quotients}
\lean{middleProducer\_allRight\_forces\_carry\_lt\_tail}{HalfCylinderMiddleCarryLowerBound.lean:1814}
\lean{middleProducer\_allRight\_forces\_rational\_skip}{HalfCylinderMiddleCarryLowerBound.lean:1888}
\end{thm}
\leannote{Lean: \lproof{Erdos249257/HalfCylinderMiddleCarryLowerBound.lean}{1814}{middleProducer_allRight_forces_carry_lt_tail}, \lproof{Erdos249257/HalfCylinderMiddleCarryLowerBound.lean}{1888}{middleProducer_allRight_forces_rational_skip}, \lproof{Erdos249257/HalfCylinderProducerLowerBound.lean}{150}{producerCarry_insert_seamBelowSupport_eq_middleCoordinate}, \lproof{Erdos249257/GenericTailOrbitRigidity.lean}{78}{binaryCoeffTail_nonneg}, and 1 further declaration in the \hyperref[sec:coverage]{coverage section}.}

\begin{thm}[A conditional two-sided dyadic bound]
\label{thm:two-sided-dyadic}
Assume the following two conditions for every row $s\ge5$:
\[
 \begin{aligned}
 \mathrm M\text{ at }s&\quad\Longrightarrow\quad
 4\,\mathrm{rem}(s)-p_s^--4\notin\{-3,-2,-1\},\\
 \mathrm R\text{ at }s\text{ and }o_s\le2^s
 &\quad\Longrightarrow\quad 4o_s+p_s^+\le2^{s+2}.
 \end{aligned}
\]
Then at every row $s\ge5$,
\[
 \mathrm{rem}(s)\le2^s\quad\text{or}\quad o_s\le2^s.
\]
The first hypothesis concerns only middle transitions.  The second
concerns only right transitions at which the old overshoot is at most
$2^s$.  Neither condition is asserted here for the whole sequence.
\par\ev{Lean}\scale{uniform}\coord{integer quotients}
\lean{SeamTwoSidedDyadicCellEscape.twoSided}{HalfCylinderMiddleCarryLowerBound.lean:4448}
\lean{seamTwoSidedDyadicAt\_five}{HalfCylinderMiddleCarryLowerBound.lean:4431}
\end{thm}
\begin{proof}
At row $5$, $D_5=\{2,4\}$ and $\mathrm{rem}(5)=71$, whereas
$B_5=\{2,3\}$ and $o_5=7\le32$.  For the induction step use
Theorem~\ref{thm:dynamics}.  The upper branch gives
$\mathrm{rem}(s+1)\le2^{s+1}$.  On a middle branch put
$c=4\,\mathrm{rem}(s)-p_s^--4$.  If $c\le-4$, the new remainder
is $2^{s+1}+c+4\le2^{s+1}$.  Otherwise the hypothesis forces
$c\ge0$.  Adding the terminal weight to the corrected lower word
then exceeds the target by $2^{s+1}-c>0$, so the new minimal
overshoot is at most $2^{s+1}$.

On a right branch, an old remainder at most $2^s$ gives a new
remainder at most $2^{s+1}$ by the recurrence.  In the other case
the inductive hypothesis gives $o_s\le2^s$; the corrected upper
word exceeds the new target by $4o_s+p_s^+-2^{s+1}$, which is
positive and at most $2^{s+1}$ by the second assumption.
\end{proof}
This conclusion is a disjunction.  It does not by itself give an
upper bound for the remainder on rows where only the overshoot is
small, and therefore is not by itself a membership criterion.

\begin{thm}[Upper-reset dyadic-band escape, checked for $13\le d\le30$]
\label{thm:upper-reset-band}
\emph{SeamUpperResetDyadicBandEscape} requires the following at every
actual upper reset $d\ge13$: for every $0\le j\le d$, the
reset charge avoids a linear-width band immediately below the dyadic power
$2^{d-j+1}$ ($2^{d-j+1}<\mathrm{resetCharge}$ or
$\mathrm{resetCharge}+2(d+j)\le 2^{d-j+1}$). Granted this, $1/2\in
\ensuremath{\mathcal A}$. The linked proof verifies this condition for
$13\le d\le30$, using exact successor remainders at rows $14$--$31$
(for example, $\mathrm{rem}(14)=392$ and
$\mathrm{rem}(31)=4187487147$). This finite verification does not
supply the hypothesis for every $d\ge13$.
\par\ev{Lean}\scale{uniform}\coord{dyadic-boundary}
\lean{half\_mem\_mersenneAchievementSet\_of\_upperResetDyadicBandEscape}{HalfCylinderMiddleCarryLowerBound.lean:4790}
\lean{seamUpperResetDyadicBandEscape\_through\_thirty}{HalfCylinderUpperResetBandCertificates.lean:78}
\end{thm}

\begin{thm}[M\"obius-centred carry nonnegativity below $1/2$]
\label{thm:mobius-centred-nonneg}
If $1\notin A$ and $X_A(2)<1/2$, then $C_A(N)\ge0$ for every $N$.
Indeed,
\[
 \operatorname{ihc}(A,N)=2^{N+1}(1/2-X_A(2))+
 \sum_{j>N+1}c_A(j)2^{N+1-j}>0.
\]
The carry is an integer, so it is at least $1$; subtracting $1$ gives
the asserted nonnegativity of the centred carry.  The strict inequality
and integrality are both needed in this argument.  The statement is
about the divisor counts $c_A$, not a new irrationality criterion.
\par\ev{Lean}\scale{uniform}\coord{divisor counts and finite sums}
\lean{mobiusCenteredHalfCarry\_nonneg\_of\_supportSeries\_lt\_half}{HalfCylinderFinalMiddleCellEscape.lean:94}
\lean{integerHalfCarry\_eq\_scaled\_residual\_add\_tail}{HalfCarryReachability.lean:871}
\end{thm}

\begin{thm}[A square-root bound on the greedy carry]
\label{thm:sqrt-bound-route}
For the greedy support $G=G_{1/2}$, the centred carry is nonnegative.
If in addition
\[
 C_G(N)\le2\sqrt N+4\qquad\text{for every }N\ge0,
\]
then $G$ has infinitely many skipped indices and $X_G(2)=1/2$.
The unresolved part is the upper bound along this particular orbit.
The carry identity explains its strength: a positive gap
$1/2-X_G(2)$ would contribute a term of order $2^N$, which cannot satisfy
a square-root bound.  This is distinct from the reset-deviation
hypothesis of Theorem~\ref{thm:upper-reset-band}; no equivalence between
the two hypotheses is asserted.
\par\ev{Lean}\scale{uniform}\coord{greedy recurrence}
\lean{greedy\_half\_infinite\_of\_mobiusCenteredHalfCarry\_sqrtBound}{HalfCarryReachability.lean:834}
\lean{infinite\_support\_half\_of\_mobiusCenteredHalfCarry\_sqrtBound}{HalfCarryReachability.lean:817}
\end{thm}
\leannote{Lean: \lproof{Erdos249257/HalfCarryReachability.lean}{919}{greedy_mobiusCenteredHalfCarry_nonneg}, \lproof{Erdos249257/HalfCarryReachability.lean}{834}{greedy_half_infinite_of_mobiusCenteredHalfCarry_sqrtBound}, \lproof{Erdos249257/HalfCarryReachability.lean}{817}{infinite_support_half_of_mobiusCenteredHalfCarry_sqrtBound}.}

\begin{thm}[Equivalent sign and vanishing conditions]
\label{thm:frozen-margin}
For the prefix $D=G\cap\{1,\ldots,k\}$ define
\[
 F_k(J)=\sum_{i=1}^{J}c_D(k+1+i)2^{J-i}-2^J C_D(k).
\]
This is a finite integer expression: it compares the next $J$ divisor
counts of the fixed prefix with its centred carry.  Let $s_n$ denote the
integer greedy remainder at depth $2n$, with weights $q(2n,d)$ for
$2\le d<n$ and target $2^{2n-1}-2^n$.
The following conditions each imply $1/2\in\mathcal A$:
\begin{enumerate}[label=(\roman*)]
\item For every skipped rank $n\ge3$, $F_{n-1}(n)\ge0$.
\item For every skipped rank $n\ge3$ at which the real and integer
      greedy words agree on $\{2,\ldots,n-1\}$, $s_n=0$.
\item For every skipped rank $n\ge3$, $B(2n)<s_n$.
\end{enumerate}
Conditions (i) and (ii) are equivalent.  Condition (iii) implies them;
no converse or strict separation is established here.  These hypotheses
must hold at every indicated skipped rank, not merely through a finite
computed range.
\par\ev{Lean}\scale{uniform}\coord{integer quotients}
\lean{half\_mem\_mersenneAchievementSet\_of\_skippedFullShellNonnegative}{HalfCylinderFullShellSeamBridge.lean:633}
\lean{skippedSeamAlignmentZero\_iff\_skippedFullShellNonnegative}{HalfCylinderFullShellSeamBridge.lean:671}
\lean{half\_mem\_mersenneAchievementSet\_of\_skippedSeamEscape}{HalfCylinderFullShellSeamBridge.lean:722}
\end{thm}

\begin{thm}[Erd\H{o}s--Borwein full-support irrationality, unconditional]
\label{thm:full-support-catalogue}
For every integer $b\ge2$, the sum $\sum_{n\ge1}(b^n-1)^{-1}$
is irrational.  This is Erd\H{o}s's theorem~\cite{erdos1948}; at base
$2$ the sum is the Erd\H{o}s--Borwein constant.  The result concerns
full support, not all its infinite subsets.
\par\ev{Lean}\scale{uniform}\coord{divisor counts and finite sums}
\lean{irrational\_erdosSum\_full\_support}{CertificateKernel.lean:8328}
\lean{irrational\_erdosBorwein\_series}{CertificateKernel.lean:8335}
\end{thm}

\begin{thm}[Pairwise-coprime support irrationality, Erd\H{o}s 1968]
\label{thm:pairwise-coprime}
For every integer $b\ge2$ and every infinite pairwise-coprime support
$A\subseteq\Npos$ with summable reciprocals, $\sum_{a\in A} 1/(b^a-1)$ is irrational.
\par\ev{Lean}\scale{uniform}\coord{divisor counts and finite sums}
\lean{irrational\_erdosSupportSeries\_pairwise\_coprime}{CertificateKernel.lean:10776}
\end{thm}

\begin{thm}[Irrationality from divisible coefficient blocks]
\label{thm:weighted-coeff-engine}
Let $b\ge2$ be an integer and let $c:\mathbb N\to\mathbb N$ satisfy
$c(n)\le n$ for every $n$.  Suppose that for every integer $q\ge1$
there are nonnegative integers $N,K,L,C$, with $K\le L$, such that
\[
 \begin{aligned}
 b^r&\mid c(N+r) &&(1\le r\le K),\\
 \sum_{r=K+1}^{L}c(N+r)b^{L-r}&\le C,\\
 c(N+L+1+t)&>0 &&\text{for some integer }t\ge0,\\
 q(C+N+L+2)&<b^L.
 \end{aligned}
\]
Then $\sum_{n\ge1}c(n)b^{-n}$ is irrational.
\par\ev{Lean}\scale{uniform}\coord{binary digits}
\lean{irrational\_coeff\_series\_of\_weighted\_coeff\_block\_certificates}{CertificateKernel.lean:8665}
\end{thm}
\begin{proof}
Write $x=\sum_{n\ge1}c(n)b^{-n}$.  Multiplication by $b^N$
makes the terms through $N$ integral.  The first condition also
makes the terms at $N+1,\ldots,N+K$ integral.  Subtracting their
sum leaves, for some integer $z$,
\[
 0<b^Nx-z=\sum_{r>K}\frac{c(N+r)}{b^r}
 \le\frac{C+N+L+2}{b^L}<\frac1q.
\]
Positivity follows from the specified nonzero coefficient beyond $L$.
The middle sum contributes at most $C/b^L$, while $c(N+r)\le N+r$
gives
\[
 \sum_{r>L}\frac{c(N+r)}{b^r}
 \le b^{-L}\left(\frac{N+L}{b-1}+\frac{b}{(b-1)^2}\right)
 \le\frac{N+L+2}{b^L}.
\]
If $x=p/q$ were rational, the positive difference $b^Nx-z$
could not be smaller than $1/q$.
\end{proof}

A certificate at one precision is not enough: the same coefficient
sequence must satisfy these conditions for every $q$.  The positive-tail
condition is essential; without it the zero sequence would satisfy the
other conditions at every precision.  The full-support theorem supplies
such certificates for $c=\tau$.  Their existence for $c=\varphi$ is not
proved here.  This is a limitation of this particular sufficient
criterion, not a comparison of every approach to Problem~249.


\begin{thm}[Irrationality from a gap beyond the preceding least common multiple]
\label{thm:lcm-gap-engine}
For an integer base $b\ge2$ and strictly increasing support $a:\mathbb N\to\mathbb N$ with $a(0)\ge
1$: if $a(k)-\mathrm{lcm}(a(0),\dots,a(k-1))\to\infty$, then $\sum'_k 1/(b^{a(k)}-1)$ is
irrational. The base may vary over the integers, but the denominators remain
$b^{a(k)}-1$.  Applying the argument to another denominator sequence
would require its own divisibility and tail estimates.
\par\ev{Lean}\scale{uniform}\coord{lcm-gap}
\lean{irrational\_erdosSum\_of\_lcm\_gap}{CertificateKernel.lean:5883}
\end{thm}

\begin{thm}[Factorial-support and $2^k$-support instances]
\label{thm:factorial-twopow-support}
For every integer $b\ge2$, both
\[
 \sum_{k\ge0}\frac1{b^{(k+1)!}-1}
 \quad\text{and}\quad
 \sum_{k\ge0}\frac1{b^{2^k}-1}
\]
are irrational.  These are instances of the preceding theorem.  For the
factorial support the preceding least common multiple is $k!$ when
$k\ge1$, so the gap is $k\,k!$.  For the powers of two it is
$2^{k-1}$, so the gap is $2^{k-1}$.  Both tend to infinity.
The conclusion concerns these two supports, not every infinite support.
\par\ev{Lean}\scale{uniform}\coord{lcm-gap}
\lean{irrational\_erdosSum\_factorial\_support}{CertificateKernel.lean:6035}
\lean{irrational\_erdosSum\_two\_pow\_support}{CertificateKernel.lean:6059}
\end{thm}

\begin{thm}[Multiples-support irrationality via dilation]
\label{thm:multiples-support}
For integers $b\ge2$ and $d\ge1$,
\[
 X_{d\Npos}(b)=\sum_{k\ge1}\frac1{(b^d)^k-1}
\]
is irrational.  This is the full-support result
(Theorem~\ref{thm:full-support-catalogue}) at the integer base $b^d$;
no new support argument is required.
\par\ev{Lean}\scale{uniform}\coord{dilation}
\lean{irrational\_erdosSupportSeries\_multiples}{CertificateKernel.lean:9103}
\lean{erdosSupportSeries\_multiples\_eq\_pow\_base\_full\_support}{CertificateKernel.lean:9054}
\end{thm}

Luca and Tachiya proved irrationality for every purely periodic integer weight that is not
identically zero~\cite[Theorem~A, p.~139]{lucatachiya2017}.  The periodic, eventually periodic,
residue-class and odd-support statements below follow from their theorem; the declarations are
formal proofs of these cases, and the mechanism described is that of the certificate proof.

\begin{thm}[Periodic-support irrationality]
\label{thm:periodic-support}
Let $b\ge2$ and $m\ge1$ be integers.  Suppose
$A\subseteq\Npos$ is nonempty and satisfies
$n\in A$ if and only if $n+m\in A$ for every $n\ge1$.
Then $X_A(b)$ is irrational.

The indicator of $A$ is a nonzero purely periodic integer weight,
so this follows from Luca and Tachiya's theorem stated above.
The linked declaration supplies a separate formal proof.
At $m=1$ the only nonempty periodic support is full support;
residue classes and unions of residue classes give the other
immediate examples.
\par\ev{Lean}\scale{uniform}\coord{periodic-sieve}
\href{https://github.com/wcook04/plectis-erdos/blob/7f79e63d0b36b5b4f0b47b6368342b4a50824f4e/lean/ErdosProblems/Erdos257/PaperCompleteR20/PositivePeriodicSupport.lean\#L49}{\textsc{Lean~source}}
\end{thm}

\begin{thm}[Eventually-periodic support irrationality]
\label{thm:eventually-periodic}
Let $b\ge2$ and $m\ge1$ be integers.  An infinite support whose
membership is $m$-periodic from some threshold $N_0$ onward has
irrational $X_A(b)$.
\par\ev{Lean}\scale{uniform}\coord{periodic-sieve}
\lean{irrational\_erdosSupportSeries\_eventuallyPeriodic}{CertificateKernel.lean:11604}
\end{thm}

\begin{thm}[Residue-class and odd-support irrationality]
\label{thm:residue-odd}
For integers $b\ge2$, $m\ge1$ and a residue $c$,
$\sum_{\substack{n\ge1\\n\equiv c\pmod m}}(b^n-1)^{-1}$ is
irrational. Specializing $m=2,c=1$: $\sum_{n\text{ odd}} 1/(b^n-1)$ is irrational for
every $b\ge 2$, the case treated explicitly in~\cite[Example~2, p.~140]{lucatachiya2017}.
\par\ev{Lean}\scale{uniform}\coord{periodic-sieve}
\lean{irrational\_erdosSupportSeries\_residueClass}{CertificateKernel.lean:11672}
\lean{irrational\_erdosSupportSeries\_odd}{CertificateKernel.lean:11686}
\end{thm}

\paragraph{Comparison with the known theorem.} Luca and Tachiya proved irrationality for every
nonzero purely periodic integer weight, of any signs~\cite[Theorem~A, p.~139]{lucatachiya2017}, so
the terminating branch below occurs only for the zero weight.  The theorem records the dichotomy
reached by the certificate method, together with its one-sided closures.

\begin{thm}[A formal dichotomy for signed periodic weights]
\label{thm:signed-periodic}
Let $b\ge2$ and $m\ge1$ be integers, and let
$w:\Npos\to\Z$ be $m$-periodic.  Put
\[
 x=\sum_{a\ge1}\frac{w(a)}{b^a-1},\qquad
 c_w(n)=\sum_{d\mid n}w(d).
\]
The cited formal argument gives the dichotomy that $x$ is
irrational or $b^kx\in\Z$ for some integer $k\ge0$.
It excludes the latter alternative if $c_w$ has one sign
throughout and is nonzero at arbitrarily large indices.
These sign assumptions concern $c_w$, not $w$.
As noted above, Luca and Tachiya's theorem already excludes
the terminating alternative for every nonzero periodic $w$,
without either sign restriction.  The zero weight gives $x=0$.
\par\ev{Lean}\scale{uniform}\coord{signed-divisor-calculus}
\lean{irrational\_or\_bpow\_mul\_eq\_intCast\_intWeightedErdosSeries\_periodic}{CertificateKernel.lean:14175}
\lean{irrational\_intWeightedErdosSeries\_periodic\_of\_coeff\_nonneg\_of\_frequently\_ne\_zero}{CertificateKernel.lean:14583}
\lean{irrational\_intWeightedErdosSeries\_periodic\_of\_coeff\_nonpos\_of\_frequently\_ne\_zero}{CertificateKernel.lean:14643}
\end{thm}

For the next two results write $M_r=2^r-1$ and, for a positive
squarefree integer $r$, set
\[
 A_r=\sum_{d\mid r}\mu(d)\frac rd\frac{M_r}{M_d}\in\Z,
 \qquad B(r)=\frac{A_r}{M_r}
           =\sum_{d\mid r}\frac{\mu(d)(r/d)}{2^d-1}.
\]
Every quotient in $A_r$ is integral because $d\mid r$.
The source names $A_r$ the M\"obius numerator; $B(r)$ is a finite
rational sum, not the full-support Lambert series.

\begin{thm}[Denominators of finite Mersenne sums]
\label{thm:mersenne-channel-survival}
Let $t\ge1$ and $h\ge1$ be integers, and let $r\ge1$ be
squarefree with every prime factor at most $t$.  Let $P$ be a set of
prime divisors of $r$ such that $t<2p$ for each $p\in P$, and put
$C=\prod_{p\in P}(2^p-1)$.  Then
\[
 \frac{C}{\gcd(C,h)}\ \bigm|\ \operatorname{den}\bigl(hB(r)\bigr).
\]
In particular, $C$ divides this reduced denominator when
$\gcd(C,h)=1$.  The factors in $C$ are pairwise coprime, since
$\gcd(2^p-1,2^q-1)=2^{\gcd(p,q)}-1=1$ for distinct primes.
The assertion concerns this finite signed sum; it does not replace
the approximation hypotheses in Theorem~\ref{thm:full-support-catalogue}.
\par\ev{Lean}\scale{uniform}\coord{cyclotomic}
\lean{upperHalfChannel\_survivorProduct\_dvd\_den}{MersenneShadowCyclotomicNoncollapse.lean:766}
\lean{upperHalfChannel\_product\_dvd\_den\_of\_coprime\_scale}{MersenneShadowCyclotomicNoncollapse.lean:795}
\end{thm}
\leannote{Lean: \lproof{ErdosProblems/Erdos257/PaperCompleteR21/MersenneChannelSurvivalAllHeights.lean}{210}{paper_mersenne_channel_survival}, \lproof{ErdosProblems/Erdos257/PaperCompleteR21/MersenneChannelSurvivalAllHeights.lean}{234}{paper_mersenne_channel_survival_of_coprime_scale}, \lproof{ErdosProblems/Erdos257/PaperCompleteR21/MersenneChannelSurvivalAllHeights.lean}{248}{paper_channel_factor_gcd_eq_one}, \lproof{ErdosProblems/Erdos257/PaperCompleteR21/MersenneChannelSurvivalAllHeights.lean}{260}{paper_channel_factors_pairwise_coprime}, and 5 further declarations in the \hyperref[sec:coverage]{coverage section}.}

\begin{thm}[Exponential growth of the reduced denominators]
\label{thm:mersenne-channel-growth}
For an integer $t\ge5$, put
\[
 H_t=\operatorname{lcm}(1,\ldots,t),\quad
 r_t=\prod_{p\le t}p,\quad h_t=H_t/r_t,\quad
 D_t=\operatorname{den}\bigl(h_tB(r_t)\bigr),
\]
where the product is over primes.  Then $D_t\ge2^{t/2}$.
Its exact denominator is
\[
 D_t=\frac{2^{r_t}-1}{\gcd(2^{r_t}-1,h_tA_{r_t})}.
\]
\par\ev{Lean}\ev{Math}\scale{uniform}\coord{cyclotomic}
\lean{upperHalfMersenneProduct\_lower\_bound}{MersenneShadowDenominatorGrowth.lean:60}
\lean{lcmHeight\_scaledMobiusShadow\_den\_exact}{MersenneShadowDenominatorGrowth.lean:147}
\end{thm}
\leannote{Lean: \lproof{Erdos249257/MersenneShadowDenominatorGrowth.lean}{85}{lcmHeight_scaledMobiusShadow_den_lower_bound}, \lproof{Erdos249257/MersenneShadowDenominatorGrowth.lean}{60}{upperHalfMersenneProduct_lower_bound}, \lproof{Erdos249257/MersenneShadowDenominatorGrowth.lean}{147}{lcmHeight_scaledMobiusShadow_den_exact}.}
\begin{proof}
Apply Bertrand's postulate~\cite[\S\S4--5, pp.~197--198]{erdos1932bertrand} to $m=\lceil t/2\rceil\ge3$.
There is a prime $m<p<2m$, hence $p\le t$ and $p-1\ge t/2$.
Every prime divisor $\ell$ of $2^p-1$ has multiplicative order
$p$ for $2$ modulo $\ell$.  Since $p$ is odd, $2p\mid\ell-1$;
thus $\ell>t$.  All prime factors of $h_t$ are at most $t$, so
Theorem~\ref{thm:mersenne-channel-survival} gives
\[
 2^{t/2}\le2^{p-1}<2^p-1\le D_t.
\]
The exact expression for $D_t$ is the usual reduction of the
fraction $h_tA_{r_t}/(2^{r_t}-1)$.
\end{proof}
The first source directly uses natural-number division $t/2$,
that is, $\lfloor t/2\rfloor$.  The ceiling choice in the ordinary
argument supplies the displayed real-exponent bound also when $t$
is odd.  The second source gives an equivalent exact denominator
formula in its numerator notation.  Neither denominator growth nor
that identity supplies an approximation-error estimate.

\subsection{Exact identities and reductions}

\begin{obs}[A finite half-value approximation and the first correction tail]
\label{obs:half-regression}
The finite prefix has value
\[
 \frac12-X_{\{2,3,6,7\}}(2)=\frac1{16002}.
\]
The excess of the full tail after rank $1$ over its dyadic value is
\[
 R_1-\frac12
 =\sum_{n\ge2}\left(\frac1{2^n-1}-2^{-n}\right)
 \le\frac1{12}+\frac1{28}=\frac5{42}.
\]
Indeed, expand each summand as $4^{-n}+\sum_{j\ge3}2^{-jn}$;
for $n\ge2$, the second sum is at most $2\,8^{-n}$.  Thus
$\tfrac12-(R_1-\tfrac12)\ge\tfrac8{21}>\tfrac38$.
The first identity is an exact rational check, whereas the two tail
inequalities use a convergent-series estimate.
\par\ev{Lean}\scale{fixed}\coord{exact sums and tail estimates}
\lean{half\_sub\_four\_term\_prefix\_eq}{HalfCutLocator.lean:664}
\lean{mersenneCorrectionTail\_one\_le}{HalfCutLocator.lean:670}
\end{obs}

\begin{obs}[Period-4 sign weight vanishes on a residue class]
\label{obs:period-four-obstruction}
Let $w(n)$ have successive values $1,0,-1,0$, repeated
with period $4$.  For every $n\equiv3\pmod4$,
\[
 c_w(n)=\sum_{d\mid n}w(d)=0.
\]
Indeed, $n$ is odd, and paired divisors $d,n/d$ have opposite
values of $w$ because their product is $3$ modulo $4$.
There are no fixed pairs: such an $n$ is not a square.
Thus an argument requiring a nonzero coefficient at every
chosen residue must check that requirement; periodicity alone
does not provide it.  This cancellation does not conflict with
irrationality of the full weighted series, which follows from
Luca and Tachiya's theorem~\cite[Theorem~A, p.~139]{lucatachiya2017}.
\par\ev{Lean}\scale{fixed}\coord{signed-divisor-calculus}
\lean{intWeightedCoeff\_periodFourSignWeight\_eq\_zero\_of\_mod\_four\_eq\_three}{CertificateKernel.lean:14226}
\end{obs}

\begin{thm}[The Lambert-series identity for the M\"obius function]
\label{thm:mobius-lambert-identity}
The absolutely convergent signed series satisfies
\[
 \sum_{d\ge1}\frac{\mu(d)}{2^d-1}=\frac12.
\]
This classical identity is also recorded in
\cite[Example~1.1, p.~4]{duverneytachiya}.
Expand each denominator geometrically and interchange the
absolutely convergent sums.  The coefficient of $2^{-n}$ is
$\sum_{d\mid n}\mu(d)$, equal to $1$ for $n=1$ and $0$
otherwise.  This proves the identity.
Its coefficients are $-1,0,1$, not indicators of a support;
the identity therefore does not represent $1/2$ as a subseries
with all coefficients in $\{0,1\}$.
\par\ev{Cited}\scale{fixed}\coord{divisor counts and finite sums}
\lean{tsum\_moebius\_div\_two\_pow\_sub\_one\_eq\_half}{MersenneLambertLadder.lean:587}
\end{thm}

\begin{cor}[Negative-M\"obius Boolean support overshoots $1/2$]
\label{cor:negative-mobius-overshoot}
Let $N=\{d\ge2:\mu(d)=-1\}$.  Isolating the $d=1$
term in Theorem~\ref{thm:mobius-lambert-identity} gives
\[
 X_N(2)=\frac12+
       \sum_{\substack{d\ge2\\\mu(d)=1}}\frac1{2^d-1}
       \ge\frac12+\frac1{63}>\frac12,
\]
since $\mu(6)=1$.  This rules out the particular candidate $N$.
It does not rule out other infinite supports representing $1/2$.
\par\ev{Lean}\scale{fixed}\coord{divisor counts and finite sums}
\lean{half\_lt\_tsum\_negativeMobius}{MobiusSignSupportNoGo.lean:164}
\lean{tsum\_negativeMobius\_eq\_half\_add\_positiveMobiusTail}{MobiusSignSupportNoGo.lean:111}
\end{cor}

\begin{obs}[Dyadic safety need not survive a take]
\label{obs:skip-block-non-invariant}
For a residual $p/(2D)$, taking rank $b$ subtracts $w_b$ and gives
\[
 \frac{p}{2D}-\frac1{2^b-1}
 =\frac{p(2^b-1)-2D}{2D(2^b-1)}.
\]
The residual $7/34$ is at most $2^{-2}$.  Taking rank $3$ leaves
$15/238$, which is strictly between $2^{-4}$ and $w_4=1/15$.
Thus rank $4$ is skipped but fails the dyadic safety test.  The displayed
numerator has increased from $7$ to $15$; that increase does not preserve
safety.

The same failure occurs with an initial denominator containing all
Mersenne denominators through rank $7$.  With
$D=\operatorname{lcm}(3,7,15,31,63,127)=1240155$, one has
\[
 \frac{12149}{2D}-\frac1{255}
 =\frac{41179}{2\cdot21082635}.
\]
The initial residual is at most $2^{-7}$.  After this rank-$8$ take,
ranks $9$ and $10$ are skipped, but the new residual exceeds $2^{-10}$.
The exact numerator of that dyadic excess is
$2^9\cdot41179-21082635=1013>0$.
These are counterexamples to the stated local safety implication,
not assertions that either state occurs on the real half-greedy orbit.
\par\ev{Cert}\scale{fixed}\coord{integer quotients}
\lean{safeBracket\_and\_numeratorMonotone\_not\_inductive\_fixture}{DyadicPrefixCompression.lean:622}
\lean{lcmSaturated\_safeBracket\_not\_inductive\_fixture}{DyadicPrefixCompression.lean:647}
\end{obs}

\begin{thm}[The half-skip dichotomy via Erd\H{o}s--Borwein irrationality]
\label{thm:half-skip-dichotomy}
The half target satisfies
\[
 \frac12\in\mathcal A\quad\Longleftrightarrow\quad
 \mathbb N_{>0}\smallsetminus G\text{ is infinite},
\]
where $G$ is the greedy support for $1/2$.  In the forward direction,
finitely many skipped exponents would express the full Mersenne sum as
$1/2$ plus a finite rational sum, contradicting Erd\H{o}s's full-support
irrationality theorem.  The reverse implication uses the greedy
tail criterion: once a remainder exceeds the whole available tail,
every later exponent is selected.  Infinitely many skips exclude that
failure.  The equivalence does not establish that infinitely many skips
actually occur.
\par\ev{Lean}\scale{cofinal}\coord{greedy recurrence}
\lean{half\_mem\_mersenneAchievementSet\_iff\_greedySkippedSupport\_infinite}{GreedyAchievementSet.lean:2583}
\lean{irrational\_erdosBorweinMersenneConstant}{GreedyAchievementSet.lean:2469}
\end{thm}

\begin{thm}[Equivalent descriptions of half-membership]
\label{thm:nine-way-hub}
Each of the following is equivalent to $1/2\in\mathcal A$:
\begin{enumerate}[label=(\arabic*),nosep]
\item The integer-greedy sequence is not eventually always on branch $\mathrm R$.
\item There are arbitrarily large $s\ge5$ with $s\notin D_{s+1}$.
\item Upper or middle transitions occur at arbitrarily large rows.
\item Some sequence $s_j\to\infty$ satisfies $s_j\notin D_{s_j+1}$ for every $j$.
\item There are rows $s_j\to\infty$ and omitted ranks
$d_j\in\{2,\ldots,s_j-1\}\smallsetminus D_{s_j}$ with $d_j\to\infty$.
\item The positive ranks omitted by the real half-greedy rule form an infinite set.
\item The real half-greedy rule has no last omitted positive rank.
\end{enumerate}
\par\ev{Lean}\scale{cofinal}\coord{integer quotients}
\lean{half\_mem\_mersenneAchievementSet\_iff\_not\_seamGreedyEventuallyRight}{HalfCylinderHalfMembershipClassification.lean:112}
\lean{half\_mem\_mersenneAchievementSet\_iff\_exists\_unboundedSkippedRanksAlong}{HalfCylinderHalfMembershipClassification.lean:213}
\lean{half\_mem\_mersenneAchievementSet\_iff\_no\_lastHalfGreedySkip}{HalfCylinderHalfMembershipClassification.lean:235}
\end{thm}
\leannote{Lean: \lproof{Erdos249257/HalfCylinderHalfMembershipClassification.lean}{112}{half_mem_mersenneAchievementSet_iff_not_seamGreedyEventuallyRight}, \lproof{Erdos249257/HalfCylinderHalfMembershipClassification.lean}{126}{half_mem_mersenneAchievementSet_iff_unboundedTerminalFalse}, \lproof{Erdos249257/HalfCylinderHalfMembershipClassification.lean}{156}{half_mem_mersenneAchievementSet_iff_unboundedUpperOrMiddle}, \lproof{Erdos249257/HalfCylinderHalfMembershipClassification.lean}{203}{half_mem_mersenneAchievementSet_iff_cofinalTerminalFalse}, and 3 further declarations in the \hyperref[sec:coverage]{coverage section}.}
Items (1)--(4) differ only by the branch classification and the choice
of a sequence of rows.  Item (5) requires the omitted ranks themselves
to tend to infinity, not merely their containing rows.  The cited
identification with the real greedy support supplies (6)--(7).
These equivalences do not prove that any of the conditions holds.

\begin{lem}[The values on either side of an eventual right continuation]
\label{lem:eventually-right-impossible}
If the seam eventually always extends ``true'' (right branch) from some row $S$ on with a
fixed lower prefix $u$, the resulting cofinite-support value stays strictly below $1/2$
(\emph{prefix\_add\_mersenneTail\_lt\_half\_of\_eventually\_right}); the matching
alternative ``upper competitor'' word gives a strict excess \emph{above} $1/2$
(\emph{half\_lt\_upper\_competitor\_of\_eventually\_right}). Neither of these two cofinite continuations represents $1/2$.
This does not rule out an eventually-right integer orbit; it describes
the two values in that case.  The inequalities are used in
Theorem~\ref{thm:final-middle-cell}.
\par\ev{Lean}\scale{cofinal}\coord{integer quotients}
\lean{prefix\_add\_mersenneTail\_lt\_half\_of\_eventually\_right}{HalfCylinderFatalGapRightTail.lean:402}
\lean{half\_lt\_upper\_competitor\_of\_eventually\_right}{HalfCylinderFatalGapRightTail.lean:627}
\end{lem}

\begin{lem}[Each Mersenne weight exceeds its remaining tail]
\label{lem:mersenne-tail-weight}
For $n\ge1$, let $w_n=(2^n-1)^{-1}$ and
$R_n=\sum_{j>n}w_j$.  Then
\[
 R_n=w_{n+1}+R_{n+1},\qquad
 2^{-n}<R_n\le2w_{n+1}<w_n.
\]
For the upper bound, compare each tail term with
$2^{1-j}w_{n+1}$ at index $n+j$, $j\ge1$, and sum the
geometric series.  The last inequality follows by comparing
$2/(2^{n+1}-1)$ with $1/(2^n-1)$.
This strict term-versus-tail inequality is the separation hypothesis
used in the greedy arguments.

It is not available for the totient weights $\varphi(n)/2^n$:
the terms at indices $4$ and $5$ already sum to
$2/16+4/32=1/4=\varphi(3)/2^3$, and the remaining tail is
positive.  A general strict-tail argument therefore cannot be
transferred to that sequence without a different hypothesis or
proof.
\par\ev{Lean}\scale{uniform}\coord{divisor counts and finite sums}
\lean{mersenneTail\_lt\_weight}{GreedyAchievementSet.lean:180}
\lean{mersenneTail\_le\_two\_mul\_weight}{GreedyAchievementSet.lean:155}
\end{lem}
\leannote{Lean: \lproof{Erdos249257/GreedyAchievementSet.lean}{114}{mersenneTail_eq_weight_add}, \lproof{Erdos249257/GreedyAchievementSet.lean}{1226}{halfTwoChannelCap_lt_mersenneTail}, \lproof{Erdos249257/GreedyAchievementSet.lean}{155}{mersenneTail_le_two_mul_weight}, \lproof{Erdos249257/GreedyAchievementSet.lean}{125}{two_mul_mersenneWeight_succ_lt}, and 1 further declaration in the \hyperref[sec:coverage]{coverage section}.}

\begin{thm}[Greedy survival and membership]
\label{thm:greedy-survival-catalogue}
With the greedy remainder $r_n(x)$ and complete tail $R_n$,
\[
 x\in\mathcal A\quad\Longleftrightarrow\quad
 x\ge0\ \text{and}\ r_n(x)\le R_n\text{ for every }n\ge0.
\]
If all inequalities hold, the nonnegative remainders tend to zero
because $R_n\to0$, so the greedy partial sums converge to $x$.
Conversely, strict tail domination $w_n>R_n$ forces the greedy
choices in any representation and therefore gives every inequality.
The argument applies to positive summable weights satisfying that
separation hypothesis; it is not a consequence of summability alone.
\par\ev{Lean}\scale{uniform}\coord{greedy recurrence}
\lean{mem\_mersenneAchievementSet\_iff\_greedy\_survival}{GreedyAchievementSet.lean:1458}
\end{thm}

\begin{lem}[The two next-prefix intervals and their gap]
\label{lem:rank-step-trichotomy}
Let $d\ge0$ be an integer, let $u\subseteq\{1,\ldots,d\}$, and suppose the target $t$ lies
in the interval
$[X_u(2),X_u(2)+R_d]$.
At depth $d+1$, either $t$ lies in the lower interval
\[
 [X_u(2),X_u(2)+R_{d+1}],
\]
or in the upper interval
\[
 [X_u(2)+w_{d+1},X_u(2)+w_{d+1}+R_{d+1}],
\]
or in the open gap between them.  These alternatives are disjoint
because $R_{d+1}<w_{d+1}$, and exhaustive because
$R_d=w_{d+1}+R_{d+1}$.
For $t=1/2$, Lemma~\ref{lem:half-endpoint-kills} also excludes
the interval endpoints, so all relevant comparisons are strict.
Endpoint exclusion is not what makes the two child intervals
disjoint; the strict-tail inequality does that.
\par\ev{Lean}\scale{uniform}\coord{dedekind-cut}
\lean{isStraddlePrefix\_step\_trichotomy}{HalfCutLocator.lean:205}
\lean{IsStraddlePrefix.half\_step\_forced}{HalfCutLocator.lean:301}
\end{lem}

\begin{lem}[A greedy gap excludes every representation]
\label{lem:fatal-gap-exclusion}
Let $d\ge0$ be an integer and $u\subseteq\{1,\ldots,d\}$ a finite prefix.  If
\[
 X_u(2)+R_{d+1}<t<X_u(2)+w_{d+1},
\]
then no support agreeing with $u$ through rank $d$ represents $t$:
omitting $d+1$ leaves value at most the lower endpoint, while
including it gives value at least the upper endpoint.
Other length-$d$ prefixes have disjoint containing intervals by
the first-difference argument using $w_n>R_n$.  Since the displayed
gap lies inside the interval for $u$, none of those prefixes can
represent $t$ either.  Thus a certified strict gap excludes every
representation, not just one proposed continuation.
\par\ev{Lean}\scale{uniform}\coord{dedekind-cut}
\lean{positiveMersenneSupportValue\_ne\_of\_rank\_gap}{HalfCutLocator.lean:380}
\end{lem}
\leannote{Lean: \lproof{ErdosProblems/Erdos257/PaperCompleteR21/GreedyGapCriteria.lean}{174}{fatal_gap_excludes_every_representation}, \lproof{ErdosProblems/Erdos257/PaperCompleteR21/GreedyGapCriteria.lean}{114}{fatal_gap_endpoint_bounds}, \lproof{ErdosProblems/Erdos257/PaperCompleteR21/GreedyGapCriteria.lean}{144}{depth_prefix_interval_disjoint}, \lproof{ErdosProblems/Erdos257/PaperCompleteR21/GreedyGapCriteria.lean}{162}{fatal_gap_within_prefix_interval}.}

\begin{lem}[No finite Mersenne sum equals one half]
\label{lem:half-endpoint-kills}
For every finite $u\subseteq\mathbb N_{>0}$,
$X_u(2)\ne1/2$: its reduced denominator is odd.
Also $X_u(2)+R_d\ne1/2$ for every $d\ge0$, since $R_d$
is the irrational full Mersenne sum minus a finite rational sum.
These observations exclude equality at the finite-prefix and
complete-tail endpoints of Lemma~\ref{lem:rank-step-trichotomy}.
\par\ev{Lean}\scale{uniform}\coord{parity-irrationality}
\lean{positiveMersenneSupportValue\_coe\_finset\_ne\_half}{HalfCutLocator.lean:243}
\lean{half\_ne\_coe\_finset\_add\_mersenneTail}{HalfCutLocator.lean:263}
\end{lem}

\begin{lem}[Straddle words are canonical: they agree with the greedy prefix]
\label{lem:straddle-agrees-greedy}
Let $u\subseteq\{1,\ldots,d\}$ satisfy
$X_u(2)\le1/2\le X_u(2)+R_d$.  Then
\[
 u=G\cap\{1,\ldots,d\},
\]
where $G$ is the real greedy support for $1/2$.
At the first disagreement, if $u$ takes rank $k$ and the greedy rule
skips it, then $X_u(2)>1/2$.  In the opposite case, the greedy
remainder before $k$ is at least $w_k$, whereas the remaining
contribution allowed by $u$, including $R_d$, is at most
$R_k<w_k$.  Both cases contradict the displayed interval.  Thus there is at most one such word
at each depth; the lemma does not assert its existence at every depth.
\par\ev{Lean}\scale{uniform}\coord{greedy recurrence}
\lean{IsStraddlePrefix.half\_agrees\_greedy}{HalfCutLocator.lean:442}
\lean{IsStraddlePrefix.erase\_top}{HalfCutLocator.lean:334}
\end{lem}

\begin{thm}[Last-skip iff local fatality: a pointwise criterion]
\label{thm:last-skip-iff-fatal}
Keep $G$ for the selected half-greedy support, $r_M(1/2)$ for its
remainder and $R_M$ for the complete tail after rank $M$.
A positive rank $M$ is the last omitted rank if and only if
$M\notin G$ and $r_M(1/2)>R_M$.  Consequently,
\[
 \frac12\in\mathcal A
 \quad\Longleftrightarrow\quad
 r_M(1/2)\le R_M\quad\text{for every positive }M\notin G.
\]
Thus it suffices to check the tail inequality at the ranks actually
omitted by the greedy rule.  This is an equivalence on that fixed orbit;
it does not establish the inequality at its unboundedly many possible
omitted ranks.
\par\ev{Lean}\scale{uniform}\coord{greedy recurrence}
\lean{half\_mem\_iff\_every\_actual\_skip\_survives}{HalfCylinderFixedTailSocket.lean:73}
\lean{isLastHalfGreedySkip\_iff\_skip\_and\_fatal}{HalfCylinderFixedTailSocket.lean:22}
\end{thm}

\begin{lem}[Seam upper-or-middle classification]
\label{lem:seam-upper-or-middle}
At a row $s\ge5$, the terminal rank is omitted in the next row,
$s\notin D_{s+1}$, if and only if either
\[
 4o_s+p_s^+\le2^{s+1},
\]
or this inequality fails and
\[
 4\,\mathrm{rem}(s)+2^{s+1}-p_s^-<2^{s+2}+4.
\]
These are exactly branches $\mathrm U$ and $\mathrm M$ of
Theorem~\ref{thm:dynamics}.  On either branch the new terminal weight
does not fit; on $\mathrm R$ it does.  This is an exact
classification, not a conjecture inferred from observed branch labels.
\par\ev{Lean}\scale{uniform}\coord{integer quotients}
\lean{seamGreedy\_terminal\_false\_iff\_upperOrMiddle}{HalfCylinderHalfMembershipClassification.lean:57}
\end{lem}

\begin{lem}[The exact gap at the largest omitted rank]
\label{lem:largest-false-rank-algebra}
Let $2\le d<s$ and $2s<3d$.  Write
$W_s(E)=\sum_{e\in E}\lfloor4^s/(2^e-1)\rfloor$ for a finite
set $E\subseteq\{2,\ldots,s-1\}$.  For
$u\subseteq\{2,\ldots,d-1\}$, put
$E_-=u\cup\{d+1,\ldots,s-1\}$ and $E_+=u\cup\{d\}$.
Then
\[
 3W_s(E_-)+3\cdot2^{s+1}+2\cdot4^{s-d}+4=3W_s(E_+).
\]
The correction is independent of the common prefix $u$.
For the integer-greedy rows $s\ge5$, a right transition preserves the
largest omitted rank, whereas an upper or middle transition makes
$s$ the largest omitted rank in the next row.
\par\ev{Lean}\scale{uniform}\coord{integer quotients}
\lean{three\_mul\_largestSkipLowerWeight\_add\_exactLateGap\_eq\_upperWeight}{HalfCylinderLargestSkipGap.lean:259}
\lean{IsLargestFalseRank.seamGreedyWord\_succ\_of\_rightBranch}{HalfCylinderLargestSkipGap.lean:339}
\end{lem}
\leannote{Lean: \lproof{ErdosProblems/Erdos257/PaperCompleteR21/SeamRowGapAndCarry.lean}{39}{largest_false_rank_algebra}, \lproof{ErdosProblems/Erdos257/PaperCompleteR21/SeamRowGapAndCarry.lean}{98}{paper_largest_false_rank_algebra}.}
\begin{proof}
For $d\le e<s$, the inequality $2s<3d$ gives
\[
 \left\lfloor\frac{4^s}{2^e-1}\right\rfloor
 =2^{2s-e}+4^{s-e}.
\]
Subtract the sum over $d<e<s$ from the term at $d$.  The powers of
$2$ leave $2^{s+1}$; the powers of $4$ leave
$(2\cdot4^{s-d}+4)/3$.  The common prefix $u$ cancels.
The branch statements follow from the support updates in
Theorem~\ref{thm:dynamics}: the right branch appends $s$ to the lower
word, and the other two branches omit $s$.
\end{proof}
For example, $s=5$, $d=4$ and $u=\{2\}$ give
$W_s(E_-)=341$ and $W_s(E_+)=409$; the displayed identity is
$3\cdot341+204=3\cdot409$.  Both the factor $3$ on the right and
the restriction $2s<3d$ are essential to this formula.


\begin{thm}[Reduction to the nearest dyadic boundary]
\label{thm:critical-dyadic-band}
Let $d,E$ be nonnegative integers with $E\le2^{d+1}$, and let $j_*$
be the largest $j\in\{0,\ldots,d\}$ for which $E\le2^{d-j+1}$.
The set is nonempty, so $j_*$ is well defined, including when $E=0$.
Then
\[
 \begin{aligned}
 &\forall j\in\{0,\ldots,d\},\quad
  2^{d-j+1}<E\ \text{or}\ E+2(d+j)\le2^{d-j+1}\\
 &\hspace{35mm}\Longleftrightarrow\quad
 E+2(d+j_*)\le2^{d-j_*+1}.
 \end{aligned}
\]
For $j>j_*$ the first alternative holds by maximality.  For $j\le j_*$,
the left side of the required inequality increases with $j$ and its
right side decreases, so the condition at $j_*$ implies all the others.
This is an elementary reduction from $d+1$ inequalities to one. Specialized to the seam reset charge
(where the actual upper-reset condition gives $E\le2^{d+1}$,
as explained in Proposition~\ref{record:257bm-c12}), the
reduced hypothesis \emph{SeamUpperResetCriticalBandEscape} is proved logically equivalent to
Theorem~\ref{thm:upper-reset-band}'s band-avoidance hypothesis.
\par\ev{Lean}\scale{uniform}\coord{dyadic-boundary}
\lean{dyadicBandEscape\_iff\_exists\_critical}{HalfUpperResetCriticalBand.lean:108}
\lean{seamUpperResetCriticalBandEscape\_iff}{HalfUpperResetCriticalBand.lean:883}
\end{thm}
\leannote{Lean: \lproof{Erdos249257/HalfUpperResetCriticalBand.lean}{46}{exists_criticalDyadicBandIndex}, \lproof{Erdos249257/HalfUpperResetCriticalBand.lean}{108}{dyadicBandEscape_iff_exists_critical}, \lproof{Erdos249257/HalfUpperResetCriticalBand.lean}{883}{seamUpperResetCriticalBandEscape_iff}.}

\par
\noindent\begin{minipage}{\linewidth}
\begin{thm}[Excluding minus three at a final middle transition]
\label{thm:final-middle-cell}
Suppose that row $D\ge13$ is a middle transition in
Theorem~\ref{thm:dynamics} and every transition at a row $s\ge D+1$
is right.  Then
\[
 C_D:=4\,\mathrm{rem}(D)-p_D^--4\ne-3.
\]
This excludes one particular value under the stated tail assumption;
it neither excludes every final middle transition nor proves the
all-middle-row hypothesis of Theorem~\ref{thm:two-sided-dyadic}.
\par\ev{Lean}\scale{uniform}\coord{integer quotients}
\lean{finalMiddleCell\_neg\_three\_not\_last}{HalfCylinderFinalMiddleCellEscape.lean:587}
\lean{mobiusCenteredHalfCarry\_add\_two}{HalfCylinderFinalMiddleCellEscape.lean:39}
\end{thm}
\end{minipage}\par
To see the contradiction, complete $D_D$ by all ranks greater than $D$.
The resulting support $B=D_D\cup\{D+1,D+2,\ldots\}$ has
$X_B(2)<1/2$ under the all-right assumption.  The exact carry identity
in the cited proof gives $C_B(2D+1)=C_D+3$.  If $C_D=-3$, this is zero.
The recurrence
$C_B(N+1)=2C_B(N)+1-c_B(N+2)$ and nonnegativity from
Theorem~\ref{thm:mobius-centred-nonneg} first force $C_B(2D+2)=0$:
indeed $2D+3\in B$, so $c_B(2D+3)\ge1$.
Next $D+2,2D+4\in B$ imply $c_B(2D+4)\ge2$, giving
$C_B(2D+3)<0$, a contradiction.  The stronger conclusion
$C_D\ge-2$ is proved in Theorem~\ref{thm:cd-neg3-impossible}.


\begin{lem}[The signed position of a skipped greedy prefix]
\label{lem:skipped-endpoint-trichotomy}
Let $s\ge5$ be omitted by the real greedy support $G$, and put
\[
 H_s=G\cap\{2,\ldots,s-1\},\qquad
 f_s=\sum_{d\in H_s}\left\lfloor\frac{4^s}{2^d-1}\right\rfloor-T_s.
\]
Exactly one of the following holds:
\[
 \begin{array}{lll}
 f_s<0:& H_s=D_s,& f_s=-\mathrm{rem}(s),\quad \mathrm{rem}(s)\ge1;\\
 f_s=0:& H_s=D_s,& \mathrm{rem}(s)=0;\\
 f_s>0:& H_s=B_s,& f_s=o_s.
 \end{array}
\]
Thus the real greedy prefix is one of the two adjacent integer words.
This classifies its signed distance $f_s$ from $T_s$; it does not
exclude the negative case.  Nor does it identify $f_s$ with the
different coordinate $C_s=4\mathrm{rem}(s)-p_s^- -4$ used in the
final-middle argument.
\par\ev{Lean}\scale{uniform}\coord{integer quotients}
\lean{halfGreedy\_skipped\_endpoint\_trichotomy}{HalfCylinderSkippedEndpointClassifier.lean:246}
\end{lem}

\begin{lem}[Spacing of reverse-carry words]
\label{lem:reverse-carry-word}
For $i=1,2$, let integer sequences $a_i,b_i,u_i$ satisfy
\[
 b_i(m)+2u_i(m)=a_i(m)+u_i(m+1).
\]
Let $k,L\ge0$ be integers.  Suppose that at index $k$ the coefficients agree and
$b_1(k)-b_2(k)=1$.  Suppose also that both coefficients and bits
agree at indices $k+1,\ldots,k+L$.  Then
\[
 u_1(k+L+1)-u_2(k+L+1)
   =2^L\bigl(2(u_1(k)-u_2(k))+1\bigr).
\]
The factor in parentheses is an odd integer.  If the two terminal
carries have absolute values at most $B_1,B_2$, respectively, then
$2^L\le B_1+B_2$.  For a common bound $B$ the conclusion is
$2^L\le2B$, not $2^L\le B$.
The formula follows by subtracting the two recurrences: the
initial bit difference gives the odd factor, and each subsequent
agreement doubles it.  An application must supply the stated
agreements and terminal bounds.
\par\ev{Lean}\scale{uniform}\coord{binary digits}
\lean{overlappingReverseCarryWords\_carryDifference\_eq\_twoPow\_mul\_odd}{HalfTrappingReturnCarry.lean:124}
\lean{overlappingReverseCarryWords\_twoPow\_le\_realBound}{HalfTrappingReturnCarry.lean:191}
\end{lem}
\leannote{Lean: \lproof{ErdosProblems/Erdos257/PaperCompleteR21/SeamRowGapAndCarry.lean}{240}{paper_reverse_carry_word}, \lproof{ErdosProblems/Erdos257/PaperCompleteR21/SeamRowGapAndCarry.lean}{209}{reverse_carry_word_common_bound_sharp}.}

\begin{lem}[Linear functionals factoring through one value]
\label{lem:linear-channel-nogo}
Let $V$ be a vector space over $\mathbb Q$, let
$\mathrm{ev}:V\to\mathbb Q$ be linear, and choose $e\in V$ with
$\mathrm{ev}(e)=1$.  Suppose that each linear functional
$\ell_j:V\to\mathbb Q$ vanishes on $\ker(\mathrm{ev})$.
Then $v-\mathrm{ev}(v)e\in\ker(\mathrm{ev})$ gives
$\ell_j(v)=\ell_j(e)\mathrm{ev}(v)$.
Consequently any finite evaluation matrix $(\ell_j(v_i))_{i,j}$ is an
outer product and has rank at most one.  Every square minor of order
at least two therefore vanishes.  This excludes determinant arguments
formed from these particular functionals, not determinant methods
with additional independent information.
\par\ev{Lean}\scale{uniform}\coord{linear-algebra}
\lean{relationInvariantLinearChannels\_det\_eq\_zero}{HalfTrappingReturnCarry.lean:42}
\end{lem}
\leannote{Lean: \lproof{Erdos249257/AdelicHeightObstruction.lean}{120}{linearDescender_eq_smul_eval}, \lproof{Erdos249257/HalfTrappingReturnCarry.lean}{42}{relationInvariantLinearChannels_det_eq_zero}.}

\begin{thm}[Two-thirds band: exact localisation of post-take skip-unsafety]
\label{thm:two-thirds-band}
Write a positive residual as $1/R$.  A skipped rank $k$ passes
the sufficient dyadic test precisely when $R\ge2^k$.
Suppose a weight at rank $b$ is taken without exhausting the
residual: with $q=2^b-1$, assume $0<R<q$.
The new reciprocal residual is $Rq/(q-R)$.
If the next take is at rank $c\ge b+2$, put $m=2^{c-1}$.
The last skipped rank is dyadically unsafe exactly when
\[
 m-1<\frac{Rq}{q-R}<m,
 \quad\text{equivalently}\quad
 \frac{q(m-1)}{q+m-1}<R<\frac{qm}{q+m}.
\]
Clearing the positive denominators proves the equivalence.
The interval has width $q^2/((q+m)(q+m-1))$.

For a single skipped rank, $c=b+2$ and $m=2q+2$.
The width is then $q^2/((3q+1)(3q+2))<1/9$, and the
interval lies inside $2q<3R<2q+2/3$.
In this single-skip case, the cited corollaries exclude integral
$R$ and show that a reduced pre-take residual $p/(2D)$ with
$p,D,q$ odd can be unsafe only if $p\ge7$.
The latter conclusion uses the divisibility by $4$ of
$6D-2pq$ under the unsafe-band inequalities.

These are local transition statements, not bounds on the actual
half-greedy orbit.  Dyadic safety is only sufficient for the
current tail test.  The odd-coprime data $(p,D,b)=(17,41,3)$
give an unsafe single-skip example; no conclusion that the
half-greedy orbit avoids such data follows from the band calculation.
\par\ev{Lean}\scale{uniform}\coord{rational-band}
\lean{postTakeUnsafeAt\_iff\_band}{HalfGreedyTwoThirdsBand.lean:88}
\lean{seven\_le\_of\_intBand\_odd}{HalfGreedyTwoThirdsBand.lean:231}
\lean{not\_twoThirdsBand\_of\_int}{HalfGreedyTwoThirdsBand.lean:185}
\end{thm}
\leannote{Lean: \lproof{ErdosProblems/Erdos257/PaperCompleteR21/PostTakeBandLocalisation.lean}{96}{paper_two_thirds_band}.}

\begin{thm}[A weaker sufficient test against the remaining tail]
\label{thm:sharp-fatal-gap}
Let $k,u,L$ be positive integers, and put
$a=2L-(2^k-1)u$.  Suppose $a>0$, equivalently that
$\rho=u/(2L)<w_k$ and the greedy rule skips weight $w_k$.
The dyadic sufficient test $\rho\le2^{-k}$ is equivalent to $u\le a$.
The weaker sufficient condition
\[
 2u\le3a
\]
ensures $\rho<R_k$, by comparison with
$2^{-k}+(3\cdot4^k)^{-1}+(7\cdot8^k)^{-1}<R_k$.
This conclusion excludes a fatal tail-mass deficit at the current step;
it does not assert that $\rho$ is representable by the remaining
weights, or that all future greedy steps survive.

The containment of these sufficient conditions is strict even for
valid rational data: $(k,u,L,a)=(2,7,13,5)$ gives
$1/4<7/26<R_2$.  In contrast, the scalar pair $(u,a)=(3,2)$ in the
linked inequality lemma does not arise from integral $L$ under the
present parity relation.  Unit numerators are nonfatal at such a
skipped step, since then $a\ge1$.  Conversely, a fatal step requires
$3a<2u$, hence $u\ge2$, or $u\ge3$ when $u$ is odd.
\par\ev{Lean}\scale{uniform}\coord{divisor counts and finite sums}
\lean{skipSafe\_of\_two\_mul\_le\_three\_mul}{HalfGreedyFatalGap.lean:107}
\lean{sharp\_strictly\_stronger}{HalfGreedyFatalGap.lean:200}
\lean{three\_le\_of\_fatal\_of\_odd}{HalfGreedyFatalGap.lean:173}
\end{thm}
\leannote{Lean: \lproof{ErdosProblems/Erdos257/PaperCompleteR21/GreedyGapCriteria.lean}{299}{paper_sharp_fatal_gap}.}

\begin{lem}[Summability of one gap length per level]
\label{lem:gap-mass-summability}
For $n\ge1$, put $g_n=w_n-R_n>0$.  The sum of one gap
length per level satisfies, for $N\ge0$,
\[
 \sum_{n>N}g_n\le\frac29\,4^{-N}+\frac37\,8^{-N}.
\]
This follows by summing the per-level upper bounds geometrically.
It is not the measure of the union of all gaps: level $n$ has
$2^{n-1}$ translated gaps of length $g_n$.  The corresponding weighted
sum and its geometric interpretation are given after
Theorem~\ref{record:257hg-i2}.  Neither estimate decides membership
of a specified point such as $1/2$.
\par\ev{Lean}\scale{uniform}\coord{divisor counts and finite sums}
\lean{mersenneGap\_tail\_le}{HalfGapMass.lean:83}
\lean{tendsto\_mersenneGap\_tail\_zero}{HalfGapMass.lean:104}
\end{lem}
\leannote{Lean: \lproof{Erdos249257/GreedyAchievementSet.lean}{2346}{mersenneGap_pos}, \lproof{Erdos249257/HalfGapMass.lean}{77}{summable_mersenneGap_succ}, \lproof{Erdos249257/HalfGapMass.lean}{83}{mersenneGap_tail_le}, \lproof{Erdos249257/HalfGapMass.lean}{104}{tendsto_mersenneGap_tail_zero}.}

\begin{lem}[The effect of adding one divisor]
\label{lem:half-divisor-unit-drop}
Let $N\ge0$ and let finite supports $D_0,D_1$ differ only by
$D_1=D_0\cup\{N+1\}$, with $N+1\notin D_0$.
For every positive integer $m$,
\[
 c_{D_1}(m)-c_{D_0}(m)=\mathbf1_{N+1\mid m}.
\]
In particular, the difference is $1$ at $m=2(N+1)$.
This follows directly from the definition of a divisor count.
Its use in another coefficient sequence requires proving that
sequence has the same support-incidence description.
\par\ev{Lean}\scale{uniform}\coord{divisor-incidence}
\lean{supportCoeff\_extend\_true\_eq\_false\_add\_one\_at\_double}{HalfDivisorUnitDrop.lean:20}
\end{lem}
\leannote{Lean: \lproof{Erdos249257/HalfCylinderIntegerGreedy.lean}{882}{supportCoeff_insert_eq_add_indicator}, \lproof{Erdos249257/HalfDivisorUnitDrop.lean}{20}{supportCoeff_extend_true_eq_false_add_one_at_double}.}

\begin{thm}[Uniqueness of an integer recurrence with a vanishing scaled limit]
\label{thm:tempered-orbit-rigidity}
For any nonnegative-integer coefficient sequence $c:\mathbb N\to\mathbb N$ with $c(n)\le
n$: the binary coefficient series $X_c = \sum'_{n\ge 1} c(n)/2^n$ is rational iff there
exists a positive integer multiplier $v$ and an integer orbit $u:\mathbb N\to\mathbb Z$
satisfying the exact carry recurrence $u(N{+}1)=2u(N)-v\cdot c(N{+}1)$ together with the
condition $u(N)/2^N\to0$.  Every integer sequence satisfying both conditions obeys $u(N) = v\cdot T_c(N)$ exactly, where $T_c(N)=\sum_{j\ge 1} c(N{+}j)/2^j$ is
the scaled tail.  For each fixed $v$ there is at most one such sequence. Positivity of the orbit alone
is deliberately \emph{not} used as an equivalent criterion: a homogeneous $2^N$-scaled
perturbation can be added to any orbit without breaking the recurrence, so the limit condition cannot be dropped.  The choices $c=c_A$ for
Problem~257 and $c=\varphi$ for Problem~249 both satisfy $0\le c(n)\le n$.
The telescoping argument is given again, with an example, in
Theorem~\ref{record:257bm-i-t7}.
\par\ev{Lean}\scale{uniform}\coord{binary digits}
\lean{binaryCoeffSeries\_rational\_iff\_exists\_temperedBinaryOrbit}{GenericTailOrbitRigidity.lean:426}
\lean{temperedBinaryOrbit\_eq\_scaledTail}{GenericTailOrbitRigidity.lean:339}
\end{thm}

\begin{lem}[Changing finitely many support elements]
\label{lem:tail-transfer}
Let $b\ge2$ be an integer, and let $A,B\subseteq\mathbb N_{\ge1}$.
If $X_A(b)$ is irrational and $A,B$ have finite symmetric difference, then $X_B(b)$ is irrational too: the two sums differ by a finite sum of
rational numbers.  This observation supplies the finite modifications in
Theorem~\ref{thm:eventually-periodic}.
\par\ev{Lean}\scale{uniform}\coord{tail-transfer}
\lean{irrational\_erdosSupportSeries\_tail\_of\_irrational}{CertificateKernel.lean:9476}
\end{lem}
\leannote{Lean: \lproof{Erdos249257/CertificateKernel.lean}{9467}{irrational_erdosSupportSeries_of_tail}, \lproof{Erdos249257/CertificateKernel.lean}{9476}{irrational_erdosSupportSeries_tail_of_irrational}.}

\begin{lem}[An integer test for the interval between the dyadic and Mersenne weights]
\label{lem:dyadic-excess-reformulation}
For positive integers $p,L$ and $n\ge0$, write the residual as $p/(2L)$
and put $E=2^np-L$.  Clearing positive denominators gives
\[
 \frac{p}{2L}\le\frac1{2^{n+1}}\quad\Longleftrightarrow\quad E\le0,
 \qquad
 \frac1{2^{n+1}}<\frac{p}{2L}<\frac1{2^{n+1}-1}
 \quad\Longleftrightarrow\quad 0<E<\frac p2.
\]
The first test concerns the dyadic bound, not the greedy selection
threshold itself.  The second identifies precisely the interval in
which a Mersenne weight is skipped but the dyadic safety test fails.
\par\ev{Lean}\scale{uniform}\coord{integer quotients}
\lean{divInt\_le\_nextDyadic\_iff\_excess\_nonpos}{DyadicPrefixCompression.lean:198}
\lean{greedyHalf\_mem\_nextMersenneDyadicSliver\_iff\_excess}{DyadicPrefixCompression.lean:1044}
\end{lem}

\begin{lem}[The odd denominator survives dyadic subtraction]
\label{lem:denominator-sandwich}
Let $r/D$ be a reduced fraction with $D>0$ odd, let $p\in\mathbb Z$,
and let $c\ge0$ be an integer.  The reduced denominator of
$p/2^c-r/D$ is divisible by $D$ and divides $2^cD$.
Indeed, its unreduced numerator $pD-2^cr$ is coprime to $D$,
since $\gcd(r,D)=\gcd(2^c,D)=1$.  Reduction can therefore remove
only powers of $2$ from the displayed denominator.  No property
of Mersenne weights is used.
\par\ev{Lean}\scale{uniform}\coord{denominator-algebra}
\lean{dyadicResidual\_denominator\_sandwich}{DyadicPrefixCompression.lean:118}
\end{lem}

\begin{lem}[A guaranteed divisor of a reduced denominator]
\label{lem:denominator-survival}
Let $a\in\mathbb Z$, $D\in\mathbb N_{>0}$ and $m,C,h\in\mathbb N$.
Write $\operatorname{den}(x)$ for the positive reduced denominator of
a rational number $x$.  If $m\mid D$ and $\gcd(m,|a|)=1$, then
$m\mid\operatorname{den}(a/D)$.  If $C\mid D$ and $\gcd(C,|a|)=1$, then
\[
 \frac{C}{\gcd(C,h)}\mid\operatorname{den}(ha/D).
\]
Thus scaling may remove part of the guaranteed divisor.
Indeed, reduction divides $D$ by $\gcd(D,|a|)$, so no prime power in $m$
can be lost. After scaling, a prime $p\mid C$ can lose at most $v_p(h)$
powers from the guaranteed divisor. This is a lower bound on the reduced
denominator, not an equality: for $D=12$, $C=6$, $a=1$, $h=2$, it guarantees
the divisor $3$, while the reduced denominator is $6$.
The lemma feeds Theorem~\ref{thm:mersenne-channel-survival} directly.
\par\ev{Lean}\scale{uniform}\coord{denominator-algebra}
\lean{divisor\_dvd\_divInt\_den}{RationalDenominatorSurvival.lean:17}
\lean{survivingDivisor\_dvd\_scaled\_divInt\_den}{RationalDenominatorSurvival.lean:38}
\end{lem}

\begin{lem}[Commuting prime-power differences]
\label{lem:mixed-prime-power-layer}
For a function $g:\mathbb N_{>0}\to\mathbb Z$, a positive integer $p$
and an integer $e\ge1$, define
\[
 (\Delta_{p,e}g)(n)=g(p^en)-g(p^{e-1}n).
\]
These operators commute.  If $p\ne q$ are primes, $e,f\ge1$,
$A\subseteq\mathbb N_{>0}$, $n\ge1$ and $\gcd(n,pq)=1$, then
\[
 (\Delta_{q,f}\Delta_{p,e}c_A)(n)
   =\sum_{d\mid n}\mathbf1_A(p^eq^fd).
\]
Thus the mixed difference counts precisely the support elements
whose $p$- and $q$-adic exponents are $e$ and $f$ in this divisor sum.
\par\ev{Lean}\scale{uniform}\coord{p-adic}
\lean{mixedPrimePowerLayerTwo\_supportCoeffInt}{MaximalOmegaLayer.lean:39}
\lean{primePowerLayer\_comm}{MaximalOmegaLayer.lean:29}
\end{lem}
\begin{proof}
Expanding either order gives the same four terms:
\[
 \begin{aligned}
 &g(p^eq^fn)-g(p^{e-1}q^fn)\\
 &\qquad{}-g(p^eq^{f-1}n)+g(p^{e-1}q^{f-1}n).
 \end{aligned}
\]
This commutation needs no primality or coprimality.  For $g=c_A$,
write each divisor of $p^eq^fn$ uniquely as $p^iq^jd$, with
$d\mid n$, $0\le i\le e$ and $0\le j\le f$.  The two differences
cancel all terms except $i=e$, $j=f$.
\end{proof}
For $A=\{12\}$, $p=2$, $e=2$, $q=3$, $f=1$, $n=1$, the
four counts are $1,0,0,0$, so the mixed difference is $1$.
This coefficient identity alone does not transfer rationality through
the subsequence operation $g(n)\mapsto g(pn)$.


\begin{defn}[The Mersenne achievement set and the support-series object]
\label{defn:achievement-set}
For a support $A\subseteq\mathbb N_{\ge1}$, its Mersenne subsum is
\[
 x_A=\sum_{n\ge1}\frac{\mathbf1_A(n)}{2^n-1},
 \qquad
 \mathcal A=\{x_A:A\subseteq\mathbb N_{\ge1}\}.
\]
Finite, empty, and infinite supports are all allowed in $\mathcal A$.
More generally, for an integer base $b\ge2$,
$X_A(b)=\sum_{n\ge1}\mathbf1_A(n)/(b^n-1)$ and
$c_A(n)=\sum_{d\mid n}\mathbf1_A(d)=(\mathbf1_A*1)(n)$ for $n\ge1$,
where $*$ denotes Dirichlet convolution.  Full support gives the series
in Theorem~\ref{thm:full-support-catalogue}.  The linked definitions
implement these positive-index sums.
\par\ev{Lean}\scale{n/a}\coord{binary digits}
\lean{mersenneAchievementSet}{GreedyAchievementSet.lean:571}
\lean{erdosSupportSeries}{CertificateKernel.lean:8897}
\lean{supportCoeff}{CertificateKernel.lean:8856}
\end{defn}

\begin{prop}[Achievement-set topology: compact, closed, perfect, measure exactly one]
\label{prop:achievement-set-topology}
$\ensuremath{\mathcal A}$ is compact (continuous image of the binary-sequence
Cantor space $\mathbb N\to\mathrm{Fin}\,2$ under the product topology, via
$\mathrm{positiveMersenneDigitValue}$), hence closed; it is also perfect, totally
disconnected, and nowhere dense, with Lebesgue measure exactly $1$. The
compactness/closedness argument (binary coding $\to$ Cantor space $\to$ continuous image)
is a fully generic technique for characterizing the achievement set of \emph{any}
absolutely convergent digit-weighted series, not specific to Mersenne denominators ; 
reusable for a $\varphi(n)/2^n$ subsum set after checking summability.
No separation of successive weights is needed for compactness or
closedness; the stronger topological conclusions require their own
hypotheses. Closedness alone is what powers every ``limit of a sequence of
achieved points is achieved'' argument in this catalogue (e.g.\
Theorem~\ref{thm:seam-limit}, Theorem~\ref{thm:straddle-closed-set}).
\par\ev{Lean}\scale{n/a}\coord{binary digits}
\lean{isCompact\_mersenneAchievementSet}{GreedyAchievementSet.lean:656}
\lean{isClosed\_mersenneAchievementSet}{GreedyAchievementSet.lean:660}
\lean{volume\_mersenneAchievementSet}{GreedyAchievementSet.lean:996}
\end{prop}
\leannote{Lean: \lproof{ErdosProblems/Erdos257/PaperCompleteR21/AchievementSetTopologyAndFiniteHalf.lean}{27}{paper_achievement_set_topology}.}

\begin{thm}[Nonmembership of one half and a finite fatal gap]
\label{thm:master-dichotomy}
The value $1/2$ is not in $\mathcal A$ if and only if there are
an integer $d\ge0$ and a finite set $u\subseteq\{1,\ldots,d\}$
such that
\[
 X_u(2)+R_{d+1}<\frac12<X_u(2)+w_{d+1}.
\]
One implication is Lemma~\ref{lem:fatal-gap-exclusion}.  For the
other, follow the greedy prefixes: nonmembership forces a first
failure of the tail inequality, hence one of these gaps.
The endpoint equalities are excluded by
Lemma~\ref{lem:half-endpoint-kills}.

A finite prefix and strict gap provide a nonmembership witness.
The tail inequality can be certified by rational truncation bounds
as in Theorem~\ref{thm:one-sided}.  Membership, in contrast, asserts
that no such witness exists; failure to find one in a finite search
is not a proof.  The reduction uses $w_n>R_n$, not compactness
alone.  The totient weights fail this hypothesis, as the example
after Lemma~\ref{lem:mersenne-tail-weight} shows.
\par\ev{Lean}\scale{n/a}\coord{dedekind-cut}
\lean{half\_mem\_mersenneAchievementSet\_iff\_no\_existsFatalHalfGap}{HalfCutLocator.lean:654}
\lean{half\_mem\_mersenneAchievementSet\_or\_exists\_fatal\_gap}{HalfCutLocator.lean:623}
\end{thm}

\begin{defn}[Separated integer values and bounded perturbations]
\label{defn:perturbed-family}
Let $\alpha$ be a set, let $s:\alpha\to\N$ be injective, and let
$p:\alpha\to\N$.  Suppose integers $g\ge1$ and $B\ge0$ satisfy
\[
 s(x)<s(y)\ \Longrightarrow\ s(x)+g\le s(y),
 \qquad 0\le p(x)\le B<3g.
\]
The updated values are $t(x)=4s(x)+p(x)$.
For a capacity $C\ge0$, an adjacent pair consists of $x_-,x_+\in\alpha$
with $s(x_-)$ the largest value at most $C$ and $s(x_+)$ the
smallest value greater than $C$.  Existence of both values is part
of the assumption; it is not guaranteed for an arbitrary family.
Set
\[
 r=C-s(x_-),\qquad o=s(x_+)-C,\qquad
 p_-=p(x_-),\quad p_+=p(x_+).
\]
The inequality $B<3g$ preserves the order of the updated values.
The next result uses the stronger bound $B<g$ to ensure that the
selected updated value is admissible at capacity $4C+g$.
The linked structures encode these data.
\par\ev{Lean}\scale{n/a}\coord{abstract-perturbed-greedy}
\lean{PerturbedFamily}{HalfCylinderIntegerGreedy.lean:1286}
\lean{AdjacentCut}{HalfCylinderIntegerGreedy.lean:1320}
\end{defn}

\begin{thm}[Perturbed-family maximality and the three-branch recurrence]
\label{thm:perturbed-family-maximality}
Use the family and adjacent pair of
Definition~\ref{defn:perturbed-family}, and assume in addition that
$B<g$.  At capacity $C'=4C+g$, the largest admissible updated value
is $t(x_+)$ if $4o+p_+\le g$, and $t(x_-)$ otherwise.
After this choice, apply the take-if-possible rule to an additional
weight $W=2g+4$.  The resulting remainder is
\[
 \begin{cases}
 g-4o-p_+,&4o+p_+\le g,\\
 4r+g-p_-,&4o+p_+>g\ \text{and}\ 4r+g-p_-<W,\\
 4r-g-p_--4,&4o+p_+>g\ \text{and}\ 4r+g-p_-\ge W.
 \end{cases}
\]
\par\ev{Lean}\scale{n/a}\coord{abstract-perturbed-greedy}
\lean{prefixChoice\_maximal}{HalfCylinderIntegerGreedy.lean:1390}
\lean{nextRemainder\_trichotomy}{HalfCylinderIntegerGreedy.lean:1463}
\end{thm}

\begin{proof}
The lower candidate satisfies $t(x_-)\le4C+B<C'$, and
$t(x_+)\le C'$ is exactly the test $4o+p_+\le g$.
Order preservation excludes earlier values from improving the
selected one.  Any value after $x_+$ has
$s(x)\ge s(x_+)+g>C+g$, and hence $t(x)>C'$.
Thus no other candidate is missed.

On the first branch the remainder is at most $g<W$, so the extra
weight is not taken.  The other two branches compare the lower
candidate's remainder $4r+g-p_-$ with $W$.
All displayed differences are nonnegative on their respective branches.
\end{proof}

The theorem describes a two-stage rule, not automatically the maximum
of $\{t(x)+\varepsilon W:x\in\alpha,\ \varepsilon\in\{0,1\}\}$
below $C'$.  For example, take old values $4,5,7$, $g=1$, $B=0$,
$p=0$ and $C=6$.  Then $C'=25$ and $W=6$.  The rule chooses
$20$ and cannot add $6$, but $16+6=22$ is a larger admissible
extended value.  A sufficient extra hypothesis for global maximality
is $4g-B\ge W$: every earlier updated value plus $W$ is then at
most the selected updated value.  In the Mersenne application
$g=2^{s+1}$, $B=2(s-2)$ and $s\ge5$, so this inequality holds.
Thus the concrete recurrence is unaffected, but an application to
another weight system must check this separation as well.

The hypothesis $B<g$ is additional; it does not follow from
$B<3g$ in Definition~\ref{defn:perturbed-family}.
The weaker bound alone does not ensure admissibility of the chosen value.
Take $g=2$, $B=3$, old values $0,2,4$, perturbations $0,3,0$, and $C=2$.
The rule selects the old value $2$, whose update is $11>C'=10$,
although the updated value $0$ is admissible.
This example explains why the stronger hypothesis must accompany
the maximality statement.
% ; - part a257_p1b ; -
\subsection{Consequence theorems}

This subsection concerns finite sets $D$ whose floor-quotient sums nearly
represent $1/2$.  Its main implication is elementary to state: if such
finite sums approximate $1/2$ at unbounded depths with an error tending to
zero, compactness supplies a representing support.  The difficult input
is the existence of the finite sets, not the final compactness argument.
Definitions and conditional implications are separated below.  A
hypothesis marked \ev{Open} is not established merely because the
implication using it has a formal proof.

\subsubsection*{Finite quotients and their remainders}

\begin{defn}[An exact finite quotient sum]\label{record:257bm-d1}
For an integer $n\ge1$, an exact row at depth $n$ is a set
$D\subseteq\{2,\ldots,n\}$ satisfying
\[
 Q(D,n)=\sum_{d\in D}\left\lfloor\frac{2^n}{2^d-1}\right\rfloor
       =2^{n-1}-1.
\]
The word ``exact'' refers to this integer equality, not to
$X_D(2)=1/2$, which is impossible for finite $D$.
The formal definition also allows $n=0$ using truncated
natural-number subtraction.  That degenerate case is immaterial
to existence at unbounded depths.
\par\noindent\ev{Lean} \scale{n/a} \coord{divisor counts and finite sums} \lean{ExactLocalMersenneHalfRow}{BooleanMobiusCofinalExactRows.lean:31}
\end{defn}

\begin{defn}[Floor quotients and the signed endpoint expression]\label{record:257bm-d2}
\[
\ensuremath{q}(M,d) := \left\lfloor \frac{2^M}{2^d-1} \right\rfloor,\qquad
\ensuremath{Q}(D,M) := \sum_{d\in D} \ensuremath{q}(M,d),
\]
The signed endpoint expression has three arguments:
\[
 \ensuremath{H}(D,k,n)
 =2\,\ensuremath{S}(D,k,n-1)+1-c_D(n),
 \qquad c_D(n)=\#\{d\in D:d\mid n\}.
\]
The suffix in this formula is defined next.  The quotients are integers;
this final expression is allowed to be negative.
\par\noindent\ev{Lean} \scale{n/a} \coord{divisor counts and finite sums} \lean{localPrefixQuotient}{BooleanMobiusLocalRepair.lean:78} \lean{localMersenneQuotient}{BooleanMobiusLocalRepair.lean:20} \lean{localRepairInteger}{BooleanMobiusLocalRepair.lean:123}
\end{defn}

\begin{defn}[The remaining binary suffix]\label{record:257bm-d3}
\[
\ensuremath{S}(D,k,M)
 :=\max\{2^{M-k}-\ensuremath{Q}(D,M)-1,0\}
 \qquad(M\ge k),
\]
the nonnegative remainder after subtracting the quotient contributions
from the nonterminating binary expansion of $2^{-k}$ through place $M$.
The formal definition uses truncated natural-number subtraction; in the
applications where the displayed difference is nonnegative, the maximum
may be omitted.  The formal convention for $M<k$ also truncates $M-k$
to zero, a case not used in these capacity estimates.
\par\noindent\ev{Lean} \scale{n/a} \coord{binary digits} \lean{localBinarySuffix}{BooleanMobiusLocalRepair.lean:102}
\end{defn}

\begin{defn}[Exact quotient sums at unbounded depths]\label{record:257bm-d4}
The cofinal exact-row condition says that for every integer $N$
there are $n\ge\max\{N,1\}$ and
$D\subseteq\{2,\ldots,n\}$ with
\[
 Q(D,n)=2^{n-1}-1.
\]
No agreement is required between supports chosen at different depths.
This removes a compatibility requirement from the data to be supplied,
not from the underlying membership question:
Proposition~\ref{prop:collapsed-list} identifies this cofinal
existence statement with half-membership.  A failed construction
of compatible rows therefore does not refute the condition,
and the condition itself is not proved here.
\par\noindent\ev{Open} \scale{cofinal} \coord{divisor counts and finite sums} \lean{CofinalExactLocalMersenneHalfRows}{BooleanMobiusCofinalExactRows.lean:38}
\end{defn}

\subsubsection*{Two implications yielding half-membership}

\begin{thm}[Compactness from exact finite sums]\label{record:257bm-c1}
Suppose the cofinal exact-row condition in
Definition~\ref{record:257bm-d4} holds.  Then $1/2\in\mathcal A$.
For each $N$, choose an exact row $D_N$ at depth
$n_N\ge\max\{N,1\}$.  Proposition~\ref{record:257bm-i9} gives
\[
 \left|X_{D_N}(2)-\frac12\right|
       \le\frac{n_N+1}{2^{n_N}}\longrightarrow0.
\]
Every $X_{D_N}(2)$ belongs to the closed set $\mathcal A$, so
its limit does too.  No finite support represents $1/2$;
the representing support is therefore infinite.
The proof does not require the chosen rows to be nested.
The unproved input is their existence at unbounded depths.
\par\noindent\ev{Lean} \scale{cofinal} \coord{topological-achievement-set} \lean{half\_mem\_mersenneAchievementSet\_of\_cofinalExactLocalRows}{BooleanMobiusCofinalExactRows.lean:71}
\end{thm}
\leannote{Lean: \lproof{Erdos249257/BooleanMobiusCofinalExactRows.lean}{58}{abs_exactLocalMersenneRowValue_sub_half_le}, \lproof{Erdos249257/BooleanMobiusCofinalExactRows.lean}{71}{half_mem_mersenneAchievementSet_of_cofinalExactLocalRows}, \lproof{Erdos249257/HalfCarryReachability.lean}{589}{finite_boolSupport_ne_half}.}

\begin{thm}[Infinitely many greedy skips suffice]\label{record:257bm-c2}
Suppose that the greedy remainder for $1/2$ satisfies
\[
 \forall N\ \exists c\ge\max\{N,4\}:\qquad
 0<r_{c-1}(1/2)<w_c.
\]
Then exact quotient rows exist at unbounded depths, and $1/2\in\mathcal A$.
A single such skip at $c$ gives an exact row at depth $2c-2$ without
an additional capacity hypothesis.  Taking arbitrarily large $c$
therefore supplies the rows needed for the compactness implication
in Theorem~\ref{record:257bm-c1}.
The existence of arbitrarily late positive skips is the hypothesis,
not an established property of the greedy sequence.  It is equivalent
to half-membership and remains unproved here.
\par\noindent\ev{Lean} \scale{cofinal} \coord{greedy recurrence} \lean{cofinalExactLocalMersenneHalfRows\_of\_positiveHalfGreedySkips}{BooleanMobiusSkipRowCofinal.lean:84} \lean{half\_mem\_mersenneAchievementSet\_of\_positiveHalfGreedySkips}{BooleanMobiusSkipRowCofinal.lean:97} \lean{exactLocalMersenneHalfRow\_of\_positiveHalfGreedySkip}{BooleanMobiusSkipRowCofinal.lean:55}
\end{thm}
\leannote{Lean: \lproof{Erdos249257/BooleanMobiusSkipRowCofinal.lean}{55}{exactLocalMersenneHalfRow_of_positiveHalfGreedySkip}, \lproof{Erdos249257/BooleanMobiusSkipRowCofinal.lean}{84}{cofinalExactLocalMersenneHalfRows_of_positiveHalfGreedySkips}, \lproof{Erdos249257/BooleanMobiusSkipRowCofinal.lean}{97}{half_mem_mersenneAchievementSet_of_positiveHalfGreedySkips}, \lproof{Erdos249257/BooleanMobiusSkipRowCofinal.lean}{110}{cofinalPositiveHalfGreedySkips_iff_half_mem}.}

\subsubsection*{Compatible finite approximations}

\begin{thm}[A compatible family of finite supports]\label{record:257bm-c3}
Let $b_{n,d}\in\{0,1\}$ satisfy $b_{n+1,d}=b_{n,d}$ whenever
$2d\le n$.  Define
\[
 E_n=\{d:2\le d\le n,\ b_{n,d}=1\},\qquad
 D_n=E_n\cap\{2,\ldots,\lfloor n/2\rfloor\}.
\]
The four conditions on the finite rows, required for every $n\ge2$, are
\[
 \begin{aligned}
 2^{\max\{c_{D_n}(n)-1,0\}}-1&\le S(D_n,1,n-1),\\
 \sum_{\substack{d\in E_n\\d>\lfloor n/2\rfloor}}2^{n-d}
     &=H(D_n,1,n),\\
 H(D_n,1,n)&<2^{n-\lfloor n/2\rfloor},\\
 Q(E_n,n)&=2^{n-1}-1.
 \end{aligned}
\]
The first condition implies that the integer in the next two lines is
nonnegative.  The second specifies the binary value of the upper half of
the row, the third bounds that value by its available number of bits, and
the fourth is the exact quotient equation.  The maximum in the first line
records natural-number subtraction at the case $c_{D_n}(n)=0$.

These are the finite hypotheses of the compatible-limit construction.
The bit condition fixes coordinate $d$ from row $2d$ onward.  Unlike the
independent finite supports in Definition~\ref{record:257bm-d4}, the rows
therefore have a prescribed common limit.  This extra structure is not,
by itself, a proof of strictness between the corresponding existence
statements.
\par\noindent\ev{Open} \scale{uniform} \coord{binary digits} \lean{GlobalBooleanMobiusRepairFeasible}{BooleanMobiusGlobalRepair.lean:174} \lean{GlobalEndpointExponentialBound}{BooleanMobiusGlobalRepair.lean:139} \lean{bit\_stable}{BooleanMobiusGlobalRepair.lean:82}
\end{thm}
\leannote{Lean: \lproof{ErdosProblems/Erdos257/PaperCompleteR21/CompatibleFiniteRowFamily.lean}{95}{paper_compatible_bit_stable}, \lproof{ErdosProblems/Erdos257/PaperCompleteR21/CompatibleFiniteRowFamily.lean}{111}{paper_compatible_finite_row_conditions}, \lproof{ErdosProblems/Erdos257/PaperCompleteR21/CompatibleFiniteRowFamily.lean}{146}{paper_compatible_first_condition_gives_nonneg}, \lproof{ErdosProblems/Erdos257/PaperCompleteR21/CompatibleFiniteRowFamily.lean}{158}{paper_compatible_rows_agree_with_limit}.}

\subsubsection*{A quotient bound at a crossing}

\begin{thm}[A sufficient quotient bound at a crossing]\label{record:257bm-c4}
Consider the following condition on finite sets $D$ and integers $c\ge4$:
\[
 \begin{gathered}
 D\subseteq\{2,\ldots,c-1\},\qquad
 X_D(2)<\frac12<X_D(2)+w_c\\[2pt]
 \Longrightarrow\qquad Q(D\cup\{c\},2c-2)\ge2^{2c-3}.
 \end{gathered}
\]
It requires the quotient inequality at every strict crossing of one half
by an added weight.  Theorem~\ref{record:257bm-i5} identifies this
inequality with the binary bound needed in the finite construction.
If the condition holds for all such $D,c$, the induction below produces
exact sums at unbounded depths and hence an infinite support of value
$1/2$.  The universal crossing condition itself is unproved.
\par\noindent\ev{Open} \scale{uniform} \coord{binary digits} \lean{SkippedCoreCriticalQuotientSupply}{BooleanMobiusCriticalCapacityCofinal.lean:30}
\end{thm}
\leannote{Lean: \lproof{Erdos249257/BooleanMobiusCriticalCapacityCofinal.lean}{1309}{cofinalExactLocalMersenneHalfRows_of_criticalQuotientSupply}, \lproof{Erdos249257/BooleanMobiusCriticalCapacityCofinal.lean}{1321}{half_mem_mersenneAchievementSet_of_criticalQuotientSupply}, \lproof{Erdos249257/HalfCarryReachability.lean}{589}{finite_boolSupport_ne_half}.}

\begin{thm}[Induction from the depth-six example]\label{record:257bm-c5}
An already-formalised induction from the endpoint-six seed (Definition~\ref{record:257bm-i-seed}) consumes Theorem~\ref{record:257bm-c4} at
every step: each \texttt{ProtectedExactLocalMersenneRow} either doubles below half
(unconditional) or recycles at its first crossing rank $e > $ cutoff, giving endpoint
$2e-2 > $ previous endpoint; \texttt{protection} (endpoint $< 2\cdot$cutoff, new ranks $>$
cutoff) is exactly what converts the non-growing recycle endpoint of the bare dichotomy
(Theorem~\ref{record:257bm-k1}) into strict progress. The below-half branch never fires twice from the seed
arithmetic, so the supply is needed at essentially every step. Ends literally at
$\mathrm{CofinalExactLocalMersenneHalfRows}$.
\par\noindent\ev{Lean} \scale{uniform} \coord{binary digits} \lean{cofinalExactLocalMersenneHalfRows\_of\_criticalQuotientSupply}{BooleanMobiusCriticalCapacityCofinal.lean:1309} \lean{ProtectedExactLocalMersenneRow}{BooleanMobiusCriticalCapacityCofinal.lean:1087} \lean{exists\_laterProtectedExactLocalMersenneRow}{BooleanMobiusCriticalCapacityCofinal.lean:1139}
\end{thm}

\begin{thm}[The quotient condition on greedy prefixes]\label{record:257bm-c6}
Let $D_c=G\cap\{2,\ldots,c-1\}$.  The condition on the actual greedy
prefixes is
\[
 \begin{gathered}
 c\ge4,\quad c\notin G\\
 \Longrightarrow\qquad Q(D_c\cup\{c\},2c-2)\ge2^{2c-3}.
 \end{gathered}
\]
It asks for the quotient bound at every skipped rank of the specified
half-greedy orbit.  This restricts the crossing test of
Theorem~\ref{record:257bm-c4} to its actual greedy prefixes.  The finite
uniqueness results justify that restriction where the cited induction
uses it; they do not prove the displayed inequality.
\par\noindent\ev{Open} \scale{uniform} \coord{greedy recurrence} \lean{HalfGreedySkippedCriticalQuotientSupply}{BooleanMobiusCriticalCapacityCofinal.lean:178}
\end{thm}
\leannote{Lean: \lproof{Erdos249257/BooleanMobiusCriticalCapacityCofinal.lean}{1068}{skippedCoreCriticalQuotientSupply_iff_halfGreedySkipped}.}

\begin{thm}[Two consecutive skips]\label{record:257bm-c6a}
If $c\ge6$ and the real half-greedy rule skips both $c$ and $c+1$,
put $D=G\cap\{2,\ldots,c-1\}$.  Then
\[
 S(D,1,2c-3)<2^{c-3}.
\]
The one-step quotient recurrence gives
$S(D,1,2c-2)<2^{c-2}$, so
Theorem~\ref{record:257bm-c7} supplies an exact row at depth $2c-2$
whose new ranks are all greater than $c$.
The hypothesis concerns this pair of skipped ranks; no unbounded
sequence of such pairs is proved here.  For the stated precritical-suffix
condition at every skipped rank, the remaining tests are the
skip-then-take cases, with ranks $4$ and $5$ handled separately in
the cited proof.
\par\noindent\ev{Lean} \scale{uniform} \coord{greedy recurrence} \lean{halfGreedy\_precriticalSuffix\_lt\_of\_next\_skip}{BooleanMobiusCriticalCapacityCofinal.lean:682} \lean{halfGreedySkippedCriticalQuotientSupply\_of\_precriticalSuffix}{BooleanMobiusCriticalCapacityCofinal.lean:1001} \lean{halfGreedySkippedPrecriticalSuffixSupply\_iff\_preTake}{BooleanMobiusCriticalCapacityCofinal.lean:838}
\end{thm}
\begin{proof}
The two skips imply $0<\rho=1/2-X_D(2)<w_{c+1}$.
Put $M=2c-3$.  The fractional-part error in $Q(D,M)$ is less
than $c-2$, because $D\subseteq\{2,\ldots,c-1\}$.  Therefore
\[
 S(D,1,M)
 <\frac{2^{2c-3}}{2^{c+1}-1}+c-3
 <2^{c-4}+c-2\le2^{c-3},
\]
where the last inequality holds for $c\ge6$.
The suffix is nonnegative because $X_D(2)<1/2$, so its next-step
identity is $S(D,1,M+1)=2S(D,1,M)+1-c_D(M+1)$.
Integrality gives the required strict bound at $M+1$.
\end{proof}


\begin{thm}[A later skip after a selected block]\label{record:257bm-c6b}
Suppose the real half-greedy rule skips rank $c$, takes
$c+1,\ldots,c+t-1$, and skips $c+t$.  If
\[
 0<t\le c-3,\qquad c-2\le2^{c-t-3},
\]
then the precritical suffix bound at $c$ holds.  It yields sharp
capacity and hence an exact row at depth $2c-2$.
The case $t=1$ includes Theorem~\ref{record:257bm-c6a}; the proof
also permits longer selected blocks.  The arithmetic condition is
precisely $t\le c-3-\lceil\log_2(c-2)\rceil$.
Its role in the proof is to make the dyadic allowance $2^{c-t-3}$
cover the bound $|D|\le c-2$ for the earlier selected support.

To obtain cofinal exact rows by this result, the stated gap condition
must hold at cofinally many skipped ranks $c$, not just at one rank
or throughout a finite sample.  A finite empirical skip frequency
alone gives no such pointwise gap bound.
\par\noindent\ev{Lean} \scale{uniform} \coord{greedy recurrence} \lean{halfGreedy\_precriticalSuffix\_lt\_of\_future\_skip\_after\_takenBlock}{BooleanMobiusCriticalCapacityCofinal.lean:603} \lean{precriticalCrossingTax\_of\_futureThreshold}{BooleanMobiusCriticalCapacityCofinal.lean:472} \lean{sub\_two\_le\_two\_pow\_sub\_four}{BooleanMobiusCriticalCapacityCofinal.lean:432}
\end{thm}

\subsubsection*{Filling the remaining binary positions}

\begin{thm}[Filling the remaining binary positions]\label{record:257bm-c7}
For any $c\ge4$ and any finite $D\subseteq[2,c)$ with $\mathrm{value}(D)<1/2$ and the sharp
capacity
\[
\ensuremath{S}(D,1,2c-2) < 2^{c-2},
\]
there is an exact row $E$ at endpoint $2c-2$ with $D\subseteq E$ and every new rank strictly
above $c$. The implication requires the displayed capacity bound,
but neither a crossing condition nor a specified real greedy prefix.
The crossing hypothesis used in Theorem~\ref{record:257bm-c4} is
one way to seek that input, not an additional premise of this filling
result.  No strict comparison between the corresponding cofinal
existence statements is asserted.
\par\noindent\ev{Lean} \scale{uniform} \coord{binary digits} \lean{exists\_exactRowStrictUpperFill\_of\_skippedCoreSharpCapacity}{BooleanMobiusSkipRow.lean:238}
\end{thm}

\begin{thm}[An exact sum at depth $2c-2$]\label{record:257bm-c8}
Under the hypotheses of Theorem~\ref{record:257bm-c7}, there is a
set $E\subseteq\{2,\ldots,2c-2\}$ with
\[
 Q(E,2c-2)=2^{2c-3}-1.
\]
This conclusion forgets the additional support-extension information
in that theorem.  It is the form used when only the existence of an
exact row is needed.
\par\noindent\ev{Lean} \scale{uniform} \coord{binary digits} \lean{exactLocalMersenneHalfRow\_two\_mul\_sub\_two\_of\_skippedCoreSharpCapacity}{BooleanMobiusSkipRow.lean:332}
\end{thm}

\begin{prop}[Quotients in the upper half of the index range]\label{record:257bm-c9}
For integers $d\ge2$ and $d\le M<2d$,
\[
 q(M,d)=\left\lfloor\frac{2^M}{2^d-1}\right\rfloor=2^{M-d}.
\]
Indeed, the quotient is $2^{M-d}+2^{M-d}/(2^d-1)$, whose second
term is strictly between zero and one.  In
Theorem~\ref{record:257bm-c7}, the available ranks
$d=c+1,\ldots,2c-2$ therefore supply the binary weights
$2^{c-3},\ldots,1$.  They represent every integer from $0$ to
$2^{c-2}-1$, which explains both the capacity bound and the
choice of terminal depth.
\par\noindent\ev{Lean} \scale{uniform} \coord{binary digits} \lean{localMersenneQuotient\_eq\_two\_pow\_sub\_of\_half\_lt}{BooleanMobiusLocalRepair.lean:26}
\end{prop}
\leannote{Lean: \lproof{Erdos249257/BooleanMobiusLocalRepair.lean}{26}{localMersenneQuotient_eq_two_pow_sub_of_half_lt}, \lproof{Erdos249257/BooleanMobiusLocalRepair.lean}{525}{exists_boolean_word_of_lt_two_pow}.}

\begin{thm}[A finite sum from a skipped prefix]\label{record:257bm-c10}
Let $c\ge4$ and $D\subseteq\{2,\ldots,c-1\}$ satisfy
\[
 0<\frac12-X_D(2)<\frac1{2^c-1}.
\]
Then there is $E$ with $D\subseteq E\subseteq\{2,\ldots,2c-2\}$ and
$Q(E,2c-2)=2^{2c-3}-1$.  No separate sharp-capacity assumption is
required.  The binary completion in the linked proof uses ranks
$c,\ldots,2c-2$; it may therefore insert $c$.
Theorem~\ref{record:257bm-c7} instead assumes the sharper
$(c-2)$-bit bound so that every added rank is strictly greater than $c$.

\par\noindent\ev{Lean} \scale{bounded} \coord{binary digits} \lean{exactLocalMersenneHalfRow\_two\_mul\_sub\_two\_of\_skippedCore}{BooleanMobiusSkipRow.lean:149}
\end{thm}
\leannote{Lean: \lproof{ErdosProblems/Erdos257/PaperCompleteR21/MersenneQuotientRowRecurrences.lean}{264}{paper_exact_row_from_skipped_prefix}.}

For a finite $F\subseteq\{2,\ldots,n\}$ with $X_F(2)>1/2$, take
$c$ to be its first crossing rank and $D=F\cap\{2,\ldots,c-1\}$.
Then $c\ge4$ because $1/3+1/7<1/2$, and the displayed hypotheses
hold: equality at the lower endpoint is excluded by odd denominators.
This is the first-crossing application in
\texttt{exists\_skippedCoreExactRow\_of\_value\_above} and
\texttt{exactLocalMersenneHalfRow\_of\_first\_localMersenne\_crossing}
(\texttt{BooleanMobiusExactRowCrossing.lean:155,\,191}).

\subsubsection*{Reducing the dyadic-boundary checks}

\begin{thm}[Testing one nearest dyadic boundary]\label{record:257bm-c11}
Let $d,E$ be nonnegative integers with $E\le2^{d+1}$, and let
\[
 j_*:=\max\{0\le j\le d:E\le2^{d-j+1}\}.
\]
The boundary $2^{d-j_*+1}$ is the smallest power at least $E$ among
the finite list $2,4,\ldots,2^{d+1}$; in particular $j_*=d$ when $E\le2$.
The following conditions are equivalent:
\[
 \begin{gathered}
 \text{for every }0\le j\le d,\quad
 2^{d-j+1}<E\ \text{or}\ E+2(d+j)\le2^{d-j+1};\\
 E+2(d+j_*)\le2^{d-j_*+1}.
 \end{gathered}
\]
The first condition is $\mathrm{DyadicBandEscape}(d,E)$ and
$j_*$ satisfies $\mathrm{CriticalDyadicBandIndex}(d,E,j_*)$.
\par\noindent\ev{Lean} \scale{uniform} \coord{dyadic-boundary} \lean{CriticalDyadicBandIndex}{HalfUpperResetCriticalBand.lean:33} \lean{dyadicBandEscape\_iff\_exists\_critical}{HalfUpperResetCriticalBand.lean:108} \lean{seamUpperResetCriticalBandEscape\_iff}{HalfUpperResetCriticalBand.lean:883} \lean{half\_mem\_mersenneAchievementSet\_of\_upperResetCriticalBandEscape}{HalfUpperResetCriticalBand.lean:896}
\end{thm}
\leannote{Lean: \lproof{ErdosProblems/Erdos257/PaperCompleteR21/DyadicBandAndTwoSidedBounds.lean}{26}{paper_critical_dyadic_band_index_unique}, \lproof{ErdosProblems/Erdos257/PaperCompleteR21/DyadicBandAndTwoSidedBounds.lean}{52}{paper_critical_dyadic_boundary_is_smallest}, \lproof{ErdosProblems/Erdos257/PaperCompleteR21/DyadicBandAndTwoSidedBounds.lean}{71}{paper_critical_dyadic_band_index_eq_top_of_le_two}, \lproof{ErdosProblems/Erdos257/PaperCompleteR21/DyadicBandAndTwoSidedBounds.lean}{89}{paper_dyadic_band_escape_iff_single_test}.}
\begin{proof}
For $j>j_*$ the boundary is below $E$; for $j<j_*$ the boundary
is larger and the required width $2(d+j)$ is smaller.  This proves
sufficiency of the single inequality; necessity is its instance
at $j_*$.
\end{proof}

For example, $d=5$, $E=20$ gives $j_*=1$ and
$20+12=32$, so equality at the allowed edge satisfies every band condition.

At an actual upper reset, Proposition~\ref{record:257bm-c12}
supplies $E\le2^{d+1}$.  The equivalence replaces all its local
band tests by one test.  The arithmetic inequality at every required
reset remains an unproved hypothesis of the membership implication.

\begin{prop}[An upper bound for the reset expression]\label{record:257bm-c12}
Let $d\ge5$ be an actual upper-reset row.  Write $E_d$ for its
reset charge, the sum of four times the adjacent upper overshoot and
the corresponding nonnegative correction term.  The upper-reset identity is
\[
 \mathrm{rem}(d+1)+E_d=2^{d+1}.
\]
The nonnegativity of the successor remainder therefore gives
$E_d\le2^{d+1}$, the side condition needed for
Theorem~\ref{record:257bm-c11}.  This argument uses the upper-reset
branch assumption.  It does not require the separate, conditional
two-sided bound of Theorem~\ref{thm:two-sided-dyadic}, and does not
assert the charge bound at arbitrary rows.
\par\noindent\ev{Lean} \scale{uniform} \coord{dyadic-boundary} \lean{seamUpperResetCharge\_le}{HalfUpperResetCriticalBand.lean:131}
\end{prop}
\leannote{Lean: \lproof{Erdos249257/HalfCylinderMiddleCarryLowerBound.lean}{3542}{seamUpperBranch_remainder_add_resetCharge_eq}, \lproof{Erdos249257/HalfUpperResetCriticalBand.lean}{131}{seamUpperResetCharge_le}.}

\subsubsection*{Nonnegativity of the omitted-tail margin}

The next three statements compare the conditions of
Theorem~\ref{thm:frozen-margin}.  The first two are equivalent; the third
is sufficient for them.

\begin{defn}[A nonnegativity condition at skipped ranks]\label{record:257bm-c13}
\[
\forall\ \text{skipped rank } n\ge3,\quad 0 \le \ensuremath{F}(n-1,n).
\]
This condition is not established here.  Its usefulness is the following implication:
assuming it, \texttt{half\_mem\_mersenneAchievementSet\_of\_skippedFullShellNonnegative} proves
$1/2\in\ensuremath{\mathcal A}$.
\par\noindent\ev{Open} \scale{uniform} \coord{frozen-margin} \lean{HalfGreedySkippedFullShellNonnegative}{HalfCylinderFullShellSeamBridge.lean:604} \lean{half\_mem\_mersenneAchievementSet\_of\_skippedFullShellNonnegative}{HalfCylinderFullShellSeamBridge.lean:633}
\end{defn}

\begin{thm}[An equivalent vanishing condition]\label{record:257bm-c14}
\[
\begin{gathered}
 \forall\ \text{skipped rank }n\ge3\text{ whose actual word equals the seam-greedy word},\\
 \text{the integer remainder at }n\text{ is zero}.
\end{gathered}
\]
This condition is equivalent to Definition~\ref{record:257bm-c13}.
At an actual skipped rank, the source proves that
$F(n-1,n)<0$ holds exactly when the real prefix agrees with the
integer greedy word and its integer remainder is positive.
On agreement, $F(n-1,n)$ is the negative of that nonnegative
remainder.  Thus nonnegativity of the margin at every skip is
equivalent to vanishing of the remainder at every aligned skip.
The equivalence is proved; neither condition is established for all
required ranks.  Its formal statement is
\texttt{skippedSeamAlignmentZero\_iff\_skippedFullShellNonnegative}.
\par\noindent\ev{Open} \scale{uniform} \coord{frozen-margin} \lean{HalfGreedySkippedSeamAlignmentZero}{HalfCylinderFullShellSeamBridge.lean:593} \lean{skippedSeamAlignmentZero\_iff\_skippedFullShellNonnegative}{HalfCylinderFullShellSeamBridge.lean:671}
\end{thm}
\leannote{Lean: \lproof{Erdos249257/HalfCylinderFullShellSeamBridge.lean}{569}{skipped_fullShell_neg_iff_alignment_and_seamRemainder_pos}, \lproof{Erdos249257/HalfCylinderFullShellSeamBridge.lean}{531}{greedyHalfFrozenMargin_fullShell_eq_neg_seamRemainder_of_alignment}, \lproof{Erdos249257/HalfCylinderFullShellSeamBridge.lean}{671}{skippedSeamAlignmentZero_iff_skippedFullShellNonnegative}.}

\begin{thm}[A sufficient lower bound for the integer remainder]\label{record:257bm-c15}
\[
\forall\ \text{skipped rank }n\ge3,\quad B(2n) < \mathrm{rem}(n),
\]
a sufficient condition for the nonnegativity condition above.
Indeed, a negative margin at an actual skip would force
$1\le\mathrm{rem}(n)\le B(2n)$, contradicting this bound.
This proves the implication, not a converse or strict separation of the conditions.
\par\noindent\ev{Open} \scale{uniform} \coord{frozen-margin} \lean{HalfGreedySkippedSeamEscape}{HalfCylinderFullShellSeamBridge.lean:687} \lean{governedFrozenMarginProducer\_of\_skippedSeamEscape}{HalfCylinderFullShellSeamBridge.lean:695}
\end{thm}
\leannote{Lean: \lproof{ErdosProblems/Erdos257/PaperCompleteR21/SeamEscapeAndTerminalStrip.lean}{36}{paper_seam_escape_forces_remainder_band}, \lproof{ErdosProblems/Erdos257/PaperCompleteR21/SeamEscapeAndTerminalStrip.lean}{47}{paper_seam_escape_implies_full_shell_nonnegative}, \lproof{ErdosProblems/Erdos257/PaperCompleteR21/SeamEscapeAndTerminalStrip.lean}{61}{paper_seam_escape_implies_half_membership}.}

\subsubsection*{Square-root bounds for half-carries}

\begin{thm}[A square-root carry bound implies half-membership]\label{record:257rig-c16}
Write $G$ for the greedy support and $C_G(N)$ for its centred carry. If
\[
\forall N,\qquad C_G(N) \le 2\sqrt N + 4,
\]
then the greedy selected support $G$ is infinite and $X_G(2)=1/2$.
Nonnegativity of $C_G$ is unconditional (see
Theorem~\ref{thm:mobius-centred-nonneg} and the greedy case above); only the upper
$2\sqrt N+4$ bound remains open. This is a different hypothesis from the dyadic-band
condition and the largest-skip condition; logical independence is not
asserted.
\par\noindent\ev{Open} \scale{cofinal} \coord{sqrt-carry-growth} \lean{infinite\_support\_half\_of\_mobiusCenteredHalfCarry\_sqrtBound}{HalfCarryReachability.lean:817} \lean{greedy\_half\_infinite\_of\_mobiusCenteredHalfCarry\_sqrtBound}{HalfCarryReachability.lean:834}
\end{thm}
\leannote{Lean: \lproof{Erdos249257/HalfCarryReachability.lean}{919}{greedy_mobiusCenteredHalfCarry_nonneg}, \lproof{Erdos249257/HalfCarryReachability.lean}{953}{greedy_half_infinite_of_mobiusCenteredHalfCarry_upperBound}.}

\begin{thm}[Terminal bounds at unbounded depths]\label{record:257rig-c17}
Suppose that for every $N\ge0$ there are $M\ge\max\{N,1\}$ and
$D\subseteq\{2,\ldots,M\}$ such that
\[
 |\operatorname{ihc}(D,M-1)|\le B(M)
       =2\lfloor\sqrt M\rfloor+4.
\]
Then some infinite $A\subseteq\Npos$ satisfies $X_A(2)=1/2$.
No compatibility between different finite supports or bound on their
earlier carries is required.

These are not the exact-row conditions of
Definition~\ref{record:257bm-d4}.  For such a finite support,
\[
 \operatorname{ihc}(D,M-1)=2^{M-1}-Q(D,M),
\]
so an exact row has carry $1$, whereas the displayed bound allows
several positive and negative values.  At $M=6$, for example,
$D=\{2,3\}$ has $Q(D,6)=30$ and carry $2$: it satisfies the
terminal bound but is not an exact row.
Both cofinal existence statements nevertheless imply, and are
implied by, half-membership, as explained in
Proposition~\ref{prop:collapsed-list}.  The difference between their
finite data must not be confused with a strict logical weakening
of the membership problem.
\par\noindent\ev{Open} \scale{cofinal} \coord{sqrt-carry-growth} \lean{HalfCarryCofinalTerminalOnlyStrip}{TerminalOnlyCofinal.lean:34} \lean{exists\_infinite\_support\_half\_of\_cofinalTerminalOnlyStrip}{TerminalOnlyCofinal.lean:192}
\end{thm}
\leannote{Lean: \lproof{ErdosProblems/Erdos257/PaperCompleteR21/TerminalStripExactRowGap.lean}{44}{paper_terminal_strip_forces_half_membership}, \lproof{ErdosProblems/Erdos257/PaperCompleteR21/TerminalStripExactRowGap.lean}{88}{paper_integerHalfCarry_eq_two_pow_sub_localPrefixQuotient}, \lproof{ErdosProblems/Erdos257/PaperCompleteR21/TerminalStripExactRowGap.lean}{106}{paper_exact_row_integerHalfCarry_eq_one}, \lproof{ErdosProblems/Erdos257/PaperCompleteR21/TerminalStripExactRowGap.lean}{122}{paper_terminal_strip_witness_six}, and 1 further declaration in the \hyperref[sec:coverage]{coverage section}.}

\begin{thm}[Cofinal returns of the greedy carry]\label{record:257rig-c18}
Suppose that for every $N\ge0$ there is $M\ge N$ such that
\[
 \operatorname{ihc}(G,M)\le B(M+1)
       =2\lfloor\sqrt{M+1}\rfloor+4.
\]
Then $G$ is infinite and $X_G(2)=1/2$.  Unlike
Theorem~\ref{record:257rig-c16}, this hypothesis bounds the actual greedy
carry only at unboundedly many indices.  It still concerns the same fixed
support $G$, not independently chosen finite supports.
The last source below compares this carry with the carry obtained by
fixing a finite prefix. Its comparison includes a contribution from
later omitted ranks. No equivalence with the finite-support hypothesis
of Theorem~\ref{record:257bm-c7} follows without controlling that
additional contribution.
\par\noindent\ev{Open} \scale{cofinal} \coord{sqrt-carry-growth} \lean{GreedyHalfCarryCofinalStripReturn}{CofinalStripReturn.lean:76} \lean{greedy\_half\_infinite\_of\_cofinalStripReturn}{CofinalStripReturn.lean:130} \lean{halfGreedy\_precriticalSuffix\_lt\_iff\_futureSkipCoverage}{BooleanMobiusCriticalCapacityCofinal.lean:949}
\end{thm}

\subsubsection*{A margin at an unspecified later horizon}

\begin{thm}[An eventual nonnegative margin suffices]\label{record:257bm-c19}
Fix a positive depth $k$ and put $D=G\cap\{2,\ldots,k\}$.
The following conditions are equivalent:
\[
 r_k(1/2)<2^{-(k+1)}
 \quad\Longleftrightarrow\quad
 F_k(J)\ge0\ \text{for some }J\ge0.
\]
Here $F_k(J)$ is the integer expression defined in
Theorem~\ref{thm:frozen-margin}.  Its normalised value is
\[
 2^{-J}F_k(J)=\sum_{i=1}^{J}c_D(k+1+i)2^{-i}-C_D(k).
\]
This normalised expression, not necessarily $F_k(J)$ itself, is
nondecreasing in $J$.  Its limit is
$1-2^{k+1}r_k(1/2)$, by the finite-support Lambert identity.
This is positive exactly under the strict dyadic inequality on the left.  At positive $k$, oddness of the
reduced excess numerator excludes equality at zero.

The result supplies some finite horizon, not the horizon $J=c-3$
required in the crossing application.  More quantitatively, write
$\eta=1-2^{k+1}r_k(1/2)>0$.  Since
$\sum_{r\ge1}c_D(m+r)2^{-r}\le m+2$, the omitted tail shows that
$(k+J+3)2^{-J}<\eta$ suffices for $F_k(J)>0$.
The value $\eta$ is rational for the finite support $D$, so this is an
effective sufficient test.  Proving it at the prescribed horizon is a
separate arithmetic obligation.
\par\noindent\ev{Lean} \scale{uniform} \coord{frozen-margin} \lean{exists\_greedyHalfFrozenMargin\_nonneg\_of\_excess\_neg}{HalfCylinderFiniteShadow.lean:1198} \lean{exists\_greedyHalfFrozenMargin\_nonneg\_iff\_excess\_neg}{HalfCylinderFiniteShadow.lean:1229}
\end{thm}
\leannote{Lean: \lproof{ErdosProblems/Erdos257/PaperCompleteR21/EventualNonnegativeMargin.lean}{56}{paper_halfGreedyPrefixSupport_eq_greedy_inter_Icc}, \lproof{ErdosProblems/Erdos257/PaperCompleteR21/EventualNonnegativeMargin.lean}{131}{paper_eventual_nonnegative_margin_equivalence}, \lproof{ErdosProblems/Erdos257/PaperCompleteR21/EventualNonnegativeMargin.lean}{165}{paper_frozen_margin_normalised_value}, \lproof{ErdosProblems/Erdos257/PaperCompleteR21/EventualNonnegativeMargin.lean}{180}{paper_frozen_margin_normalised_monotone}, and 7 further declarations in the \hyperref[sec:coverage]{coverage section}.}

\subsubsection*{Equivalent conditions for infinitely many greedy skips}

\begin{thm}[Half-membership and infinitely many skips]\label{record:257bm-c20}
\[
(1/2:\mathbb{R})\in\ensuremath{\mathcal A} \;\Longleftrightarrow\; (\mathrm{greedyMersenneSkippedSupport}(1/2)).\mathrm{Infinite}.
\]
Odd-denominator parity supplies the positivity clause in Theorem~\ref{record:257bm-c2} automatically. Thus its
cofinal-positive-skip hypothesis is exactly the assertion that the greedy skipped support is
infinite, which the displayed theorem identifies with half-membership. See Observation~\ref{record:257bm-k5} for the
reason this restatement does not itself prove membership.
\par\noindent\ev{Lean} \scale{cofinal} \coord{topological-achievement-set} \lean{half\_mem\_mersenneAchievementSet\_iff\_greedySkippedSupport\_infinite}{GreedyAchievementSet.lean:2583} \lean{mem\_mersenneAchievementSet\_of\_greedySkippedSupport\_infinite}{GreedyAchievementSet.lean:1528}
\end{thm}

\subsection{Invariants}

This subsection collects identities and bounds used in the preceding
criteria.  Each statement retains its own hypotheses; a bound for one
branch is not an invariant of the entire recurrence.

\subsubsection*{One-step quotient identities}

\begin{thm}[The next floor quotient]\label{record:257bm-i1a}
For $d\ge2$ and $M\ge0$, the floor quotient satisfies
\[
 q(M+1,d)=2q(M,d)+\mathbf1_{d\mid M+1}.
\]
Indeed, $q(M,d)=\sum_{j=1}^{\lfloor M/d\rfloor}2^{M-jd}$.
Increasing $M$ doubles each existing term and adds $1$ exactly when
$M+1$ is divisible by $d$.  Thus a quotient doubles or doubles plus
one; it does not remain unchanged except when both values are zero.
Summing over a fixed finite support gives the corresponding update
with added term $c_D(M+1)$.
\par\noindent\ev{Lean} \scale{uniform} \coord{binary digits} \lean{localMersenneQuotient\_endpoint\_succ}{BooleanMobiusExactTransition.lean:29}
\end{thm}
\leannote{Lean: \lproof{ErdosProblems/Erdos257/PaperCompleteR21/MersenneQuotientRowRecurrences.lean}{37}{paper_next_floor_quotient}, \lproof{ErdosProblems/Erdos257/PaperCompleteR21/MersenneQuotientRowRecurrences.lean}{74}{paper_next_floor_quotient_no_fixed_point}, \lproof{ErdosProblems/Erdos257/PaperCompleteR21/MersenneQuotientRowRecurrences.lean}{82}{paper_next_quotient_sum}, \lproof{ErdosProblems/Erdos257/PaperCompleteR21/MersenneQuotientRowRecurrences.lean}{44}{paper_floor_quotient_geometric_sum}.}

\begin{thm}[The next quotient sum]\label{record:257bm-i1b}
For a fixed finite set $D\subseteq\{2,3,\ldots\}$ and $M\ge0$,
\[
 Q(D,M+1)=2Q(D,M)+c_D(M+1).
\]
This is the sum of Theorem~\ref{record:257bm-i1a} over $d\in D$.
Keeping $D$ fixed is essential: changing the support between steps
adds a separate difference of quotient sums.  The exclusion of
$d=1$ is also essential for the displayed correction term, since
$q(M,1)=2^M$ doubles without an added unit.
\par\noindent\ev{Lean} \scale{uniform} \coord{binary digits} \lean{localPrefixQuotient\_succ}{BooleanMobiusExactTransition.lean:131}
\end{thm}

\begin{thm}[The signed endpoint recurrence]\label{record:257bm-i1c}
The signed next-step expression is
$H(D,k,n)=2S(D,k,n-1)+1-c_D(n)$.  The last term counts the selected
exponents dividing the new endpoint. This is the same recurrence shape as the
generic tempered-orbit recurrence $u(N+1)=2u(N)-v\cdot c(N+1)$ from the public
\texttt{GenericTailOrbitRigidity.lean} T7 criterion (Theorem~\ref{record:257bm-i-t7} below), specialised to Mersenne
local repair. A matching recurrence alone does not imply rationality:
Theorem~\ref{record:257bm-i-t7} also requires one fixed coefficient
sequence, integer states and a vanishing scaled limit.  Those additional
conditions must be checked before that criterion can be applied to these
finite-row quantities.
\par\noindent\ev{Lean} \scale{uniform} \coord{binary digits} \lean{localEndpointDefect\_succ}{BooleanMobiusExactTransition.lean:189} \lean{localRepairInteger\_eq\_localEndpointDefect\_succ}{BooleanMobiusExactTransition.lean:203}
\end{thm}
\leannote{Lean: \lproof{ErdosProblems/Erdos257/PaperCompleteR21/MersenneQuotientRowRecurrences.lean}{92}{paper_signed_endpoint_recurrence}, \lproof{ErdosProblems/Erdos257/PaperCompleteR21/MersenneQuotientRowRecurrences.lean}{99}{paper_endpoint_term_counts_divisors}, \lproof{ErdosProblems/Erdos257/PaperCompleteR21/MersenneQuotientRowRecurrences.lean}{106}{paper_signed_endpoint_defect_succ}, \lproof{ErdosProblems/Erdos257/PaperCompleteR21/MersenneQuotientRowRecurrences.lean}{114}{paper_repair_integer_eq_endpoint_defect}.}

\subsubsection*{Bounds for the remaining binary positions}

\begin{defn}[The sharper binary bound]\label{record:257bm-i-cap}
Let $c\ge4$, $D\subseteq\{2,\ldots,c-1\}$, and $X_D(2)<1/2$.
The sharper bound is the condition
\[
 S(D,1,2c-2)<2^{c-2}.
\]
It permits a binary completion using ranks $c+1,\ldots,2c-2$,
with rank $c$ omitted, as explained below.
\par\noindent\ev{Lean} \scale{n/a} \coord{binary digits} \lean{localBinarySuffix\_two\_mul\_sub\_two\_lt\_criticalCapacity\_iff}{BooleanMobiusSkippedCoreCriticalCapacity.lean:22}
\end{defn}

\begin{thm}[Equivalent forms of the sharper bound]\label{record:257bm-i5}
Under the hypotheses of Definition~\ref{record:257bm-i-cap}:
\[
\ensuremath{S}(D,1,2c-2) < 2^{c-2} \;\Longleftrightarrow\; 2^{(2c-2)-1} \le \ensuremath{Q}(\mathrm{insert}\ c\ D,\ 2c-2).
\]
The threshold on the right is one greater than the exact-row target
$2^{2c-3}-1$.  Thus this is a test for omitting rank $c$, not for
including it in the completed row.
\par\noindent\ev{Lean} \scale{uniform} \coord{binary digits} \lean{localBinarySuffix\_two\_mul\_sub\_two\_lt\_criticalCapacity\_iff}{BooleanMobiusSkippedCoreCriticalCapacity.lean:22}
\end{thm}

\begin{proof}
Put $M=2c-2$.  Since $X_D(2)<1/2$, one has
$Q(D,M)<2^{M-1}$, and the integer remainder is
$S(D,1,M)=2^{M-1}-Q(D,M)-1\ge0$ without truncation.
The geometric quotient identity gives
$q(M,d)=2^{M-d}$ for $c\le d\le M$, since $2d>M$.
In particular $q(M,c)=2^{c-2}$, so
\[
 S(D,1,M)<2^{c-2}
 \quad\Longleftrightarrow\quad
 Q(D\cup\{c\},M)\ge2^{M-1}.
\]
For the completion, use the ranks $c+1,\ldots,M$ instead.
Their quotient weights are $2^{c-3},\ldots,1$ and represent every
integer from $0$ to $2^{c-2}-1$.  A subset of them therefore has
quotient sum $S(D,1,M)$ and completes $D$ to the exact target
$2^{M-1}-1$ while omitting $c$.  These finite identities prove
the equivalence and explain the sharper bound.  They do not assert
that a completed row stays below one half in real value.
\end{proof}

For example, $D=\{2,3\}$ and $c=5$ give $M=8$, $Q(D,8)=121$
and remainder $6<8$.  Inserting rank $5$ adds $8$, giving $129$
rather than the exact target $127$.  The correct completion adds
ranks $6$ and $7$, whose quotients are $4$ and $2$.  \ev{Math}

\begin{thm}[An unconditional bound with one extra bit]\label{record:257bm-i6}
Unconditionally, for every $c\ge4$ and every below-half core $D\subseteq[2,c)$ with deficit
$<\ensuremath{w}(c)$:
\[
\ensuremath{S}(D,1,2c-2) < 2^{c-1}.
\]
The bound holds uniformly in $c$, but is twice the threshold needed
in Theorem~\ref{record:257bm-i5} to apply Theorem~\ref{record:257bm-c7}.
Before using $|D|\le c-2\le2^{c-2}$, the proof gives the sharper
additive estimate
$\ensuremath{S}(D,1,2c-2)<2^{c-2}+|D|$.
Thus the critical-capacity inequality would follow by excluding the
integer band $[2^{c-2},\,2^{c-2}+c-3]$, which contains $c-2$ integers. This is the same band
shape as Theorem~\ref{record:257bm-c11}'s dyadic-band condition.
Its width is linear in $c$, whereas the square-root reset condition
discussed in Section~\ref{sec:o4} has width $2^{(r+5)/2}$ in a
different parameter $r$.  Comparing these widths alone proves no implication
between the two hypotheses.
\par\noindent\ev{Lean} \scale{uniform} \coord{binary digits} \lean{localBinarySuffix\_two\_mul\_sub\_two\_lt\_upperHalfCapacity}{BooleanMobiusSkippedCoreExactRow.lean:123} \lean{exists\_upperHalfBooleanWord\_of\_skippedCore}{BooleanMobiusSkippedCoreExactRow.lean:228}
\end{thm}
\leannote{Lean: \lproof{ErdosProblems/Erdos257/PaperCompleteR21/MersenneQuotientRowRecurrences.lean}{126}{paper_unconditional_bound_one_extra_bit}, \lproof{ErdosProblems/Erdos257/PaperCompleteR21/MersenneQuotientRowRecurrences.lean}{141}{paper_sharper_additive_estimate}, \lproof{ErdosProblems/Erdos257/PaperCompleteR21/MersenneQuotientRowRecurrences.lean}{236}{paper_capacity_band_exclusion}.}

\subsubsection*{Doubling the endpoint}

\begin{thm}[The binary bound after doubling]\label{record:257bm-i7}
Let $n\ge6$ and $D\subseteq\{2,\ldots,n\}$ satisfy
$2\in D$, $Q(D,n)=2^{n-1}-1$, and $X_D(2)<1/2$.  Then
\[
 S(D,1,2n-1)<2^{n-1}.
\]
The missing quotient can therefore be filled using only the ranks
$n+1,\ldots,2n-1$, without changing $D$.
\par\noindent\ev{Lean} \scale{uniform} \coord{binary digits} \lean{localBinarySuffix\_two\_mul\_sub\_one\_lt\_upperWindow\_of\_exact\_below}{BooleanMobiusExactRowDoubling.lean:44}
\end{thm}
\leannote{Lean: \lproof{Erdos249257/BooleanMobiusExactRowDoubling.lean}{44}{localBinarySuffix_two_mul_sub_one_lt_upperWindow_of_exact_below}, \lproof{Erdos249257/BooleanMobiusExactRowDoubling.lean}{185}{exists_exactRowStrictUpperExtension_two_mul_sub_one_of_exact_below}.}

\begin{proof}
Write $F_M=\sum_{d\in D}\{2^M/(2^d-1)\}$ and $P=2^{n-1}$.
The exact row at $n$ gives
$2^n(1/2-X_D(2))=1-F_n$.
At $M=2n-1$ this becomes
\[
 S(D,1,2n-1)=P(1-F_n)+F_{2n-1}-1.
\]
The below-half assumption makes this an ordinary integer difference,
so no truncation is needed.  Since $2\in D$, one has $F_n\ge1/3$;
also $F_{2n-1}\le|D|\le n-1$.  Therefore
\[
 S(D,1,2n-1)\le\frac23P+n-2<P,
\]
using $3(n-2)<2^{n-1}$ for $n\ge6$.
The ranks $n+1,\ldots,2n-1$ supply exactly $n-1$ binary digits,
with total capacity $2^{n-1}-1$.  The binary expansion of the
nonnegative missing quotient gives the required extension.
\end{proof}

For example, $n=6$ and $D=\{2,3,6\}$ give $X_D(2)=31/63$,
$Q(D,11)=1006$ and missing quotient $1023-1006=17$.
The ranks $7$ and $11$ supply $16+1$.
The resulting exact row $E=\{2,3,6,7,11\}$ has
$X_E(2)-1/2=13955/32756094>0$.
Thus the extension is exact at its integer target, but need not
remain below half in real value.

\begin{prop}[Every exact sum contains exponent two]\label{record:257bm-i-rank2}
Let $n\ge3$ and let $D\subseteq\{2,\ldots,n\}$ satisfy
$Q(D,n)=2^{n-1}-1$.  Then $2\in D$.
Thus the rank-two hypothesis in Theorem~\ref{record:257bm-i7}
is automatic for its exact rows.  The interval restriction on
$D$ is part of the assertion, not an assumption about arbitrary
finite quotient sums.
\par\noindent\ev{Lean} \scale{uniform} \coord{binary digits} \lean{two\_mem\_of\_exact\_localMersenneQuotient}{BooleanMobiusExactRowRankTwo.lean:23}
\end{prop}

\subsubsection*{The depth-six example}

\begin{defn}[An exact sum at depth six]\label{record:257bm-i-seed}
$\mathrm{exactRowSixSupport} := \{2,3,6\}$ satisfies
\[
\ensuremath{Q}(\{2,3,6\},6) = 2^5-1 = 31,\qquad \ensuremath{V}(\{2,3,6\}) < 1/2,
\]
hence $\mathrm{ExactLocalMersenneHalfRow}(6)$. Closed numeric fixture, \texttt{norm\_num}-checked.
This is the chosen seed for the $n\ge6$ transition theorem; the induction
starts at this depth.  It is not the first
exact quotient row without that restriction: at $n=2$, $D=\{2\}$ already
gives $Q(D,2)=1=2^{2-1}-1$.
\par\noindent\ev{Cert} \scale{fixed} \coord{divisor counts and finite sums} \lean{exactRowSixSupport}{BooleanMobiusExactRowSeed.lean:14} \lean{exactLocalMersenneHalfRow\_six}{BooleanMobiusExactRowSeed.lean:34}
\end{defn}

\subsubsection*{Approximation error and finite denominators}

\begin{prop}[Error in the finite subseries value]\label{record:257bm-i9}
An exact row at endpoint $n$ has real value within $O(n/2^n)$ of $1/2$: the quantitative bound
$|y_n-1/2|\le(n+1)/2^n$ consumed by Theorem~\ref{record:257bm-c1} to turn a sequence of exact rows tending to depth
infinity into a sequence of values tending to $1/2$.
\par\noindent\ev{Lean} \scale{uniform} \coord{divisor counts and finite sums} \lean{abs\_localMersennePrefixValue\_sub\_half\_le}{BooleanMobiusGlobalRepair.lean:243}
\end{prop}

\begin{prop}[No finite support has value one half]\label{record:257bm-i10}
A finite sum of reciprocals of odd integers has odd denominator in
lowest terms, so it cannot equal $1/2$. Applied to the denominators
$2^a-1$, this proves that a support representing $1/2$ must be infinite,
as used in Theorem~\ref{record:257bm-c1}.
\par\noindent\ev{Cited} \scale{uniform} \coord{p-adic} \lean{finite\_boolSupport\_ne\_half}{HalfCarryReachability.lean:589}
\end{prop}
\leannote{Lean: \lproof{ErdosProblems/Erdos257/PaperCompleteR21/OddReciprocalDenominators.lean}{54}{paper_finite_sum_inv_odd_den_odd}, \lproof{ErdosProblems/Erdos257/PaperCompleteR21/OddReciprocalDenominators.lean}{73}{paper_finite_sum_inv_odd_ne_half}, \lproof{ErdosProblems/Erdos257/PaperCompleteR21/OddReciprocalDenominators.lean}{94}{paper_finiteErdosSum_den_odd}, \lproof{ErdosProblems/Erdos257/PaperCompleteR21/OddReciprocalDenominators.lean}{114}{paper_finite_support_series_ne_half}, and 1 further declaration in the \hyperref[sec:coverage]{coverage section}.}

\subsubsection*{Uniqueness for gap-dominated finite weights}

\begin{thm}[Small remainder forces the greedy word]\label{record:257bm-i11a}
Let $g\ge1$ and let $w_1,\ldots,w_m$ be nonnegative integer weights
satisfying
\[
 w_i\ge g+\sum_{j>i}w_j\qquad(1\le i\le m).
\]
Fix an integer capacity $C\ge0$. The greedy word $\gamma$ is obtained
by taking $w_i$ exactly when it does not exceed the remaining capacity.
For an admissible Boolean word $\varepsilon$, meaning
$\sum_i\varepsilon_iw_i\le C$, one has
\[
 C-\sum_i\varepsilon_iw_i<g
 \quad\Longleftrightarrow\quad
 \varepsilon=\gamma\ \text{and}\ 
 C-\sum_i\gamma_iw_i<g.
\]
Indeed, two different words have sums separated by at least $g$, as
seen at their first differing index; the greedy sum is the largest
admissible sum. Thus there is at most one admissible word with remainder
below $g$, and any such word is greedy. The greedy word alone need not
have remainder below $g$: for weights $(5)$, gap $g=2$ and capacity
$C=3$, its remainder is $3$.
\par\noindent\ev{Lean} \scale{uniform} \coord{greedy recurrence} \lean{remainder\_lt\_gap\_iff\_eq\_integerGreedyBits}{BooleanMobiusGreedyReduction.lean:918}
\end{thm}

\begin{prop}[Gap domination of the integer weights]\label{record:257bm-i11b}
For integers $1\le d\le R\le M$, the quotient weights satisfy
\[
 \left\lfloor\frac{2^M}{2^d-1}\right\rfloor
 \ge 2^{M-R}+
 \sum_{j=d+1}^{R}\left\lfloor\frac{2^M}{2^j-1}\right\rfloor.
\]
Thus the finite weight list at ranks $2,\ldots,R$ satisfies the
hypothesis of Theorem~\ref{record:257bm-i11a} with $g=2^{M-R}$.
At $M=2R-1$ this gap is $2^{R-1}$, and at $M=2R$ it is $2^R$.
The inequality is unconditional; it does not assert that the greedy
remainder is smaller than the gap.
\par\noindent\ev{Lean} \scale{uniform} \coord{greedy recurrence} \lean{localMersenneWeightsFrom\_gapDominates}{BooleanMobiusGreedyReduction.lean:684} \lean{localMersenneWeights\_gapDominates\_even}{BooleanMobiusGreedyReduction.lean:846} \lean{localMersenneWeights\_gapDominates\_odd}{BooleanMobiusGreedyReduction.lean:855}
\end{prop}

\begin{prop}[Uniqueness of the finite greedy representation]\label{record:257bm-i11c}
Let $1\le M$, $0\le R\le M$, and $0\le a<2^{M-R}$. If
\[
 \sum_{d=2}^{R}\varepsilon_d
       \left\lfloor\frac{2^M}{2^d-1}\right\rfloor+a
       =2^{M-1}-1,\qquad \varepsilon_d\in\{0,1\},
\]
then $\varepsilon$ is the greedy word for capacity $2^{M-1}-1$
and $a$ is its remainder. This follows by applying
Theorem~\ref{record:257bm-i11a} with the preceding gap inequality.
The statement identifies any such representation; it does not prove
that a representation with $a<2^{M-R}$ exists. The degenerate formal
case $M=0$ has empty word and zero target, using truncated natural-number
subtraction.
\par\noindent\ev{Lean} \scale{uniform} \coord{greedy recurrence} \lean{localMersenneHalfTarget\_lower\_word\_eq\_greedy\_and\_remainder\_eq}{BooleanMobiusGreedyReduction.lean:997}
\end{prop}

\subsubsection*{A geometric form of the quotient condition}

\begin{thm}[A division-free form of the binary bound]\label{record:257bm-i12}
For $M\ge0$ and $d\ge2$,
\[
 \left\lfloor\frac{2^M}{2^d-1}\right\rfloor
   =\sum_{j=1}^{\lfloor M/d\rfloor}2^{M-jd},
\]
where an empty sum is zero. To see this, write $M=qd+r$ with
$0\le r<d$ and expand the finite geometric sum. The remaining fraction
is $2^r/(2^d-1)$, which lies strictly between $0$ and $1$.
Substitution in Theorem~\ref{record:257bm-i5} expresses its quotient
condition as a finite sum of powers of $2$. This is an exact rewriting,
not a weaker hypothesis or a new existence result.
\par\noindent\ev{Lean} \scale{uniform} \coord{geometric-series} \lean{localMersenneQuotient\_eq\_geometric}{BooleanMobiusCriticalCapacityGeometric.lean:27} \lean{localBinarySuffix\_two\_mul\_sub\_two\_lt\_criticalCapacity\_iff\_geometric}{BooleanMobiusCriticalCapacityGeometric.lean:196} \lean{localBinarySuffix\_two\_mul\_sub\_two\_lt\_criticalCapacity\_iff\_geometricCore}{BooleanMobiusCriticalCapacityGeometric.lean:208}
\end{thm}

\subsubsection*{Bounds for a general coefficient sequence}

\begin{thm}[An upper bound for a binary coefficient tail]\label{record:257bm-i2}
Let $c:\N\to\N$ satisfy $c(n)\le n$ for every $n$.  Then
\[
 \sum_{r\ge1}c(N+r)2^{-r}\le N+2\qquad(N\ge0).
\]
Indeed, $c(N+r)\le N+r$, while $\sum_{r\ge1}2^{-r}=1$ and
$\sum_{r\ge1}r2^{-r}=2$.  Thus the scaled tail is $O(N)$, and in
particular $o(2^N)$, as required in Theorem~\ref{record:257bm-i-t7}.
The coefficient bound also holds for Euler's totient function.
\par\noindent\ev{Cited} \scale{uniform} \coord{binary digits} \lean{binaryCoeffTail\_le}{GenericTailOrbitRigidity.lean:85}
\end{thm}
\leannote{Lean: \lproof{Erdos249257/GenericTailOrbitRigidity.lean}{85}{binaryCoeffTail_le}, \lproof{Erdos249257/GenericTailOrbitRigidity.lean}{91}{binaryCoeffTail_div_pow_tendsto_zero}.}

\begin{thm}[Rationality through an integer recurrence]\label{record:257bm-i-t7}
Let $c:\N\to\N$ satisfy $c(n)\le n$ for every $n$.  Then
$\sum_{n\ge1}c(n)2^{-n}$ is rational if and only if there exist a positive
integer $v$ and a sequence $u:\N\to\Z$ such that
\[
 u(N+1)=2u(N)-v c(N+1)\quad(N\ge0),\qquad
 \frac{u(N)}{2^N}\longrightarrow0.
\]
The limit condition is $u(N)=o(2^N)$; it does not assume that $u$ is bounded.
No divisor-count hypothesis is imposed on $c$.  Thus the ordinary statement
also applies to $c=\varphi$; the identification with the separate formal
\#249 development, including its indexing conventions, is not asserted here.
\par\noindent\ev{Lean} \scale{uniform} \coord{binary digits} \lean{IsTemperedBinaryOrbit}{GenericTailOrbitRigidity.lean:67} \lean{binaryCoeffSeries}{GenericTailOrbitRigidity.lean:37} \lean{binaryCoeffSeries\_rational\_iff\_exists\_temperedBinaryOrbit}{GenericTailOrbitRigidity.lean:426}
\end{thm}
\leannote{Lean: \lproof{Erdos249257/GenericTailOrbitRigidity.lean}{435}{not_irrational_binaryCoeffSeries_iff_exists_temperedBinaryOrbit}.}

The equivalence follows by telescoping the recurrence:
\[
 \frac{u(N)}{2^N}=u(0)-v\sum_{n=1}^N\frac{c(n)}{2^n}.
\]
The limit condition gives the rational value $u(0)/v$.  Conversely, if the
series equals $p/v$, define
$u(N)=2^Np-v\sum_{n=1}^N2^{N-n}c(n)$.  This is an integer and equals
$v\sum_{r\ge1}c(N+r)2^{-r}$, so the preceding tail bound proves the limit.
This argument also explains why an arbitrary integer solution of the
recurrence is insufficient: its homogeneous part can grow like $2^N$.

For example, the finite support $A=\{2\}$ has
$c_A(n)=\mathbf1_{2\mid n}$ and $X_A(2)=1/3$.  With $v=3$, the sequence
$u(N)=1$ for even $N$ and $u(N)=2$ for odd $N$ satisfies the recurrence.
The criterion therefore includes finite supports; it does not provide the
infinite support needed to contradict Problem~257.

\begin{thm}[M\"obius inversion of the divisor counts]\label{record:257bm-i-mob}
For a set $A\subseteq\N$ and each positive integer $n$,
\[
 \sum_{d\mid n}\mu(d)c_A(n/d)=\mathbf1_A(n),
\]
where $\mu$ is the M\"obius function and $c_A(n)$ counts the positive
elements of $A$ dividing $n$.  Thus the divisor counts determine the
positive support exactly.  In particular their M\"obius transform takes
only the values $0$ and $1$.  The identity does not depend on a base.
\par\noindent\ev{Lean} \scale{uniform} \coord{mobius-inversion} \lean{moebius\_mul\_supportCoeffAF}{BooleanMobiusCarry.lean:95} \lean{mobius\_supportCoeff\_boolean}{BooleanMobiusCarry.lean:118} \lean{card\_divisors\_le\_two\_mul\_sqrt}{BooleanMobiusCarry.lean:209}
\end{thm}
\leannote{Lean: \lproof{Erdos249257/BooleanMobiusCarry.lean}{95}{moebius_mul_supportCoeffAF}, \lproof{Erdos249257/BooleanMobiusCarry.lean}{112}{mobius_supportCoeff_eq_one_iff}, \lproof{Erdos249257/BooleanMobiusCarry.lean}{118}{mobius_supportCoeff_boolean}.}

\begin{thm}[The support series as a coefficient series]\label{record:257bm-i-bridge}
For $A\subseteq\Npos$,
\[
 X_A(2)=\sum_{a\in A}\frac1{2^a-1}
       =\sum_{n\ge1}\frac{c_A(n)}{2^n}.
\]
Expanding each denominator as a geometric series and interchanging
nonnegative sums gives the identity.  Since $c_A(n)\le\tau(n)\le n$,
Theorem~\ref{record:257bm-i-t7} applies.  Together with
Theorem~\ref{record:257bm-i-mob}, it expresses rationality through an
integer recurrence whose coefficients all come from the same support.
\par\noindent\ev{Lean} \scale{uniform} \coord{divisor counts and finite sums} \lean{erdosSupportSeries\_two\_eq\_binaryCoeffSeries}{BooleanMobiusCarry.lean:377} \lean{erdosSupportSeries\_rational\_iff\_exists\_temperedCarry}{BooleanMobiusCarry.lean:384}
\end{thm}
\leannote{Lean: \lproof{Erdos249257/BooleanMobiusCarry.lean}{377}{erdosSupportSeries_two_eq_binaryCoeffSeries}, \lproof{Erdos249257/CertificateKernel.lean}{8860}{supportCoeff_le_card_divisors}, \lproof{Erdos249257/CertificateKernel.lean}{8868}{supportCoeff_le_self}, \lproof{Erdos249257/BooleanMobiusCarry.lean}{384}{erdosSupportSeries_rational_iff_exists_temperedCarry}.}

\subsubsection*{Restrictions on a support with rational value}

The next results impose separate restrictions on rational-valued supports.
Their sources are \texttt{RationalSupportCarrySkeleton.lean},
\texttt{SublogDivisorCoverage.lean} and \texttt{MaximalOmegaLayer.lean}.
The related strict-gap result is Proposition~\ref{record:257bm-i10}.

\begin{thm}[A restriction on dyadic rational values]\label{record:257rig-i2}
If an infinite support $A\subseteq\Npos$ has $X_A(2)=p/2^c$ for
integers $p$ and $c\ge0$, then $\sum_{a\in A}1/a$ either diverges or
converges to a value greater than $1$. The cited proof averages the
shifted integer recurrence and uses a common multiple of two distinct
support elements. Thus a convergent reciprocal sum of at most $1$ is
excluded. This is a separate necessary condition: the reciprocal-summable
criterion proved earlier already excludes every support with a convergent
reciprocal sum, not just those whose sum is at most $1$.
\par\noindent\ev{Lean} \scale{uniform} \coord{cesaro-tail} \lean{dyadic\_support\_fraction\_reciprocalMass\_diverges\_or\_gt\_one}{RationalSupportCarrySkeleton.lean:2210} \lean{one\_lt\_reciprocalMass\_of\_dyadic\_support\_fraction\_of\_two\_pos\_mem}{RationalSupportCarrySkeleton.lean:2124}
\end{thm}

\begin{thm}[Unboundedness of a positive shifted recurrence]\label{record:257rig-i3}
Let $A\subseteq\Npos$ be infinite, fix a shift $c\ge0$ and a positive
integer $v$, and let $u$ be a positive integer sequence satisfying
\[
 u(n+1)+v c_A(c+n+1)=2u(n)\qquad(n\ge0).
\]
Then $u$ is unbounded. Indeed, under a proposed bound $u(n)\le B$,
choose $2B+1$ distinct support elements and a common multiple larger
than $c$. At the corresponding index the divisor count is at least
$2B+1$, contradicting the recurrence. This does not conflict with
Theorem~\ref{record:257bm-i-t7}, whose condition is $u(n)/2^n\to0$,
not boundedness. An unbounded sequence can still satisfy that limit.
\par\noindent\ev{Lean} \scale{uniform} \coord{common-multiple-forcing} \lean{shifted\_state\_unbounded\_of\_infinite\_support}{RationalSupportCarrySkeleton.lean:2327} \lean{exists\_unbounded\_shifted\_odd\_tail\_nat\_state\_of\_support\_fraction}{RationalSupportCarrySkeleton.lean:2383} \lean{one\_add\_mul\_card\_le\_two\_mul\_shifted\_state}{RationalSupportCarrySkeleton.lean:2237}
\end{thm}

\begin{thm}[Bounds for intervals with zero divisor counts]\label{record:257rig-i4a}
Suppose that $X_A(2)=p/(2^cv)$ for an infinite positive support $A$
and an odd positive integer $v$. For every $\varepsilon>0$ there is
$B$ such that, for $N\ge1$, a run of $h$ zero divisor counts starting
after $c+N$ satisfies $h\le\varepsilon\log_2N+B$.
In fact this conclusion holds without rationality: fix any $a\in A$.
Every $a$ consecutive positive integers include a multiple of $a$, where
$c_A$ is positive. Hence $h\le a-1$, so one can take $B=a-1$.
The zero-run conclusion therefore imposes no additional restriction on
a fixed nonempty support, regardless of the tail estimates used in the
linked proof.
\par\noindent\ev{Lean} \scale{uniform} \coord{divisor-envelope} \lean{supportCoeffZeroWindow\_length\_le\_eps\_logb}{SublogDivisorCoverage.lean:435}
\end{thm}
\leannote{Lean: \lproof{ErdosProblems/Erdos257/PaperCompleteR21/SkipSafetyAndDivisorZeroRuns.lean}{219}{paper_zero_run_le_eps_logb}, \lproof{ErdosProblems/Erdos257/PaperCompleteR21/SkipSafetyAndDivisorZeroRuns.lean}{237}{paper_zero_run_le_of_mem}.}

\begin{prop}[A subpower bound for divisor counts]\label{record:257rig-i4b}
For positive integers $n,k$,
$\tau(n)^k \le (k^{2^k})^k\cdot n$, or equivalently
$\tau(n)\le k^{2^k}n^{1/k}$. The constant absorbs the finitely many primes
below $2^k$; for larger primes, $(\nu+1)^k\le p^\nu$ controls each factor
of the divisor product. This estimate is independent of the support
problem. The elementary zero-run bound above already follows from a
single positive support element, without this estimate or a recurrence.
\par\noindent\ev{Lean} \scale{uniform} \coord{divisor-envelope} \lean{card\_divisors\_pow\_le\_divisorSubpowerConst\_pow\_mul}{SublogDivisorCoverage.lean:107}
\end{prop}
\leannote{Lean: \lproof{Erdos249257/SublogDivisorCoverage.lean}{107}{card_divisors_pow_le_divisorSubpowerConst_pow_mul}, \lproof{Erdos249257/SublogDivisorCoverage.lean}{142}{card_divisors_le_divisorSubpowerConst_mul_rpow}.}

\begin{prop}[Two distinct prime-power differences commute]\label{record:257rig-i5}
Lemma~\ref{lem:mixed-prime-power-layer} gives the four-term expansion
and its divisor-count interpretation.  Commutation is elementary
for arbitrary positive multipliers.  The extraction formula uses
distinct primes, \emph{positive} exponents $e,f$ and $\gcd(n,pq)=1$.
The example $A=\{12\}$ shows the extracted coefficient explicitly.
No rationality statement about multiplicative subsequences follows
from commutation alone.
\par\noindent\ev{Lean} \scale{uniform} \coord{mobius-inversion} \lean{mixedPrimePowerLayerTwo\_supportCoeffInt}{MaximalOmegaLayer.lean:39} \lean{primePowerLayer\_comm}{MaximalOmegaLayer.lean:29}
\end{prop}
\leannote{Lean: \lproof{Erdos249257/MaximalOmegaLayer.lean}{29}{primePowerLayer_comm}, \lproof{Erdos249257/MaximalOmegaLayer.lean}{39}{mixedPrimePowerLayerTwo_supportCoeffInt}, \lproof{Erdos249257/MaximalOmegaLayer.lean}{65}{mixedPrimePowerLayerTwo_twelve_fixture}.}

\subsubsection*{The first crossing of the half-value}

\begin{defn}[The first crossing rank]\label{record:257bm-i-cross}
\[
\mathrm{localMersenneCrossingRanks}(E) \;:=\; E.\mathrm{filter}\big(c \mapsto \tfrac12 < \ensuremath{V}(E.\mathrm{filter}(\cdot\le c))\big),
\]
the ranks at which $E$'s inclusive ordered prefix is already above one half.
\par\noindent\ev{Lean} \scale{n/a} \coord{binary digits} \lean{localMersenneCrossingRanks}{BooleanMobiusExactRowCrossing.lean:24}
\end{defn}

\begin{thm}[Existence of a first crossing]\label{record:257bm-i-cross2}
Let $E\subseteq\{2,3,\ldots\}$ be finite and suppose $X_E(2)>1/2$.
There is a least $c\in E$ such that
\[
 X_{E\cap[2,c)}(2)<\frac12
       <X_{E\cap[2,c]}(2),\qquad c\ge4.
\]
Positivity of the summands makes the partial sums increasing.  Their
odd reduced denominators exclude equality with $1/2$, so the first
crossing is strict on both sides.  Finally
$w_2+w_3=10/21<1/2$, which excludes a crossing before rank $4$.
This is the finite crossing data used in the later exact-row
constructions; the lower bound on $c$ and strictness are specific
to these weights.
\par\noindent\ev{Lean} \scale{uniform} \coord{binary digits} \lean{exists\_first\_localMersenne\_crossing}{BooleanMobiusExactRowCrossing.lean:31}
\end{thm}

\subsubsection*{Centred carries at skipped ranks}

\begin{thm}[Nonnegative centred carries below the half-value]\label{record:257hg-i6}
If $1\notin A$ and $X_A(2)<1/2$, then
\[
 C_A(N)=\operatorname{ihc}(A,N)-1\ge0\qquad(N\ge0).
\]
Indeed, Lemma~\ref{lem:collapse-mech} gives
\[
 \operatorname{ihc}(A,N)
 =2^{N+1}\bigl(1/2-X_A(2)\bigr)
     +\sum_{r\ge1}c_A(N+1+r)2^{-r}>0.
\]
The first term is strictly positive and the tail is nonnegative.
Since the half-carry is an integer, it is at least $1$; subtracting
$1$ proves the centred bound.  Nonnegativity alone, without this
strictness and integrality step, would only give $C_A(N)\ge-1$.
The argument supplies the lower bound used with the conditional
upper bound in Theorem~\ref{record:257rig-c16}.
\par\noindent\ev{Lean} \scale{uniform} \coord{mobius-centred-carry} \lean{integerHalfCarry\_eq\_scaled\_residual\_add\_tail}{HalfCarryReachability.lean:871} \lean{mobiusCenteredHalfCarry\_nonneg\_of\_supportSeries\_lt\_half}{HalfCylinderFinalMiddleCellEscape.lean:94}
\end{thm}

\begin{thm}[Three possibilities at a skipped endpoint]\label{record:257hg-i7}
For a rank $s\ge5$ omitted by the real greedy support, let $H_s$
and $f_s$ be as in Lemma~\ref{lem:skipped-endpoint-trichotomy}.
That lemma identifies the actual prefix as $D_s$ when $f_s\le0$
and as $B_s$ when $f_s>0$, with the exact remainder or overshoot
in each case.  To obtain the nonnegative-margin condition in
Definition~\ref{record:257bm-c13}, the negative case would still have
to be excluded at every required skipped rank.  The trichotomy
itself does not exclude it.  In particular, the value $-3$
excluded for $C_D$ under the all-right-tail hypothesis of
Theorem~\ref{thm:final-middle-cell} is not an exclusion for $f_s$:
the coordinates and the hypotheses are different.
\par\noindent\ev{Lean} \scale{uniform} \coord{frozen-margin} \lean{halfGreedy\_skipped\_endpoint\_trichotomy}{HalfCylinderSkippedEndpointClassifier.lean:246}
\end{thm}

\subsubsection*{Gap lengths and the measure of their union}

\begin{defn}[The gap after a Mersenne weight]\label{record:257hg-i1}
For $n\ge1$, let $R_n=\sum_{j>n}w_j$ and $g_n=w_n-R_n>0$.
The quantity $g_n$ is the length of each gap created at level $n$, not
the sum of all gap lengths created at that level.
\par\noindent\ev{Lean} \scale{n/a} \coord{divisor counts and finite sums} \lean{mersenneGap}{GreedyAchievementSet.lean:2337} \lean{mersenneGap\_pos}{GreedyAchievementSet.lean:2346}
\end{defn}

\begin{thm}[The unweighted sum of remaining gap lengths]\label{record:257hg-i2}
For every $N\ge0$, the unweighted sum of gap lengths satisfies
\[
 \sum_{n>N}g_n\le\frac29\,4^{-N}+\frac37\,8^{-N},
\]
which tends to zero.  The estimate follows by summing the per-level
bound; it counts one length per level.  It does not, by itself, bound
the measure of the union of all gaps, because there are $2^{n-1}$
disjoint translated gaps at level $n$.
\par\noindent\ev{Lean} \scale{uniform} \coord{summed-gap-mass} \lean{summable\_mersenneGap\_shift}{HalfGapMass.lean:69} \lean{mersenneGap\_tail\_le}{HalfGapMass.lean:83} \lean{tendsto\_mersenneGap\_tail\_zero}{HalfGapMass.lean:104} \lean{mersenneGap\_le}{HalfGapMass.lean:40}
\end{thm}

The measure of the union of gaps created after level $N$ is instead
\[
 \sum_{n>N}2^{n-1}g_n=2^N R_N-1
 \le\frac23\,2^{-N}\qquad(N\ge1).
\]
Indeed, the finite sum from $N+1$ to $M$ telescopes to
$2^N R_N-2^M R_M$, and $2^M R_M\to1$.  The final bound uses
$R_N\le2^{-N}+(2/3)4^{-N}$.  Disjointness follows from the strict-tail
construction of the achievement set.  This is an ordinary geometric
deduction, not the statement of the preceding linked gap-sum lemma.
A small union measure still gives no answer about a prescribed point.
\ev{Math}

\subsubsection*{A sufficient test for a safe skip}

\begin{defn}[Three geometric terms in the tail lower bound]\label{record:257hg-i3}
$\mathrm{mersenneTailLB3}(k) := 1/2^k + 1/(3\cdot4^k) + 1/(7\cdot8^k) < \ensuremath{R}(k)$
;  a strict rational lower bound obtained by truncating the Lambert-type series at three terms.
\par\noindent\ev{Lean} \scale{n/a} \coord{lambert-bound} \lean{mersenneTailLB3}{HalfGreedyFatalGap.lean:56} \lean{three\_div\_lt\_mersenneTailLB3}{HalfGreedyFatalGap.lean:70} \lean{mersenneTailLB3\_lt\_mersenneTail}{HalfGreedyFatalGap.lean:216}
\end{defn}

\begin{thm}[A sufficient inequality for a safe skip]\label{record:257hg-i4}
Let $k,u,L\ge1$ be integers, set $a=2L-(2^k-1)u$, and
suppose $a>0$.  For the skipped rational remainder $\rho=u/(2L)$,
$\rho\le2^{-k}$ is equivalent to $u\le a$, whereas
\[
 2u\le3a\quad\Longrightarrow\quad\rho<R_k.
\]
The latter implication uses the three-term lower bound above.  More
explicitly, with $t=2^k\ge2$,
\[
 \left(\frac1t+\frac1{3t^2}+\frac1{7t^3}\right)
       -\frac3{3t-1}
 =\frac{2t-3}{21t^3(3t-1)}>0.
\]
The inequality $2u\le3a$ gives $\rho\le3/(3t-1)$, proving the
claim.  The first two terms alone would not prove this comparison.
This is a sufficient exclusion of the current tail-mass deficit, not
an exact test for membership in the remaining achievement set.

The sufficient condition is weaker than $u\le a$.  A realizable example
is $(k,u,L,a)=(2,7,13,5)$; it passes $2u\le3a$ but not $u\le a$.
The linked scalar example $(u,a)=(3,2)$ proves strict containment of the
inequalities, but is not integral rational data for this substitution.
For comparison, the exact mass threshold is
$a/u\ge R_k^{-1}-(2^k-1)$; its right-hand side lies strictly between
$0$ and $2/3$.  None of these comparisons asserts that the half-greedy
orbit satisfies the sufficient inequality at every skipped rank.
\par\noindent\ev{Lean} \scale{uniform} \coord{lambert-bound} \lean{skipSafe\_of\_two\_mul\_le\_three\_mul}{HalfGreedyFatalGap.lean:107} \lean{sharp\_of\_dyadic}{HalfGreedyFatalGap.lean:197} \lean{sharp\_strictly\_stronger}{HalfGreedyFatalGap.lean:200} \lean{skipSafe\_actualTail\_of\_two\_mul\_le\_three\_mul}{HalfGreedyFatalGap.lean:235}
\end{thm}
\leannote{Lean: \lproof{ErdosProblems/Erdos257/PaperCompleteR21/SkipSafetyAndDivisorZeroRuns.lean}{30}{paper_dyadic_skip_test_iff}, \lproof{ErdosProblems/Erdos257/PaperCompleteR21/SkipSafetyAndDivisorZeroRuns.lean}{54}{paper_sharp_skip_safe_lb3}, \lproof{ErdosProblems/Erdos257/PaperCompleteR21/SkipSafetyAndDivisorZeroRuns.lean}{62}{paper_sharp_skip_safe_actual_tail}, \lproof{ErdosProblems/Erdos257/PaperCompleteR21/SkipSafetyAndDivisorZeroRuns.lean}{71}{paper_three_channel_margin_identity}, and 5 further declarations in the \hyperref[sec:coverage]{coverage section}.}

\begin{prop}[Two unconditional safety cases]\label{record:257hg-i5}
Under the positive-integer and skipped-step hypotheses of
Theorem~\ref{record:257hg-i4}, $u=1$ implies $a\ge1$, so the current
remainder is less than $R_k$.  Conversely, a fatal tail-mass deficit
forces $3a<2u$, hence $u\ge2$, or $u\ge3$ when $u$ is odd.
These are statements about the current step, not infinite survival. The exact declaration names are linked
below.
\par\noindent\ev{Lean} \scale{uniform} \coord{lambert-bound} \lean{unitNumerator\_skipSafe}{HalfGreedyFatalGap.lean:135} \lean{two\_le\_of\_fatal}{HalfGreedyFatalGap.lean:161} \lean{three\_le\_of\_fatal\_of\_odd}{HalfGreedyFatalGap.lean:173} \lean{unitNumerator\_skipSafe\_actualTail}{HalfGreedyFatalGap.lean:246}
\end{prop}

\subsubsection*{Two-sided bounds near dyadic endpoints}

\begin{thm}[Two-sided dyadic bounds]\label{record:257bm-i13}
Under the local hypothesis excluding the specified three middle cells and imposing the right-pulse
bound, induction (base case at row 5 verified by \texttt{decide}) propagates it to the universal
two-sided bound
\[
\forall s\ge5,\quad \min(\mathrm{rem}(s),\ \mathrm{overshoot}(s)) \;\le\; 2^s.
\]
This conditional bound says that at least one of the remainder and
the adjacent overshoot is at most $2^s$.  It does not assert that both
are at most $2^s$; that stronger assertion would replace the minimum
by a maximum.  It is not needed for the charge bound at an actual upper reset:
Proposition~\ref{record:257bm-c12} obtains that bound directly from the
upper-reset identity.  Applying the inductive argument to another
recurrence would require its own transition and separation estimates.
\par\noindent\ev{Lean} \scale{uniform} \coord{dyadic-scale} \lean{SeamTwoSidedDyadicCellEscape}{HalfCylinderMiddleCarryLowerBound.lean:4374} \lean{SeamTwoSidedDyadicCellEscape.step}{HalfCylinderMiddleCarryLowerBound.lean:4398} \lean{SeamTwoSidedDyadicCellEscape.twoSided}{HalfCylinderMiddleCarryLowerBound.lean:4448}
\end{thm}
\leannote{Lean: \lproof{ErdosProblems/Erdos257/PaperCompleteR21/DyadicBandAndTwoSidedBounds.lean}{105}{paper_two_sided_dyadic_bound}.}

\begin{prop}[A finite band check for $13\le d\le30$]\label{record:257bm-i14}
For every actual upper-reset index $13\le d\le30$ and every
$0\le j\le d$, the linked certificate verifies
\[
2^{d-j+1}<\mathrm{resetCharge}
\quad\text{or}\quad
\mathrm{resetCharge}+2(d+j)\le2^{d-j+1}.
\]
It computes the successor remainder at each row $d+1$, from
$\mathrm{rem}(14)=392$ through $\mathrm{rem}(31)=4187487147$.
This proves the required condition on a finite
range, not the universal condition for all $d\ge13$. A proof for
unbounded reset indices must supply a further argument; repeatedly
extending the table alone does not establish that statement.
\par\noindent\ev{Cert} \scale{fixed} \coord{dyadic-boundary} \lean{SeamUpperResetDyadicBandEscape}{HalfCylinderMiddleCarryLowerBound.lean:4566} \lean{seamUpperResetDyadicBandEscape\_through\_thirty}{HalfCylinderUpperResetBandCertificates.lean:78} \lean{half\_mem\_mersenneAchievementSet\_of\_upperResetDyadicBandEscape}{HalfCylinderMiddleCarryLowerBound.lean:4790}
\end{prop}
\leannote{Lean: \lproof{ErdosProblems/Erdos257/PaperCompleteR21/UpperResetBandCertificate.lean}{29}{paper_finite_band_check_thirteen_to_thirty}, \lproof{ErdosProblems/Erdos257/PaperCompleteR21/UpperResetBandCertificate.lean}{44}{paper_successor_remainders_fourteen_through_thirtyone}, \lproof{ErdosProblems/Erdos257/PaperCompleteR21/UpperResetBandCertificate.lean}{77}{paper_universal_band_condition_unfolded}, \lproof{ErdosProblems/Erdos257/PaperCompleteR21/UpperResetBandCertificate.lean}{70}{paper_universal_band_condition_would_close_half}.}

\subsubsection*{General identities for perturbed greedy recurrences}

\begin{thm}[A general perturbed greedy recurrence]\label{record:257bm-i15}
Let the old integer values be separated by at least $g\ge1$,
and let the update be $t(x)=4s(x)+p(x)$ with
$0\le p(x)\le B<g$.  Suppose $x_-$ and $x_+$ are the adjacent
old values on either side of capacity $C$.
Theorem~\ref{thm:perturbed-family-maximality} gives the largest
admissible updated value at capacity $4C+g$ and the exact three-branch
remainder after testing the extra weight $2g+4$.

The underlying structure in
Definition~\ref{defn:perturbed-family} assumes only $B<3g$;
that suffices for order preservation, but the cited maximality
application has the additional hypothesis $B<g$.
The counterexample after Theorem~\ref{thm:perturbed-family-maximality}
shows what can fail without it.  For a maximum over both choices of the extra weight, rather than
the stated two-stage rule, one also needs a separation argument;
$4g-B\ge2g+4$ suffices, as shown after that theorem.
\par\noindent\ev{Lean} \scale{uniform} \coord{abstract-perturbed-greedy} \lean{PerturbedFamily}{HalfCylinderIntegerGreedy.lean:1286} \lean{prefixChoice\_maximal}{HalfCylinderIntegerGreedy.lean:1390} \lean{nextRemainder\_trichotomy}{HalfCylinderIntegerGreedy.lean:1463}
\end{thm}
\leannote{Lean: \lproof{ErdosProblems/Erdos257/PaperCompleteR21/UpperResetBandCertificate.lean}{97}{paper_perturbed_order_preservation}, \lproof{ErdosProblems/Erdos257/PaperCompleteR21/UpperResetBandCertificate.lean}{106}{paper_perturbed_prefixChoice_maximal}, \lproof{ErdosProblems/Erdos257/PaperCompleteR21/UpperResetBandCertificate.lean}{124}{paper_perturbed_nextRemainder_three_branches}, \lproof{ErdosProblems/Erdos257/PaperCompleteR21/UpperResetBandCertificate.lean}{140}{paper_perturbed_separation_global_maximality}, and 3 further declarations in the \hyperref[sec:coverage]{coverage section}.}

\begin{thm}[Spacing of compatible reverse-carry words]\label{record:257bm-i16}
Consider two integer carry recurrences
$b_i(m)+2u_i(m)=a_i(m)+u_i(m+1)$, $i=1,2$.
If $a_1(k)=a_2(k)$ and the output bits at $k$ are $1$ and $0$,
then the carry difference at $k+1$ is odd.  If the coefficients
and bits agree for the following $L$ positions, the difference
at $k+L+1$ is $2^L$ times that odd integer.
Consequently absolute terminal bounds $B_1,B_2$ give
$2^L\le B_1+B_2$.
Lemma~\ref{lem:reverse-carry-word} states the exact identity and
proves it by subtraction.  The linked formal statements encode
these agreements; they do not supply terminal bounds for an
unrelated digit system.
\par\noindent\ev{Lean} \scale{uniform} \coord{reverse-carry-word} \lean{ReverseCarryWord}{HalfTrappingReturnCarry.lean:70} \lean{overlappingReverseCarryWords\_carryDifference\_eq\_twoPow\_mul\_odd}{HalfTrappingReturnCarry.lean:124} \lean{overlappingReverseCarryWords\_twoPow\_le\_realBound}{HalfTrappingReturnCarry.lean:191} \lean{overlappingMidpointReturns\_twoPow\_le\_realBound}{HalfTrappingReturnCarry.lean:244}
\end{thm}
\leannote{Lean: \lproof{Erdos249257/HalfTrappingReturnCarry.lean}{124}{overlappingReverseCarryWords_carryDifference_eq_twoPow_mul_odd}, \lproof{Erdos249257/HalfTrappingReturnCarry.lean}{191}{overlappingReverseCarryWords_twoPow_le_realBound}, \lproof{Erdos249257/HalfTrappingReturnCarry.lean}{244}{overlappingMidpointReturns_twoPow_le_realBound}, \lproof{ErdosProblems/Erdos257/PaperCompleteR21/SeamRowGapAndCarry.lean}{240}{paper_reverse_carry_word}.}

\begin{prop}[Changing one support bit at a doubled rank]\label{record:257bm-i17}
Let two supports agree except that the second includes $N+1$
while the first does not.  At the argument $2(N+1)$ their divisor
counts differ by exactly $1$, because $N+1$ divides $2(N+1)$.
More generally the difference at a positive argument $m$ is
$\mathbf1_{N+1\mid m}$.  The linked statements apply this identity
to the finite supports in their hypotheses.  No conclusion about
a different coefficient sequence follows without identifying its
own support change.
\par\noindent\ev{Lean} \scale{uniform} \coord{half-divisor} \lean{supportCoeff\_extend\_true\_eq\_false\_add\_one\_at\_double}{HalfDivisorUnitDrop.lean:20} \lean{supportCoeff\_boundaryPair\_unitDrop\_at\_double}{HalfDivisorUnitDrop.lean:35}
\end{prop}
\leannote{Lean: \lproof{Erdos249257/HalfCylinderFiniteShadow.lean}{642}{supportCoeff_insert_divisor}, \lproof{Erdos249257/HalfDivisorUnitDrop.lean}{20}{supportCoeff_extend_true_eq_false_add_one_at_double}, \lproof{Erdos249257/HalfDivisorUnitDrop.lean}{35}{supportCoeff_boundaryPair_unitDrop_at_double}.}

\subsection{Obstructions and countermodels}

The results below test particular proposed implications.  A countermodel
to an abstract transition rule shows that the rule needs an additional
hypothesis; it does not disprove a more structured arithmetic
argument.  Each entry states the information retained by its model
and the conclusion that this information fails to imply.

\begin{thm}[A bounded model of the doubling-or-return alternative]\label{record:257bm-k1}
The dichotomy $\mathrm{ExactLocalMersenneHalfRow}(2n{-}1) \vee \exists c,\,4\le c\le n\wedge
\mathrm{ExactLocalMersenneHalfRow}(2c{-}2)$ (proved for $n\ge6$ at Theorem~\ref{record:257bm-k-dich} below) is not by
itself enough for cofinality. Countermodel: $\mathrm{boundedDoubleOrRecycleModel}(n) := (n=6)$
satisfies exactly the same two-branch transition shape (seed at 6, and
\texttt{boundedDoubleOrRecycleModel\_transition} reproduces the $\vee$ shape by always taking
the recycle branch with $c=4$, conclusion back at $2\cdot4-2=6$), yet
\[
\neg\big(\forall N,\ \exists n\ge N,\ \mathrm{boundedDoubleOrRecycleModel}(n)\big) \qquad \text{(not cofinal ;  only ever true at } n=6\text{)}.
\]
Packaged existentially as \texttt{exists\_seeded\_bounded\_double\_or\_recycle\_model}. The
double-or-recycle transition shape, even together with an endpoint-six seed, is logically
insufficient to conclude cofinal exact rows. A genuinely new progress input (strict endpoint
growth, or an independent cofinality argument) is required. The countermodel retains only the seed and this transition rule.  It
therefore tests those premises, not additional arithmetic information
that a more structured induction may use.
\par\noindent\ev{Lean} \scale{fixed} \coord{meta-logical} \lean{boundedDoubleOrRecycleModel}{BooleanMobiusExactRowDichotomy.lean:58} \lean{boundedDoubleOrRecycleModel\_transition}{BooleanMobiusExactRowDichotomy.lean:64} \lean{boundedDoubleOrRecycleModel\_not\_cofinal}{BooleanMobiusExactRowDichotomy.lean:79} \lean{exists\_seeded\_bounded\_double\_or\_recycle\_model}{BooleanMobiusExactRowDichotomy.lean:91}
\end{thm}
\leannote{Lean: \lproof{ErdosProblems/Erdos257/PaperCompleteR21/ExactRowDichotomyCountermodels.lean}{84}{paper_bounded_double_or_recycle_countermodel}, \lproof{ErdosProblems/Erdos257/PaperCompleteR21/ExactRowDichotomyCountermodels.lean}{102}{paper_exists_seeded_bounded_double_or_recycle_model}.}

\begin{thm}[Doubling or returning to an earlier depth]\label{record:257bm-k-dich}
For $n\ge6$, $\mathrm{ExactLocalMersenneHalfRow}(n)$:
\[
\mathrm{ExactLocalMersenneHalfRow}(2n-1) \;\vee\; \exists c,\ 4\le c\le n \wedge \mathrm{ExactLocalMersenneHalfRow}(2c-2).
\]
Every exact row of depth at least $6$ has at least one of the two
asserted continuations.  The proof separates the below-half case from
the above-half case, but the two existential conclusions need not be
exclusive.  The returned depth $2c-2$ is not claimed to exceed $n$. The finite Mersenne value cannot equal $1/2$, by
Observation~\ref{record:257bm-k11}
(\texttt{finiteErdosSum\_den\_odd}); this says nothing about equality
of the returned and original depths.  For example, at $n=6$, the
support $\{2,3,6\}$ is an exact row, and so is
$\{2,3,6,7,11\}$ at depth $11$.  The second conclusion also holds
with $c=4$, since $2c-2=6$.  Thus both conclusions hold in this example,
and the second allows no depth increase.  What is
missing, per Theorem~\ref{record:257bm-k1}, is not another dichotomy but strict endpoint progress in the recycle
branch, or a proof the below-half branch recurs ;  re-deriving this disjunction without an additional progress argument
would not establish cofinality.
\par\noindent\ev{Lean} \scale{bounded} \coord{divisor counts and finite sums} \lean{exactLocalMersenneHalfRow\_double\_or\_recycle}{BooleanMobiusExactRowDichotomy.lean:26}
\end{thm}
\leannote{Lean: \lproof{ErdosProblems/Erdos257/PaperCompleteR21/ExactRowDichotomyCountermodels.lean}{33}{paper_exact_row_double_or_recycle}, \lproof{ErdosProblems/Erdos257/PaperCompleteR21/ExactRowDichotomyCountermodels.lean}{41}{paper_finite_row_value_ne_half}, \lproof{ErdosProblems/Erdos257/PaperCompleteR21/ExactRowDichotomyCountermodels.lean}{55}{paper_exact_row_example_six_and_eleven}.}

\begin{prop}[A sufficient fractional-mass bound need not hold]\label{record:257bm-k2}
Take $D=\{2,3\}$ and $c=5$.  Direct calculation gives
\[
 X_D(2)=\frac{10}{21}<\frac12
 <\frac{331}{651}=X_{D\cup\{5\}}(2),\qquad S(D,1,8)=6<8.
\]
However, the combined fractional mass is
\[
 \sum_{d\in D\cup\{5\}}\frac{2^8\bmod(2^d-1)}{2^d-1}
 =\frac{757}{651}>1.
\]
The linked \texttt{norm\_num} proof checks this example.  Thus the
combined fractional-mass bound by $1$, which suffices for the sharp
capacity estimate, is not necessary for that estimate.
Thus the fractional-mass bound cannot hold at every real crossing core.
A proof using it would need to handle the exceptional cores separately;
the example does not rule out all uses of fractional-mass estimates.
Theorem~\ref{record:257rig-k6} identifies the support at a critical
crossing with a real greedy prefix, but does not prove the missing bound
on those prefixes.
\par\noindent\ev{Cert} \scale{fixed} \coord{divisor counts and finite sums} \lean{splitFractionMass\_one\_bound\_not\_necessary\_fixture}{BooleanMobiusSkippedCoreCriticalCapacity.lean:183}
\end{prop}
\leannote{Lean: \lproof{ErdosProblems/Erdos257/PaperCompleteR21/ExactRowDichotomyCountermodels.lean}{118}{paper_fractional_mass_bound_not_necessary}, \lproof{ErdosProblems/Erdos257/PaperCompleteR21/ExactRowDichotomyCountermodels.lean}{139}{paper_fractional_mass_bound_suffices_for_sharp_capacity}.}

\begin{obs}[A doubled finite sum can exceed one half]\label{record:257bm-k3}
Theorem~\ref{record:257bm-i7}'s doubling extension is not preserved iteratively: below-halfness is not preserved
under one application of doubling, and demonstrably fails at the very first step. From the
identity $\mathrm{value}(E)-\tfrac12 = \big(\mathrm{localFractionMass}(E,2n{-}1)-1\big)/2^{2n-1}$,
at the seed $D=\{2,3,6\}$ (Definition~\ref{record:257bm-i-seed}) with $M=11$ the core alone contributes
\[
\frac{2^{11}\bmod3}{3} + \frac{2^{11}\bmod7}{7} + \frac{2^{11}\bmod63}{63} = \frac23+\frac47+\frac{32}{63} \approx 1.746 > 1,
\]
so the doubled row is strictly above half. Iterating the cheap (doubling) arm is therefore
impossible from the only recorded seed; the chain is forced into the capacity-gated recycle arm
at step one. What would close the doubling route instead: cofinally many BELOW-HALF exact rows,
$\forall N\ \exists n\ge N\ \exists D$ exact at $n$ with $\mathrm{localFractionMass}(D,n)<1$ ;  a
single such row at each of cofinally many $n$ would make doubling unnecessary; a
self-reproducing one would close everything.
\par\noindent\ev{Cert} \scale{fixed} \coord{divisor counts and finite sums} \lean{exists\_exactRowStrictUpperExtension\_two\_mul\_sub\_one\_of\_exact\_below}{BooleanMobiusExactRowDoubling.lean:185}
\end{obs}

\begin{thm}[The returning endpoint need not be larger]\label{record:257bm-k4}
Theorem~\ref{record:257bm-c10} (\texttt{exists\_skippedCoreExactRow\_of\_value\_above}) is unconditional, but its
witness is $\exists c\le n$, not $\exists c$ large: $2c-2$ may be $\le n$, so the endpoint need
not grow. Theorem~\ref{record:257bm-k1}'s \texttt{exists\_seeded\_bounded\_double\_or\_recycle\_model} is the
explicit falsifier of the naive hope that growth comes for free:
$\mathrm{boundedDoubleOrRecycleModel}(n):=(n=6)$ satisfies the same transition schema plus a
seed and is not cofinal. Growth is recovered only inside
$\mathrm{ProtectedExactLocalMersenneRow}$ (Theorem~\ref{record:257bm-c5}), whose invariants
$\mathrm{endpoint}<2\cdot\mathrm{cutoff}$ and $\mathrm{new\_above\_cutoff}$ force $c>\mathrm{cutoff}$
hence $2c-2>\mathrm{endpoint}$ ;  and maintaining those invariants is precisely what needs the
strict-upper (sharp capacity) fill of Theorem~\ref{record:257bm-c7} rather than this general recycling theorem.
\par\noindent\ev{Lean} \scale{bounded} \coord{binary digits} \lean{exists\_skippedCoreExactRow\_of\_value\_above}{BooleanMobiusExactRowCrossing.lean:155}
\end{thm}
\leannote{Lean: \lproof{ErdosProblems/Erdos257/PaperCompleteR21/ExactRowDichotomyCountermodels.lean}{163}{paper_skipped_core_recycling_witness_bounded}, \lproof{ErdosProblems/Erdos257/PaperCompleteR21/ExactRowDichotomyCountermodels.lean}{172}{paper_returning_endpoint_may_fail_to_grow}, \lproof{ErdosProblems/Erdos257/PaperCompleteR21/ExactRowDichotomyCountermodels.lean}{102}{paper_exists_seeded_bounded_double_or_recycle_model}, \lproof{ErdosProblems/Erdos257/PaperCompleteR21/ExactRowDichotomyCountermodels.lean}{184}{paper_protected_row_crossing_beyond_cutoff}, and 1 further declaration in the \hyperref[sec:coverage]{coverage section}.}

\begin{obs}[The skip condition is equivalent to membership]\label{record:257bm-k5}
The positivity conjunct of Theorem~\ref{record:257bm-c2}'s hypothesis,
$0<\ensuremath{r}(1/2,c{-}1)$, is unconditionally dischargeable
(\texttt{localMersennePrefixValue\_halfGreedy\_lt\_half}, itself just the odd-denominator parity
fact of Observation~\ref{record:257bm-k11}, together with \texttt{greedyMersenneRemainderRat\_eq\_sub\_finiteErdosSum}). Hence
\[
\mathrm{CofinalPositiveHalfGreedySkips} \;\Longleftrightarrow\; (\mathrm{greedyMersenneSkippedSupport}(1/2)).\mathrm{Infinite},
\]
which Theorem~\ref{record:257bm-c20} (A6) proves equivalent to $1/2\in\ensuremath{\mathcal A}$. Proving Theorem~\ref{record:257bm-c2}'s hypothesis would therefore settle
the half-target question.  This equivalence does not by itself make the
hypothesis easier, but it also does not preclude a useful argument in
greedy coordinates.  Theorem~\ref{record:257bm-c7} offers a different
sufficient condition whose finite supports need not be greedy prefixes.
\par\noindent\ev{Lean} \scale{cofinal} \coord{greedy recurrence} \lean{localMersennePrefixValue\_halfGreedy\_lt\_half}{BooleanMobiusCriticalCapacityCofinal.lean:138}
\end{obs}

\begin{thm}[Uniqueness at a critical crossing]\label{record:257rig-k6}
For $c\ge4$, $D$ bounded $[2,c)$, below-half, with genuine crossing deficit
$\tfrac12-\mathrm{value}(D) < \ensuremath{w}(c)$:
\[
D \;=\; \mathrm{halfGreedyPrefixSupport}(c-1).
\]
Thus $D=G\cap\{1,\ldots,c-1\}$: at a critical crossing the support is
fixed by $c$.  The proof uses the strict inequality between each Mersenne
weight and the sum of all later weights, as in
Theorem~\ref{record:257bm-i11a}.  Consequently, the support quantifier in
Theorem~\ref{record:257bm-c4} does not allow arbitrary choices of $D$.
This reduces that hypothesis to a statement about the actual greedy
sequence.  It does not rule out proving a new bound for that sequence.
\par\noindent\ev{Lean} \scale{uniform} \coord{greedy recurrence} \lean{eq\_halfGreedyPrefixSupport\_of\_critical\_crossing}{BooleanMobiusCriticalCapacityCofinal.lean:50}
\end{thm}
\leannote{Lean: \lproof{ErdosProblems/Erdos257/PaperCompleteR21/ExactRowDichotomyCountermodels.lean}{231}{paper_critical_crossing_support_is_greedy_prefix}.}

\begin{obs}[Different requirements on finite witnesses]\label{record:257rig-k7}
Theorem~\ref{record:257rig-c17} (the cofinal terminal-strip condition) tolerates any support inside the depth-$M$
window but demands the carry inside a $\sim2\sqrt M$ strip; Theorem~\ref{record:257bm-c7} tolerates a carry up to
$2^{c-2}$ (exponentially looser) but demands the support be confined to ranks $<c$, because the
exact fill encodes the residue in the pure upper window where the Mersenne quotient is exactly
$2^{M-d}$ (Proposition~\ref{record:257bm-c9}). The two conditions therefore differ in the permitted support as well as
the size of the carry.  If the square-root witnesses could also be chosen
with support in $[2,c)$ at depth $2c-3$, then
$2\sqrt{2c}+4<2^{c-2}$ for $c\ge8$ would give the numerical bound needed
for the fill.  The required support restriction is not proved here.
Comparing the displayed bounds alone does not order the two existence
hypotheses.
\par\noindent\ev{Lean} \scale{cofinal} \coord{sqrt-carry-growth} \lean{HalfCarryCofinalTerminalOnlyStrip}{TerminalOnlyCofinal.lean:34}
\end{obs}

\begin{obs}[An unsafe example in the two-thirds interval]\label{record:257hg-k8}
Read in full alongside \texttt{HalfGreedyFatalGap} for scope contrast: writing a greedy residual
as $\mathrm{rem}=1/R$, a skip at rank $k$ is dyadically safe iff $R\ge2^k$. Two unconditional
corollaries hold (an integral reciprocal $R$ is never unsafe, \texttt{not\_twoThirdsBand\_of\_int};
an unsafe run with odd $p,D,q$ forces $p\ge7$, \texttt{seven\_le\_of\_intBand\_odd}) ;  but the
file's own docstring is explicit that dyadic safety is sufficient, not necessary, for the true
greedy process to survive (the true test uses the actual tail $\sum_{j>k}\mathrm{weight}(j)$,
which exceeds $2^{-k}$), and nothing in this file asserts the actual half-greedy orbit avoids
the band: $(p,D,b)=(17,41,3)$ is an explicit odd-coprime example that IS dyadically unsafe. This
sharpens directly into Theorem~\ref{record:257hg-i4}'s $2u\le3a$ criterion, which narrows exactly this gap for the
single-skip case ;  but the general question of whether the fatal region overlaps a reachable
band under the actual greedy trajectory is settled by neither file.
\par\noindent\ev{Lean} \scale{n/a} \coord{rational-band} \lean{TwoThirdsBand}{HalfGreedyTwoThirdsBand.lean:120} \lean{not\_twoThirdsBand\_of\_int}{HalfGreedyTwoThirdsBand.lean:185} \lean{seven\_le\_of\_intBand\_odd}{HalfGreedyTwoThirdsBand.lean:231}
\end{obs}

\begin{thm}[Vanishing of the specified linear-channel determinant]\label{record:257bm-k9}
Let $V$ be a vector space over $\Q$, let $e:V\to\Q$ be linear, and
let $(\ell_j)_{j\in\iota}$ be a finite family of linear functionals
vanishing on $\ker e$.  For any vectors $(v_i)_{i\in\iota}$, the matrix
$(\ell_j(v_i))_{i,j\in\iota}$ has rank at most one, so every square
minor of size at least two vanishes.
Indeed, the functionals descend to $V/\ker e$, which has dimension at
most one.  This elementary linear-algebra argument applies at every
matrix size.  It does not cover additional functionals that fail to
vanish on $\ker e$.
\par\noindent\ev{Lean} \scale{uniform} \coord{linear-algebra} \lean{relationInvariantLinearChannels\_det\_eq\_zero}{HalfTrappingReturnCarry.lean:42}
\end{thm}
\leannote{Lean: \lproof{ErdosProblems/Erdos257/PaperCompleteR21/LinearChannelAndMiddleCellExclusion.lean}{37}{paper_relationInvariant_channels_rank_le_one}, \lproof{ErdosProblems/Erdos257/PaperCompleteR21/LinearChannelAndMiddleCellExclusion.lean}{63}{paper_relationInvariant_channels_det_eq_zero}.}

\begin{obs}[The additive error in the binary bound]\label{record:257bm-k10}
Theorem~\ref{record:257bm-i6}'s $2^{c-1}$ bound and Theorem~\ref{record:257bm-i5}'s sharp $2^{c-2}$ requirement differ by one bit, and the
proof of the loose bound shows exactly where: it derives $A<2^{c-2}+|D|$ and discards
$|D|\le c-2\le2^{c-2}$. Under this additive upper bound, the remaining assertion is that the
integer suffix avoids $[2^{c-2},\,2^{c-2}+c-3]$.  This identifies the
unresolved interval; it does not prescribe how to exclude it.  A sharper
estimate, an arithmetic restriction on the suffix, or another argument
could in principle establish the required inequality.
\par\noindent\ev{Lean} \scale{uniform} \coord{binary digits} \lean{localBinarySuffix\_two\_mul\_sub\_two\_lt\_upperHalfCapacity}{BooleanMobiusSkippedCoreExactRow.lean:123}
\end{obs}

\begin{obs}[Source status of finite-denominator parity]\label{record:257bm-k11}
For finite $D\subseteq\Npos$, each denominator $2^d-1$ is odd.
The sum $V(D)$ can therefore be written over an odd common denominator;
its reduced denominator, being a divisor of that denominator, is odd.
In particular, $V(D)\ne1/2$.  This elementary argument supplies the
parity fact used in Proposition~\ref{record:257bm-k2} and
Proposition~\ref{record:257bm-i10}.  The linked declaration is retained
for formal-source comparison; no new kernel replay is asserted.
\par\noindent\ev{Cited} \scale{uniform} \coord{p-adic} \lean{finiteErdosSum\_den\_odd}{DyadicPrefixCompression.lean:1513}
\end{obs}

\begin{thm}[Excluding the cell with value minus three]\label{record:257hg-k12}
At a middle row $D\ge13$ followed only by right transitions,
$4\,\mathrm{rem}(D)-p_D^--4\ne-3$, by
Theorem~\ref{thm:final-middle-cell}.  Its proof uses the nonnegative
centred carry for the completed support, not a finite search.
The values $-2,-1$ remain among the three exceptional negative cells.
Excluding them under this extra tail assumption would still not
exclude nonnegative values of the coordinate, or establish the
all-middle-row and right-branch hypotheses in
Theorem~\ref{thm:two-sided-dyadic}.  The complete remaining tail
inequality is stated in Remark~\ref{rem:tail-dominance-open}.
\par\noindent\ev{Lean} \scale{uniform} \coord{mobius-centred-carry} \lean{finalMiddleCell\_neg\_three\_not\_last}{HalfCylinderFinalMiddleCellEscape.lean:587} \lean{mobiusCenteredHalfCarry\_add\_two}{HalfCylinderFinalMiddleCellEscape.lean:39} \lean{cofiniteRightTail\_ne\_zero\_centeredEndpoint}{HalfCylinderFinalMiddleCellEscape.lean:547}
\end{thm}
\leannote{Lean: \lproof{ErdosProblems/Erdos257/PaperCompleteR21/LinearChannelAndMiddleCellExclusion.lean}{100}{paper_final_middle_cell_ne_neg_three}, \lproof{ErdosProblems/Erdos257/PaperCompleteR21/LinearChannelAndMiddleCellExclusion.lean}{118}{paper_final_middle_cell_remaining_negative_values}, \lproof{ErdosProblems/Erdos257/PaperCompleteR21/LinearChannelAndMiddleCellExclusion.lean}{155}{paper_mobiusCenteredHalfCarry_add_two}, \lproof{ErdosProblems/Erdos257/PaperCompleteR21/LinearChannelAndMiddleCellExclusion.lean}{163}{paper_cofiniteRightTail_ne_zero_centeredEndpoint}.}
% ; - part a257_newdecls ; -
\subsection{Further finite and conditional results}

This subsection records further finite identities and conditional
membership results from the supplied source files.  They concern the
middle transition, finite families with a common prefix, and rational
targets other than $1/2$.  Their hypotheses are stated separately from
the existence assertions still needed to apply them.

\subsubsection{The integer update in a middle interval}

The following identities explain the two-sided bound and the reset-band
condition used above.  They also distinguish the signed middle coordinate
$C_s=4\,\mathrm{rem}(s)-p_s^- -4$ from the nonnegative label in the
next estimate.

\begin{prop}[The unsafe middle range is exactly three integers]
For $s\ge5$, write $R=\mathrm{rem}(s)$ and $P=p_s^-$.  Since
$4R-P-4$ is an integer,
\[
\neg(4R-P-4\le-4\ \vee\ 0\le 4R-P-4)\ \Longleftrightarrow\ 4R-P-4\in\{-3,-2,-1\}
\]
\end{prop}
\leannote{Lean: \lproof{ErdosProblems/Erdos257/PaperCompleteR21/LinearChannelAndMiddleCellExclusion.lean}{83}{paper_unsafe_middle_range_is_three_integers}.}
(\lean{seamMiddleCoordinate\_not\_safe\_iff\_three\_cells}{HalfCylinderMiddleCarryLowerBound.lean:4125},
\ev{Lean}, \scale{uniform}, \coord{divisor counts and finite sums}, proved by \texttt{omega} from the integer
definitions). The complement of the two safe integer ranges is precisely
$\{-3,-2,-1\}$.  The conclusion is only a classification; excluding
any of these values requires an additional argument.

For a row $s\ge5$, set
\[
 E_s=\sum_{d\in D_s}\left\{\frac{4^s}{2^d-1}\right\},\qquad
 Z_s=\mathrm{rem}(s)-E_s,\qquad
 c=4\,\mathrm{rem}(s)-p_s^-=C_s+4.
\]
Here braces denote fractional part.  Assume $c\ge0$, as in the linked
statements below.  In particular, the exceptional signed values
$C_s=-3,-2,-1$ correspond to $c=1,2,3$, not to negative labels.
The floor-quotient recurrence gives $p_s^-\le4E_s$, and hence
\[
 4Z_s\le c.
\]
This is \lean{seamMiddleCell\_four\_mul\_floorZ\_le}{HalfCylinderMiddleCarryLowerBound.lean:4142}
(\ev{Lean}).  On a middle transition $D_{s+1}=D_s$.  Writing
$\rho=1/2-X_{D_s}(2)$, the exact quotient identity therefore gives
\[
 4^{s+1}\rho=2^{s+2}+4Z_s\le2^{s+2}+c,
\]
as in
\lean{seamMiddleCell\_scaledRationalRemainder\_le}{HalfCylinderMiddleCarryLowerBound.lean:4167}
(\ev{Lean}).

For $C_s=-3$, the correct label is $c=1$.  This estimate does have a
local consequence: since
$4^{s+1}R_s>2^{s+2}+4/3$, it gives $\rho<R_s$.
Thus the completed support $D_s\cup\{s+1,s+2,\ldots\}$ has value
greater than $1/2$.  It does not prove that every later greedy step
survives.  For $s\ge13$, under the additional all-right-tail hypothesis, the same
completed support would have value less than $1/2$, giving the
contradiction in
\lean{finalMiddleCell\_neg\_three\_not\_last}{HalfCylinderFinalMiddleCellEscape.lean:587}.
That tail hypothesis must not be dropped.  For $c=2,3$, the displayed
upper bound alone no longer gives $\rho<R_s$.

A parallel iff for the right branch: an overshoot witness $\le2^s$ produces a pulse leak above
$2^{s+2}$ exactly when the overshoot sits within an explicit distance $k$ of $2^s$ with
$4k<\mathrm{abovePulse}$ (\lean{seamRightPulseLeak\_iff\_distanceCell}{HalfCylinderMiddleCarryLowerBound.lean:4202},
\ev{Lean}). A new one-coordinate normal form, $\mathtt{seamPureHalfPrefixRemainder}\,s :=
\mathtt{integerGreedyRemainder}(\mathtt{seamWeights}\,s)(2^{2s-1})$
(\lean{seamPureHalfPrefixRemainder}{HalfCylinderMiddleCarryLowerBound.lean:4232}), has an exact
case-split formula (\lean{seamPureHalfPrefixRemainder\_eq\_ite}{HalfCylinderMiddleCarryLowerBound.lean:4238},
\ev{Lean}) whose equivalence with the two-sided bound is stated explicitly in
\lean{seamPureHalfPrefixRemainder\_le\_iff\_twoSided}{HalfCylinderMiddleCarryLowerBound.lean:4356} (\ev{Lean}).

\begin{obs}[The reset-band hypothesis has a certified base case]
\texttt{SeamUpperResetDyadicBandEscape} (O5's second hypothesis, discussed after \S\ref{sec:o4} above) is
reduced by three new lemmas to a quarter-scale overshoot exclusion: the iff
\lean{seamUpperReset\_band\_iff\_successorRemainder\_avoids}{HalfCylinderMiddleCarryLowerBound.lean:4469}
(\ev{Lean}) rewrites the reset-charge band as a successor-remainder avoidance; the identity
\lean{upperResetCharge\_dyadicDistance\_add\_abovePulse}{HalfCylinderMiddleCarryLowerBound.lean:4490}
(\ev{Lean}) rescales it by a factor of $4$; and
\lean{seamUpperReset\_band\_of\_overshoot\_band}{HalfCylinderMiddleCarryLowerBound.lean:4514} (\ev{Lean})
derives the band from a quarter-scale overshoot exclusion plus the pulse bound $\mathrm{abovePulse}\le
2(d-2)$. Row $d=13$, the first row of the late regime, is then proved unconditionally to escape
\emph{every} band simultaneously: its successor remainder is exactly $392$
(\lean{seamUpperResetDyadicBandEscape\_at\_thirteen}{HalfCylinderMiddleCarryLowerBound.lean:4547},
\ev{Lean}, \ev{Cert} for the numeral $392$). This is one certified row of a hypothesis that is $\forall
d\ge13$; it is evidence the base case is not itself an obstruction, nothing more.
\end{obs}

\subsubsection{Finite approximations without compatibility}
For a finite support $D\subseteq\{2,\ldots,M\}$, the integer
\[
 K_D(M)=2^{M-1}-\sum_{j=2}^{M}2^{M-j}\#\{d\in D:d\mid j\}
\]
is the half-carry at the terminal rank.  It is
$\operatorname{ihc}(D,M-1)$ in the notation used earlier.  The condition
below asks for $M_j\to\infty$ and finite sets $D_j$ with
$|K_{D_j}(M_j)|/2^{M_j}\to0$.  It does not require the sets to agree or
bound the carries at earlier ranks.
The estimate
\[
 \left|\sum_{d\in D_j}\frac1{2^d-1}-\frac12\right|
 \le\frac{|K_{D_j}(M_j)|+2\sqrt{M_j}+4}{2^{M_j}}
\]
explains the hypothesis: it makes the finite values converge to $1/2$.
Compactness then gives membership, but not the finite approximants needed
to apply the theorem.  A square-root terminal bound is sufficient and
stronger as a bound on these particular witnesses; no claim of a strict
logical separation between the corresponding existence statements follows.


Here is the finite family used in the cited construction.  Put
$B(m)=2\lfloor\sqrt m\rfloor+4$.  A depth-$M$ family with a shared
prefix through $K\le M$ consists of sets
$A_1,\ldots,A_{B(M)}\subseteq\{2,\ldots,M\}$ such that
\[
 1\le K_{A_k}(m)\le B(m)\quad(1\le m\le M),
 \qquad K_{A_k}(M)=k,
\]
all $A_k$ agree through $K$, and one integer $E\ge B(M)$ satisfies
\[
 \sum_{j=K+1}^{M}1_{A_k}(j)2^{M-j}+k=E
 \quad(1\le k\le B(M)).
\]
The suffixes therefore have consecutive binary values $E-k$.
The following statements explain how these particular families supply
finite approximations; they do not construct them at arbitrary depths.

\begin{prop}
Every shared-prefix family just defined at depth $N$ contains a finite
support $A$ with $|K_A(N)|\le B(N)$
(\lean{CylinderStage.halfTerminalOnlyStripWitness}{SuffixCylinderTerminalOnlyBridge.lean:63}, \ev{Lean},
\coord{support rigidity, half-carry recurrence}), and at a feedback row this survives \emph{both} outputs of
the corpus's total feedback theorem; full-cylinder advance or a localized one-hole seam, which
can delete at most one of carries $3,4$
(\lean{ErdosProblems.Erdos257.PaperCompleteR21.paper\_shared\_prefix\_family\_strip\_witness\_after\_feedback\_all\_depths}{lean/ErdosProblems/Erdos257/PaperCompleteR21/FeedbackRowStripWitnessAllDepths.lean:104},
\ev{Lean}).
\end{prop}
\leannote{Lean: \lproof{ErdosProblems/Erdos257/PaperCompleteR21/SharedPrefixFamiliesAndMeasureDichotomy.lean}{37}{paper_shared_prefix_family_contains_strip_witness}, \lproof{ErdosProblems/Erdos257/PaperCompleteR21/FeedbackRowStripWitnessAllDepths.lean}{104}{paper_shared_prefix_family_strip_witness_after_feedback_all_depths}, \lproof{ErdosProblems/Erdos257/PaperCompleteR21/FeedbackRowStripWitnessAllDepths.lean}{84}{paper_feedback_row_total_dichotomy}, \lproof{ErdosProblems/Erdos257/PaperCompleteR21/FeedbackRowStripWitnessAllDepths.lean}{74}{halfTerminalOnlyStripWitness_of_feedbackAdvance}, and 3 further declarations in the \hyperref[sec:coverage]{coverage section}.}
Four concrete certified instances follow at depths $51$--$54$
(\lean{fullCylinderStage51\_halfTerminalOnlyStripWitness}{SuffixCylinderTerminalOnlyBridge.lean:231},
\lean{fullCylinderStage54\_halfTerminalOnlyStripWitness}{SuffixCylinderTerminalOnlyBridge.lean:257},
\ev{Lean}, \ev{Cert} for the underlying stage), which is a finite witness chain, not a cofinal one.

\begin{obs}[Shared-prefix families at unbounded depths]
The conditional theorem
\[
\left(\begin{gathered}\text{such shared-prefix families exist}\\\text{at arbitrarily large depths }M\end{gathered}\right)
\ \Longrightarrow\ \exists A,\ A.\mathrm{Infinite}\wedge X_{A}(2)=1/2
\]
is proved unconditionally
(\lean{exists\_infinite\_support\_half\_of\_cofinalCylinderStages}{SuffixCylinderTerminalOnlyBridge.lean:277},
\ev{Lean}), with a positive-support refinement using the same odd-denominator lemma
\lean{finite\_boolSupport\_ne\_half}{HalfCarryReachability.lean:589} already used for O2 above
(\lean{exists\_infinite\_positive\_support\_half\_of\_cofinalCylinderStages}{SuffixCylinderTerminalOnlyBridge.lean:287},
\ev{Lean}). The hypothesis requires these finite families at arbitrarily large
depths and is \ev{Open}.  The terminal bound then implies half-membership.
This is a sufficient condition; no converse or strict comparison with
other sufficient conditions is established here.
\end{obs}

\subsubsection{Removing one case from a sufficient tail condition}

Five declarations extend the already-cited \texttt{SeamMiddleProducerSqrtEscape} hypothesis with a
reduction in the cases that remain to be checked, using the proved
exclusion of the $-3$ cell.  No strict separation of the existence
conditions is inferred from dropping that case.

\begin{defn}
$\mathrm{SeamMiddleProducerTailEscape}$ is the exact hypothesis that at every genuine middle
transition ($s\ge13$, no successor carry, below the terminal-weight threshold) the terminal-augmented
binary coefficient tail is strictly below the signed carry
(\lean{SeamMiddleProducerTailEscape}{HalfCylinderLastProducerContradiction.lean:34}).
$\mathrm{SeamMiddleProducerTailEscapeExceptNegThree}$ is the same hypothesis restricted to cells
$\ne-3$ (\lean{SeamMiddleProducerTailEscapeExceptNegThree}{HalfCylinderLastProducerContradiction.lean:54}).
\end{defn}
The reduction $\mathrm{TailEscape}\Rightarrow\mathrm{ExceptNegThree}$ is trivial by hypothesis-dropping
(\lean{SeamMiddleProducerTailEscape.toExceptNegThree}{HalfCylinderLastProducerContradiction.lean:76},
\ev{Lean}); the substance is in the converse direction of the implication. The reduced hypothesis, with cell
$-3$ discharged by the already-proved
\lean{middleProducer\_neg\_three\_not\_last}{HalfCylinderLastProducerContradiction.lean:364}, already
suffices to close HALF:
\[
\mathrm{SeamMiddleProducerTailEscapeExceptNegThree}\ \Longrightarrow\ (1/2:\mathbb{R})\in\mathcal{A}
\]
(\lean{half\_mem\_mersenneAchievementSet\_of\_middleProducerTailEscapeExceptNegThree}{HalfCylinderLastProducerContradiction.lean:436},
\ev{Lean}), with the unreduced hypothesis and the already-cited square-root majorant both routing through
it (\lean{half\_mem\_mersenneAchievementSet\_of\_middleProducerTailEscape}{HalfCylinderLastProducerContradiction.lean:451},
\ev{Lean}). So the chain now reads $\mathrm{SqrtEscape}\Rightarrow\mathrm{TailEscape}\Rightarrow
\mathrm{ExceptNegThree}\Rightarrow\mathrm{HALF}$, a chain of sufficient implications.  The displayed arrows do not prove
that any implication is strict; all three antecedents remain \ev{Open}.

\subsubsection{An eventual nonnegative margin}

The finite-horizon condition can be stated without a separate private
predicate.  Recall the integer margin $F_k(J)$ from
Theorem~\ref{thm:frozen-margin} and the equivalence in
\lean{exists\_greedyHalfFrozenMargin\_nonneg\_iff\_excess\_neg}{HalfCylinderFiniteShadow.lean:1229}.

\begin{defn}[A nonnegative margin by the next depth]
At every skipped step $k+1$, require
\[
 \exists J\in\{0,\ldots,k+1\},\qquad F_k(J)\ge0.
\]
This is the hypothesis encoded by
\lean{HalfGreedyGovernedFrozenMarginProducer}{HalfCylinderFiniteShadow.lean:1267}.
\end{defn}
It implies $1/2\in\mathcal A$, as stated in
\lean{half\_mem\_mersenneAchievementSet\_of\_governedFrozenMarginProducer}{HalfCylinderFiniteShadow.lean:1275}
and deduced through
\lean{half\_mem\_mersenneAchievementSet\_of\_skipped\_dyadicCap}{GreedyAchievementSet.lean:1444}
(\ev{Lean}).  Indeed, the finite coefficient window is at most its
infinite tail, so $F_k(J)\ge0$ implies
$r_k(1/2)\le2^{-(k+1)}$ at that skip.  This is below $R_{k+1}$.
At $k=0$ the remainder is exactly $1/2$ and the same inequality
holds; the strict version in Theorem~\ref{record:257bm-c19}
requires positive $k$.

The bound $J\le k+1$ is \emph{not used} in this implication.
Allowing any finite $J$ at each skip already suffices; requiring
first passage by $k+1$ is a stronger finite-horizon demand, not a
necessary ingredient of the membership proof.  Neither all-skip
existence condition is established here, and no equivalence between
the bounded and unbounded horizons is asserted. \ev{Open}

\subsubsection{Rewriting the signed carry recurrence}

Numeral-adjacent boundary words have more structure than the previously used
coefficient bound.  If the lower suffix numeral is one larger than the upper
suffix numeral, there is a unique carry pivot $d$: rank $d$ enters the lower
support, every rank $e$ with $d<e\le N$ flips from lower-unselected to
upper-selected, and all ranks below $d$ agree
(\lean{exists\_carry\_pivot\_of\_wordSuffixNumeral\_succ}{SuffixCylinderCarryPivot.lean:141},
\ev{Lean}, \scale{uniform}, \coord{support rigidity, half-carry recurrence}).  Hence for
every positive coefficient row $m\ge1$ the boundary discrepancy is the exact divisor-incidence
identity
\[
 c_{\rm low}(m)+\#\{e:d<e\le N,\ e\mid m\}
 =c_{\rm up}(m)+\mathbf 1_{d\mid m},
\]
not merely an inequality
(\lean{supportCoeff\_add\_deepDivisorCount\_eq\_of\_carry\_pivot}{SuffixCylinderCarryPivot.lean:163},
\ev{Lean}).

The stage-level theorem consumes adjacent profiled prefixes directly
(\lean{profiledGapStage\_carry\_pivot\_of\_hasAdjacentPrefixes}{SuffixCylinderCarryPivot.lean:228},
\ev{Lean}).  This is a local normal form for the two-sheet boundary.  It does
not produce adjacent stages cofinally and therefore does not close HALF or
universal \#257.

\subsubsection{Further restrictions on represented values}

This module works in the new \texttt{ErdosProblems.Erdos257} namespace and states in its own header
that its results ``concern achievement-set geometry and do not settle the universal irrationality
problem.'' For an arbitrary support restriction $J\subseteq\mathbb{N}$, define
$\mathtt{supportedMersenneAchievementSet}\,J$ as the range of Mersenne digit values supported on $J$
(\lean{supportedMersenneAchievementSet}{ErdosProblems/Erdos257/MersenneSubseriesRigidity.lean:76}). It is compact
(\lean{isCompact\_supportedMersenneAchievementSet}{ErdosProblems/Erdos257/MersenneSubseriesRigidity.lean:90}, \ev{Lean}),
nowhere dense for every $J$ (\lean{isNowhereDense\_supportedMersenneAchievementSet}{ErdosProblems/Erdos257/MersenneSubseriesRigidity.lean:112},
\ev{Lean}), and perfect; compact with no isolated points; whenever $J$ is infinite
(\lean{perfect\_supportedMersenneAchievementSet}{ErdosProblems/Erdos257/MersenneSubseriesRigidity.lean:167}, \ev{Lean}), via a
Cantor-space no-isolated-point argument on the coding space itself
(\lean{preperfect\_supportedMersenneDigitSet}{ErdosProblems/Erdos257/MersenneSubseriesRigidity.lean:120}, \ev{Lean}).

\begin{prop}[Exact Lebesgue-measure dichotomy]
Either $J=F^c$ for a finite $F$, and $\mathrm{vol}(\mathtt{supportedMersenneAchievementSet}\,J)=2^{-|F|}$
exactly, or $J^c$ is infinite and the volume is exactly $0$.
\end{prop}
\leannote{Lean: \lproof{ErdosProblems/Erdos257/PaperCompleteR21/SharedPrefixFamiliesAndMeasureDichotomy.lean}{68}{paper_volume_supportedMersenneAchievementSet_dichotomy}.}
(\lean{volume\_supportedMersenneAchievementSet\_dichotomy}{ErdosProblems/Erdos257/MersenneSubseriesRigidity.lean:397}, \ev{Lean},
\scale{uniform}, \coord{divisor counts and finite sums}). The finite case is an exact doubling induction over inserted
coordinates (\lean{volume\_supportedMersenneAchievementSet\_insert}{ErdosProblems/Erdos257/MersenneSubseriesRigidity.lean:286},
\ev{Lean}, using disjointness of the two coordinate-split faces,
\lean{disjoint\_supportedMersenneAchievementSet\_translate}{ErdosProblems/Erdos257/MersenneSubseriesRigidity.lean:262}, \ev{Lean},
which is exactly where unique binary coding is spent). None of this module touches irrationality of any
subseries value; it is a self-contained measure-and-topology classification of achievement sets, orthogonal
to O1--O5. The dichotomy is the case of retained coordinates of Hornich's strict-tail measure
formula~\cite[Theorem~4(1), p.~9]{nitecki2013}: the retained Mersenne weights still exceed their
tails, and the Mersenne tail estimate supplies the values $2^{-|F|}$ and $0$.

\subsubsection{Conditional counterexamples at rational targets}

This short module restates $\mathrm{U257}$ in the new namespace
(\lean{UniversalMersenneSubseriesIrrationality}{ErdosProblems/Erdos257/HalfCounterexampleFrontier.lean:26}) and gives it a
problem-centric interface to the legacy half-terminal required input already used for O2 above: any
finite supports with vanishing scaled terminal carry witness gives an infinite $A$ with
$X_{A}(2)=1/2$
(\lean{exists\_rational\_half\_counterexample\_of\_terminalScaledVanishing}{ErdosProblems/Erdos257/HalfCounterexampleFrontier.lean:31},
\ev{Lean}), hence refutes $\mathrm{U257}$ outright
(\lean{not\_universal\_of\_terminalScaledVanishing}{ErdosProblems/Erdos257/HalfCounterexampleFrontier.lean:39}, \ev{Lean}). The
module's own docstring states plainly that the required input sequence itself ``has only been checked to a
large finite cutoff'': this is a restatement of the already-\ev{Open} O2 required input obligation in a new
namespace, not a new result, and it duplicates \S O2's own logic (any HALF witness is $\neg\mathrm{U257}$)
rather than extending it.

The module now also proves an unconditional finite-support obstruction at
the second rational target $1/21$:
\[
  F\subseteq\{n:n\ge2\}\quad\Longrightarrow\quad
  \sum_{n\in F}\frac1{2^n-1}\ne\frac1{21}
\]
for every finite $F$
(\lean{finiteErdosSum\_ne\_one\_div\_twenty\_one}{ErdosProblems/Erdos257/HalfCounterexampleFrontier.lean:61}, \ev{Lean},
\scale{uniform}, \coord{finite-lcm}).  The denominator-order identity forces
$\operatorname{lcm}(F)=6$, after which the eight subsets of $\{2,3,6\}$ are exact arithmetic.
This removes finite termination for that target; it does not establish achievement-set membership.

\subsubsection{Multiplicity of the stated two--three representations}

This module records the exact arithmetic of the primitive cone
$2p+3q=n$, $p,q>0$, $\gcd(p,q)=1$
(\lean{IsPrimitive23Solution}{Primitive23Multiplicity.lean:20}).  It proves that every $n\ge11$
has such a representation
(\lean{exists\_primitive23\_solution\_of\_eleven\_le}{Primitive23Multiplicity.lean:52},
\ev{Lean}, \scale{uniform}, \coord{finite-lcm}), while rank $10$ has none
(\lean{no\_primitive23\_solution\_ten}{Primitive23Multiplicity.lean:26}, \ev{Lean}).
Rank $11$ already has two distinct witnesses
(\lean{exists\_two\_primitive23\_solutions\_eleven}{Primitive23Multiplicity.lean:38}, \ev{Lean}),
and so does every rank $10k$ with $k\ge2$
(\lean{exists\_two\_primitive23\_solutions\_mul\_ten}{Primitive23Multiplicity.lean:86},
\ev{Lean}, \scale{uniform}, \coord{finite-lcm}).  These theorems are deliberately recorded as
route diagnostics rather than as a counterexample: the file does not prove that the associated
integer-multiplicity expansion of $1/21$ Booleanises.  The recurring multiplicities are precisely
the unresolved collision data, not zero--one support digits.

\subsubsection{Integer-quotient tests for one twenty-first}

The later quotient-greedy analysis removes the diffuse alternatives.  Let
\texttt{TwentyOneFatalAlignedBranch} denote the explicit branch consisting
of a fatal greedy gap, finite skipped support (hence cofinite selection),
eventual alignment between the quotient rows and the rational greedy prefix,
and eventual visits to every doubling block.  Then
\[
 \frac1{21}\in\ensuremath{\mathcal A}
 \quad\Longleftrightarrow\quad
 \neg\,\mathtt{TwentyOneFatalAlignedBranch}.
\]
\lean{one\_div\_twenty\_one\_mem\_iff\_not\_fatalAlignedBranch}{TwentyOneQuotientGreedy.lean:3507}
\quad\ev{Lean}\quad\scale{uniform}\quad\coord{greedy recurrence}.

On the fatal branch the canonical quotient remainder $s_R$ is eventually
strictly above the closed capacity $2^R$
(\lean{twentyOneFatalAlignedBranch\_eventually\_strict\_supercapacity}{TwentyOneQuotientGreedy.lean:5625}),
and thereafter the support appends $R+1$ while $s_R$ follows one explicit
affine recurrence
(\lean{twentyOneFatalAlignedBranch\_eventually\_affine\_supercapacity}{TwentyOneQuotientGreedy.lean:5658}).
Conversely, closed rows $s_R\le2^R$ along any unbounded sequence decay after
normalisation and put $1/21$ in the achievement set
(\lean{twentyOneCofinalEvenQuotientGreedyDecay\_of\_closedRows}{TwentyOneQuotientGreedy.lean:5554};
\lean{one\_div\_twenty\_one\_mem\_mersenneAchievementSet\_of\_cofinalGreedyDecay}{TwentyOneQuotientGreedy.lean:5458}).

This is the exact current contribution: membership has been reduced to
excluding one named permanent affine-supercapacity regime.  No theorem
excludes that regime, forces unbounded closed returns, or proves membership.
Combined with the finite-support obstruction above, membership would provide
an infinite-support rational counterexample to universal \#257; its present
status remains \ev{Open}.
% ; - part a257_p2 ; -
\clearpage
\part*{IV. Scales, conditional interfaces and unresolved inputs}
\addcontentsline{toc}{part}{IV. Scales, conditional interfaces and unresolved inputs}
\section{Finite ranges and unbounded parameters}
\label{sec:scale-ladder}

Every claim in this programme carries a \emph{scale} tag alongside its evidence band. The two
axes are independent and both matter to a reader deciding whether a result can be used as a
premise. Evidence labels record the type of support supplied with a claim.
\ev{Lean} refers to the cited formal source and its recorded build
status, not a new kernel run in this revision.  A newly cited release
file marked \texttt{release\_only} in the packet is described as
\ev{Lean source}; that label does not promote it to the packet's
checked-build class.  \ev{Cert} denotes an exact finite computation;
the surrounding text distinguishes fresh checks from inherited
reports.  \ev{Math} denotes an ordinary mathematical argument, not
independent human certification.  \ev{Cited} invokes the stated
external result, and \ev{Open} marks an unproved hypothesis.

The scale records the quantified range.  A \scale{fixed} claim concerns a
specified value; \scale{bounded} denotes either a finite tested range or a
witness bounded in terms of the input.  A \scale{uniform} assertion applies
to every index in its stated range, \emph{with its hypotheses retained}; it
need not be unconditional.  A \scale{cofinal} assertion has the form
$\forall N\,\exists n\ge N,\ P(n)$.  Compatibility between its witnesses is
an additional requirement only when explicitly imposed.  Definitions and
scale-free statements are marked \scale{n/a}.

The important distinction is between obtaining a witness in a finite
range and proving the quantified assertion used by a membership test.
For example, half-membership follows from unbounded successful exact
rows or from infinitely many greedy skips.  A different sufficient
condition asks for a bound at every reset beyond its threshold.
These statements have different quantifiers, even when the same
transition formulas occur in their proofs.
Each row must retain its predicate, quantified variable, and every premise. An unconditional
eventual assertion $\forall n\ge N_0,\,P(n)$ does give cofinal witnesses to that same predicate,
by taking $n=\max(N,N_0)$. A uniform conditional implication $H(n)\Rightarrow C(n)$ supplies no
witness until $H(n)$ is established, and a checked finite band supplies no assertion beyond its
endpoints. These are different logical shapes and are recorded separately below.

Table~\ref{tab:scale-ladder} separates established finite or quantified
results from unproved membership hypotheses and the proved implications
that use them.  Its lower part therefore contains proved theorems as well
as open assumptions.  The scale alone determines neither evidence nor
truth status; the stated predicate and premises do.

\begingroup
\begin{landscape}
\fontsize{8}{9.8}\selectfont
\setlength{\tabcolsep}{3pt}
\begin{longtable}{@{}p{0.42\linewidth}p{0.11\linewidth}p{0.29\linewidth}p{0.12\linewidth}@{}}
\toprule
\papertablehead
\textbf{Result} & \textbf{Scale} & \textbf{Quantifier prefix} & \textbf{Coordinate} \\
\midrule
\endfirsthead
\toprule
\papertablehead
\textbf{Result} & \textbf{Scale} & \textbf{Quantifier prefix} & \textbf{Coordinate} \\
\midrule
\endhead
\bottomrule
\endfoot
\papertablehead\multicolumn{4}{l}{\emph{Fixed ;  one concrete witness or numeral}} \\
Seed for continuation from $n\ge6$: support $\{2,3,6\}$ satisfies
$\ensuremath{Q}=2^5-1$. \lean{exactLocalMersenneHalfRow\_six}{BooleanMobiusExactRowSeed.lean:34}
\ev{Cert}
& \scale{fixed} & $n=6$ & \coord{divisor counts and finite sums} \\
One-point counter-model: the double-or-recycle transition shape alone, even seeded, is not
cofinal. \lean{boundedDoubleOrRecycleModel\_not\_cofinal}{BooleanMobiusExactRowDichotomy.lean:79}
\ev{Lean}
& \scale{fixed} & $n=6$ (free-standing model) & \coord{meta-logical} \\
The combined-residue-mass $\le 1$ sufficient test for sharp capacity is not necessary: fixture
$D=\{2,3\}$, $c=5$. \lean{splitFractionMass\_one\_bound\_not\_necessary\_fixture}{BooleanMobiusSkippedCoreCriticalCapacity.lean:183}
\ev{Cert}
& \scale{fixed} & $D=\{2,3\},\ c=5$ & \coord{divisor counts and finite sums} \\
Exact residual fixture: $\tfrac12-W(\{2,3,6,7\})=\tfrac1{16002}$.
\lean{half\_sub\_four\_term\_prefix\_eq}{HalfCutLocator.lean:664} \ev{Cert}
& \scale{fixed} & none (numeric instance) & \coord{regression-fixture} \\
Closed-form Möbius--Lambert value $\sum_{d\ge1}\mu(d)/(2^d-1)=1/2$ (imported identity).
\lean{tsum\_moebius\_div\_two\_pow\_sub\_one\_eq\_half}{MersenneLambertLadder.lean:587} \ev{Cited}
& \scale{fixed} & none (closed-form value) & \coord{divisor counts and finite sums} \\
The one natural negative-Möbius candidate support overshoots: $1/2 <
\sum_{d\in N}1/(2^d-1)$ for $N=\{d\ge2:\mu(d)=-1\}$.
\lean{half\_lt\_tsum\_negativeMobius}{MobiusSignSupportNoGo.lean:164} \ev{Lean}
& \scale{fixed} & one fixed candidate support $N$ & \coord{divisor counts and finite sums} \\
\papertablehead\multicolumn{4}{l}{\emph{Bounded ;  a witness exists, controlled by the input}} \\
From an exact row, either repeated doubling is available or an exact row is obtained at some $2c-2$ with $c\le n$; these alternatives need not be disjoint.
\lean{exactLocalMersenneHalfRow\_double\_or\_recycle}{BooleanMobiusExactRowDichotomy.lean:26}
\ev{Lean}
& \scale{bounded} & $\forall n\ge6$, given a row at $n$: $(\ldots)\lor\exists c\le n(\ldots)$
& \coord{divisor counts and finite sums} \\
Every above-half bounded support recycles into a genuine exact row at some $2c-2$ with $c\le n$.
\lean{exists\_skippedCoreExactRow\_of\_value\_above}{BooleanMobiusExactRowCrossing.lean:155} \ev{Lean}
& \scale{bounded} & $E\subseteq[2,n]$, $X_E(2)>1/2$: $\exists c\le n$ & \coord{binary digits} \\
(\#249 cross-argument check) $\sum\varphi(n)/2^n$ matches no rational whose denominator divides
$2^{14}(2^h-1)$, for some $1\le h\le16$.
\lean{totient\_series\_ne\_rat\_of\_den\_dvd\_upto\_sixteen}{CarrySurvivorExtinction.lean:587} \ev{Cert}
& \scale{bounded} & $1\le h\le16$ (explicit bound) & \coord{binary digits} \\
Truncation-rung ladder: $\mathrm{HalfRung}(J)$ proved for every $3\le J\le22$ via the finite
decision procedure (\S\ref{sec:four-coordinates}, Theorem~\ref{thm:tr-finite-decision}).
(source: \texttt{erdos257\_truncation\_rung\_ladder\_2026\_07\_24.md}, \S5) \ev{Math}+\ev{Cert}
& \scale{bounded} & $3\le J\le22$ (finite table, exhaustive) & \coord{truncated greedy sums} \\
\papertablehead\multicolumn{4}{l}{\emph{Uniform assertions and conditional implications; premises retained}} \\
Master achievement/greedy-survival criterion.
\lean{mem\_mersenneAchievementSet\_iff\_greedy\_survival}{GreedyAchievementSet.lean:1458} \ev{Lean}
& \scale{uniform} & $\forall x\ge0$, membership iff survival at every rank & \coord{greedy recurrence} \\
Membership iff every actually-skipped rank survives.
\lean{no\_lastHalfGreedySkip\_iff\_every\_skip\_survives}{HalfCylinderFixedTailSocket.lean:57} \ev{Lean}
& \scale{uniform} & $\forall M\in$ skipped support & \coord{greedy recurrence} \\
Every crossing core is forced to be the canonical rational half-greedy prefix.
\lean{eq\_halfGreedyPrefixSupport\_of\_critical\_crossing}{BooleanMobiusCriticalCapacityCofinal.lean:50}
\ev{Lean}
& \scale{uniform} & $c\ge4$, $D\subseteq[2,c)$, $0<1/2-X_D(2)<w_c$ & \coord{greedy recurrence} \\
Largest-skip conditional implication used in the reset argument.
\lean{half\_mem\_mersenneAchievementSet\_of\_largestSkipLateStepSocket}{HalfCylinderLargestSkipInduction.lean:167}
\ev{Lean}
& \scale{uniform} & $\forall s\ge14$, given the hypothesis & \coord{integer quotients} \\
Band-avoidance certified at reset indices $13\le r\le30$, using
successor remainders at indices $14,\ldots,31$.  This is a finite range
of the conditional hypothesis. \lean{seamUpperResetDyadicBandEscape\_through\_thirty}{HalfCylinderUpperResetBandCertificates.lean:78}
\ev{Cert}
& \scale{bounded} & $13\le r\le30$, exhaustive & \coord{integer quotients} \\
Critical-band-index reduction: a $(d{+}1)$-way check collapses to one nearest-boundary check.
\lean{half\_mem\_mersenneAchievementSet\_of\_upperResetCriticalBandEscape}{HalfUpperResetCriticalBand.lean:896}
\ev{Lean}
& \scale{uniform} & $\forall d,E$ with $E\le2^{d+1}$ & \coord{dyadic-boundary} \\
Parity obstruction: no finite support has value $1/2$. \lean{finite\_boolSupport\_ne\_half}{HalfCarryReachability.lean:589}
\ev{Lean}
& \scale{uniform} & $\forall A$ finite, $0\notin A$ & \coord{denominator-parity} \\
Full-support irrationality at every integer base.
\lean{irrational\_erdosSum\_full\_support}{CertificateKernel.lean:8328}; base two:
\lean{irrational\_erdosBorwein\_series}{CertificateKernel.lean:8335} \ev{Lean}
& \scale{uniform} & $\forall b\ge2$ & \coord{divisor counts and finite sums} \\
Erdős 1968: any infinite, pairwise-coprime, reciprocal-summable index set gives an irrational
support series. \lean{irrational\_erdosSupportSeries\_pairwise\_coprime}{CertificateKernel.lean:10776}
\ev{Lean}
& \scale{uniform} & $\forall b\ge2\ \forall A$ (infinite, coprime, summable) & \coord{divisor counts and finite sums} \\
Generic block-certificate engine (problem-agnostic): weighted-coefficient series is irrational
given a per-precision divisibility/window certificate.
\lean{irrational\_coeff\_series\_of\_weighted\_coeff\_block\_certificates}{CertificateKernel.lean:8665}
\ev{Lean}
& \scale{uniform} & $\forall b\ge2\ \forall c\,(\forall m,c(m)\le m)$, given $\forall q\,\exists N,K,L,C$
& \coord{binary digits} \\
Finite greedy-prefix denominators are odd (source of the parity boundary exclusion,
\S\ref{sec:four-coordinates}). \lean{halfGreedyPrefixDenominator\_odd}{DyadicPrefixCompression.lean:1521}
\ev{Lean}
& \scale{uniform} & $\forall n$ & \coord{denominator-parity} \\
\papertablehead\multicolumn{4}{c}{\textbf{Unproved hypotheses and the proved implications that use them}} \\
\papertablehead\multicolumn{4}{l}{\emph{Unbounded-witness and all-reset conditions have different quantifiers}} \\
The exact-row existence condition: exact rows recur at arbitrarily large endpoints, with no coherence
required between witnesses. \lean{CofinalExactLocalMersenneHalfRows}{BooleanMobiusCofinalExactRows.lean:38}
\ev{Open}
& \scale{cofinal} & $\forall N\ \exists n\ge N,\ \mathrm{ExactLocalMersenneHalfRow}(n)$
& \coord{divisor counts and finite sums} \\
The preceding cofinal exact-row condition implies $1/2\in\mathcal A$ by compactness. \lean{half\_mem\_mersenneAchievementSet\_of\_cofinalExactLocalRows}{BooleanMobiusCofinalExactRows.lean:71}
\ev{Lean} (theorem proved; \emph{hypothesis} open)
& \scale{cofinal} & given $\forall N\,\exists n\ge N(\ldots)$ & \coord{topological-achievement-set} \\
An equivalent membership condition: the canonical rational half-greedy orbit itself has infinitely many
positive skip events. \lean{CofinalPositiveHalfGreedySkips}{BooleanMobiusSkipRowCofinal.lean:22}
\ev{Open}
& \scale{cofinal} & $\forall N\ \exists c\ge\max(N,4),\ 0<\mathrm{rem}(c{-}1)<w_c$
& \coord{greedy recurrence} \\
Infinitely many positive greedy skips imply $1/2\in\mathcal A$.
\lean{half\_mem\_mersenneAchievementSet\_of\_positiveHalfGreedySkips}{BooleanMobiusSkipRowCofinal.lean:97}
\ev{Lean} (theorem proved; hypothesis open)
& \scale{cofinal} & given $\forall N\,\exists c\ge N(\ldots)$ & \coord{greedy recurrence} \\
Master half-branch dichotomy: membership iff the greedy skip set is infinite.
\lean{half\_mem\_mersenneAchievementSet\_iff\_greedySkippedSupport\_infinite}{GreedyAchievementSet.lean:2583}
\ev{Lean} (iff proved; both sides cofinal-shaped)
& \scale{cofinal} & iff $(\mathrm{skippedSupport}).\mathrm{Infinite}$ & \coord{greedy recurrence} \\
Seven-way classification hub: membership iff the seam word is not eventually right, iff
unbounded terminal-false, iff cofinally many skipped ranks, \emph{etc.}
\lean{half\_mem\_mersenneAchievementSet\_iff\_not\_seamGreedyEventuallyRight}{HalfCylinderHalfMembershipClassification.lean:112}
\ev{Lean}
& \scale{cofinal} & every right-hand side is $\exists^\infty$/Unbounded-shaped & \coord{integer quotients} \\
Ladder dichotomy: either infinitely many rungs survive (\#257 false, explicit limit support) or an
eventual bad-rank obstruction closes only this required input. (source: \texttt{erdos257\_truncation\_rung\_ladder\_2026\_07\_24.md},
Cor.~8) \ev{Math}
& \scale{cofinal} & open branch (a): $\exists^\infty J,\ \mathrm{HalfRung}(J)$ & \coord{truncated greedy sums} \\
Reset square-root escape, a sufficient all-reset hypothesis.
The release source gives the conditional implication; the run-length
comparison is discussed separately in Section~\ref{sec:o4}.
\ev{Open}
& \scale{all resets} & $\forall r\ge10$ that are resets,
$|\mathrm{rem}(r+1)-2^{r+1}|>2^{(r+5)/2}$
& \coord{integer quotients} \\
\caption{Quantifiers and evidence for the stated results.  The labels
separate finite tests, uniform implications, unbounded witness
conditions and all-reset hypotheses.  A theorem conditional on an
unproved hypothesis is not a proof of that hypothesis.}
\label{tab:scale-ladder}
\end{longtable}
\end{landscape}
\endgroup

\begin{rem}[Finite checks and quantified conclusions]
\label{rem:promotion-audit}
The supplied audit examined eighteen possible implications from recorded
results to the required unbounded-witness conditions.  It reported no
such implication.  This is a report about those attempts, not a theorem
excluding other arguments and not a fresh kernel replay.

For example, \lean{seamUpperResetDyadicBandEscape\_through\_thirty}{HalfCylinderUpperResetBandCertificates.lean:78}
checks reset indices $13$ through $30$; it has no conclusion beyond that
range.  Likewise, a bound depending on the depth need not produce a
witness at a larger depth.  Conversely, a proved eventual assertion of
the \emph{same} witness predicate immediately gives cofinal witnesses.
The issue is the precise premise and conclusion, not the name of its
scale or the tactic used in a finite proof.
\end{rem}

\section{Approximations to the greedy orbit}
\label{sec:four-coordinates}

\subsection*{Four reported calculations}

Fix the weights $x_n:=1/(2^n-1)$, $n\ge2$, and the target $1/2$. The single object under study
throughout this programme is the greedy orbit for this target: a residual $\rho$ starts at $1/2$,
and at rank $k$ the orbit \emph{takes} iff $\rho\ge x_k$. A skip at rank $k$ is \emph{safe} iff
$\rho\le T_{k+1}:=\sum_{j>k}x_j$ and \emph{fatal} iff $\rho\in(T_{k+1},x_k)$ ;  the fatal interval is
exactly the gap no later tail can repair. The four calculations below concern this target, but their finite
parameters and certification records differ.  Agreement with the real
greedy prefix requires a convergence argument for the approximants used
in each calculation; agreement over a tested range is not that argument.

\begin{itemize}
\item \textbf{Truncation rung $J$.} Replace the full weight $x_n$ by the $J$-term truncation
$w_n^{(J)}:=\sum_{q=1}^{J}2^{-qn}$ and ask whether some Boolean support hits $1/2$ exactly under
the truncated weights. Certified for $3\le J\le22$ (Table~\ref{tab:scale-ladder}, bounded row;
Theorem~\ref{thm:tr-witness-exclusion}--\ref{thm:tr-finite-decision} below).
\item \textbf{Seam row $s$.} The base-4 seam recursion tracking the deviation $\Delta_s:=\mathrm{rem}(s)-2^s$
of the integer greedy orbit from its dyadic target. Certified rows $6$ through $2500$ for the
reset-crossing classification (\S\ref{sec:scale-ladder}, uniform rows; source:
\texttt{erdos257\_reset\_crossing\_unification\_2026\_07\_24.md}, \S8), and the underlying real
orbit separately certified to row $200{,}000$ (source: \texttt{erdos257\_sqrt\_escape\_digit\_reduction\_2026\_07\_24.md},
\S3): exactly one TARGET violation in that range, at $r=7$.
\item \textbf{Integer margin $m_c$.} An advisory Type~B scan of the integer-greedy margin out to
$5\times10^5$, tracking how far the certified region sits from the danger threshold.
\item \textbf{Sharp tail margin.} The rank-by-rank take/skip margin against the exact tail, checked
for ranks $2$ through $3000$: $1497$ takes against $1502$ skips, zero fatal skips anywhere, minimum
safe-skip margin $+2.9922$ bits at rank $5$, and the smallest
reported take-margin $+0.0340$ bits at rank $7$, where the remainder
$1/126$ exceeds the weight $1/127$.  Here the margins are the base-two
logarithms of the corresponding positive slack multiplied by $4^k$ (source:
\texttt{erdos257\_reset\_crossing\_unification\_2026\_07\_24.md}, \S8).
\end{itemize}

\begin{prop}[Stability of each fixed greedy prefix]
\label{prop:one-orbit}
Let $t_j\to1/2$ and, for each fixed $n\ge2$, let
$v_n^{(j)}\to x_n=(2^n-1)^{-1}$, with $v_n^{(j)}>0$.
Apply the greedy rule with target $t_j$ and weights $v_n^{(j)}$ in
increasing order of $n$, through depths $m_j\to\infty$.
For every fixed depth $K$, the decisions at ranks $2,\ldots,K$
eventually agree with those of the real half-greedy rule.
This assertion concerns finite prefixes, not survival at all ranks.
\end{prop}
\leannote{Lean: \lproof{ErdosProblems/Erdos257/PaperCompleteR21/GreedyOrbitNoTies.lean}{272}{paper_one_orbit_stability}, \lproof{ErdosProblems/Erdos257/PaperCompleteR21/GreedyOrbitNoTies.lean}{182}{approx_orbit_induction}, \lproof{ErdosProblems/Erdos257/PaperCompleteR21/GreedyOrbitNoTies.lean}{166}{tailGreedyRemainder_mersenne}.}
\begin{proof}
At rank $2$, the target and the weight converge to their real counterparts.
Their comparison is strict by Lemma~\ref{lem:no-ties}, so its sign is
eventually unchanged.  Inductively, once all earlier decisions agree,
the approximate remainder is the target minus a fixed finite sum of
convergent weights.  It therefore converges to the real remainder.
The next comparison is again strict, so the next decision eventually
agrees.  A finite induction through $K$ proves the claim.
\end{proof}
\begin{lem}[No-ties lemma]
\label{lem:no-ties}
At every rank $k$ of the full greedy orbit for target $1/2$, both defining comparisons are strict:
$\rho\ne x_k$ and $\rho\ne T_{k+1}$.
\end{lem}
\leannote{Lean: \lproof{ErdosProblems/Erdos257/PaperCompleteR21/GreedyOrbitNoTies.lean}{126}{paper_no_ties}, \lproof{ErdosProblems/Erdos257/PaperCompleteR21/GreedyOrbitNoTies.lean}{71}{paper_no_ties_take}, \lproof{ErdosProblems/Erdos257/PaperCompleteR21/GreedyOrbitNoTies.lean}{114}{paper_no_ties_skip}, \lproof{ErdosProblems/Erdos257/PaperCompleteR21/GreedyOrbitNoTies.lean}{44}{irrational_mersenneTail}.}
\begin{proof}
The two strict comparisons have different reasons: parity excludes
equality with a weight, and irrationality excludes equality with the
complete tail.

\emph{Take boundary excluded (parity).} Write a finite Mersenne sum as $p/q=\sum_{d\in F}x_d$ in lowest
terms.  Its denominator $q$ is odd, because every $2^d-1$ is odd
(equivalently, \lean{halfGreedyPrefixDenominator\_odd}{DyadicPrefixCompression.lean:1521}, itself
traced to \lean{finiteErdosSum\_den\_odd}{DyadicPrefixCompression.lean:1513}). Hence the residual
after any finite prefix is $\rho=\tfrac12-\tfrac pq=\tfrac{q-2p}{2q}$: its reduced denominator
is \emph{even}. Every weight $x_k=1/(2^k-1)$ has \emph{odd} denominator. An even-denominator
rational can never equal an odd-denominator rational, so $\rho=x_k$ is impossible at any rank ; 
the take comparison is never a tie.

\emph{Skip-safety boundary excluded (irrationality).} The tail past rank $k$ is
$T_{k+1}=E-1-\sum_{2\le j\le k}x_j$, where $E=\sum_{n\ge1}1/(2^n-1)$ is the Erdős--Borwein
constant: a rational shift of $E$. Erdős's own 1948 theorem gives $E$ irrational
(\lean{irrational\_erdosBorwein\_series}{CertificateKernel.lean:8335}, instantiated at $b=2$), so
$T_{k+1}$ is irrational for every $k$. The residual $\rho$ arising from any finite Boolean prefix
is rational. An irrational number never equals a rational one, so $\rho=T_{k+1}$ is impossible at
any rank ;  the skip-safety comparison is never a tie either.

Both boundaries of the fatal interval $(T_{k+1},x_k)$ are therefore approached only strictly, at
every rank, unconditionally.
\end{proof}
\begin{rem}
The hypotheses of Proposition~\ref{prop:one-orbit} can be checked for
two concrete approximations.  The truncated geometric weights
$\sum_{q=1}^{J}2^{-qn}$ increase to $x_n$.  For the integer quotient
rows, division by $4^s$ gives weights
$4^{-s}\lfloor4^s/(2^n-1)\rfloor\to x_n$ and target
$4^{-s}(2^{2s-1}-2^s)=1/2-2^{-s}\to1/2$.
Hence both greedy rules agree with each fixed real prefix once their
parameters are large enough.  Convergent upper and lower bounds for
$R_k$ likewise determine the strict tail comparison at a fixed rank.
No conclusion about the advisory margin scan follows without specifying
its approximants and proving their convergence.  Fixed-prefix agreement is not uniform agreement through a growing
depth.  For the quotient rows, however, summability bounds the unmatched
tail by $R_K$ and does give convergence of the coded values, as proved
in Theorem~\ref{thm:seam-limit}.  It identifies their limit as
$X_G(2)$; it does not establish that this value is $1/2$.
\end{rem}

\subsection*{Successive tail truncations}

Finite truncations of the geometric expansion give rational weights
and permit a finite decision procedure for each fixed truncation
order $J$.  The bounds below explain why only finitely many ranks
need to be tested.  Fix $J\ge2$ and set
\[
\begin{aligned}
w_n^{(J)}&:=\sum_{q=1}^{J}2^{-qn}, &
T_{n+1}^{(J)}&:=\sum_{q=1}^{J}\frac{2^{-qn}}{2^q-1},\\
\mathrm{HalfRung}(J)&:\Longleftrightarrow
\exists A\subseteq\{2,3,\ldots\},&
\sum_{n\in A}w_n^{(J)}&=\tfrac12.
\end{aligned}
\]
The arguments below are ordinary mathematics.  The accompanying exact
integer computation rechecks the finite instances $3\le J\le22$; neither
the arguments nor that computation are presented as a Lean verification.

\begin{lem}[Forced greedy]
\label{lem:tr-forced-greedy}
For every $J\ge2$: $w_n^{(J)}>T_{n+1}^{(J)}$ for all $n$ (the $q=1$ terms agree exactly, and every
$q\ge2$ tail term is strictly smaller than the corresponding weight term). Consequently the greedy support is the unique candidate support, and
$\mathrm{HalfRung}(J)$ holds iff the greedy orbit for $1/2$ under weights $w_n^{(J)}$ never lands in
a fatal interval $(T_{n+1}^{(J)},w_n^{(J)})$. Rank $1$ is always a safe skip; ranks $2$ and $3$ are
always takes.
\end{lem}
\leannote{Lean: \lproof{ErdosProblems/Erdos257/PaperCompleteR21/TruncatedRungGreedyDecision.lean}{486}{paper_forced_greedy_unique_support_and_criterion}, \lproof{ErdosProblems/Erdos257/PaperCompleteR21/TruncatedRungGreedyDecision.lean}{567}{paper_forced_greedy_low_ranks}.}

\begin{lem}[Parity forces infinite support]
\label{lem:tr-parity}
For every $J\ge2$, no finite $A\subseteq\{2,3,\ldots\}$ attains
$\mathrm{HalfRung}(J)$.
\end{lem}
\leannote{Lean: \lproof{ErdosProblems/Erdos257/PaperCompleteR21/TruncatedRungWitnessHorizon.lean}{176}{paper_parity_excludes_finite_support}.}
\begin{proof}
The empty set has value zero.  For a nonempty finite set, put $m=\max A$.
After multiplication by $2^{Jm}$, the $m$th summand is odd and every earlier
summand is even, whereas the target $2^{Jm-1}$ is even.  This parity
contradiction excludes finite support.
\end{proof}

\begin{thm}[Witness exclusion]
\label{thm:tr-witness-exclusion}
Define the misalignment mass
$$\mu_J(M):=\sum_{q=2}^{J}\frac{2^{M\bmod q}}{2^q-1}.$$
Let $J\ge3$, $n\ge4$, and suppose some $M\in[n,2n-2]$ has $\mu_J(M)\le\tfrac{11}{15}$. Then no
Boolean prefix $D\subseteq\{2,\ldots,n-1\}$ satisfies
$$T_{n+1}^{(J)}<\tfrac12-\sum_{d\in D}w_d^{(J)}<w_n^{(J)};$$
This excludes the specified gap for every Boolean prefix, not only the
greedy prefix.
\end{thm}
\leannote{Lean: \lproof{ErdosProblems/Erdos257/PaperCompleteR21/TruncatedRungWitnessHorizon.lean}{250}{paper_witness_exclusion}.}
\begin{proof}
Suppose such a prefix exists and write
$\rho=1/2-\sum_{d\in D}w_d^{(J)}$.  Since
$T_{n+1}^{(J)}\ge2^{-n}$ and $n\le M\le2n-2$,
\[
 0<2^M\rho-2^{M-n}
   <\sum_{q=2}^{J}2^{M-qn}
   \le\frac{2^{M-2n}}{1-2^{-n}}\le\frac4{15}.
\]
Here $n\ge4$ gives the last inequality.  Split each scaled summand
$2^{M-qd}$ according to whether $qd\le M$.  The terms with
$qd\le M$ are integers.  The other terms have total
\[
 F=\sum_{d\in D}\sum_{\substack{1\le q\le J\\qd>M}}2^{M-qd}
 \le\sum_{q=2}^{J}\sum_{d>\lfloor M/q\rfloor}2^{M-qd}
 =\mu_J(M)\le\frac{11}{15}.
\]
There is no $q=1$ contribution to $F$, since $d<n\le M$.
Thus $2^M\rho+F-2^{M-n}$ is an integer strictly between zero and one,
a contradiction.  The same calculation also applies when $J=2$.
\end{proof}

\begin{cor}[Half-LCM horizon]
\label{cor:tr-half-lcm}
Let $J\ge2$ and $L_J:=\mathrm{lcm}(2,3,\ldots,J)$.  For every
$n\ge\max\{4,L_J/2+1\}$ there is an $M\in[n,2n-2]$ with
$\mu_J(M)<11/15$.  Thus witness exclusion leaves only the finite
window $[4,L_J/2]$ to check.
\end{cor}
\leannote{Lean: \lproof{ErdosProblems/Erdos257/PaperCompleteR21/TruncatedRungWitnessHorizon.lean}{467}{paper_half_lcm_horizon}.}
\begin{proof}
Let $M$ be the least multiple of $L_J$ with $M\ge n$.
If $M=L_J$, then $M\le2n-2$ is the assumed inequality.
If $M=kL_J$ with $k\ge2$, then $n\ge(k-1)L_J+1$ gives
$2n-2\ge2(k-1)L_J\ge M$.  At this witness,
\[
 \mu_J(M)=\sum_{q=2}^{J}\frac1{2^q-1}
 <\sum_{q=2}^{\infty}\frac{4}{3\cdot2^q}
 =\frac23<\frac{11}{15}.
\]
The bound uses $2^q-1\ge3\cdot2^{q-2}$, with strict inequality for
$q>2$, so no numerical approximation to the full series is needed.
Apply Theorem~\ref{thm:tr-witness-exclusion}; its proof also covers
$J=2$.  Ranks $2$ and $3$ are takes, as already noted.
\end{proof}

\begin{lem}[Mod-12 filter, $J\ge7$]
\label{lem:tr-mod12}
For $J\ge7$, the inequality $\mu_J(M)\le11/15$ implies $12\mid M$.
Thus only multiples of $12$ need be tested as witnesses in each interval
$[n,2n-2]$; no potentially uncovered rank is discarded.
\end{lem}
\leannote{Lean: \lproof{ErdosProblems/Erdos257/PaperCompleteR21/TruncatedRungWitnessHorizon.lean}{544}{paper_mod_twelve_filter}.}
\begin{proof}
If $M$ is odd, the $q=2,3$ terms give at least $2/3+1/7>11/15$.
If $M\equiv2\pmod4$, the $q=2,3,4$ terms give at least
$1/3+1/7+4/15=26/35>11/15$.  Finally, if $4\mid M$ but $3\nmid M$,
the $q=2,3,4,6$ terms give at least
$1/3+2/7+1/15+4/63=236/315>11/15$.
Each case contradicts the assumed upper bound.
\end{proof}

\begin{thm}[Finite decision procedure]
\label{thm:tr-finite-decision}
Call $n\in[4,L_J/2]$ \textbf{bad} if no $M\in[n,2n-2]$ has $\mu_J(M)\le\tfrac{11}{15}$, and set
$B(J):=\max(\mathrm{bad}\cup\{3\})$. Then $\mathrm{HalfRung}(J)$ holds iff the greedy orbit for
$1/2$ under weights $w_n^{(J)}$ survives every rank from $2$ through
$B(J)$.  This is a finite exact decision procedure.
\end{thm}
\leannote{Lean: \lproof{ErdosProblems/Erdos257/PaperCompleteR21/TruncatedRungGreedyDecision.lean}{896}{paper_rung_finite_decision}.}
\begin{proof}
Survival at all ranks plainly implies survival through $B(J)$.
Conversely, every later rank at most $L_J/2$ has a witness by the
definition of $B(J)$, and every larger rank has one by
Corollary~\ref{cor:tr-half-lcm}.  Witness exclusion rules out a fatal gap
at all these ranks.  The strict-tail criterion then gives membership.
For $J=2$, the same proof works because $\mu_2(M)\le2/3$ for every $M$;
the window is empty and $B(2)=3$.
\end{proof}

\begin{rem}[Verified truncations and the unbounded-index question]
The finite decision procedure verifies $\mathrm{HalfRung}(J)$ for every
$3\le J\le22$.  The supplied checker uses exact integers: it verifies
witness intervals covering the ranks beyond $B(J)$, proves that the
listed bad ranks have no witness, and tests the greedy remainders through
$B(J)$.  In particular, it reproduces $B(19)=1008$ and
$B(20)=B(21)=B(22)=1530$.
The certificate data and an independent verifier are in
\nolinkurl{audit/final-proof-reading/truncation/}.  These are finite
computer-assisted proofs using the ordinary lemmas above, not Lean
proofs and not a test of the full weights through all ranks.

If $A$ represents $1/2$ with the $J$-truncated weights, then
\[
 0\le X_A(2)-\tfrac12
 \le\sum_{q>J}\frac{2^{-2q}}{1-2^{-q}}
 \le2\sum_{q>J}4^{-q}=\frac23\,4^{-J}.
\]
Consequently, such representations for an unbounded sequence of $J$
would put $1/2$ in the closed set $\mathcal A$.  Its representing support
would be infinite, since finite Mersenne sums have odd denominator
(Lemma~\ref{lem:no-ties}); Lemma~\ref{lem:tr-parity} separately excludes
finite support for each truncation.  The unbounded-index assertion is
still unproved.  Twenty verified values of $J$ do not establish it.
\end{rem}
% ; - part a257_p3 ; -
\section{Which hypotheses remain unproved}

This section compares the stated results with the remaining hypotheses. Parts I--II assembled the corpus; here every
result in it is tested against three exact open obligations and rendered as a usable premise:
what it yields, the exact mismatch against the obligation, a classified gap kind
(\texttt{hypothesis\_strength}, a different parameter range, \texttt{quantifier\_order},
a different expression of the same condition, or \texttt{multiple}), and the exact auxiliary statement that would
close it. Nothing here decides \#249 or \#257; the corpus's own moderator standard applies
throughout: a conditional reduction is a result with a hypothesis,
not a proof of that hypothesis, and a finite verified list is not a
cofinal supply.

\subsection{The three \#257 obligations, stated exactly}

\begin{defn}[257-reset]
For every upper or middle reset row $r\ge10$, require
$\Delta_{r+1}^2>2^{r+5}$, where
$\Delta_{r+1}=\mathrm{rem}(r+1)-2^{r+1}$.
The finite range $10\le r<31$ must be covered as well as the tail.
Section~\ref{sec:o4} gives the exact conditional implication from this
hypothesis.  A proposed bound on right-run lengths is a different
statement until the pulse corrections and endpoint conventions are
matched.  Finite observations establish neither failure of every
constant run bound nor logarithmic asymptotic growth.
\end{defn}

\begin{defn}[257-cofinal-rows]
$\forall N\ \exists n \ge N$ with \lean{ExactLocalMersenneHalfRow}{BooleanMobiusCofinalExactRows.lean:31} $n$
;  i.e.\ cofinally many finite prefixes whose local Mersenne quotient lands exactly on the
half-value residue class. Consumed unconditionally by
\lean{half\_mem\_mersenneAchievementSet\_of\_cofinalExactLocalRows}{BooleanMobiusCofinalExactRows.lean:71}.
No coherence between the witnessing rows is demanded and none should be assumed: the rows may
be, and in the corpus's own words, mutually incompatible ;  only their endpoint lengths need
tend to infinity. Requiring compatibility would add a condition on the witnesses.
Whether that condition can be obtained from this existence statement
needs a separate extraction argument; it is not a strict logical
comparison merely because the witnesses are described differently.
\end{defn}

\begin{defn}[257-universal]
The original question is the base-$2$ statement:
$\sum_{n\in A}(2^n-1)^{-1}$ is irrational for every infinite
$A\subseteq\mathbb N_{\ge1}$.  Asking the same for every integer
base $b\ge2$ is a stronger all-base extension, not the definition of
Problem~\#257 itself.
\end{defn}

\subsection{Reading a row}

Each entry identifies the available result, the assumptions needed to
apply it, and the further estimate still required.  A source link supports
the stated declaration, not every proposed consequence in the discussion.
Ordinary deductions and finite computations are identified separately.

\subsection{Reset estimates and their missing hypotheses}

\begin{obs}[Remainders along a right-branch run] \ev{Lean}
\lean{upperOrMiddle\_within\_of\_two\_pow\_deficit}{HalfCylinderResetDeficitEscape.lean:261}.
For integers $s\ge5$ and $k\ge0$, the cited theorem assumes
$\mathrm{rem}(s)\le2^s$ and
$2^k(2^s-\mathrm{rem}(s))\ge2^s$ and gives an upper or middle
transition at some $t\in[s,s+k]$.  Its literal contrapositive concerns
\emph{all} $k+1$ transition indices in that closed interval.

For a run of exactly $k\ge1$ right transitions, at indices
$s,\ldots,s+k-1$, a direct recurrence gives the sharper conventionally
indexed statement.  Put $D_j=2^{s+j}-\mathrm{rem}(s+j)$ and
$p_j=p^-_{s+j}\ge0$.  The right-branch identity yields
\[
 D_{j+1}=4D_j+p_j+4,\qquad
 4^kD_0<D_k\le2^{s+k}.
\]
The strict inequality uses $k\ge1$, and the last bound uses
$\mathrm{rem}(s+k)\ge0$.  Hence $D_0<2^{s-k}$.
In particular, $D_0\ge1$ excludes $k\ge s$ consecutive right
transitions.  A stronger lower bound on $D_0$ shortens that run, but
must be established at the starting row of the run.  Applying this
argument after a reset requires the post-reset index and the
\emph{deficit-side} condition; it does not prove the absolute reset
bound on both signs.
\end{obs}

\begin{obs}[The affine excess recurrence] \ev{Lean}
\label{obs:affine-excess-recurrence}
\lean{affineRightExcess\_exactIterate}{HalfCylinderMiddleCarryLowerBound.lean:3320},
\lean{affineRightExcess\_charge\_lt\_scaledInitial}{HalfCylinderMiddleCarryLowerBound.lean:3350},
\lean{eventualRightTail\_excess\_exactIterate}{HalfCylinderMiddleCarryLowerBound.lean:3361}
(recurrence \lean{rightBranch\_excess\_succ\_eq}{HalfCylinderMiddleCarryLowerBound.lean:2497}).
For $k$ right transitions from row $S\ge5$, let
$X_j=\mathrm{rem}(S+j)-2^{S+j}$ and $p_j=p^-_{S+j}$.
The cited identity is
\[
 X_0=\frac{X_k}{4^k}+\sum_{j=0}^{k-1}\frac{p_j+4}{4^{j+1}}.
\]
The source also gives $0\le p_j\le2(S+j-2)$.  Thus the following
consequences are ordinary deductions from the identity:
\[
 \sum_{j=0}^{k-1}\frac{p_j+4}{4^{j+1}}
 \le\sum_{j\ge0}\frac{2(S+j)}{4^{j+1}}
 =\frac{2S}{3}+\frac29.
\]
If the terminal row satisfies $\mathrm{rem}(S+k)<2^{S+k+1}$, then
\[
 X_0<2^{S-k}+\frac{2S}{3}+\frac29.
\]
Only this terminal bound is needed; no bound at every intermediate row
is required.  If instead the largest omitted rank $d$ at the terminal
row satisfies $2(S+k)<3d$, the next entry gives the different bound
\[
 X_0<2^{S-k}+\frac{2\,4^{S+k-d}+4}{3\,4^k}
                +\frac{2S}{3}+\frac29.
\]
The extra term must be retained.  Lateness of an omitted rank is not
known to be equivalent to $\mathrm{rem}(s)<2^{s+1}$; the cited estimate
alone does not imply that sharper ceiling.  Neither upper bound supplies
the missing lower bound on an actual reset deviation.
\end{obs}

\begin{obs}[A bound at later rows] \ev{Lean}
\lean{three\_mul\_remainder\_lt\_exactLateGap}{HalfCylinderResetDeficitEscape.lean:117}.
For $s\ge5$, if $d$ is the largest omitted rank and $2s<3d$,
the cited theorem gives
\[
 \mathrm{rem}(s)<2^{s+1}+\frac{2\,4^{s-d}+4}{3}.
\]
Since $s-d<s/3$, its additive error is at most
$(2\cdot2^{2s/3}+4)/3$, hence $O(2^{2s/3})$.
The undivided error $2\,4^{s-d}+4$ need not be less than
$2^{2s/3}$: already $s=7,d=5$ gives $36>2^{14/3}$.
These numbers test the numerical estimate, not an asserted pair on the
greedy orbit.  The theorem supplies an upper bound only under the
largest-omitted-rank and lateness assumptions.  When it is used in the
preceding recurrence, its additive error remains part of the result.
\end{obs}

\begin{obs}[Exclusion of equality at a branch boundary] \ev{Lean}
\lean{seamGreedyFloorZ\_ne\_takeThreshold}{HalfCylinderMiddleCarryLowerBound.lean:1665}, with
\lean{halfGreedy\_skipped\_endpoint\_trichotomy}{HalfCylinderSkippedEndpointClassifier.lean:246}.
The first cited theorem excludes equality between the corrected
real-valued quotient expression and the terminal \emph{take threshold}:
such equality would give a finite Mersenne sum equal to $1/2$, contrary
to denominator parity.  The second classifies the signed integer
position of a real greedy prefix at a skipped rank; its statement
explicitly permits the zero-position case.

Neither statement identifies its boundary with
$\mathrm{rem}(r)=2^r$.  They therefore do not, by themselves, prove
$|\mathrm{rem}(r)-2^r|\ge1$.  Integrality gives that inequality only
after nonvanishing of this particular deviation has been proved.
An estimate of the form $2^{r/2+o(r)}$ would not settle the required
constant either: $2^{r/2}/r$ has that form but is smaller than
$2^{(r+5)/2}$.  The reset hypothesis requires its stated constant,
strict inequality and post-reset index, not merely an exponent
approaching $1/2$.
\end{obs}

\begin{obs}[A sharper test for a fatal greedy gap] \ev{Lean}
\lean{skipSafe\_of\_two\_mul\_le\_three\_mul}{Erdos249257/HalfGreedyFatalGap.lean:107},
\lean{three\_mul\_lt\_two\_mul\_of\_fatal}{Erdos249257/HalfGreedyFatalGap.lean:146},
\lean{two\_le\_of\_fatal}{Erdos249257/HalfGreedyFatalGap.lean:161},
\lean{three\_le\_of\_fatal\_of\_odd}{Erdos249257/HalfGreedyFatalGap.lean:173}.
Let $k\ge1$ be a skipped rank and write its positive rational
remainder in lowest terms as $\rho=u/(2L)<w_k$, with $u,L$ odd and
positive.  Set $a=2L-(2^k-1)u>0$.  The cited sufficient condition
$2u\le3a$ guarantees $\rho<R_k$ and improves the simpler condition
$u\le a$.  Its contrapositive gives $3a<2u$ at a fatal skip,
and hence $u\ge3$ when $u$ is odd.

This is not the exact threshold.  With
$\theta_k=R_k^{-1}-(2^k-1)$, direct rearrangement gives
\[
 \rho<R_k\quad\Longleftrightarrow\quad u\theta_k<a.
\]
The tail estimate gives $\theta_k<2/3$ and $\theta_k\to2/3$;
$2u\le3a$ is therefore sufficient, not necessary.
For example, at $k=2$ the local rational state $\rho=77/282$ has
$u=77$, $L=141$, $a=51$, so $2u>3a$, yet
$\sum_{j=3}^{12}(2^j-1)^{-1}>77/282$ proves $\rho<R_2$.
This is a local state, not an asserted point of the half-greedy orbit.
Improving a bounded tail ratio alone does not provide a lower bound
growing exponentially in $k$ on an arithmetic numerator; such a bound
would require additional information about the actual prefix.
\end{obs}

\begin{obs}[Exclusion of a linear-width band at rows 13--30] \ev{Lean}
\lean{seamUpperResetDyadicBandEscape\_through\_thirty}{HalfCylinderUpperResetBandCertificates.lean:78}.
\textbf{Conclusion.} At every upper reset row $d$, $13 \le d \le 30$, the reset charge $E_d =
4\cdot\mathrm{overshoot}_d + \mathrm{abovePulse}_d$ avoids the linear-width band
$(2^{d-j+1}-2(d+j),\ 2^{d-j+1}]$ for every $j \le d$ ;  the finite segment of
\lean{SeamUpperResetDyadicBandEscape}{HalfCylinderMiddleCarryLowerBound.lean:4566}, which by
\lean{half\_mem\_mersenneAchievementSet\_of\_upperResetDyadicBandEscape}{HalfCylinderMiddleCarryLowerBound.lean:4790}
delivers the whole \#257-false payoff.
\textbf{Parameter range.} proved at 18 explicit rows by \texttt{interval\_cases}
plus per-row \texttt{decide+kernel}; the obligation needs all $d \ge 13$. The excluded band has width $O(d)$, whereas the reset condition uses
a width of order $2^{(r+5)/2}$.  This is a comparison of numerical scales,
not a proof that one hypothesis on the orbit implies the other.
\textbf{Missing step.} A structural lower bound on the distance from $E_d$ to the nearest dyadic
power at or above it, of width $2(d+j_0)$ at the single critical index $j_0$ ;  already collapsed
from $\forall j$ to one inequality by
\lean{HalfUpperResetCriticalBand.dyadicBandEscape\_iff\_exists\_critical}{HalfUpperResetCriticalBand.lean:108}.
Only $O(d)$ separation is required, not sqrt-of-scale.
\end{obs}

\begin{obs}[Middle predecessors and right-branch runs] \ev{Lean}
\lean{seamMiddleBranch\_nextRemainder\_ge\_row}{HalfCylinderMiddleCarryLowerBound.lean:2655},
\lean{seamMiddleBranch\_nextRemainder\_ge\_expBarrier}{HalfCylinderMiddleCarryLowerBound.lean:2680},
\lean{SeamUpperResetDyadicBandEscape.remainder\_ge\_row}{HalfCylinderMiddleCarryLowerBound.lean:4641}.
\textbf{Conclusion.} Unconditional, uniform in $s\ge5$: a middle (M) branch forces $s+1 \le
\mathrm{rem}(s+1)$, using no lower bound on the source remainder; with the row-scale source
bound this upgrades to $2^{s+1}+(s+1) \le \mathrm{rem}(s+1)$. The minimal-counterexample argument
at line 3606 then shows any row whose remainder is smaller than its index must have an upper-reset (U) ancestor, since M-ancestors
and R-runs are both excluded.
\textbf{Scope of the hypothesis.} two of the three ancestor channels (M, R) are
closed unconditionally at all scales; only the U channel remains, and it needs only the
$O(d)$-wide band escape of F5b above ;  sqrt-of-dyadic-scale escape is being demanded by the
257-reset statement where linear-width escape at U-resets alone suffices for this downstream
route.
\textbf{Missing step.} The single remaining channel ;  $\forall d \ge 13$ with successor carries, the
reset charge avoids the width-$2(d+j_0)$ band at its critical dyadic index. Then
\texttt{remainder\_ge\_row}, hence \texttt{SeamMiddleProducerRowEscape}, hence $1/2 \in
\ensuremath{\mathcal A}$, all already in Lean.
\end{obs}

\begin{obs}[A linear lower bound at middle transitions] \ev{Lean}
\lean{SeamMiddleProducerRowEscape}{HalfCylinderMiddleCarryLowerBound.lean:1769},
\lean{half\_mem\_mersenneAchievementSet\_of\_middleProducerRowEscape}{HalfCylinderMiddleCarryLowerBound.lean:2371}
(via \texttt{.toCardEscape}, \texttt{.toTailEscape}, \lean{}{HalfCylinderMiddleCarryLowerBound.lean:1510}).
\textbf{Conclusion.} If at every late middle row ($s \ge 13$, no carry, middle branch)
$s \le \mathrm{rem}(s)$, then $1/2 \in \ensuremath{\mathcal A}$. Module docstring: ``far
weaker than an exponential bound, but stronger than the already exposed square-root-scale carry
target.''
\textbf{Scope of the hypothesis.} this hypothesis asks for a \emph{linear-in-row}
lower bound on the remainder itself; the obligation supplies an \emph{exponential} lower bound
on the deviation from $2^s$ at reset rows. For unrestricted numerical values, a large absolute deviation from
$2^s$ permits $\mathrm{rem}(s)=0$, which violates the linear lower bound.
That observation is not a counterexample within the actual recurrence,
so incomparability of the two orbit hypotheses is not established.
\textbf{Missing step.} $\forall s \ge 13$ at a middle transition, $s \le \mathrm{rem}(s)$.
The upper-reset band condition described above is one sufficient input
to \texttt{remainder\_ge\_row}; no converse is asserted.
\end{obs}

\begin{obs}[A square-root lower bound in the row index] \ev{Lean}
\lean{SeamMiddleProducerSqrtEscape}{HalfCylinderLastProducerContradiction.lean:84},
its implication to the exact tail bound (\texttt{.toTailEscape}, line~98),
\lean{half\_mem\_mersenneAchievementSet\_of\_middleProducerSqrtEscape}{HalfCylinderLastProducerContradiction.lean:458}.
\textbf{Conclusion.} If at every late middle row $2\sqrt{2s+2}+4 < \mathrm{producerCarry}(s)$, then
$1/2 \in \ensuremath{\mathcal A}$; \texttt{producerCarry} equals the middle coordinate
$4\,\mathrm{rem}-\mathrm{belowPulse}-4$
(\lean{}{HalfCylinderProducerLowerBound.lean:150}), so the requirement is essentially
$\mathrm{rem}(s) \gtrsim (\sqrt{s}+\mathrm{belowPulse})/4$.
\textbf{Comparison.} this is ``square root'' of the \emph{row index}
$s$ (a supportCoeff tail majorant), not square root of the \emph{dyadic scale} $2^{(r+5)/2}$. The
two sqrt's in this record are exponentially far apart and easy to conflate; this hypothesis is row-index bound.  The cited conditional theorem uses it to rule out a
last non-right transition.  An implication between the different orbit
hypotheses would require a further argument.
\textbf{Missing step.} $\forall s \ge 13$ at a middle transition, $4\,\mathrm{rem}(s) -
\mathrm{belowPulse}(s) - 4 > 2\sqrt{2s+2}+4$ ;  a row-scale, not dyadic-scale, lower bound on the
middle coordinate.
\end{obs}

\begin{obs}[The remaining interval of possible integer values] \ev{Lean}
\lean{upperProducer\_not\_last}{HalfCylinderLastProducerContradiction.lean:299},
\lean{middleProducer\_neg\_three\_not\_last}{HalfCylinderLastProducerContradiction.lean:364},
\lean{finalMiddleCell\_neg\_three\_not\_last}{HalfCylinderFinalMiddleCellEscape.lean:587},
\lean{middleProducer\_allRight\_landingExcess\_window}{HalfCylinderMiddleCarryLowerBound.lean:2547},
\lean{middleThenAllRight\_landingExcess\_two\_le}{HalfCylinderMiddleCarryLowerBound.lean:2319}.
\textbf{Conclusion.} Under the hypothesis of a last non-right transition
at row $D\ge13$, the upper case is excluded.  In the middle case put
\[
 X_{D+1}:=\mathrm{rem}(D+1)-2^{D+1}
          =4\,\mathrm{rem}(D)-p_D^-=C_D+4.
\]
The cited bounds give $2\le X_{D+1}<2\sqrt{2D+2}+8$.
This permits nonnegative $C_D$ as well as the negative exceptional
values $-2,-1$; the latter do not exhaust the remaining interval.
\textbf{Scope of the hypothesis.} the residual bad set is an explicit
$O(\sqrt D)$-wide integer window at each candidate last non-right transition row, not an interval at dyadic
scale. The reset hypothesis would exclude this window.  Its width is much
smaller than $2^{(r+5)/2}$, but width alone does not measure the difficulty
of proving that the actual sequence avoids it.
\textbf{Missing step.} Exclude every integer
\[
 2\le X_{D+1}\le\bigl\lceil2\sqrt{2D+2}\bigr\rceil+7
\]
at a hypothetical last middle transition $D\ge13$.
The ceiling is necessary when converting the strict real upper bound
to an integer interval.  Excluding only $C_D=-2,-1$ does not suffice.
\end{obs}

\begin{obs}[Excluding three integer values] \ev{Lean}
\lean{SeamTwoSidedDyadicCellEscape}{HalfCylinderMiddleCarryLowerBound.lean:4374},
\lean{SeamTwoSidedDyadicCellEscape.step}{HalfCylinderMiddleCarryLowerBound.lean:4398},
\lean{SeamTwoSidedDyadicCellEscape.twoSided}{HalfCylinderMiddleCarryLowerBound.lean:4448}.
Theorem~\ref{thm:two-sided-dyadic} states the two needed hypotheses
and their row ranges explicitly.  Excluding $-3,-2,-1$ is required
only on middle branches; the overshoot inequality is required only
on right branches where $o_s\le2^s$.  The resulting invariant bounds
at least one of the two distances from the target, not necessarily
the lower remainder.  The exclusion of $-3$ in the special
last-middle-then-all-right configuration does not establish the
unrestricted middle-branch hypothesis.  Both universal assumptions
remain to be proved for the actual sequence.
\end{obs}

\begin{obs}[An upper bound on the remainder] \ev{Lean}
\lean{SeamGreedyRemainderGapBound}{HalfCylinderSeamLimit.lean:192},
\lean{half\_mem\_mersenneAchievementSet\_of\_seamRemainderGapBound}{HalfCylinderSeamLimit.lean:303}
and the cofinal implication \lean{}{HalfCylinderSeamLimit.lean:181}.
The estimate $\mathrm{rem}(s)<2^{s+1}$ for every $s\ge6$ would
give $\mathrm{rem}(s)/4^s\to0$ and hence $1/2\in\mathcal A$.
Theorem~\ref{thm:seam-limit} identifies the normalized limit
unconditionally; requiring its value to be zero, even only along
one unbounded sequence, is equivalent to membership.  No strict
separation of these conditions on the actual sequence is proved.
The upper estimate is also the missing ceiling in the excess-side
right-run argument.  The existing bound on rows satisfying $2s<3d$,
where $d$ is the largest omitted rank, still requires that row
condition wherever it is applied.  The finite checks do not supply
it at every later row.
\end{obs}

\begin{obs}[Reported computation through row 200{,}000] \ev{Cert}
\textbf{Conclusion.} Exact integer arithmetic to row 200{,}000 ;  branch counts $R\!:\!M\!:\!U =
100197\!:\!49899\!:\!49898$; the obligation \emph{fails only at} $r=7$; maximum R-run $19$ at row
$158{,}096$ against the required threshold $\approx (r-3)/2 \approx 79{,}000$ (margin
$\sim\!4000\times$). Independently: 1209 resets checked to row 2500, zero classification
anomalies, exactly two crossing cells ever ($s{=}7,d{=}5$ and $s{=}10,d{=}7$), both below hypothesis
scope.
\textbf{Parameter range.} fixed-scale exhaustive verification versus a statement
at every reset $r \ge 31$. This is the canonical difference between a finite and an unbounded parameter range this programme warns against ;  a
finite verified list is not a cofinal supply, however large the margin.
\textbf{Missing step.} Nothing computational; only a theorem. The certificate's value is fixing the
true constant ($r \ge 31$, sole violation $r=7$) so the theorem's scope can be stated exactly.
\end{obs}

\begin{obs}[Structure of resets after long runs] \ev{Math}
\textbf{Conclusion.} At a dangerous reset $r$ ($|\mathrm{rem}(r{+}1)-2^{r+1}| \le 2^{(r+5)/2}$) whose
preceding reset $r_0$ has pure R-run length $L' = r-r_0-1 > (r+5)/4$, the deviation
$w_{r_0+1}$ is pinned to at most two adjacent integers determined by the divisor-pulse stream
alone; a chain of $n$ dangerous resets forces $n{-}1$ nested exact integer-coincidence towers.
\textbf{Scope of the hypothesis.} applies only to dangerous resets preceded by a
long pure R-run ($L' > (r+5)/4$) ;  an isolated dangerous reset after a short run is untouched.
Pinning to two integers is not yet a contradiction: no theorem excludes those two values. Prose,
not Lean.
\textbf{Missing step.} (i) a congruence or pulse-stream obstruction showing the two pinned integer
values are unattainable, and (ii) removal of the long-preceding-run hypothesis, or a proof that
dangerous resets must enter chains to which the exclusion applies.
Excluding the two pinned values addresses only resets satisfying the
long-run hypothesis; covering all remaining resets is a separate
requirement.  Neither step alone proves a global exclusion.
\end{obs}

\begin{obs}[The sign condition and growth of the margin] \ev{Cert}
\textbf{Conclusion.} At all 1209 observed resets, M $\Rightarrow$ deviation $>0$ and U $\Rightarrow$
deviation $<0$, with zero exceptions; the margin over the $2^{(r+5)/2}$ threshold grows from
1.11 bits at $r=14$ to $\sim$1246 bits at $r=2500$. The record isolates this as a separate,
possibly more tractable target than full sqrt-escape.
\textbf{Different restrictions.} The upper sign follows from the
positive reset charge: $w_{r+1}=-(4o_r+p_r^+)<0$.
The middle sign is equivalent to $4\mathrm{rem}(r)>p_r$ and needs
further information about the orbit.  Neither sign implies a
square-root lower bound on the absolute deviation.
\textbf{Missing step.} Prove that middle inequality and then a
quantitative margin on each branch.  The exact branch calculation
is given in Section~\ref{sec:invent-R1}; it is not a substitute for
the magnitude estimate.
\end{obs}

\begin{obs}[Comparison of finite truncation lengths] \ev{Cert}
The exact finite procedure in
Theorem~\ref{thm:tr-finite-decision} verifies the twenty truncations
$3\le J\le22$.  The proof combines finite integer certificates with a
witness bound valid beyond the finite search horizon; it does not infer
an infinite statement from a long finite run.  These certificates are
reproduced in \nolinkurl{audit/final-proof-reading/truncation/}.
Representations for an unbounded sequence of $J$ would yield a
representation of $1/2$ for the full weights by the compactness estimate
following that theorem.  That unbounded-index premise remains unproved.
The size of the least-common-multiple horizon affects this finite
algorithm, but is not an impossibility theorem for other arguments.
\end{obs}

\subsection{Finite quotient sums and unbounded depths}

\begin{obs}[Exact rows from a sharp capacity bound] \ev{Lean}
\lean{exists\_exactRowStrictUpperFill\_of\_skippedCoreSharpCapacity}{BooleanMobiusSkipRow.lean:238},
\lean{exactLocalMersenneHalfRow\_two\_mul\_sub\_two\_of\_skippedCoreSharpCapacity}{BooleanMobiusSkipRow.lean:332}.
\textbf{Conclusion.} For any $c \ge 4$ and any finite $D \subseteq [2,c)$ with
$\ensuremath{V}\,D < 1/2$ and $\ensuremath{S}\,D\,1\,(2c{-}2) <
2^{c-2}$: an exact row $E$ at endpoint $2c{-}2$ with $D \subseteq E$, every new rank strictly
above $c$.  No additional crossing condition or identification with a greedy prefix is needed.
\textbf{Scope of the hypothesis.} needs a per-$c$ input
$\ensuremath{S}\,D\,1\,(2c{-}2) < 2^{c-2}$, and the corpus only ever supplies that
input through the crossing-deficit route, which
\lean{eq\_halfGreedyPrefixSupport\_of\_critical\_crossing}{BooleanMobiusCriticalCapacityCofinal.lean:50}
then identifies the core with
$D=\mathrm{halfGreedyPrefixSupport}(c{-}1)$, under its critical-crossing
premises.  The fill implication does not require that identification;
this difference in its formulation is not a proof of strictness
between the associated cofinal existence statements.
\textbf{Closes it ($\star$):} $\forall N\ \exists c \ge N,\ \exists$ finite $D \subseteq [2,c)$
with $\ensuremath{V}\,D<1/2$ and $\ensuremath{S}\,D\,1\,(2c{-}2) <
2^{c-2}$ ;  equivalently $\ensuremath{C}\,{\uparrow}D\,(2c{-}3) < 2^{c-2}$. No
coherence between the $D$'s, no requirement that $D$ be a greedy prefix.
\end{obs}

\begin{obs}[A conditional induction for unbounded depths] \ev{Lean}
\lean{SkippedCoreCriticalQuotientSupply}{BooleanMobiusCriticalCapacityCofinal.lean:30} $\to$
\lean{exists\_laterProtectedExactLocalMersenneRow}{BooleanMobiusCriticalCapacityCofinal.lean:1139}
$\to$
\lean{cofinalExactLocalMersenneHalfRows\_of\_criticalQuotientSupply}{BooleanMobiusCriticalCapacityCofinal.lean:1309}
(\lean{ProtectedExactLocalMersenneRow}{BooleanMobiusCriticalCapacityCofinal.lean:1087}; seed
line 1023; strict progress \lean{}{BooleanMobiusExactRowCrossing.lean:233}).
\textbf{Conclusion.} Under the sharp-capacity supply hypothesis below,
the induction starts from the depth-six seed and constructs exact rows
at unbounded depths.  A protected row is either doubled while its real
value stays below $1/2$, or extended using its first crossing rank $e$.
In the latter case protection gives $2e-2$ strictly larger than the old
endpoint.
\textbf{Scope of the hypothesis.} the recycle branch consumes
\texttt{SkippedCoreCriticalQuotientSupply} whenever recycling is required.
Protection (endpoint $<2\cdot$cutoff and new ranks $>$cutoff) makes that
extension strictly increase the endpoint.  The induction establishes the
conditional implication, not the supply hypothesis.
\textbf{Missing step.} The sharp inequality at one crossing rank per step, canonical form
\lean{HalfGreedySkippedCriticalQuotientSupply}{BooleanMobiusCriticalCapacityCofinal.lean:178}:
$\forall c \ge 4$ skipped by the rational half-greedy orbit, $2^{(2c-2)-1} \le
\ensuremath{Q}(\mathrm{insert}\,c\,(\mathrm{halfGreedyPrefixSupport}(c{-}1)))\,(2c{-}2)$.
\end{obs}

\begin{obs}[Unconditional capacity, off by one bit] \ev{Lean}
\lean{localBinarySuffix\_two\_mul\_sub\_two\_lt\_upperHalfCapacity}{BooleanMobiusSkippedCoreExactRow.lean:123}
(word form line 228).
\textbf{Conclusion.} Unconditionally, $\forall c \ge 4$ and every below-half core $D \subseteq [2,c)$
with deficit $< \ensuremath{w}\,c$: $\ensuremath{S}\,D\,1\,(2c{-}2) <
2^{c-1}$. Uniform in $c$, no table, no case split.
\textbf{Parameter range.} one bit short. The strict-upper fill needs
$<2^{c-2}$, this gives $<2^{c-1}$. The proof derives $A < 2^{c-2}+|D|$ then discards
$|D| \le c-2 \le 2^{c-2}$ ;  the loss is purely additive. The true unproved statement is only that
$A$ avoids the linear-width band $[2^{c-2},\,2^{c-2}+c-3]$, width $c{-}2$ inside a range of size
$2^{c-2}$.
\textbf{Missing step.} A linear-width dyadic-band anti-concentration: at cofinally many crossing
ranks $c$, $\ensuremath{S}\,D\,1\,(2c{-}2) \notin [2^{c-2},\,2^{c-2}+c-3]$. Same band
shape as the earlier linear-band conditions.  The reset condition
uses a different coordinate and parameter, with width $2^{(r+5)/2}$;
no implication between the hypotheses follows merely from this
comparison of widths.
\end{obs}

\begin{obs}[Doubling extension, fails at the only recorded seed] \ev{Lean}
\lean{exists\_exactRowStrictUpperExtension\_two\_mul\_sub\_one\_of\_exact\_below}{BooleanMobiusExactRowDoubling.lean:185},
line 296; rank-two discharge
\lean{two\_mem\_of\_exact\_localMersenneQuotient}{BooleanMobiusExactRowRankTwo.lean:23}, line 108.
\textbf{Conclusion.} $\forall n \ge 6$: a below-half exact row at $n$ produces an exact row at
$2n{-}1$, support-extending, all new ranks $>n$. The only unconditional endpoint-doubling
required input in the corpus.
\textbf{Scope of the hypothesis.} below-halfness is not preserved and
demonstrably fails at the first step. For the seed $D=\{2,3,6\}$ at $M=11$ the core alone
contributes $2/3+4/7+32/63 \approx 1.746 > 1$, so the doubled row is strictly \emph{above} half.
Iterating the cheap arm is therefore impossible from the only recorded seed; the chain is forced
into the capacity-gated recycle arm at step one.
\textbf{Missing step.} Cofinally many below-half exact rows ;  $\forall N\ \exists n \ge N\ \exists D$
exact at $n$ with $\mathrm{localFractionMass}\,D\,n < 1$. A single such row at each of cofinally
many $n$ makes doubling unnecessary; a self-reproducing one closes everything.
\end{obs}

\begin{obs}[Consecutive-skip precritical bound] \ev{Lean}
\lean{halfGreedy\_precriticalSuffix\_lt\_of\_next\_skip}{BooleanMobiusCriticalCapacityCofinal.lean:682}
(feeds \lean{halfGreedySkippedCriticalQuotientSupply\_of\_precriticalSuffix}{BooleanMobiusCriticalCapacityCofinal.lean:1001},
then lines 970/986).
\textbf{Conclusion.} Unconditional: if $c \ge 6$ and both rank $c$ and rank $c{+}1$ are skipped by
the rational half-greedy orbit, then $\ensuremath{S}(\mathrm{halfGreedyPrefixSupport}
(c{-}1))\,1\,(2c{-}3) < 2^{c-3}$, doubling into sharp capacity at $c$, hence an exact row at
$2c{-}2$. No supply hypothesis anywhere.
\textbf{Scope of the hypothesis.} needs two \emph{consecutive} skipped ranks.
The residual open set is exactly the skip-then-take rows:
\lean{halfGreedySkippedPrecriticalSuffixSupply\_iff\_preTake}{BooleanMobiusCriticalCapacityCofinal.lean:838}
proves the whole hypothesis equivalent to
\lean{HalfGreedyPreTakePrecriticalSuffixSupply}{BooleanMobiusCriticalCapacityCofinal.lean:200},
with $c=4,5$ already discharged by \texttt{decide}.
\textbf{Missing step.} Cofinally many $c \ge 6$ at which two consecutive ranks $c,c{+}1$ are both
skipped by the rational half-greedy orbit ;  or dually, the pre-take residual at cofinally many
skipped-then-taken ranks. Either one closes the obligation with no further input.
\end{obs}

\begin{obs}[Bounded-gap precritical bound] \ev{Lean}
\lean{halfGreedy\_precriticalSuffix\_lt\_of\_future\_skip\_after\_takenBlock}{BooleanMobiusCriticalCapacityCofinal.lean:603}
(with \lean{precriticalCrossingTax\_of\_futureThreshold}{BooleanMobiusCriticalCapacityCofinal.lean:472},
\lean{sub\_two\_le\_two\_pow\_sub\_four}{BooleanMobiusCriticalCapacityCofinal.lean:432}).
\textbf{Conclusion.} Unconditional and uniform in \emph{both} parameters: skip at $c$, takes at
$c{+}1,\dots,c{+}t{-}1$, skip at $c{+}t$, with $0<t\le c-3$ and $c-2\le 2^{c-t-3}$
$\Rightarrow$ precritical suffix bound at $c$ $\Rightarrow$ sharp capacity $\Rightarrow$ exact
row at $2c{-}2$. Handles any skip gap $t \lesssim c-3-\log_2(c-2)$, not just $t=1$.
\textbf{Scope of the hypothesis.} the only missing input is an orbit-level
skip-gap bound ;  that the next skipped rank after $c$ occurs at $c{+}t$ with $c-2 \le
2^{c-t-3}$. The proof is fully uniform in $(c,t)$, every constant explicit, no table.
\textbf{Missing step.} For cofinally many skipped ranks $c$ of the rational half-greedy orbit, the
next skipped rank is at most $c+(c-3-\lceil\log_2(c-2)\rceil)$.
The finite skip-density observation does not imply this pointwise gap
bound at unbounded ranks.  Establishing the latter would supply the
hypothesis of the stated conditional argument.
\end{obs}

\begin{obs}[Frozen margin, ineffective horizon] \ev{Lean}
\lean{exists\_greedyHalfFrozenMargin\_nonneg\_of\_excess\_neg}{HalfCylinderFiniteShadow.lean:1198},
\lean{exists\_greedyHalfFrozenMargin\_nonneg\_iff\_excess\_neg}{HalfCylinderFiniteShadow.lean:1229}
(margin def line 1021, recurrence line 1036, monotonicity line 1094).
\textbf{Conclusion.} At every dyadically safe rank $k>0$, $\ensuremath{F}\,k\,J \ge
0$ for \emph{some} horizon $J$ ;  the finite coefficient window eventually covers the centred
carry. Nonnegativity is monotone up in $J$.
\textbf{Parameter range.} the horizon comes from a non-effective limit argument
(\texttt{tendsto\_finiteCoeffWindow\_atTop} plus an existence extraction), so $J$ is unbounded.
The exact-row route needs the crossing by the specific horizon $J=c-3$
(\lean{halfGreedy\_precriticalSuffix\_lt\_iff\_frozenMargin\_nonneg}{BooleanMobiusCriticalCapacityCofinal.lean:866},
$k=c-1$). All ingredients to make it effective already exist:
\lean{binaryCoeffTail\_eq\_finiteCoeffWindow\_add\_shiftedTail}{HalfCylinderFiniteShadow.lean:203} gives
$\mathrm{gap}=\mathrm{tail}(k{+}1{+}J)/2^J$, and
\lean{binaryCoeffTail\_le}{GenericTailOrbitRigidity.lean:85} gives $\mathrm{tail}(m)\le m+2$.
\textbf{Missing step.} Margin$(k,J)\ge0$ as soon as $(k{+}J{+}3)/2^J <
|\mathrm{halfGreedyNextDyadicExcessNumerator}\,k|/\mathrm{halfGreedyPrefixDenominator}\,k$.
Setting $J=k-2$ turns the hypothesis into one explicit dyadic-safety margin
$|E_k|/D_k > (2k{+}1)/2^{k-2}$ at every skipped rank. Pure bookkeeping over existing lemmas.
\end{obs}

\begin{obs}[A nonnegative margin at an unspecified horizon] \ev{Lean}
\lean{HalfGreedySkippedFullShellNonnegative}{HalfCylinderFullShellSeamBridge.lean:604},
\lean{half\_mem\_mersenneAchievementSet\_of\_skippedFullShellNonnegative}{HalfCylinderFullShellSeamBridge.lean:633},
\lean{HalfGreedySkippedSeamEscape}{HalfCylinderFullShellSeamBridge.lean:687}.
\textbf{Open property and proved conditional implication:} \texttt{HalfGreedySkippedFullShellNonnegative} asks
that every skipped rank $n\ge3$ satisfy
$0\le\ensuremath{F}(n{-}1)\,n$. Assuming this property,
\texttt{half\_mem\_mersenneAchievementSet\_of\_skippedFullShellNonnegative} proves
$1/2\in\ensuremath{\mathcal A}$.
\textbf{Order of quantifiers.} exactly three doubling steps of horizon short.
The exact-row hypothesis needs margin$(k,k{-}2)\ge0$; this gives margin$(k,k{+}1)\ge0$. By the
succ-recurrence, margin$(k,k{+}1)=8\cdot$margin$(k,k{-}2)+4\,\mathrm{sc}(2k)+2\,\mathrm{sc}(2k{+}1)
+\mathrm{sc}(2k{+}2)$, sc a divisor-count supportCoeff $\le\tau\le 2\sqrt{2k+2}$. Since
nonnegativity is monotone \emph{up} in $J$, the cited statement is the weaker one and cannot
supply the obligation as it stands.
\textbf{Missing step.} A horizon-tightening ;  at every skipped rank, margin$(k,k{+}1) \ge
4\,\mathrm{sc}(2k)+2\,\mathrm{sc}(2k{+}1)+\mathrm{sc}(2k{+}2)$: the full-shell margin must exceed
an explicit $O(\sqrt k)$ divisor-count quantity. Same $\sqrt{}$-scale currency as the LPC-sqrt row
above.
\end{obs}

\begin{obs}[Cofinal positive skips \emph{is} the obligation] \ev{Lean}
\lean{CofinalPositiveHalfGreedySkips}{BooleanMobiusSkipRowCofinal.lean:22},
\lean{cofinalExactLocalMersenneHalfRows\_of\_positiveHalfGreedySkips}{BooleanMobiusSkipRowCofinal.lean:84},
line 60; A5
\lean{half\_mem\_mersenneAchievementSet\_iff\_greedySkippedSupport\_infinite}{GreedyAchievementSet.lean:2583},
A6 (\lean{}{GreedyAchievementSet.lean:1515}); positivity discharge
\lean{}{BooleanMobiusCriticalCapacityCofinal.lean:138}.
\textbf{Conclusion.} Cofinally many positive rational-greedy skips $\Rightarrow$
the existence of exact quotient sums at unbounded depths $\Rightarrow 1/2 \in
\ensuremath{\mathcal A}$. Advertised elsewhere as the shortest known path.
\textbf{Scope of the hypothesis.} it is not weaker than the problem ;  it \emph{is}
the problem. The positivity conjunct $0 < \ensuremath{r}(1/2)(c{-}1)$ is
unconditionally dischargeable by the odd-denominator parity fact, so
$\mathrm{CofinalPositiveHalfGreedySkips} \Leftrightarrow
(\mathrm{greedyMersenneSkippedSupport}(1/2))$.Infinite, which Theorem~\ref{record:257bm-c20} proves equivalent to
$1/2\in\mathcal A$.  Thus proving this hypothesis would settle the half-target
question itself.
\textbf{What remains.} One must prove that this fixed greedy sequence has
infinitely many skips.  Equivalence neither supplies that proof nor rules
out finding a useful invariant for the sequence.  The condition in
Theorem~\ref{record:257bm-c7} instead permits finite supports that are
not greedy prefixes.
\end{obs}

\begin{obs}[Terminal bounds and restrictions on finite support] \ev{Lean}
\lean{HalfCarryCofinalTerminalOnlyStrip}{TerminalOnlyCofinal.lean:34},
\lean{HalfTerminalOnlyStripWitness}{TerminalOnlyCofinal.lean:24},
\lean{exists\_infinite\_support\_half\_of\_cofinalTerminalOnlyStrip}{TerminalOnlyCofinal.lean:192}.
\textbf{Conclusion.} Cofinally many depths $M$ carrying a finite normalized word $a$ with
$|\ensuremath{\operatorname{ihc}}(\mathrm{wordSupport}\,a)(M{-}1)| \le \ensuremath{B}\,M
\Rightarrow \exists A$ infinite with $X_{A}(2)=1/2$. Same architecture
as the obligation: cofinal, no coherence, mutually incompatible witnesses allowed.

\par\medskip

\textbf{Different restrictions.} The terminal condition permits \emph{any} support inside the depth-$M$ window but
demands the carry inside a $\sim 2\sqrt M$ strip; the obligation form ($\star$) tolerates a carry
up to $2^{M/2-1}$ ;  exponentially looser ;  but demands the support be confined to ranks $<
M/2+1$, since the exact fill encodes the residue in the pure upper window where the Mersenne
quotient is exactly $2^{M-d}$. The gap is support \emph{locality}, not carry size.

\par\medskip

\textbf{Missing step.} A lower-half confinement upgrade of the strip witnesses ;  $\forall N\
\exists c\ge N$ with a strip-scale witness whose support lies in $[2,c)$ at depth $2c{-}3$. Then
$2\sqrt{2c}+4<2^{c-2}$ for $c\ge8$ gives the required numerical
capacity bound.  The confinement of the support remains an additional,
unproved condition; the comparison does not establish a simpler
existence theorem.
\end{obs}

\begin{obs}[The live greedy carry and the carry of a fixed prefix] \ev{Lean}
\lean{GreedyHalfCarryCofinalStripReturn}{CofinalStripReturn.lean:76},
\lean{greedy\_half\_infinite\_of\_cofinalStripReturn}{CofinalStripReturn.lean:130}, line 157;
\lean{infinite\_support\_half\_of\_mobiusCenteredHalfCarry\_sqrtBound}{HalfCarryReachability.lean:817},
line 805 (nonnegativity from line 842 /
\lean{}{HalfCylinderFinalMiddleCellEscape.lean:94}).
\textbf{Conclusion.} Cofinal returns of the actual greedy carry to the square-root strip (C4c), or a
pointwise $2\sqrt N+4$ bound on the Möbius-centred carry (G3), each give an infinite support with
value exactly $1/2$.
\textbf{Comparison.} these bound the carry of the \emph{infinite}
greedy support $\mathrm{greedyMersenneSupport}(1/2)$; sharp capacity is about the \emph{frozen
finite prefix} $\mathrm{halfGreedyPrefixSupport}(c{-}1)$. The two carries agree only at index $k$
(inside \lean{greedyHalfFrozenMargin\_eq\_actual\_skip\_margin}{HalfCylinderFiniteShadow.lean:1120}),
not at depth $2c{-}3$ where the fill needs it. The bridge is already recorded:
\lean{halfGreedy\_precriticalSuffix\_lt\_iff\_futureSkipCoverage}{BooleanMobiusCriticalCapacityCofinal.lean:949}:
sharp capacity $\Leftrightarrow \ensuremath{C}\,G\,(2c{-}4) \le
\mathrm{futureSkipCapacity}\,G\,c\,(c{-}3)$ ;  transport costs exactly the future-skip capacity
term.
\textbf{Missing step.} Either a strip bound stated for the frozen prefix rather than the live orbit,
or $\mathrm{futureSkipCapacity}\,G\,c\,(c{-}3) \ge \ensuremath{B}(2c{-}3)$ ;  holds as
soon as one skip occurs early in $[c{+}1,\,2c{-}4]$, since the capacity is dyadically weighted.
Same input as the bounded-gap precritical bound above.
\end{obs}

\begin{obs}[Above-support recycling, no growth guarantee] \ev{Lean}
\lean{exists\_skippedCoreExactRow\_of\_value\_above}{BooleanMobiusExactRowCrossing.lean:155},
\lean{exactLocalMersenneHalfRow\_of\_first\_localMersenne\_crossing}{BooleanMobiusExactRowCrossing.lean:191},
\lean{exactLocalMersenneHalfRow\_two\_mul\_sub\_two\_of\_skippedCore}{BooleanMobiusSkipRow.lean:149}.
\textbf{Conclusion.} Unconditional: every finite above-half support bounded by $n$ recycles through
its first crossing rank $c$ into a genuine exact row at endpoint $2c{-}2$, with $4 \le c \le n$.
No capacity hypothesis at all.
\textbf{Order of quantifiers.} the witness is $\exists c \le n$, not $\exists c$
large ;  $2c{-}2$ may be $\le n$, so the endpoint need not grow.
\lean{exists\_seeded\_bounded\_double\_or\_recycle\_model}{BooleanMobiusExactRowDichotomy.lean:91}
is the explicit falsifier ;  $\mathrm{boundedDoubleOrRecycleModel}\,n := (n{=}6)$ satisfies the
same transition schema plus a seed and is \emph{not} cofinal. Growth is recovered only inside
\texttt{ProtectedExactLocalMersenneRow}, whose invariants force $c>$cutoff, hence $2c{-}2>$
endpoint ;  and maintaining those invariants is precisely what needs the sharp-capacity fill
rather than this general recycling.
\textbf{Missing step.} A crossing-location lower bound ;  for the supports actually produced, the
first crossing rank satisfies $2c{-}2>n$. Equivalently, keep protection alive without sharp
capacity, e.g.\ a fill placing all new ranks above $c$ using only the recorded $c{-}1$-bit
capacity.
\end{obs}

\begin{obs}[Double-or-recycle dichotomy, proved insufficient] \ev{Lean}
\lean{exactLocalMersenneHalfRow\_double\_or\_recycle}{BooleanMobiusExactRowDichotomy.lean:26},
\lean{boundedDoubleOrRecycleModel\_not\_cofinal}{BooleanMobiusExactRowDichotomy.lean:79},
\lean{exists\_seeded\_bounded\_double\_or\_recycle\_model}{BooleanMobiusExactRowDichotomy.lean:91}.
\textbf{Conclusion.} $\forall n \ge 6$, an exact row at $n$ gives an exact row at $2n{-}1$ or at
$2c{-}2$ for some $4 \le c \le n$; plus a formal proof that this shape, even seeded at endpoint
six, is logically insufficient for cofinality.
\textbf{Order of quantifiers.} the dichotomy is proved by splitting on the sign
of the witness value (\texttt{finiteErdosSum\_den\_odd} excludes equality); the no-go pins the
deficit precisely ;  what is missing is not another dichotomy but strict endpoint progress in the
recycle branch, or a proof that the below-half branch recurs. Re-deriving this dichotomy without controlling repeated returns to
bounded depths would leave the same gap.
\textbf{Missing step.} Either (a) the recycle branch upgraded to $2c{-}2>n$, or (b) cofinally many
below-half exact rows so the doubling branch alone is cofinal. These are two sufficient ways to obtain progress; the countermodel
does not rule out other assumptions that also force unbounded depths.
\end{obs}

\subsection{Universal support criteria and their limits}

\begin{obs}[A sufficient block-certificate condition] \ev{Lean}
\lean{irrational\_erdosSupportSeries\_of\_weighted\_coeff\_certificates}{CertificateKernel.lean:9031}
(carry variant \lean{}{CertificateKernel.lean:9053}).
\textbf{Conclusion.} $\forall b\ge2,\ \forall A\subseteq\mathbb N$ (no infinitude, no structure
hypothesis at all): if $\forall q>0\ \exists N,K,L,C$ with $K\le L$, $b^r \mid
c_{A}(N{+}r)$ for $r\in[1,K]$, $\sum_{r=K+1}^L c_{A}(N{+}r)
b^{L-r}\le C$, $\exists t,\ 0<c_{A}(N{+}L{+}1{+}t)$, and $q(C{+}N{+}L{+}2)<b^L$
;  then $X_{A}(b)$ is irrational. The cited full-support, multiples, periodic, residue-class,
pairwise-coprime and fixed-core constructions use this criterion.  No
such factorisation is asserted for the weighted support theorem in
Section~\ref{sec:257-problem}.
\textbf{Scope of the hypothesis.} already universally quantified over $A$ ;  the
universal quantifier is \emph{not} the gap. The hypothesis \texttt{hcert} is a cofinal required input
supplied only for named classes; the theorem's own docstring: ``certificate existence for
$f_A$... for arbitrary infinite $A$ remains an open obligation.''
\textbf{Limit of this sufficient condition.} Supplying these
certificates for every infinite support would imply the universal
statement, but Proposition~\ref{prop:squarefree} rules out that supply
for the two specified certificate schemes.  The carry variant replaces
digit-wise divisibility by aggregate divisibility
$b^K\mid\sum_r c_A(N+r)b^{K-r}$; the squarefree obstruction applies to
that variant too.  Universal certificate supply is therefore not being
posed as an unresolved equivalent formulation of Problem~257.
\end{obs}

\begin{obs}[Rational targets and infinitely many greedy skips] \ev{Lean}
\lean{half\_mem\_mersenneAchievementSet\_iff\_greedySkippedSupport\_infinite}{GreedyAchievementSet.lean:2583}
(forward \lean{}{GreedyAchievementSet.lean:2482}; reverse
\lean{mem\_mersenneAchievementSet\_of\_greedySkippedSupport\_infinite}{GreedyAchievementSet.lean:1528}).
For every rational $x\ge0$,
\[
 x\in\mathcal A\quad\Longleftrightarrow\quad
 \text{the greedy algorithm for $x$ skips infinitely many ranks}.
\]
The reverse implication holds for every real $x\ge0$:
a fatal remainder $r_n(x)>R_n$ forces every subsequent weight to
be taken, because $R_n=w_{n+1}+R_{n+1}$.
Hence infinitely many skips exclude a fatal remainder, and
Theorem~\ref{thm:greedy-survival} gives membership.
For the forward implication, membership identifies the greedy
sum with $x$.  Finitely many skips would make that sum $E$ minus
a finite rational sum, which is irrational by
Erd\H{o}s's theorem~\cite{erdos1948}.  This contradicts rationality
of $x$ and proves the stated extension as an ordinary argument.

Infinitely many skips do not mean infinitely many selected ranks:
$x=1/3$ has the single-element support $\{2\}$ and skips every
other rank.  Thus the equivalence does not decide which rational
targets are represented, or provide an infinite support with rational
value.  The supplied half-target declarations remain identified above;
a new target-uniform Lean theorem is not claimed. \ev{Math}
\end{obs}

\begin{obs}[Finite sums exclude every rational with even denominator] \ev{Lean}
\lean{finite\_boolSupport\_ne\_half}{HalfCarryReachability.lean:589}; engine
\lean{finiteErdosSum\_den\_odd}{DyadicPrefixCompression.lean:1513}.
\textbf{Conclusion.} $\forall A$ finite with $0\notin A$: $X_{A}(2) \ne
1/2$ ;  any support achieving exactly $1/2$ must be infinite.
\textbf{Parameter range.} the headline is pinned to target $1/2$ while its engine
$\mathrm{Odd}((\mathrm{finiteErdosSum}\,F\,2).\mathrm{den})$ is already fully general over finite
supports.
\textbf{Sufficient specialization:} verbatim re-derivation ;  $\forall t:\mathbb Q$ with even reduced denominator,
$\forall A$ finite with $0\notin A$, $X_{A}(2) \ne (t:\mathbb R)$. No new
mathematics. Thus a represented rational with even reduced denominator would refute universal
\#257 at $b=2$. The converse is not established. If $\mathrm{FinVal}$ denotes the finite-support
values, the exact unrestricted condition is
$(\mathbb Q\cap\ensuremath{\mathcal A})\smallsetminus\mathrm{FinVal}\ne\varnothing$.
\end{obs}

\begin{obs}[The denominator in the reciprocal-mass estimate] \ev{Lean}
\lean{dyadic\_support\_fraction\_reciprocalMass\_diverges\_or\_gt\_one}{RationalSupportCarrySkeleton.lean:2210}
(line 2116); generic-$v$ sibling
\lean{shiftedBooleanCollision\_wrap\_bound\_of\_common\_multiple}{RationalSupportCarrySkeleton.lean:1968}.
\textbf{Conclusion.} For $A$ infinite with $X_{A}(2) = p/2^c$ (dyadic
rational): $\neg\mathrm{Summable}(1/a$ on $A) \vee 1 < \mathrm{reciprocalMass}\,A$.
\textbf{Comparison.} restricted to denominators $2^c$ ($v{=}1$) while
the obligation needs every rational $p/(2^c v)$. This restriction is a call-site choice, not a
proof constraint ;  every supporting lemma carries a free $\{v:\mathbb N\}\ (hv:0<v)$, and the
general-$v$ conclusion is \emph{already recorded} at line 1962 for arbitrary $(v,h,L,F)$.
\textbf{Missing step.} Instantiate the recorded generic theorem at $F=\{a,b\}$, $L=\mathrm{lcm}\,a\,b$,
$h=$ multiplicative order of $2$ mod $v$, giving $\rho(A) \ge
\mathrm{doublingWrapCount}(p,v,h)/h + 1/\mathrm{lcm}(a,b)$ for every rational value. The one
genuinely new step: for $v{>}1$ the excess can be $0$, so strictness needs $1 \le
\mathrm{doublingWrapCount}(p,v,h)/h + \mathrm{booleanCollisionSurplus}\,|F|/L$ (equivalently: the
mean least residue of $p\cdot2^n \bmod v$ is $\ge v(1-|F|/2L)$).
\end{obs}

\begin{obs}[Zero windows of a fixed support] \ev{Lean}
\lean{supportCoeffZeroWindow\_length\_le\_eps\_logb}{SublogDivisorCoverage.lean:435} (line 392).
\textbf{Conclusion.} For any $A$ with a positive element and
$X_{A}(2)=p/(2^c v)$, $v{>}0$ (any rational, not just dyadic):
$\forall\varepsilon{>}0\ \exists B\ge0\ \forall N\ge1\ \forall h$,
$\mathrm{SupportCoeffZeroWindow}\,A\,(c{+}N)\,h \to h \le \varepsilon\log_2 N + B$. The constant $B$ may depend on the fixed support and the other parameters.
More basically, if $a_0$ is any positive element of $A$, then each block
of $a_0$ consecutive integers contains a multiple of $a_0$.  Every zero
window of $c_A$ therefore has length at most $a_0-1$, whether or not
$X_A(2)$ is rational.  For a fixed nonempty support, the displayed
logarithmic bound thus gives no additional restriction.
A proposed alternative between super-logarithmic zero windows and general
block certificates cannot split this class: the first branch is already
impossible.  It would require the certificate branch for every support,
which the obstruction in Proposition~\ref{prop:squarefree} rules out for
the specified schemes.
\end{obs}

\begin{obs}[Prime support and the reciprocal-sum hypothesis] \ev{Lean}\ev{Cited}
\lean{irrational\_erdosSupportSeries\_pairwise\_coprime}{CertificateKernel.lean:10776}
(direct certificate required input
\lean{exists\_weighted\_coeff\_certificates\_supportCoeff\_pairwise\_coprime}{CertificateKernel.lean:9902}).
\textbf{Conclusion.} For every $b\ge2$, an infinite pairwise-coprime support with
summable reciprocal mass has irrational support series.
\textbf{Method boundary:} the primes satisfy the other hypotheses but
$\sum_p1/p$ diverges, so this Lean required input does not apply.  It is false,
however, to infer that the base-$2$ prime-support value is open.
Tao--Ter\"av\"ainen prove
\[
 \sum_{p\ {\rm prime}}\frac1{2^p-1}
 =\sum_{n\ge1}\frac{\omega(n)}{2^n}
\]
irrational in Theorem~1.3 of
\href{https://arxiv.org/abs/2512.01739}{their 2025 preprint}.  That result is
cited, not formalised here.
\textbf{What remains diagnostic:} replacing the global tail budget by a
windowed one would ask whether this particular certificate engine can recover
the known theorem.  It is no longer an open mathematical target or a claimed
new route for \#257.
\end{obs}

\begin{obs}[The two hypotheses of the sunflower criterion] \ev{Lean}
\lean{irrational\_erdosSupportSeries\_of\_orthogonalPetalBouquet}{SupportSunflowerDichotomy.lean:540}
(selector def line 406; carry hypothesis line 381; bouquet structure line 33).
\textbf{Conclusion.} For $A=$ (finite exceptional set dividing $Q$) $\cup \{\mathrm{core}_i \cdot
\mathrm{petal}_i\}$ with cores $\mid Q$, petals pairwise coprime and coprime to $Q$, $\sum
1/\mathrm{petal}_i<\infty$: if \texttt{SunflowerForcedSlotTailSelection} $A$ then irrational.

\par\medskip

\textbf{Different restrictions.} two gaps. (i) On the support: the bouquet class strictly
generalises pairwise-coprime but still demands globally pairwise-coprime, summable petals ;  it
inherits NM-06's exclusion of the primes and adds nothing in the divisor-dense regime. (ii) On
the selector: \texttt{SunflowerForcedSlotTailSelection} is an unproved cofinal supply ;  grep for
``dichotomy'' in the file returns only the namespace; despite the filename there is no dichotomy
theorem there.

\par\medskip

\textbf{Missing step.} For the selector, $\forall K{>}0\ \exists N$, $2^K \mid$ the carried first
block $\wedge\ \mathrm{binaryCoeffTail}(c_{A})(N{+}K)\le16$, from the
bouquet's persistent-reduced-modulus averaging ;  the file states this is what that averaging
argument is supposed to give. For the class: a bouquet-existence theorem $A$.Infinite $\to$
Nonempty(bouquet structure) is \emph{false} as stated (take $A=$ all multiples of 6 ;  no coprime
petal decomposition).  A theorem for every infinite support would
therefore require a different structural hypothesis, not merely a
proof of that false existence assertion.
\end{obs}

\begin{obs}[What eventual periodicity supplies] \ev{Lean}
\lean{irrational\_erdosSupportSeries\_eventuallyPeriodic}{CertificateKernel.lean:11604} (engine
line 11262, required input line 11072; finite-perturbation transfer lines 9139, 9148).
\textbf{Conclusion.} $\forall b\ge2\ \forall m{>}0\ \forall N_0$: $A$ infinite and $m$-periodic above
$N_0$ $\Rightarrow$ irrational. Plus: the class of supports with irrational series is closed
under finite symmetric difference in both directions. The statement follows from the periodic
theorem of Luca and Tachiya~\cite[Theorem~A, p.~139]{lucatachiya2017}.

\par\medskip

\textbf{Scope of the hypothesis.} The theorem allows every positive
period and every finite starting threshold.  It still requires an
eventually periodic support, so it does not include all aperiodic
supports covered by the other criteria in this record.  Finite
symmetric difference preserves eventual periodicity.  Thus the
finite-perturbation argument alone cannot extend this theorem to an
aperiodic tail; a further hypothesis or argument is needed.

\par\medskip

\textbf{Missing step.} A ``locally periodic at cofinally many scales'' version ;  $\forall q{>}0\
\exists N\ \exists m\le f(N)$ such that $A\cap[N,N{+}L]$ is $m$-periodic for $L$ large enough
that the periodic sieve manufactures one block certificate. The periodic required input already only
needs periodicity across the certificate window (window-local), so this is the closest genuine
widening ;  but proving cofinally many periodic windows exist for arbitrary $A$ is a new
statement, not a re-run.
\end{obs}

\begin{obs}[The reduced denominator in a gap criterion] \ev{Lean}
\lean{irrational\_erdosSum\_of\_lcm\_gap}{CertificateKernel.lean:5883} (instances
\lean{irrational\_erdosSum\_factorial\_support}{CertificateKernel.lean:6035},
\lean{irrational\_erdosSum\_two\_pow\_support}{CertificateKernel.lean:6059}).
Let $A=\{a_0<a_1<\cdots\}\subseteq\mathbb N_{>0}$, let $b\ge2$ be an integer, and
write the first $k$ terms in lowest terms as
\[
 S_k=\sum_{j<k}\frac1{b^{a_j}-1}=\frac{P_k}{Q_k},\qquad Q_k\ge1.
\]
A sufficient irrationality condition is
\[
 Q_k b^{-a_k}\longrightarrow0,
 \quad\text{equivalently}\quad a_k-\log_b Q_k\longrightarrow+\infty.
\]
The next omitted exponent is $a_k$, not the last exponent included
in $S_k$.  Since $a_{k+j}\ge a_k+j$, the positive omitted tail obeys
\[
 0<X_A(b)-S_k
 \le\left(\frac{b}{b-1}\right)^2b^{-a_k}.
\]
If $X_A(b)=p/q$ were rational, positivity would give
$qQ_k(X_A(b)-S_k)\ge1$.  The displayed tail estimate contradicts
this under the proposed condition.  This is the denominator-times-error
argument in \lean{irrational\_of\_den\_mul\_abs\_sub\_tendsto\_zero}{CertificateKernel.lean:5371},
with all indices specified.

For $L_k=\operatorname{lcm}(a_0,\ldots,a_{k-1})$ one has
$Q_k\mid b^{L_k}-1$, hence $\log_b Q_k<L_k$.
Thus the earlier condition $a_k-L_k\to\infty$ suffices.
Using actual reduced denominators can improve this estimate, but
requires an \emph{upper} bound for $Q_k$ relative to $b^{a_k}$.
A surviving-divisor \emph{lower} bound does not provide that upper
bound.  In particular, the divisibility statement in
\lean{}{RationalDenominatorSurvival.lean:17} and the lower bound in
\lean{upperHalfMersenneProduct\_lower\_bound}{MersenneShadowDenominatorGrowth.lean:60}
do not alone establish the proposed condition for an arbitrary
support.  No such condition is proved here for the remaining dense
supports. \ev{Math}
\end{obs}

\begin{obs}[Rationality expressed by the integer carry] \ev{Lean}
\lean{erdosSupportSeries\_rational\_iff\_exists\_temperedCarry}{BooleanMobiusCarry.lean:384}
(line 376); T7 \lean{binaryCoeffSeries\_rational\_iff\_exists\_temperedBinaryOrbit}{GenericTailOrbitRigidity.lean:426}.
\textbf{Conclusion.} $\forall A\subseteq\mathbb N$ (no hypothesis at all):
$\mathrm{HasRationalValue}(X_{A}(2)) \Leftrightarrow \exists q{>}0\
\exists U:\mathbb N\to\mathbb Z,\ \mathrm{IsTemperedBinaryOrbit}(c_{A})\,q\,U$.
\textbf{Comparison.} already universal over $A$ ;  the exact coordinate
in which the obligation should be stated, and the literal shared argument with \#249 (instantiate T7
at $c:=\mathrm{Nat.totient}$). It does not refute the orbit:
\lean{temperedBinaryOrbit\_eq\_scaledTail}{GenericTailOrbitRigidity.lean:339} pins the orbit
uniquely to $u(N)=qT_c(N)$, and $T_c(N)\le N{+}2$ confirms it is never ruled out on growth
grounds for any $c$ with $c(n)\le n$.
\textbf{Missing step.} A support-specific non-existence theorem ;  $\forall A$ infinite, $\forall
q{>}0$, $\neg\exists U$ with $\mathrm{IsTemperedBinaryOrbit}(c_{A})\,q\,U$.
Because T7 is an iff and the orbit is unique, this is logically equivalent to the obligation ; 
the row's value is naming the single object (the tempered carry orbit of $c_{A}$)
whose non-existence is the entire content, for every $A$ at once.
\end{obs}

\begin{obs}[Why unbounded carry alone is insufficient] \ev{Lean}
\lean{shifted\_state\_unbounded\_of\_infinite\_support}{RationalSupportCarrySkeleton.lean:2327},
\lean{exists\_unbounded\_shifted\_odd\_tail\_nat\_state\_of\_support\_fraction}{RationalSupportCarrySkeleton.lean:2383},
\lean{one\_add\_mul\_card\_le\_two\_mul\_shifted\_state}{RationalSupportCarrySkeleton.lean:2237}.
Let $A\subseteq\Npos$ be infinite, let $c\ge0$ and $v\ge1$ be integers,
and let $u(n)$ be positive integers satisfying
\[
 u(n+1)+v c_A(c+n+1)=2u(n)\qquad(n\ge0).
\]
For a finite $F\subseteq A$ and a common multiple $L$ of its
elements with \emph{$L>c$}, put $n=L-c-1\ge0$.  Directly from
the recurrence,
\[
 2u(L-c-1)=u(L-c)+v c_A(L)\ge1+v|F|.
\]
This argument uses positivity of the next state, not an assumed
tail representation or a limit condition.  A general positive
solution may also have an exponentially growing homogeneous part.
The condition $L>c$ ensures that the coefficient row really is
$L$; truncated natural-number subtraction cannot replace it.

Finite subsets $F$ of arbitrarily large size, followed by positive
common multiples exceeding $c$, prove that $u$ is unbounded.
In the tempered case the separate identity
$u(n)=v\sum_{r\ge1}c_A(c+n+r)2^{-r}$ gives
$u(n)\le v(c+n+2)$.  Unboundedness is consistent with this linear
upper bound and with $u(n)/2^n\to0$.  A contradiction would require
an incompatible upper estimate at the same chosen indices.
\ev{Math}
\end{obs}

\begin{obs}[A finite-gap criterion for every target in the convex hull] \ev{Lean}
\lean{half\_mem\_mersenneAchievementSet\_or\_exists\_fatal\_gap}{HalfCutLocator.lean:623},
\lean{half\_mem\_mersenneAchievementSet\_iff\_no\_existsFatalHalfGap}{HalfCutLocator.lean:654};
general-$A$ killer \lean{positiveMersenneSupportValue\_ne\_of\_rank\_gap}{HalfCutLocator.lean:380}.
\label{obs:general-target-gap}
For every real $x\in[0,E]$, not just rational targets, one has
\[
 x\notin\mathcal A\quad\Longleftrightarrow\quad
 \exists m\ge1\ \exists D\subseteq\{1,\ldots,m-1\},\qquad
 X_D(2)+R_m<x<X_D(2)+w_m.
\]
Here $R_m=\sum_{j>m}w_j$.  The implication to the right has a
short greedy proof.  If $x\notin\mathcal A$, choose the first $m$
with $r_m(x)>R_m$ using Theorem~\ref{thm:greedy-survival}.
Since $r_0(x)=x\le E$, we have $m\ge1$ and
$r_{m-1}(x)\le R_{m-1}=w_m+R_m$.  A take at $m$ would leave a
remainder at most $R_m$, so this step must be a skip.  Thus
$R_m<r_{m-1}(x)<w_m$, giving the displayed gap with the greedy
prefix $D$.

Conversely, the displayed gap contains no achievable value.
A support with prefix $D$ either skips $m$, giving value at most
$X_D(2)+R_m$, or takes $m$, giving value at least $X_D(2)+w_m$.
A different prefix first differs at some $j<m$; the inequality
$w_j>R_j$ orders its entire continuation interval strictly on one
side of the interval with prefix $D$.
Both gap endpoints themselves are achievable, by the supports
$D\cup\{m+1,m+2,\ldots\}$ and $D\cup\{m\}$.
Consequently no parity argument or irrationality assertion about
$E$ is needed for this ordinary general-target dichotomy.
The restriction to $[0,E]$ is essential to this formulation:
targets outside the convex hull are already excluded and need not
lie in an internal gap.  The general-target equivalence and both endpoint
constructions are Lean-checked in
\lean{paper\_general\_target\_gap}{lean/ErdosProblems/Erdos257/PaperCompleteR20/GeneralTargetGap.lean:173}.
\ev{Math}
\end{obs}

\begin{obs}[Why the signed periodic theorem still requires periodicity] \ev{Lean}
\lean{irrational\_or\_bpow\_mul\_eq\_intCast\_intWeightedErdosSeries\_periodic}{CertificateKernel.lean:14175},
\lean{irrational\_intWeightedErdosSeries\_periodic\_of\_coeff\_nonneg\_of\_frequently\_ne\_zero}{CertificateKernel.lean:14583}
(nonpos mirror line 14315).
\textbf{Conclusion.} $\forall b\ge2\ \forall m{>}0\ \forall$ $\mathbb Z$-valued $m$-periodic weight
$w$: irrational$(\sum w(a)/(b^a-1))$ $\vee$ $\exists k,z,\ b^k x=z$; the terminating branch closes
unconditionally as soon as $\mathrm{intWeightedCoeff}\,w$ is one-signed and frequently nonzero.
Luca and Tachiya's theorem already gives irrationality for every nonzero periodic
weight~\cite[Theorem~A, p.~139]{lucatachiya2017}.

\par\medskip

\textbf{Different restrictions.} two mismatches. (i) periodicity of the weight is still
required ;  this is the signed lift of NM-08, not an escape from it. (ii) the one-sidedness clause
is \emph{free} for \#257 (a 0/1 support indicator gives $c_{A}\ge0$
automatically) ;  the sign machinery buys nothing here; it was built for the signed \#249 argument.
The genuinely transferable content is the dichotomy shape itself: irrational-or-$b$-adically-
terminating, terminating branch killable by a positive far tail.

\par\medskip

\textbf{Missing step.} An aperiodic version of the dichotomy ;  for arbitrary $A$, either
$X_{A}(2)$ is irrational or $2^k\cdot$it is an integer. For arbitrary infinite $A$, exclusion of the terminating branch is itself an additional unproved
assertion; finite-support odd-denominator parity does not pass automatically to an infinite limit.
An aperiodic dichotomy together with a valid terminating-branch exclusion would close the
obligation. Periodicity is used only to manufacture the full-block certificates
(``periodicity alone manufactures the full-block certificates,'' per the file's own docstring),
so removing it is again NM-01's missing required input, restated a third time.
\end{obs}

\subsection{Three combined conditional arguments}

The first conditional implication also appears in the supplied
release source, as identified below.  The other two entries are
ordinary mathematical arguments about what the stated hypotheses
would imply.  None establishes its unproved orbit hypothesis or
decides~\#257.

\begin{thm}[A: a sufficient lower bound at every reset]
Assume that
$|\mathrm{rem}(r+1)-2^{r+1}|>2^{(r+5)/2}$ for every upper or
middle reset $r\ge10$.  Then $1/2\in\mathcal A$.

The supplied release source states this implication as
\href{https://github.com/wcook04/plectis-erdos-lean/blob/52f29ad173b04e3bac941b3663f2b9aebe5de0bb/Erdos257PeriodNoncollapse/HalfCylinderResetSqrtEscape.lean#L442}{\textsc{Lean source}}.
The packet classifies that file as \texttt{release\_only}; its proof
body was inspected for this exposition, without a new kernel replay.
Here is the role of the constants.  Put
\[
 B(r)=2^{\lfloor(r+4)/2\rfloor}+2r+3.
\]
For $r\ge10$, the source proves $B(r)^2\le2^{r+5}$.
A right branch at the first crossing of the largest omitted rank $r$
would force $|\Delta_{r+1}|\le B(r)$, by tracing the exact affine
recurrence back to that reset.  The assumed strict lower bound
excludes this crossing.  The largest-skip induction then gives
half-membership.  The hypothesis concerns every actual reset beyond
the threshold; neither a long finite check nor an asserted
run-length asymptotic supplies it.
\end{thm}
\leannote{Lean: \lproof{ErdosProblems/Erdos257/PaperCompleteR21/ResetSqrtEscapeHalfMembership.lean}{595}{paper_theoremA_half_membership}, \lproof{ErdosProblems/Erdos257/PaperCompleteR21/ResetSqrtEscapeHalfMembership.lean}{545}{paperResetSqrtEscape_iff_square}, \lproof{ErdosProblems/Erdos257/PaperCompleteR21/ResetSqrtEscapeHalfMembership.lean}{121}{paper_theoremA_crossing_bound_square_le}, \lproof{ErdosProblems/Erdos257/PaperCompleteR21/ResetSqrtEscapeHalfMembership.lean}{261}{paper_theoremA_right_branch_forces_small_deviation}, and 1 further declaration in the \hyperref[sec:coverage]{coverage section}.}

\begin{rem}[B: rigidity of dangerous resets]
A dangerous reset (one within $2^{(r+5)/2}$ of the sqrt-escape threshold) that is preceded by a
long pure R-run pins the deviation at the preceding reset to one of at most two adjacent
pulse-determined integers, where one might expect a free real number: the affine excess recurrence
(Observation~\ref{obs:affine-excess-recurrence}) is exact, so once the run length and the pulse
word are fixed, $w_{r_0+1}$ has no remaining degree of freedom.
Observation~\ref{obs:dangerous-reset-rigidity} states the hypotheses and the run-length
threshold. Consequently a chain of $n$ consecutive dangerous resets with long runs is a tower of
$n{-}1$ nested exact integer conditions, none free to vary once the pulse stream is fixed. Failure of sqrt-escape therefore
requires such an integer-coincidence tower. No probability model is supplied, so these conditions
do not by themselves make dangerous resets unlikely; the open step is excluding the two pinned
values. \ev{Math}
\end{rem}

\begin{thm}[C ;  the one-sided finite decision boundary]
The checked equivalence
\lean{mem\_mersenneAchievementSet\_iff\_greedy\_survival}{GreedyAchievementSet.lean:1458}
separates two logically different kinds of evidence. A fatal greedy gap found at a finite rank is
a finite certificate that $1/2\notin\mathcal A$. By contrast, membership requires survival at
every rank; the corpus supplies no finite certificate that the orbit survives forever and no
completion theorem turning a long surviving prefix into membership.

The reset-crossing hypotheses, truncation-rung ladder, and sharp-capacity inequalities feed
sufficient conditions for membership. They are not proved equivalent to one another,
nor is failure to find one at a given depth evidence of survival. Consequently these searches
give at most a semi-decision procedure for \emph{nonmembership of the specific value $1/2$}.
Halting does not prove the universal statement in \#257, and non-halting does not prove that
$1/2$ is represented. In particular, neither truth nor falsity of the universal problem is
shown semi-decidable here. Deeper negative search alone establishes nothing beyond the tested
finite ranks; this is the quantifier boundary recorded by SE-7/RC-7.
\end{thm}
\leannote{Lean: \lproof{ErdosProblems/Erdos257/PaperCompleteR21/CentredCompletionAndDecisionBoundary.lean}{77}{paper_one_sided_finite_decision_boundary}.}

\subsection{Four limitations of the stated implications}

The following conditions are equivalent to their target statements
within the specified scope.  An equivalence may still be useful by
exposing a quantity that is easier to estimate.  What it does not
supply is a proof of that estimate, or a strict weakening of the
logical obligation.

\begin{rem}[hbound is equivalent to the conclusion, greedy support only]
\lean{greedy\_half\_infinite\_of\_mobiusCenteredHalfCarry\_sqrtBound}{HalfCarryReachability.lean:834}
and its packaged form
\lean{greedy\_half\_infinite\_of\_mobiusCenteredHalfCarry\_upperBound}{HalfCarryReachability.lean:953}
take a hypothesis \texttt{hbound}: $\forall N,\ \ensuremath{C}\,A\,N \le
2\sqrt N + 4$. For the canonical greedy support $A=\mathrm{greedyMersenneSupport}(1/2)$
specifically, \texttt{hbound} is logically \emph{equivalent} to the theorem's own conclusion
$X_{A}(2)=1/2$ ;  not a strictly weaker sufficient condition. This follows
from the exact residual identity at line 842, the unconditional sqrt-envelope tail bound
(\lean{binaryCoeffTail\_supportCoeff\_le\_two\_sqrt\_add\_four}{BooleanMobiusCarry.lean:290}), the
unconditional tail-nonnegativity lemma
(\lean{binaryCoeffTail\_nonneg}{GenericTailOrbitRigidity.lean:78}), and the greedy-specific fact
$X_{A}(2)\le1/2$ (line 859). \textbf{The scope limit is exact and must
not be widened by inference:} this equivalence does \emph{not} extend to every $A$ with $1\notin
A$ ;  dropping the greedy-specific $\le1/2$ fact yields only $\mathrm{hbound}\Leftrightarrow
\delta\le0$ (a one-sided condition), not $\delta=0$. For the sibling hypothesis
\lean{GreedyHalfCarryCofinalStripReturn}{CofinalStripReturn.lean:76}, the former
\texttt{Real.sqrt}/\texttt{Nat.sqrt} adapter gap is repaired by restricting cofinal witness depths
to perfect squares. Any remaining gap lies in the required input hypothesis, not in rounding.
\end{rem}

\begin{rem}[What the terminal carry condition proves]
\lean{HalfCarryCofinalTerminalOnlyStrip}{TerminalOnlyCofinal.lean:34} at the original strip
constant is logically equivalent to $1/2\in\ensuremath{\mathcal A}$ after restricting
cofinal witness depths to perfect squares and using the recorded tail bound. The square-root
conversion is no longer an open adapter; the open content is the required input itself.
\end{rem}

\begin{rem}[Positive greedy skips and half-membership]
Repeated here for emphasis from the 257-cofinal-rows entry above: once the positivity conjunct of
\lean{CofinalPositiveHalfGreedySkips}{BooleanMobiusSkipRowCofinal.lean:22} is recognised as
unconditionally free (odd-numerator/even-subtrahend parity plus greedy nonnegativity, discharged
at \lean{}{BooleanMobiusCriticalCapacityCofinal.lean:138}), the hypothesis collapses to
$(\mathrm{greedyMersenneSkippedSupport}(1/2))$.Infinite, proved by A6
(\lean{}{GreedyAchievementSet.lean:2514}) to be \emph{equivalent} to $1/2\in
\ensuremath{\mathcal A}$. A proof or disproof of this one Prop settles \#257-at-1/2 and
vice versa; it is not a waypoint toward that settlement.
\end{rem}

\begin{rem}[Scope discipline for the equivalence class]
Each equivalence uses its stated domain.  The centred-carry upper
bound concerns the actual greedy support $G$, not every support omitting
$1$.  The cofinal terminal condition already works with the original
constant $4$, because perfect-square depths avoid a rounding loss.
The positive-skip equivalence is stated here at the target $1/2$;
extending it to another target requires checking the target-dependent
arguments.  None of these qualifications supplies the still-missing
unbounded sequence of witnesses.
\end{rem}
% ; - part a257_p4 ; -
\section{Hypothesis-specific obstructions}

The examples in this section test particular proposed implications
for the reset estimate.  Some give counterexamples in abstract
recurrence models; others compute an exact finite orbit or compare
two quantitative bounds.  An abstract counterexample is not a state
of the actual greedy orbit unless that correspondence is established.
Additional arithmetic hypotheses may therefore change the conclusion.

Throughout, $\mathrm{rem}(s)$ is the quotient-row remainder, $\lambda_n := \Delta_n/2^n$ the normalized seam
deviation of Theorem~\ref{thm:dynamics}, and $\mathcal A$ the Mersenne achievement set. Reset $\sqrt{}$-escape
from row 10 is the statement $|\mathrm{rem}(r+1)-2^{r+1}| > 2^{(r+5)/2}$ for every reset row $r\ge 10$
(RC-4, Part~III).  It is a sufficient input for the cited
half-membership argument, not a necessary condition for every
possible proof.

\subsection{Four tested approaches to the reset estimate}

\begin{obs}[Prime steps can occur within a bounded drift window]
\label{obs:prime-doubling}
\textbf{Route as conceived.} The universal recursion $K(M+1) = 2K(M) - (\tau(M+1)-1)$ (Part~I,
SE-1) doubles at every step and subtracts the divisor-excess $\tau(M+1)-1$. At a prime $M+1=p$,
$\tau(p)=2$ so the subtracted term is $1$; the naive hope is that primality of an index forces the
step to be "almost pure doubling," and that prime density (Chebyshev/PNT-scale) then bounds how
long a run of near-pure-doubling steps can last, which would bound the danger run-length $L_r$ in
the right-branch recurrence in Theorem~\ref{thm:dynamics}.
The upper branch has a different, adjacent-spacing correction, so
a run-length argument cannot use the discarded uniform-error
formula across resets.

\textbf{A finite test.} Use the skip set of the quotient row $n=607$
from Section~\ref{ssec:seam-model}, and define the frozen-row drift
\[
 J_{607}(M)=K(M)+\sum_{d\in\mathrm{Skip}_{607}}
                   \left\lfloor\frac{2^M}{2^d-1}\right\rfloor.
\]
Exact integer arithmetic gives, in order at $M=607,\ldots,617$,
\[
 5,\ 5,\ 4,\ 4,\ 6,\ 2,\ 3,\ 4,\ 5,\ 2,\ 3.
\]
There are eleven values and ten transitions.  The indices $607$, $613$,
and $617$ are prime.  For $M+1\ge607$ the recurrence is
\[
 J_{607}(M+1)=2J_{607}(M)
 -\bigl(\tau(M+1)-1-c_{\mathrm{Skip}_{607}}(M+1)\bigr).
\]
At a prime index the last coefficient is $1$.  The composite-index
coefficients can compensate for those prime steps, as the displayed
values illustrate.  This table is a fresh finite arithmetic check,
not a replay of the longer historical run search.

\textbf{Scope.} Primality at some indices does not force this drift to
leave a small interval immediately.  The example supplies no universal
run-length bound and no statistical conclusion about prime indices.
It also does not exclude a proof that combines prime counts with
quantitative information about the intervening divisor counts.
\coord{prime-density}
\ev{Cert}
\end{obs}

\begin{obs}[Parity and perfect squares in the recurrence]
\label{obs:parity-squares}
\textbf{Route as conceived.} $K(M)$ is even exactly when $M$ is a perfect square (a true,
provable fact). The hope was to use this parity law to forbid certain residues or run patterns in
the digit stream of $K$, since parity constraints are a classical first attack on digit-run
questions.

\textbf{What the parity calculation says.} The recurrence implies
$J(M+1)\equiv\delta(M+1)\pmod2$ for every integer forcing sequence.
For $K$, the additional arithmetic fact that $\tau(M)$ is odd exactly
at squares gives the stated parity pattern.  One should not call that
pattern a property of every forcing sequence, or say that it excludes
no residue: at each fixed index it does specify the parity.

\textbf{Scope.} Restating a parity consequence of a known recurrence
does not add a new independent constraint.  To turn it into a
run-length bound would require further information about the sizes or
arrangement of the forcing terms.  No argument here rules out using
parity together with such information.
\coord{parity-tautology}

\textbf{Use of the observation.} The parity law does constrain
$K(M)$ at each specified index and is a useful check on any proposed
run pattern.  The missing implication is from this residue information
to a quantitative bound on run lengths.  No such implication is
proved by the parity calculation alone.
\ev{Math}
\end{obs}

\begin{obs}[What the denominator estimate does not give]
\label{obs:irrationality-measure}
One proposed separation argument would prevent
$\sum_{d\in\mathrm{Skip}_n}x_d$ from approximating
$C:=E-\tfrac32=0.1066951524152917\ldots$ at scale $2^{-3n/2}$.
An irrationality-measure bound is relevant only after relating the
reduced denominator of that finite sum to $n$.

For $n\ge2$ and $\mathrm{Skip}_n\subseteq\{2,\ldots,n-1\}$, its
reduced denominator $D_n$ satisfies
\[
 \log_2 D_n\le\sum_{d\in\mathrm{Skip}_n}\log_2(2^d-1)
 \le\sum_{d\in\mathrm{Skip}_n}d
 \le\frac{n(n-1)}2-1.
\]
The middle inequality is strict when the skip set is nonempty.
This is an upper bound; it is not an asymptotic formula for $D_n$.
In particular, cardinality alone does not justify an asymptotic of
order $n^2/4$.

Zudilin proves $\mu(E)\le2.46497868\ldots$~\cite[Theorem~1, p.~154]{zudilin2004}.
Rational translation preserves the irrationality exponent, so for
any fixed $\nu>2.46497868\ldots$ there is a constant $c_\nu>0$
such that every reduced rational $p/q$, with $q\ge1$, satisfies
$|C-p/q|\ge c_\nu q^{-\nu}$.
Using only the displayed denominator estimate therefore gives
\[
 \left|C-\sum_{d\in\mathrm{Skip}_n}x_d\right|
 \ge c_\nu\,2^{-\nu(n(n-1)/2-1)}.
\]
This guaranteed lower bound is eventually smaller than $2^{-3n/2}$,
so it does not prove the desired separation.  A sharper estimate
for the actual reduced denominators, or an approximation result
using additional structure of these numerators, could change the
comparison.  Neither is excluded by this calculation.
\coord{diophantine-approximation}\ev{Math}
\end{obs}

Retaining only the parity of an input coefficient leaves considerable
freedom.  In $e'=2e+c$, every incoming integer $e$ admits the odd
choice $c=-1-2e$, which gives $e'=-1$.  The next proposition uses
the same centred-remainder construction when more unit bits and
a valuation are prescribed.

\begin{prop}[Centred completion of fixed-precision 2-adic data]
\label{prop:2adic-nogo}
Fix $u\ge1$ and a finite list $(v_i,a_i)_{0\le i<m}$, where
$v_i\ge0$ are integers and the integers $a_i$ are odd.  For every
initial integer $e_0$, there are integers $z_i,e_{i+1}$ such that
\[
 e_{i+1}=2e_i+2^{v_i}(a_i+2^u z_i),\qquad
 |e_{i+1}|\le2^{v_i+u-1}\qquad(0\le i<m).
\]
Thus one may prescribe the valuation and $u$ bits of the odd unit
of each input coefficient while keeping every \emph{successor}
state in its stated centred interval.  The initial state is arbitrary.
\coord{2-adic-valuation-unit}\ev{Lean}\scale{finite words}
\end{prop}
\leannote{Lean: \lproof{ErdosProblems/Erdos257/PaperCompleteR21/CentredCompletionAndDecisionBoundary.lean}{34}{paper_centred_completion_of_fixed_precision}.}

\begin{proof}
For a single step, put $M=2^{v_i+u}$ and $R=M/2$, and choose
\[
 e_{i+1}=\bigl((2e_i+2^{v_i}a_i+R)\bmod M\bigr)-R.
\]
The remainder convention $0\le x\bmod M<M$ gives
$-R\le e_{i+1}<R$.  Also
$e_{i+1}-2e_i\equiv2^{v_i}a_i\pmod M$, so this difference has
the required form $2^{v_i}(a_i+2^u z_i)$.  Repeating the
construction proves the finite-list statement.  It uses no
relation between the freely chosen quotients $z_i$ at different
steps.
\end{proof}

The formal-source statements are
\lean{fixedPrecisionTropicalNoGo}{TropicalCurvatureCarry.lean:137},
\lean{vu\_step\_has\_centred\_completion}{TropicalCurvatureCarry.lean:69}
and
\lean{vu\_word\_has\_prefix\_locked\_completion}{TropicalCurvatureCarry.lean:115}.
They concern this unrestricted congruence model, not coefficients
already fixed by a divisor-count sequence.  Fixed-precision data
alone cannot exclude every completion in this model.  The proposition
does not show that the completion comes from a Mersenne support,
satisfies additional arithmetic relations, or belongs to the
particular greedy orbit.  In particular, it is not a general
obstruction to every fixed-window argument for Problem~249 or~257.

These observations have separate scopes: a finite prime-window
example, a parity consequence of a known recurrence, an insufficient
denominator estimate, and a completion theorem for specified
fixed-precision symbols.  They do not exhaust the possible approaches
to the reset estimate or force every proof to use digit-carry
machinery.  The next section studies that machinery because it gives
explicit conditional implications, not because other methods have
been ruled out.

\subsection{Excluding particular final dyadic intervals}
\label{sec:branch-exclusions}

The two-sided bound
$\min(\mathrm{rem}(s),\mathrm{overshoot}(s))\le2^s$ for all $s\ge5$
is conditional.  It propagates by induction from the local hypothesis \textsf{SeamTwoSidedDyadicCellEscape}, which at every row
excludes three specified integer values on the middle branch and imposes
one further inequality on the right branch:
\[
 \begin{gathered}
 s\ge5,\quad\text{no carry and a middle transition}
 \quad\Longrightarrow\quad C_s\notin\{-3,-2,-1\};\\
 s\ge5,\quad\text{no carry and a right transition},\quad
 \mathrm{overshoot}(s)\le2^s\\
 \Longrightarrow\quad\mathrm{charge}(s)\le2^{s+2}.
 \end{gathered}
\]
Here $C_s=4\,\mathrm{rem}(s)-p_s^- -4$.  On a middle transition,
$\mathrm{rem}(s+1)-2^{s+1}=C_s+4$; on a right transition,
$\mathrm{rem}(s+1)=C_s-2^{s+1}$.  Thus $C_s$ must not be identified
with the next row's centred remainder.
\lean{Erdos249257.SeamTwoSidedDyadicCellEscape}{HalfCylinderMiddleCarryLowerBound.lean:4374}

\begin{prop}[The upper (carries) successor needs no exceptional-cell exclusion]
\label{prop:upper-unconditional}
Let $s\ge5$ and suppose the upper branch occurs at row $s$, that is,
$(\mathrm{seamAdjacentCut}\ s\ hs).\mathrm{successorCarries}$ holds.
Then
\[
  \mathrm{rem}(s+1)\le2^{s+1}.
\]
\lean{Erdos249257.seamUpperBranch\_nextRemainder\_le\_pow}{HalfCylinderMiddleCarryLowerBound.lean:3924}
Indeed, the reset identity expresses $2^{s+1}$ as
$\mathrm{rem}(s+1)$ plus a nonnegative reset charge:
\lean{Erdos249257.seamUpperBranch\_remainder\_add\_resetCharge\_eq}{HalfCylinderMiddleCarryLowerBound.lean:3542}.
Thus the upper branch needs no exceptional-value exclusion for the
next-row two-sided invariant.  This conclusion concerns that induction
step, not the absence of upper transitions on the orbit.
\coord{quotient-branch-classification}\ \ev{Lean}\scale{n/a}
\end{prop}
\leannote{Lean: \lproof{ErdosProblems/Erdos257/PaperCompleteR21/BranchCellHorizonExclusions.lean}{31}{paper_upper_branch_needs_no_exceptional_cell}.}

\begin{thm}[$C_D=-3$ is impossible at the final middle transition, $D\ge 13$]
\label{thm:cd-neg3-impossible}
Consider a hypothetical \emph{final middle transition}: a middle transition at row $D\ge 13$
($\lnot\mathrm{carries}$, and the middle-branch inequality
$4\mathrm{rem}(D)+\mathrm{gap}-\mathrm{belowPulse} < \mathrm{terminalWeight}$ holds at $D$), followed
by an all-right tail forever after ($\forall s\ge D{+}1$, $\mathrm{seamGreedyWord}(s+1) =
\mathrm{seamGreedyWord}(s).\mathrm{extend}\ \mathrm{true}$). Under these assumptions,
\[
  \mathrm{belowPulse}(D) + 2 \;\le\; 4\cdot\mathrm{rem}(D),
\]
i.e.\ $C_D \ge -2$ ;  the cell $C_D=-3$ (and every more negative value) is excluded.
\lean{Erdos249257.middleThenAllRight\_landingExcess\_two\_le}{HalfCylinderMiddleCarryLowerBound.lean:2319}

\textbf{Proof of the bound.} The all-right-tail assumption puts the
terminal-augmented finite prefix strictly above $1/2$
(\lean{Erdos249257.half\_lt\_upper\_competitor\_of\_eventually\_right}{HalfCylinderFatalGapRightTail.lean:627}).
Consequently the producer carry is below its complete incidence tail
(\lean{Erdos249257.middleProducer\_allRight\_forces\_carry\_lt\_tail}{HalfCylinderMiddleCarryLowerBound.lean:1814}),
or equivalently
$Z_D<1+1/(2^D-1)$, where
$Z_D=\mathrm{rem}(D)-\mathrm{seamWordFloorError}(D)$
(\lean{Erdos249257.middleProducer\_allRight\_forces\_floorZ\_lt\_takeThreshold}{HalfCylinderMiddleCarryLowerBound.lean:1869}).
The lower bound uses the other endpoint: under the same tail assumption,
the finite value at $D$ plus the complete Mersenne tail remains below
$1/2$.  The exact scaled remainder identity therefore gives
\[
 Z_D>4^D\sum_{e>D}\frac1{2^e-1}-2^D>\frac13.
\]
The last inequality follows by retaining the first two geometric
channels of the complete Mersenne tail.  Put
$E_D=\mathrm{seamWordFloorError}(D)$.  Pulse absorption gives
$\mathrm{belowPulse}(D)\le4E_D$, so
\[
 4\,\mathrm{rem}(D)-\mathrm{belowPulse}(D)
 \ge4(\mathrm{rem}(D)-E_D)=4Z_D>\frac43.
\]
The left side is an integer, hence is at least $2$, as required.
The upper threshold inequality alone would not give this lower bound.

\textbf{Scope.} This excludes $C_D=-3$ at a middle row $D\ge13$
under the additional all-right-tail assumption.  The general induction
requires exclusion of all three cells at every middle row, not just
under this extra hypothesis, together with its right-branch bound. The two statements share the same coordinate $C_D$
but have different hypotheses.  Theorem~\ref{thm:cd-neg3-impossible}
is universal over rows satisfying the additional all-right-tail
assumption; it does not assert that such a row exists.  The induction
hypothesis concerns every middle row, without that tail assumption.
Excluding $C_D=-3$ in the former statement therefore does not supply
the exclusion required in the latter.
\coord{final-required input; conclusion-excess}\ \textbf{Object or representation:} about the \emph{object}
under the stated hypothesis ;  a genuine arithmetic consequence of the all-right-tail assumption
via the fatal-gap orbit, not a coordinate artifact. \ev{Lean}\scale{n/a}
\end{thm}
\leannote{Lean: \lproof{ErdosProblems/Erdos257/PaperCompleteR21/BranchCellHorizonExclusions.lean}{50}{paper_final_middle_cell_at_least_neg_two}.}

\begin{cor}[Only $C_D\in\{-2,-1\}$ remain among the exceptional dyadic cells]
\label{cor:cd-remaining}
At a final middle row $D\ge13$, Theorem~\ref{thm:cd-neg3-impossible}
removes $-3$ from the exceptional set $\{-3,-2,-1\}$, leaving
$-2,-1$ \emph{within that set}.  It does not assert
$C_D\in\{-2,-1\}$: nonnegative values have not been excluded.
Proposition~\ref{prop:upper-unconditional} handles upper transitions
in the two-sided induction; it supplies no additional restriction on
a middle coordinate.  For the all-middle-row induction, the exclusion
of all three values still needs proof without an all-right-tail
assumption, together with the separate right-branch inequality.
\coord{final-required input; conclusion-excess}\ \ev{Lean}
\end{cor}
\leannote{Lean: \lproof{ErdosProblems/Erdos257/PaperCompleteR21/BranchCellHorizonExclusions.lean}{81}{paper_final_middle_cell_remaining_cells}.}

\begin{rem}[A sufficient tail inequality and its source]
\label{rem:tail-dominance-open}
Let $U_D$ be the finite set of exponents selected by the below
word at row $D$, and put $F_D=U_D\cup\{D\}$.  With
$c_{F_D}(n)=\#\{d\in F_D:d\mid n\}$, define
\[
 \Theta_D:=\sum_{j\ge1}c_{F_D}(2D+2+j)2^{-j}.
\]
This is a tail sum, not a ratio.  Since
$0\le c_{F_D}(n)\le |F_D|$, it satisfies
$0\le\Theta_D\le |F_D|$.  The integer
$C_D=4\,\mathrm{rem}(D)-\mathrm{belowPulse}(D)-4$
is the producer carry in the cited source.  Under the final-middle
and all-right-tail assumptions, that source proves $C_D<\Theta_D$:
\lean{Erdos249257.middleProducer\_allRight\_forces\_carry\_lt\_tail}{HalfCylinderMiddleCarryLowerBound.lean:1814}.

An inequality in the opposite direction would exclude such a final
middle transition.  The following sufficient condition retains the
exception already handled by Theorem~\ref{thm:cd-neg3-impossible}:
\[
 \textbf{(TD)}\qquad
 \forall D\ge13,\quad
 \bigl(D\text{ is a middle row and }C_D\ne-3\bigr)
 \ \Longrightarrow\ \Theta_D<C_D.
\]
Because $\Theta_D\ge0$, this condition in particular excludes
$C_D=-2$ and $C_D=-1$.  It requires an inequality at the other indicated
middle rows as well, so it is not merely a reformulation of excluding
two integer values.  Nor does it supply the separate right-branch
hypothesis used in the two-sided induction.

The supplied release already defines this condition as
\textsf{SeamMiddleProducerTailEscapeExceptNegThree} and proves
half-membership from it:
\lean{half\_mem\_mersenneAchievementSet\_of\_middleProducerTailEscapeExceptNegThree}{HalfCylinderLastProducerContradiction.lean:436}.
The inequality itself is not established here.  Its role is a
conditional implication, not a newly identified necessary condition.

The two more concrete sufficient bounds in
Theorem~\ref{thm:middle-producer-escape} imply the tail inequality even
without the exception.  In the source they are
\lean{Erdos249257.SeamMiddleProducerCardEscape}{HalfCylinderMiddleCarryLowerBound.lean:279}
and
\lean{Erdos249257.SeamMiddleProducerRowEscape}{HalfCylinderMiddleCarryLowerBound.lean:1769}.
The cardinality condition is
\[
 |U_D|+\mathrm{belowPulse}(D)+5<4\,\mathrm{rem}(D).
\]
It gives $|F_D|=|U_D|+1<C_D$, hence $\Theta_D<C_D$.
The simpler row condition $D\le\mathrm{rem}(D)$ is stronger:
$\mathrm{belowPulse}(D)\le2|U_D|$ and $|U_D|\le D-2$ give
$|U_D|+\mathrm{belowPulse}(D)+5\le3D-1<4D$.
The corresponding half-membership implication is
\lean{Erdos249257.half\_mem\_mersenneAchievementSet\_of\_middleProducerCardEscape}{HalfCylinderMiddleCarryLowerBound.lean:2364}.
No converse between these sufficient conditions is claimed.
\ev{Open}\scale{cofinal}
\end{rem}

\subsection{Related counterexamples and restrictions}

The next results give counterexamples to particular summaries and
estimates for integer carries.  Their hypotheses matter: a countermodel
for one summary is not a prohibition on every proof that uses carries
or on every other input.

\begin{prop}[A finite-state restriction for the stated pulse family]
\label{prop:finite-state-nogo}
For a ``balanced pulse'' family at location $m\ge2$ (radius $\rho=\lfloor(m{+}1)/2\rfloor$, parameters
$0\le r\le\rho$, moving mass between positions $m$ and $m{+}1$ with the weighted total
$2c(m)+c(m{+}1)=2\rho$ of the two coefficients fixed), if
a predecessor state is constant across the whole family, then no function
$\mathrm{decode}:\mathrm{State}\to\mathbb N$ can recover the parameter $r$ from $\mathrm{state}(r)$
for every $r$. The family has exactly $\rho+1=\lfloor(m{+}1)/2\rfloor+1$ members, so the fan-out is
unbounded in $m$. More strongly, any finite set of states carrying an exact decoder for the family
has at least $\rho+1$ elements, unbounded in $m$.
\lean{Erdos249257.balancedPulse\_no\_autonomous\_decoder}{GenericTailOrbitRigidity.lean:247}
(exact family size at
\lean{ErdosProblems.Erdos257.PaperCompleteR21.paper\_balanced\_pulse\_fanout\_is\_radius\_succ}{lean/ErdosProblems/Erdos257/PaperCompleteR21/BalancedPulseFanOutCount.lean:61}).

\textbf{Scope.} Excludes bounded-state encodings of pre-$m$ history that must distinguish every
member of the displayed balanced-pulse family, whether the coefficients are $\varphi$ or a
Möbius-support indicator. Applying it to \#249 or \#257 requires showing that the relevant orbit
realises that family, and it does not rule out every finite-state proof strategy.
\coord{binary digits, generic}\ \textbf{Object or representation:} about a \emph{representation}
class, the encodings of history that separate the balanced-pulse family; it says nothing about $C$
or $\mathcal A$ directly.
\ev{Lean}\scale{n/a}
\end{prop}
\leannote{Lean: \lproof{ErdosProblems/Erdos257/PaperCompleteR21/BalancedPulseFanOutCount.lean}{113}{paper_pulse_family_no_autonomous_decoder}, \lproof{ErdosProblems/Erdos257/PaperCompleteR21/BalancedPulseFanOutCount.lean}{61}{paper_balanced_pulse_fanout_is_radius_succ}, \lproof{ErdosProblems/Erdos257/PaperCompleteR21/BalancedPulseFanOutCount.lean}{95}{paper_balanced_pulse_fanout_unbounded_corrected}, \lproof{ErdosProblems/Erdos257/PaperCompleteR21/BalancedPulseFanOutCount.lean}{123}{paper_pulse_family_finite_state_card}, and 2 further declarations in the \hyperref[sec:coverage]{coverage section}.}

\begin{prop}[Möbius-support countermodel: the natural negative-sign candidate overshoots]
\label{prop:mobius-nogo}
The signed Lambert identity $\sum_{d\ge1}\mu(d)/(2^d-1) = 1/2$ is exact. Writing
$N:=\{d:\mu(d)=-1\}$: $\sum_{d\in N}1/(2^d-1) = 1/2 + \sum_{d\in P}1/(2^d-1)$ where $P:=\{d\ge2:
\mu(d)=1\}$, and quantitatively $1/2 + 1/63 \le \sum_{d\in N}1/(2^d-1)$ (using the first positive
tail term $d=6$, $\mu(6)=1$) ;  the negative-Möbius support strictly \emph{overshoots} $1/2$ by at
least $1/63$.
\lean{MobiusSignSupportNoGo.half\_lt\_tsum\_negativeMobius}{MobiusSignSupportNoGo.lean:164}
(exact decomposition at
\lean{MobiusSignSupportNoGo.tsum\_negativeMobius\_eq\_half\_add\_positiveMobiusTail}{MobiusSignSupportNoGo.lean:111}).

\textbf{Scope.} Rules out exactly one natural candidate infinite Boolean support (the negative-Möbius
set) as a witness for $1/2\in\mathcal A$. A route-sufficiency no-go only ;  it says nothing about
whether \emph{some other} infinite Boolean support sums to $1/2$.
\coord{divisor counts and finite sums}\ \textbf{Object or representation:} about the \emph{object} ;  a genuine
value inequality for one specific candidate set, not a coordinate artifact. Margin $1/63$ is the
concrete number any repair attempt (adding/removing finitely many elements) must close or exceed.
\ev{Lean}\scale{fixed}
\end{prop}
\leannote{Lean: \lproof{ErdosProblems/Erdos257/PaperCompleteR21/MobiusSignAndFiniteCertificates.lean}{70}{paper_mobius_support_overshoots_half}, \lproof{ErdosProblems/Erdos257/PaperCompleteR21/MobiusSignAndFiniteCertificates.lean}{48}{paper_first_positiveMobius_tail_term}.}

\begin{prop}[Finite computation and half-membership]
\label{prop:finite-boolSupport-and-onesided}
No finite positive-index Boolean support has value exactly $1/2$: the reduced denominator of any
finite Mersenne subset-sum is provably \textbf{odd} (each $2^n-1$ is odd), while $1/2$ needs an even
denominator.
\lean{finite\_boolSupport\_ne\_half}{HalfCarryReachability.lean:589}. Separately,
$\mathsf{CertifiedGreedyMersenneDeath}$ is a finite-depth certificate
of $x\notin\mathcal A$, obtained by decidable tests on rational input.
For example, the supplied source excludes $3/4$ at level $1$ with
lookahead $0$:
\lean{three\_fourths\_certifiedGreedyMersenneDeath}{GreedyAchievementSet.lean:1778}.
Failure to find such a certificate at a given depth does not establish
membership.  It records survival of that finite test, not survival
at every depth.
\lean{Erdos249257.CertifiedGreedyMersenneDeath}{GreedyAchievementSet.lean:1756}.

\textbf{Scope.} Any support representing $1/2$, if it exists, must
be infinite.  This follows from denominator parity, independently
of the certificate search.  The search supplies a different fact:
a successful exclusion certificate proves nonmembership, whereas
its absence in a bounded search does not prove membership or
rationality.  The same distinction must be checked separately for
any other certificate family.
\coord{denominator parity; finite search}
\ev{Lean}\scale{fixed}
\end{prop}
\leannote{Lean: \lproof{ErdosProblems/Erdos257/PaperCompleteR21/MobiusSignAndFiniteCertificates.lean}{101}{paper_finite_support_and_onesided_certificate}.}

\begin{prop}[Period exclusion for the totient series]
\label{prop:carry-survivor-extinction}
For $\sum_n\varphi(n)/2^n$, the cited finite test proves that
$\mathrm{totientTail}(N+h)-\mathrm{totientTail}(N)$ is not an integer
by excluding every possible integer state in a bounded range within
finitely many steps.  If the series were rational, some positive
period $h_0$ would make these tail differences integral for every
sufficiently large $N$.  Telescoping would then give integrality
also for every positive multiple $mh_0$.

Consequently, the required certificate supply is: for every $h_0\ge1$
and every lower bound $N_0$, there exist $m\ge1$, $N\ge N_0$ and a
finite test length $K$ excluding integrality for $(mh_0,N)$.  The
unbounded choice of $N$ is essential; a counterexample before the
unknown eventual threshold does not exclude an eventual period.
\lean{Erdos257PeriodNoncollapse.irrational\_totient\_series\_of\_multiple\_survivor\_supply}{CarrySurvivorExtinction.lean:491}.
It is sufficient instead to obtain such certificates for
$h=\mathrm{lcm}(1,\ldots,t)$ with both $t$ and $N$ exceeding arbitrary
prescribed lower bounds:
\lean{Erdos257PeriodNoncollapse.irrational\_totient\_series\_of\_lcm\_survivor\_supply}{CarrySurvivorExtinction.lean:540}.
The supplied finite theorem excludes the tail differences for
$1\le h\le16$ at $(N,L)=(14,9)$, and hence excludes rational values
whose reduced denominator divides $2^{14}(2^h-1)$ for one of those
$h$.  It does not establish the unbounded certificate supply.
\lean{Erdos257PeriodNoncollapse.totient\_series\_ne\_rat\_of\_den\_dvd\_upto\_sixteen}{CarrySurvivorExtinction.lean:587}.

\textbf{Scope.} The declarations in \textsf{CarrySurvivorExtinction.lean},
\textsf{AdjacentCarryTube.lean}, \textsf{AdjacentPhaseSeparation.lean}
and \textsf{TotientCarryKernelRigidity.lean} used in this comparison
concern the totient series of Problem~249.  Their carry recurrences
may suggest constructions for Problem~257, but applying them there
requires a new argument for the divisor transform of its chosen
support.  The totient-specific theorems do not supply that argument.
\coord{integer carries; totient series}
\ev{Lean}\scale{cofinal}
\end{prop}
\leannote{Lean: \lproof{ErdosProblems/Erdos257/PaperCompleteR21/TotientPeriodCertificateSupply.lean}{40}{paper_carry_survivor_extinction}, \lproof{ErdosProblems/Erdos257/PaperCompleteR21/TotientPeriodCertificateSupply.lean}{23}{paper_periodLcm_is_prefix_lcm}.}

\subsection{The scalar-localisation height obstruction}

\begin{lem}[Denominator complement survives scaling]
\label{lem:scalar-localization}
For $x:\mathbb Q$, $c:\mathbb Z$, $H:\mathbb N$: if $H\mid x.\mathrm{den}$ and
$(c\cdot x).\mathrm{den}\mid H$ ;  i.e.\ multiplying by the integer $c$ shrinks the displayed
denominator down \emph{into} $H$ ;  then the \textbf{complementary} denominator factor
$x.\mathrm{den}/H$ divides $c$:
\[
  H\mid x.\mathrm{den} \ \land\ (c\cdot x).\mathrm{den}\mid H \implies x.\mathrm{den}/H \mid |c|.
\]
\lean{Erdos249257.AdelicHeightObstruction.scalarLocalization\_complement\_dvd}{AdelicHeightObstruction.lean:23}.
Equivalently, writing $D=x.\mathrm{den}/H$, we have $D\mid c$,
so $t=c/D$ is an integer and $Hcx=t\,x.\mathrm{num}$.
Here $H>0$ follows from $H\mid x.\mathrm{den}$, since the
reduced denominator is positive.
\lean{Erdos249257.AdelicHeightObstruction.scalarLocalization\_integer\_eq\_mul\_num}{AdelicHeightObstruction.lean:56}.

\textbf{Scope.} This is elementary rational arithmetic, independent
of the problem.  If $c\ne0$, divisibility gives the size bound
$|c|\ge x.\mathrm{den}/H$.  The nonzero condition is essential:
$c=0$ clears every denominator and satisfies the divisibility
conclusion without any positive lower bound on $|c|$.  The lemma
therefore constrains denominator clearing by a nonzero bounded
multiplier, not arbitrary multiplication.
\coord{rational-denominator; height}\ \textbf{Object or representation:} about a
\emph{representation} ;  pure denominator bookkeeping for $\mathbb Q$, with zero problem-specific
content; the upstream primitive
(\lean{Erdos249257.RationalDenominatorSurvival.divisor\_dvd\_divInt\_den}{RationalDenominatorSurvival.lean:17})
is the domain-neutral extraction this lemma builds on.
\ev{Lean}\scale{n/a}
\end{lem}
\leannote{Lean: \lproof{ErdosProblems/Erdos257/PaperCompleteR21/ScalarLocalisationHeightObstruction.lean}{29}{paper_scalar_localization}, \lproof{ErdosProblems/Erdos257/PaperCompleteR21/ScalarLocalisationHeightObstruction.lean}{53}{paper_scalar_localization_size_bound}, \lproof{ErdosProblems/Erdos257/PaperCompleteR21/ScalarLocalisationHeightObstruction.lean}{62}{paper_scalar_localization_zero_degenerate}.}

\begin{cor}[Mersenne specialisation]
\label{cor:mersenne-height}
Let $x$ be a positive rational number and let $r\ge0$, $n\ge1$
be integers.  If $2^r\mid x.\mathrm{num}.\mathrm{natAbs}$ and
$x<2/(2^n-1)$, then
$2^r\cdot(2^n-1) < 2\cdot x.\mathrm{den}$ ;  a numerator $2$-power lower bound plus a Mersenne-scale
upper bound on $x$ together force a denominator lower bound, \emph{without} introducing a global
prefix LCM.
\lean{Erdos249257.AdelicHeightObstruction.positiveRat\_mersenne\_height}{AdelicHeightObstruction.lean:103}
(generic form
\lean{Erdos249257.AdelicHeightObstruction.positiveRat\_numDivisor\_mul\_lt\_two\_mul\_den}{AdelicHeightObstruction.lean:77}).
To see the inequality directly, write $x=a/b$ in lowest terms
with $a,b>0$.  The hypotheses give $a\ge2^r$ and
$a(2^n-1)<2b$.  The conclusion follows by substitution.
An application of Lemma~\ref{lem:scalar-localization} must in
addition supply its denominator-divisibility assumptions.
\coord{rational-denominator; height, Mersenne instance}\ \ev{Lean}\scale{n/a}
\end{cor}
\leannote{Lean: \lproof{ErdosProblems/Erdos257/PaperCompleteR21/ScalarLocalisationHeightObstruction.lean}{77}{paper_mersenne_height}.}

\subsection{The final-skip band formula, and what it does not show}

\begin{defn}[Upper-reset dyadic band escape]
\label{defn:band-escape}
At each upper-reset row $d\ge13$, write
$E_d=4\,\mathrm{overshoot}(d)+\mathrm{abovePulse}(d)$ for the reset
expression in the preceding transition formulas.  The condition is
\[
 \text{for every }0\le j\le d,\qquad
 E_d>2^{d-j+1}\quad\text{or}\quad
 E_d+2(d+j)\le2^{d-j+1}.
\]
Thus $E_d$ must avoid the interval
$(2^{d-j+1}-2(d+j),\,2^{d-j+1}]$ at every indicated scale.
\lean{Erdos249257.SeamUpperResetDyadicBandEscape}{HalfCylinderMiddleCarryLowerBound.lean:4566}.
In words: at every upper-reset row $d\ge13$, the reset charge must avoid a linear-width forbidden
band immediately below every dyadic threshold $2^{d-j+1}$, for every $j\le d$ simultaneously.
Granted this hypothesis,
\lean{Erdos249257.half\_mem\_mersenneAchievementSet\_of\_upperResetDyadicBandEscape}{HalfCylinderMiddleCarryLowerBound.lean:4790}
gives $1/2\in\mathcal A$ outright.
\end{defn}

\begin{prop}[Quantifier collapse: one critical index suffices, not $d{+}1$]
\label{prop:critical-band-index}
Let $d,E\in\mathbb N$ and assume $E\le2^{d+1}$. Then the purely combinatorial statement
$\mathsf{DyadicBandEscape}(d,E) \iff \exists j,\ \mathsf{CriticalDyadicBandIndex}(d,E,j)\land
E+2(d+j)\le 2^{d-j+1}$ collapses the $\forall j\in[0,d]$ band-avoidance condition (formally $d{+}1$
separate inequalities) to checking exactly \emph{one} nearest-boundary index $j$.
The range assumption ensures that a dyadic threshold lies at or above $E$.
Choose the smallest such threshold: smaller thresholds are already below
$E$, while larger thresholds have narrower forbidden bands. Without the
range assumption all bands escape automatically when $E>2^{d+1}$, but no
critical index exists; $(d,E)=(0,3)$ is the smallest example.
Specialised to the
concrete seam reset charge, $\mathsf{SeamUpperResetCriticalBandEscape}$ is proved logically
\emph{equivalent} to Definition~\ref{defn:band-escape}'s hypothesis.
\lean{Erdos249257.dyadicBandEscape\_iff\_exists\_critical}{HalfUpperResetCriticalBand.lean:108}
(seam specialisation
\lean{Erdos249257.seamUpperResetCriticalBandEscape\_iff}{HalfUpperResetCriticalBand.lean:883}).
Zero Mersenne/seam content in the core lemma ;  pure $(d,E,j)$ arithmetic over powers of 2.
\coord{dyadic-boundary, generic}\ \ev{Lean}\scale{n/a}
\end{prop}
\leannote{Lean: \lproof{ErdosProblems/Erdos257/PaperCompleteR21/CriticalDyadicBandCollapse.lean}{20}{paper_critical_band_index_collapse}.}

\begin{obs}[A verified finite range of reset indices]
\label{cert:band-through-31}
The linked certificate proves the condition in
Definition~\ref{defn:band-escape} for each reset index $13\le d\le30$.
It uses successor remainders at rows $14$--$31$, including
$\mathrm{rem}(14)=392$ and $\mathrm{rem}(31)=4{,}187{,}487{,}147$.
Thus the last remainder row is $31$, but the last reset index covered
by this theorem is $30$.
\lean{Erdos249257.seamUpperResetDyadicBandEscape\_through\_thirty}{HalfCylinderUpperResetBandCertificates.lean:78}.
\ev{Cert}\scale{bounded}
\end{obs}

\begin{rem}[What this does \emph{not} show ;  read the quantifiers exactly]
\label{rem:band-caveat}
Definition~\ref{defn:band-escape} is a $\forall d\ge13$ statement; Certificate
\ref{cert:band-through-31} is a finite verification in the cited Lean source for $13\le d\le30$. \textbf{The
finite certificate does not establish the universal hypothesis.} Nothing in the tree proves
$\mathsf{SeamUpperResetDyadicBandEscape}$ for all $d\ge13$ ;  this is exactly the unproved hypothesis named
in Part~I (SE-9) and Part~III (RC-4/D1's \textsf{LargestSkipLateStepSocket}), stated here in its own
coordinate.

\par\medskip

The finite theorem supplies no reset index beyond $30$. It therefore
does not prove that the orbit avoids the band at unboundedly many reset
indices, let alone at every reset. The universal band condition remains
an unproved sufficient hypothesis.

\par\medskip

For the branch convention of Theorem~\ref{thm:dynamics}, direct
integer recomputation at reset-test indices $6\le r\le2500$
gives $605$ middle and $604$ upper resets, hence $1209$ resets.
There are $1286$ right transitions.  The earlier reported count
$1264$ does not agree with this convention.  The recomputation
checks the sign inequalities at these resets, the square-root
bound at resets $r\ge10$, and the displayed band condition at
upper resets $r\ge13$, with no failures.  This remains a finite
ordinary computation, distinct from the cited Lean theorem for
$13\le r\le30$ and from an assertion at every later reset.
\ev{Open}\scale{cofinal}
\end{rem}
% ; - part a257_invent ; -
\section{Unproved inputs for the recorded implications}
\label{sec:invent-257}

The preceding parts prove support criteria, finite identities and
conditional implications.  They do not decide the universal problem or
either rational target.  This section states the remaining hypotheses
for particular implications and compares them with the estimates
already available.  Ordinary arguments, reported finite computations
and unproved proposals are distinguished by the source tags explained
at the start of the paper.  No single necessary obstruction is claimed
to account for every unsuccessful estimate.

For example, the first two weighted divisor counts after $M=3$
give $(\tau(4)-1)/2+(\tau(5)-1)/4=5/4$.
The finite weighted divisor sum below records the same calculation
for any starting index and length.  Its role in particular universal-support
criteria is discussed separately; it is not a proved common necessary
obstruction for the two problems.

\begin{defn}[The short-window divisor phase]
\label{defn:theta}
For $M\ge 1$ and $L\ge 1$ put
\[
\Theta_L(M)\ :=\ \sum_{i=1}^{L}\bigl(\tau(M+i)-1\bigr)\,2^{-i}\ \in\ \Q_{\ge 0},
\]
where $\tau$ is the number-of-divisors function, so that $\tau(n)-1=\#\{d\ge 2: d\mid n\}$ counts
the divisors of $n$ other than $1$.
\end{defn}

Iteration of the integer recurrence
$K(M+1)=2K(M)-(\tau(M+1)-1)$ gives the exact identity
\[
 \Theta_L(M)=K(M)-2^{-L}K(M+L)\qquad(M\ge2,\ L\ge1).
\]
Consequently $2^L\Theta_L(M)$ is an integer, and
\[
 \operatorname{dist}(\Theta_L(M),\mathbb Z)
 =2^{-L}\operatorname{dist}(K(M+L),2^L\mathbb Z).
\]
A nonzero distance is at least $2^{-L}$, but nonvanishing itself
requires an argument.  For instance $\Theta_1(3)=1$ has distance
zero.  The identity concerns $K$; the quotient deviation also
contains the skip-set correction in
Theorem~\ref{thm:master-identity}.  No positive lower bound for
this phase is asserted to be necessary for the universal support
problem.

\subsection{Reset bounds and finite weighted divisor sums}
\label{sec:invent-R1}

\subsubsection*{The statement needed}

The first two conditions below each have a conditional implication to
half-membership.  No implication between them is asserted.  The sign
condition is weaker than a quantitative margin in its conclusion,
but its relation to either membership hypothesis must still be
proved.

\emph{Form (i), the reset-deviation bound.} Write
$w_s=\mathrm{rem}(s)-2^s$.  The actual hypothesis is
\[
 \forall r\ge10\ \text{that are upper or middle resets}:\qquad
 |w_{r+1}|>2^{(r+5)/2}.
\]
Equivalently, $|\lambda_{r+1}|>2^{(3-r)/2}$ for
$\lambda_{r+1}=w_{r+1}/2^{r+1}$.  This normalisation concerns
distance of the deviation from zero, not distance of an unrelated
phase from the nearest integer.
The run-length estimate $L_r<(r-3)/2$ for resets $r\ge31$ is a
proposed comparison target.  Its exact relation to the displayed
bound requires the pulse terms and the initial range $10\le r<31$;
it is not stated as an equivalence here.  A sublinear estimate gives
an eventual ceiling, not automatically the required finite initial
cases.  The supplied computation reports a maximum run of length
$19$ through row $200{,}000$, attained at row $158{,}096$.
That observation is finite evidence, not an asymptotic theorem.

\emph{Form (ii), the Lean-facing form already proved as a conditional implication.}
\[
\mathtt{SeamUpperResetDyadicBandEscape}\ :=\ \forall d\ge 13,\ \bigl(\mathtt{seamAdjacentCut}\ d\bigr).\mathtt{successorCarries}\ \Longrightarrow\ \forall j\le d,
\]
\[
2^{\,d-j+1} < E_d\quad\lor\quad E_d + 2(d+j)\ \le\ 2^{\,d-j+1},
\qquad
E_d := 4\cdot\mathtt{overshoot}_d+\mathtt{abovePulse}_d,
\]
(\lean{SeamUpperResetDyadicBandEscape}{HalfCylinderMiddleCarryLowerBound.lean:4566}), with the
closed conditional implication
\lean{half_mem_mersenneAchievementSet_of_upperResetDyadicBandEscape}{HalfCylinderMiddleCarryLowerBound.lean:4790},
\ev{Lean}. The $\forall j\le d$ band check collapses to a single nearest-boundary index per row by
\lean{dyadicBandEscape_iff_exists_critical}{HalfUpperResetCriticalBand.lean:108}, and rows
$13\le d\le 30$ are verified in the cited source by
\lean{seamUpperResetDyadicBandEscape_through_thirty}{HalfCylinderUpperResetBandCertificates.lean:78}.

\emph{Form (iii), the sign of the reset deviation.}
\[
\forall\ \text{reset rows } r:\qquad
r \text{ is an } M\text{-reset}\Rightarrow w_{r+1}>0,
\qquad
r \text{ is a } U\text{-reset}\Rightarrow w_{r+1}<0 .
\]
The finite recomputation for this revision verifies these signs
at all $1209$ resets in rows $6$--$2500$, with $605$ middle and $604$ upper resets.
The upper sign follows directly from the exact branch identity below.
The middle sign remains an unproved assertion about the orbit; neither
sign alone gives the required lower bound on the magnitude.

\subsubsection*{Divisor residues and the changing skip set}

Expanding each divisor count separates its residue conditions.
This is a finite counting identity.  It does not prescribe whether
an eventual proof should be analytic, combinatorial, or use another
method.

\begin{lem}[Divisor-residue form of the short-window phase]
\label{lem:odometer}
For all $M,L\ge 1$,
\[
\Theta_L(M)\ =\ \sum_{d\ge 2}\ \sum_{\substack{1\le i\le L\\ i\,\equiv\,-M\ (\mathrm{mod}\ d)}} 2^{-i}
\ =\ \sum_{d\ge 2}\ 2^{-i_d(M)}\cdot\frac{1-2^{-d\,m_d}}{1-2^{-d}},
\]
where $i_d(M)\in[1,d]$ is the least $i\ge1$ with $d\mid M+i$ (so $i_d(M)$ depends only on
$M \bmod d$), $m_d := \#\{1\le i\le L: d\mid M+i\}$, and terms with $i_d(M)>L$ are empty. In
particular, for every integer $D\ge M+L$,
\[
\Theta_L(M)\ =\ \sum_{d=2}^{D}\ \sum_{i=1}^{L} 2^{-i}\,\mathbf 1_{d\mid M+i},
\]
and the right-hand side, with an arbitrary integer $x\ge0$ in place of $M$, depends on $x$ only
through its residues modulo the integers $d$ with $2\le d\le D$. The cutoff $D$ has to be fixed
before the argument varies: at $(M,L)=(1,1)$ the cutoff is $M+L=2$, the integers $1$ and $3$ have
the same residue modulo every $d$ with $2\le d\le 2$, and yet $\Theta_1(1)=1/2$ while
$\Theta_1(3)=1$.
\end{lem}
\leannote{Lean: \lproof{ErdosProblems/Erdos257/PaperCompleteR21/ShortWindowDivisorPhase.lean}{417}{theta_eq_tsum_divisorResidue}, \lproof{ErdosProblems/Erdos257/PaperCompleteR21/ShortWindowDivisorPhase.lean}{442}{theta_eq_tsum_geometricForm}, \lproof{ErdosProblems/Erdos257/PaperCompleteR21/ShortWindowDivisorPhase.lean}{368}{theta_eq_divisorResidueSum}, \lproof{ErdosProblems/Erdos257/PaperCompleteR21/ShortWindowDivisorPhase.lean}{378}{theta_eq_geometricForm}, and 13 further declarations in the \hyperref[sec:coverage]{coverage section}.}
\begin{proof}
$\tau(n)-1=\sum_{d\ge2}[\,d\mid n\,]$; exchange the order of summation and sum the geometric
progression $i_d, i_d+d, i_d+2d,\dots$ inside $[1,L]$. No $d>M+L$ divides any $M+i$ with
$1\le i\le L$, which gives the cutoff form, and each indicator $\mathbf 1_{d\mid x+i}$ depends
only on $x \bmod d$.
\end{proof}
\ev{Lean}
(\lean{ErdosProblems.Erdos257.PaperCompleteR21.ShortWindowDivisorPhase.theta\_eq\_tsum\_divisorResidue}{lean/ErdosProblems/Erdos257/PaperCompleteR21/ShortWindowDivisorPhase.lean:417},
cutoff form at
\lean{ErdosProblems.Erdos257.PaperCompleteR21.ShortWindowDivisorPhase.theta\_eq\_Psi}{lean/ErdosProblems/Erdos257/PaperCompleteR21/ShortWindowDivisorPhase.lean:339}),
and verified exactly as rationals at $(M,L)\in\{1000,10001,123456\}\times\{12,20\}$
\ev{Cert}.

A finite divisor cutoff makes the averaging question precise.  For
integers $D\ge2$ and $L\ge1$, define
\[
 \Psi_{L,D}(x)=\sum_{d=2}^{D}\sum_{i=1}^{L}2^{-i}
                  \mathbf1_{d\mid x+i},\qquad
 Q_D=\operatorname{lcm}(1,\ldots,D).
\]
This is the cutoff sum of Lemma~\ref{lem:odometer}.  It is periodic in $x$
with period $Q_D$, and
$\Theta_L(M)=\Psi_{L,D}(M)$ whenever $M\ge1$ and $D\ge M+L$.
Averaging each indicator over a complete period gives
\[
 \frac1{Q_D}\sum_{x=0}^{Q_D-1}\Psi_{L,D}(x)
   =(1-2^{-L})\sum_{d=2}^{D}\frac1d.
\]
The mean diverges as $D\to\infty$, even for fixed $L$.
Thus this family does not converge to a finite integrable observable
whose limiting mean could simply be used in an equidistribution
argument.

The same issue arises if one uses the profinite integers
$\widehat{\mathbb Z}$, whose Haar probability measure is uniform
on each finite residue quotient.  The finite-cutoff functions are
continuous there.  Their increasing limit is infinite almost
everywhere: the events $p\mid x+1$ are independent over primes by
the Chinese remainder theorem, have probabilities $1/p$, and occur
infinitely often almost surely.  Each contributes $1/2$ to the
term at offset $1$.  The positive integer inputs $M$, for which
$\Theta_L(M)$ is finite, do not justify treating this divergent
limit as a continuous or integrable function.  Only the finite
period average above is used in the present argument.

There is a related residue expression in
Theorem~\ref{thm:real-form}, but its index and divisor set differ:
\[
 \sum_{d\in D_n}\left\{\frac{4^n}{2^d-1}\right\}
 =\sum_{d\in D_n}\frac{2^{-i_d(2n)}}{1-2^{-d}}.
\]
Indeed $i_d(2n)=d-(2n\bmod d)$, with value $d$ when the
remainder is zero.  The argument is $2n$, not $n$, because the
quotient scale is $4^n=2^{2n}$.  For example, at $n=7,d=3$ the
fractional part is $4/7$; using $i_3(7)$ would instead give $2/7$.
This sum uses the selected finite set $D_n$, whereas $\Theta_L$
uses all divisors with offsets truncated at $L$.  Their common
geometric summand does not make the two sums equal.

To keep the changing support visible, put
$S_n=\sum_{d\in\mathrm{Skip}_n}\lfloor4^n/(2^d-1)\rfloor$.
The master identity $\Delta_n=K(2n)+S_n$ and the recurrence for $K$
give, for $n\ge6$ and $\ell\ge1$,
\[
 2^{-2\ell}\Delta_{n+\ell}
 =\Delta_n-\Theta_{2\ell}(2n)
       +2^{-2\ell}S_{n+\ell}-S_n.
\]
A phase estimate would have to control this additional term and
identify the required reset rows before implying the reset bound.
No such implication is proved here.  The scales also differ from
a logarithmic-window assertion: $\ell$ quotient transitions use
$2\ell$ binary steps starting at $M=2n$.  If $\ell$ is proportional
to $n$, the phase length is comparable to $M$, not to $\log M$.
A theorem for logarithmic windows would need a separate argument
connecting its range to the reset condition.

Finite residue expansions, divisor-sum estimates and probabilistic
local estimates are possible sources of further information.  None
is an input already supplied here.  An application must state the
divisor cutoff, offset length, error at the required precision, and
control on the changing skip-set sum, rather than appeal to the
name of a method or to equidistribution of a different statistic.

\subsubsection*{The specific obstacle}

\textbf{What the cited short-interval theorems control.}  Matom\"aki and
Radziwi\l{}\l{} prove that a real-valued multiplicative function bounded by $1$ has, in almost all
windows $[x,x+h]$ with $h\to\infty$, essentially its long
average~\cite[Theorem~1, pp.~1--2]{matomakiradziwill}.  Mangerel extends such bounds to a class of
divisor-bounded multiplicative
functions~\cite[Theorem~1.7, p.~7, and Corollary~1.8, pp.~7--8]{mangerel2021}, and Sun describes the
almost-all transition profile of the $k$-fold divisor function at its critical polylogarithmic
window~\cite[Theorem~1.1 and Corollary~1.2, p.~3]{sun2026}.  These results control short-interval averages.  They do not by
themselves supply a reset-deviation estimate or its unproved
weighted-phase substitute, for two reasons.

First, their conclusions hold for almost all $x$, whereas the proposed
application needs every reset row beyond its threshold.  An
exceptional set of density zero may still contain infinitely many
required indices.  No density assertion about the reset set, or
avoidance of the exceptional set by that orbit, is proved here.

Second, $\tau$ is not bounded by $1$, and the relevant statistic is the distance to $\Z$ of a
\emph{dyadically weighted} sum of $\tau$-values at a precision shrinking exponentially in the window
parameter.  The first term $(\tau(M+1)-1)/2$ has fixed weight
$1/2$, independent of the window length; an unweighted average
over that window does not by itself control this term.  An application through the phase would need both an exceptional-set
argument along the specified indices and control of the weighted
sum.  It would still have to justify the passage from that sum to
the reset deviation, including the skip-set correction.

\textbf{The required statistic is different.} A central limit theorem
for prime-factor counts, or an almost-all theorem for unweighted
interval averages, is not by itself a bound for the distance to an
integer of this dyadically weighted divisor sum.  An application must
state its precision, its dependence on the growing window length,
and its exceptional set.  The references cited above do not supply
that orbitwise estimate in the form used here.  This is a limitation
of the proposed application, not a survey claim that every relevant
local estimate is unavailable.

The countermodels in the record impose similarly specific
limitations.  They concern fixed-precision residues
(\lean{affineBinaryOrbit_mod_twoPow_eq}{GenericTailOrbitRigidity.lean:307}),
a decoder without fresh arithmetic input
(\lean{balancedPulse_no_autonomous_decoder}{GenericTailOrbitRigidity.lean:247}),
and fixed-precision signatures
(\lean{fixedPrecisionTropicalNoGo}{TropicalCurvatureCarry.lean:137}).
They do not partition all possible proof methods.  Together with the
denominator estimate, they explain why the particular shortcuts
examined here are insufficient; they do not prove that a
weighted-window method is necessary or impossible.

\subsubsection*{Possible estimates to prove}

\emph{A conditional bound along a right-branch run.}

Let $r\ge6$, $L\ge1$, and suppose that the transitions at rows
$r,\ldots,r+L-1$ all take the right branch.  Put
$w_s=\mathrm{rem}(s)-2^s$ and let $p_s$ be the lower-prefix
pulse, so $0\le p_s\le2(s-2)$.  Iterating the exact recurrence
$w_{s+1}=4w_s-4-p_s$ gives
\[
 w_{r+L}=4^Lw_r-
       \sum_{j=0}^{L-1}4^{L-1-j}(4+p_{r+j}).
\]
The reverse triangle inequality and the finite geometric sum give
\[
 |w_{r+L}|
 \ge4^L\left(|w_r|-\frac{4+2(r+L)}3\right).
\]
This estimate alone does not force a right-branch run to end.
For that conclusion one also needs an upper bound on its terminal
deviation.  More precisely, if $C>0$ and
\[
 |w_{r+L}|\le C2^{r+L},\qquad
 |w_r|>\frac{4+2(r+L)}3,
\]
then division by $4^L$ and taking logarithms give
\[
 L\le r+\log_2 C-
 \log_2\left(|w_r|-\frac{4+2(r+L)}3\right).
\]
Every term and hypothesis is explicit here.  A separate window
estimate must supply the displayed terminal upper bound; it does
not follow from the right-branch recurrence alone.  This explains
the possible connection between the initial deviation and the
run length without asserting a sharp threshold, a run-length
equivalence, or an unproved bound for all later rows.  \ev{Math}
for this conditional calculation.

\emph{An exact finite phase experiment.}

For the $200{,}000$ integers $10^6\le M<1.2\times10^6$, an
integer divisor sieve through $1{,}200{,}039$ computes
$2^{40}\Theta_{40}(M)$ exactly.  Reducing that integer modulo
$2^{40}$ gives the following counts and empirical proportions for
$\operatorname{dist}(\Theta_{40}(M),\mathbb Z)<t$.
\begin{center}
\begin{tabular}{lccccc}
\hline
$t$ & $2^{-6}$ & $2^{-8}$ & $2^{-10}$ & $2^{-12}$ & $2^{-14}$\\
\hline
Count & $7083$ & $1858$ & $491$ & $125$ & $37$\\
Proportion & $0.035415$ & $0.009290$ & $0.002455$ & $0.000625$ & $0.000185$\\
Proportion$/t$ & $2.267$ & $2.378$ & $2.514$ & $2.560$ & $3.031$\\
\hline
\end{tabular}
\end{center}
These counts were reproduced for this revision; the experiment
samples consecutive integers, not reset rows.  Its median distance
is $0.24470815549\ldots$, compared with $1/4$ for a uniform
fractional part.  A uniform fractional part has probability $2t$
of distance less than $t$, for $0<t\le1/2$, which explains the
comparison in the last row.  Agreement of a finite histogram with
that model is not a distribution theorem.

Truncation matters at the smallest tested thresholds.  With the
same starting indices, the counts below $2^{-14}$ are $22,36,37,37$
for $L=16,24,32,40$, respectively.  Thus similar central quantiles
do not justify ignoring the truncation error in a rare-event test.
More generally, for $L'\ge L\ge1$ and $B\ge0$, if
$\tau(M+i)-1\le B$ over $L<i\le L'$, then
\[
 0\le\Theta_{L'}(M)-\Theta_L(M)
 \le B(2^{-L}-2^{-L'})\le B2^{-L}.
\]
Distance to the nearest integer is $1$-Lipschitz.  A threshold
comparison is therefore guaranteed to remain unchanged when
the distance at the shorter cutoff is farther than this error
from the threshold.  This is a sufficient stability test, not a
necessary one.
The estimate uses a bound over the specified finite range, not an
unproved uniform bound on the divisor function.

A probabilistic application requires further assumptions.  For
example, if events $E_r$ in a specified probability space satisfied
$\mathbb P(E_r)\le C2^{(3-r)/2}$ for every $r\ge R$, the union
bound would give
\[
 \mathbb P\left(\bigcup_{r\ge R}E_r\right)
 \le \frac{C2^{(3-R)/2}}{1-2^{-1/2}}.
\]
Summability would also imply only finitely many failures almost
surely, not zero failures at every index.  Neither the finite
histogram nor its fit to a uniform-phase model proves these
probability bounds, identifies their events with reset failures,
or defines a probability for the actual deterministic orbit.
In particular, no membership conclusion is drawn from the fitted
small expected counts reported in the earlier record.  \ev{Cert}
for the finite counts; \ev{Math} for the elementary error and
probability inequalities; \ev{Open} for an application to resets.

\emph{The reset signs and the exact branch identities.}
Write $p_s$ for the nonnegative lower-prefix pulse and
$w_s=\mathrm{rem}(s)-2^s$.  On a middle transition, where the upper
carry does not occur, the update is
\[
 \mathrm{rem}(s+1)=4\mathrm{rem}(s)+2^{s+1}-p_s,
 \qquad
 w_{s+1}=4w_s+2^{s+2}-p_s=4\mathrm{rem}(s)-p_s.
\]
The natural-number subtraction can be read as signed subtraction
because $p_s\le2(s-2)<2^{s+1}$ for $s\ge5$.
The middle-branch condition is
$4\mathrm{rem}(s)<2^{s+1}+4+p_s$, together with failure of the
upper carry.  Thus a positive middle-reset deviation is exactly the
additional inequality $4\mathrm{rem}(s)>p_s$.
The pulse bound and branch inequality alone do not prove it:
$s=5$, $\mathrm{rem}=1$, $p=5$ satisfy both numerical bounds but give
$4\mathrm{rem}-p=-1$.  This is a counterexample to those numerical
premises, not a claim that the concrete orbit reaches that state.

On an upper transition, let $o_s>0$ be the overshoot and $p_s^+\ge0$
the upper-prefix pulse.  The carry condition pays the subtraction,
so
\[
 w_{s+1}=-(4o_s+p_s^+)<0.
\]
This sign already follows from the exact upper-branch identity;
only the middle sign needs further orbit information.  On a right
transition the corresponding formula is
$w_{s+1}=4w_s-4-p_s$.
The same middle and upper deviation identities appear in the supplied
\href{https://github.com/wcook04/plectis-erdos-lean/blob/52f29ad173b04e3bac941b3663f2b9aebe5de0bb/Erdos257PeriodNoncollapse/HalfResetSqrtEscapeScaleProducers.lean#L79-L144}{release source},
classified \texttt{release\_only} by the packet.  They separate the
sign calculation from the unproved magnitude estimate needed for square-root escape.  At the
reported middle step $s=14$, for example,
$\mathrm{rem}(14)=392$, $p_{14}=3$ and
$\mathrm{rem}(15)=34333$, giving $w_{15}=1565$.

\subsection{Progress to larger endpoint depths}
\label{sec:invent-R2}

\subsubsection*{The statement needed}

The exact-row transition
\lean{exactLocalMersenneHalfRow_double_or_recycle}{BooleanMobiusExactRowDichotomy.lean:26}
(\ev{Lean}) says every exact row $n\ge6$ either doubles to $2n-1$ or recycles to $2c-2$ for some
$4\le c\le n$. Cofinality of exact rows; which gives $\mathrm{HALF}$; needs the recycle
branch not to return to a bounded endpoint forever. Three sufficient forms:
\[
\text{(a)}\quad
\mathtt{SkippedCoreCriticalQuotientSupply}:\ \forall D\subseteq[2,c),\ 4\le c,\
v(D)<\tfrac12\ \wedge\ \tfrac12-v(D)<\ensuremath{w}\,c
\]
\[
\Longrightarrow\quad 2^{\,(2c-2)-1}\ \le\ \ensuremath{Q}\,(\{c\}\cup D)\,(2c-2),
\]
where $v(D):=\ensuremath{V}\ D$
(\lean{SkippedCoreCriticalQuotientSupply}{BooleanMobiusCriticalCapacityCofinal.lean:30});
\[
\text{(b)}\quad
\text{the crossing rank } c \text{ produced by the recycle branch does not repeat a bounded value infinitely often};
\]
\[
\text{(c)}\quad
c\ge g(n),\qquad 2g(n)-2>n
\quad\text{at every reachable recycle input }n.
\]

Unboundedness of $g$ alone would not suffice in (c).  For example,
$g(n)=\lfloor n/2\rfloor+1$ tends to infinity but allows the fixed
transition $n=6$, $c=4$, $2c-2=6$.  The displayed strict-progress
inequality excludes such cycles; the doubling branch already
increases every endpoint $n\ge6$.

\subsubsection*{What kind of object would furnish it}

The real and integer greedy rules must be distinguished here.
A real straddling word agrees with the real greedy prefix
(\lean{IsStraddlePrefix.half_agrees_greedy}{HalfCutLocator.lean:442}).
An exact local row instead satisfies
$\ensuremath{Q}\ D\ n=2^{n-1}-1$.  Although floor division is not
injective in general, the quotient weights at fixed $n$ have a gap of
at least $1$ between distinct subset sums.  A witness is therefore
unique when it exists and is recovered by the integer greedy test
(Theorem~\ref{record:257bm-i11a} and
Propositions~\ref{record:257bm-i11b}--\ref{record:257bm-i11c}).
There is no exponential family of witnesses to choose from at that
depth.  What remains to prove is that successful depths are unbounded.

One possible additional ingredient is an integer-valued function
$\Phi(n,D)\ge0$ on reachable exact rows that increases by at least one
under each permitted transition and is bounded on rows of bounded
endpoint.  Along an infinite trajectory, $\Phi$ would then tend to
infinity, forcing the endpoints to be unbounded.  Strict increase only
on recycle steps would need a separate check of what doubling does to
$\Phi$; an unspecified real-valued increasing quantity need not tend
to infinity.
This describes one possible proof mechanism, not a necessary class of
all proofs. The corpus already has the sharp
capacity test in exact form: by
\lean{localBinarySuffix_two_mul_sub_two_lt_criticalCapacity_iff}{BooleanMobiusSkippedCoreCriticalCapacity.lean:22},
form (a) is exactly the statement that the $c-2$-bit suffix capacity
$\ensuremath{S}\,D\,1\,(2c-2)<2^{\,c-2}$ is met.

\subsubsection*{The specific obstacle}

The formal counterexample concerns exactly the transition schema:
that schema \emph{plus} a seed provably does not give cofinality. The theorem
\lean{exists_seeded_bounded_double_or_recycle_model}{BooleanMobiusExactRowDichotomy.lean:91}
(\ev{Lean}) exhibits the one-point predicate $P(n):=(n=6)$, which satisfies the seed
(\lean{exactLocalMersenneHalfRow_six}{BooleanMobiusExactRowSeed.lean:34}, the explicit witness
$\{2,3,6\}$) and the exact shape of the transition, and is not cofinal. So no argument that uses
only the dichotomy and the base case can work, however cleverly the induction is arranged. Any
proof must consume arithmetic content of $\mathtt{ExactLocalMersenneHalfRow}$ that the abstract
schema does not see.

There is a second, more specific trap, and it is also formalised. The natural sufficient
condition via fractional split mass is \emph{not necessary}: the fixture
\lean{splitFractionMass_one_bound_not_necessary_fixture}{BooleanMobiusSkippedCoreCriticalCapacity.lean:183}
(\ev{Lean}, at $D=\{2,3\}$, $c=5$) exhibits a real crossing core violating it. Thus the fractional-mass bound is not necessary for an individual
crossing core.  This counterexample does not rule out proving that
sufficient bound on a more restricted family.

\subsubsection*{Possible estimates to prove}

\emph{The uniquely determined prefix at a crossing.} The identification
below is proved; the required capacity bound remains open.

Form (a) is universally quantified over cores $D$, which looks like an exponential obligation.
It is not. Its hypotheses are exactly those of
\lean{eq_halfGreedyPrefixSupport_of_critical_crossing}{BooleanMobiusCriticalCapacityCofinal.lean:50}
(\ev{Lean}), which forces $D=\mathtt{halfGreedyPrefixSupport}(c-1)$. So the $\forall D$ collapses
to a single orbit and form (a) is a statement to be checked along \emph{one} sequence indexed by
$c$, namely
\[
\forall c\ge4\ \text{skipped by the real half-greedy orbit}:\qquad
\ensuremath{S}\ \bigl(\mathtt{halfGreedyPrefixSupport}(c-1)\bigr)\ 1\ (2c-2)\ <\ 2^{\,c-2}.
\]
The crossing hypothesis has not disappeared: these are the ranks
where the real greedy rule skips the next term.  Each resulting
finite inequality is decidable.  Computing its margin can identify
examples to explain a proposed estimate, but a finite table neither
selects a proof method nor establishes a cofinal supply.

The progress-function criterion above remains unproved.  Endpoint
doubling is not a proved doubling law for the largest selected
exponent, and an unbounded function of the incoming endpoint alone
does not rule out a bounded cycle.

\subsection{Prime support and the limits of one formal argument}
\label{sec:invent-R3}

The base-$2$ statement
\[
\mathrm{Irrational}\Bigl(\sum_{p\ {\rm prime}}\frac1{2^p-1}\Bigr)
\]
is known.  Expanding geometrically and collecting divisors gives
\[
 \sum_{p\ {\rm prime}}\frac1{2^p-1}
 =\sum_{n\ge1}\frac{\omega(n)}{2^n},
\]
and Tao--Ter\"av\"ainen prove the latter series irrational in Theorem~1.3
of
\href{https://arxiv.org/abs/2512.01739}{\emph{Quantitative correlations and
some problems on prime factors of consecutive integers}}.  This is a cited
analytic result, not a theorem formalised in the pinned Lean corpus.

The existing pairwise-coprime Lean theorem does not recover it because that
theorem assumes $\sum_{a\in A}1/a<\infty$, whereas $\sum_p1/p$ diverges.  Its
generic block-certificate engine would also demand, in the prime case,
\[
  2^r\mid\omega(N+r)\qquad(1\le r\le K),
\]
a simultaneous exact divisibility pattern far stronger than the input used
by Tao--Ter\"av\"ainen.  Small computational witnesses for $K=2,3$ diagnose
that the condition is not vacuous, but they neither prove the known theorem
nor define a new open case.

Thus the contribution of this route audit is negative and
coordinate-specific: it identifies why the present formal certificate
engine fails to reproduce a theorem already available by quantitative
correlation estimates.  The earlier proposal to prove the prime-support
irrationality as new mathematics is withdrawn.  Extensions to other bases
or prime-power supports are mentioned by the cited authors with proof
modifications omitted; this paper promotes only the fully written base-$2$
theorem.

\subsection{Equivalent carry conditions at perfect-square depths}
\label{sec:invent-R4}

\subsubsection*{The statement needed}

The following five conditions are equivalent to $1/2\in\mathcal A$:
\begin{enumerate}[label=(\arabic*)]
\item $C_G(N)\le2\sqrt N+4$ for every $N\ge0$.
\item The set $\mathbb N_{>0}\smallsetminus G$ is infinite.
\item For arbitrarily large $n$ there is
$D\subseteq\{2,\ldots,n\}$ with $Q(D,n)=2^{n-1}-1$.
\item For arbitrarily large $M$,
$\operatorname{ihc}(G,M)\le B(M+1)$.
\item For arbitrarily large $M\ge1$ there is
$D\subseteq\{2,\ldots,M\}$ with
$|\operatorname{ihc}(D,M-1)|\le B(M)$.
\end{enumerate}
Here $G$ is the actual greedy support for $1/2$.
The forward directions of (2)--(5) are recorded
Lean (\lean{half_mem_mersenneAchievementSet_of_positiveHalfGreedySkips}{BooleanMobiusSkipRowCofinal.lean:97},
\lean{half_mem_mersenneAchievementSet_of_cofinalExactLocalRows}{BooleanMobiusCofinalExactRows.lean:71},
\lean{half_mem_mersenneAchievementSet_of_cofinalStripReturn}{CofinalStripReturn.lean:157},
\lean{half_mem_mersenneAchievementSet_of_cofinalTerminalOnlyStrip}{TerminalOnlyCofinal.lean:134},
all \ev{Lean}). The reverse implications are ordinary deductions in
this paper.  The perfect-square argument below supplies the terminal
bounds without changing the constant in the original statements.

\subsubsection*{The carry bound at perfect-square depths}

\begin{lem}[The terminal bound at square depths]
\label{lem:sqwitness}
Let $A\subseteq\N$ with $1\notin A$ and $X_{A}(2)=1/2$. Then for every
$k\ge1$,
\[
\bigl(\ensuremath{\operatorname{ihc}}\ A\ (k^{2}-1)\ :\ \R\bigr)
\ =\ \mathtt{binaryCoeffTail}\ (c_{A})\ (k^{2})
\ \le\ 2k+4\ =\ \ensuremath{B}\ (k^{2}).
\]
\end{lem}
\leannote{Lean: \lproof{ErdosProblems/Erdos257/PaperCompleteR21/SquareDepthAndHalfMembershipEquivalences.lean}{27}{paper_square_depth_terminal_bound}, \lproof{ErdosProblems/Erdos257/PaperCompleteR20/CarryCollapseCorrespondence.lean}{8}{square_depth_witness}.}
\begin{proof}
By \lean{integerHalfCarry_eq_scaled_residual_add_tail}{HalfCarryReachability.lean:871} (\ev{Lean}),
\[
\bigl(\ensuremath{\operatorname{ihc}}\ A\ N:\R\bigr)
= 2^{\,N+1}\Bigl(\tfrac12-X_{A}(2)\Bigr)
+\mathtt{binaryCoeffTail}\ (c_{A})\ (N+1),
\]
and the first term vanishes by hypothesis. Take $N=k^{2}-1$, so $N+1=k^{2}$. By
\lean{binaryCoeffTail_supportCoeff_le_two_sqrt_add_four}{BooleanMobiusCarry.lean:290} (\ev{Lean},
unconditional) the tail at $k^2$ is at most $2\sqrt{k^{2}}+4=2k+4$, and by
\lean{halfStripBound}{HalfCarryReachability.lean:32}, $\ensuremath{B}\ n
= 2\,\mathtt{Nat.sqrt}\ n+4$, which at $n=k^{2}$ equals $2k+4$ exactly. Non-negativity of the tail
(\lean{binaryCoeffTail_nonneg}{GenericTailOrbitRigidity.lean:78}, \ev{Lean}) gives the lower side.
\end{proof}
\ev{Math}; every input is supplied formal and the composition is four lines.

The former square-root adapter issue is resolved by restricting cofinal witness depths to perfect
squares: at $n=k^{2}$ the two roots coincide
exactly and no widening is needed. Applying Lemma~\ref{lem:sqwitness} along $M=k^{2}-1$ gives (4)
directly; taking $a$ to be the indicator of an achieving support truncated to
$\mathtt{Fin}(M+1)$ at $M=k^{2}$, and using
\lean{integerHalfCarry_inter_Iic_eq_of_succ_le}{HalfCarryReachability.lean:49} together with
$1\notin A$ (forced because $\ensuremath{w}\,1=1>1/2$), gives the witness required by
\lean{HalfTerminalOnlyStripWitness}{TerminalOnlyCofinal.lean:24} and hence (5).

\subsubsection*{What the equivalences establish}

Each biconditional is a theorem about the exact logical content of a
proposed hypothesis.  It is also a reformulation of the membership
question.  These are compatible descriptions, not competing evidence
classes. The equivalence identifies the missing content but does
not show that the reformulation cannot expose a useful invariant. In particular, (3) shows that the existence of finite witnesses at
unbounded depths, even without compatibility between those witnesses,
is equivalent to membership.  The absence of a compatibility requirement
may still be useful in constructing witnesses; equivalence alone says
nothing about which formulation is easier to prove.

\textbf{Scope of the equivalences.} (1) does \emph{not} extend to every
$A$ with $1\notin A$: without
\lean{erdosSupportSeries_greedyMersenneSupport_le}{HalfCarryReachability.lean:888} (which supplies
$\delta:=1/2-S_G\ge0$ for the greedy support specifically) one obtains only the one-sided
equivalence $\text{hbound}\iff(1/2-S_A)\le0$. The mechanism in all five cases is the same and
should be recorded once: $2^{\,N+1}\delta$ dominates any $\sqrt{N}$ envelope, so a strip
hypothesis forces $\delta=0$; and $\delta=0$ makes the unconditional tail bound supply the strip.
Membership is then read off by
\lean{mem_mersenneAchievementSet_iff_greedy_survival}{GreedyAchievementSet.lean:1458} (\ev{Lean}).

\subsection{Sparse and dense supports under the certificate criterion}
\label{sec:invent-R5}

\subsubsection*{The statement needed}

\[
 \begin{gathered}
 A\subseteq\mathbb N_{>0},\quad A\text{ infinite}\quad\Longrightarrow\\
 \text{$c_A$ has arbitrarily late super-logarithmic zero windows}\\
 \text{or the specified block certificates exist at every precision.}
 \end{gathered}
\]
For a fixed nonempty $A$, let $a_0=\min A$.  Every interval of $a_0$
consecutive integers contains a multiple of $a_0$, so $c_A$ has no zero
window longer than $a_0-1$.  The first alternative is therefore
impossible, independently of rationality.  The cited logarithmic result
\lean{supportCoeffZeroWindow_length_le_eps_logb}{SublogDivisorCoverage.lean:435}
does not create a sparse subclass here.  The separate lcm-gap theorem
\lean{irrational_erdosSum_of_lcm_gap}{CertificateKernel.lean:5883}
remains a valid sufficient criterion, but does not prove this proposed
alternative.
The proposal would consequently require the following certificates for
every infinite support:
\[
\exists N,K,L,C,\quad K\le L,\quad
\forall r\in[1,K]:\ 2^{r}\mid c_{A}(N+r),
\]
\[
\sum_{r=K+1}^{L}c_{A}(N+r)\,2^{\,L-r}\le C,
\qquad
q\,(C+N+L+2)<2^{L}.
\]

\subsubsection*{What kind of object would furnish it}

The identity $c_A=\mathbf1_A*\mathbf1$ is a Dirichlet convolution,
and M\"obius inversion gives $\mathbf1_A=\mu*c_A$.  These identities
hold for every support, but do not supply the displayed divisibility
or tail estimates.  Shared factors restrict which residue prescriptions
can be realised; they do not make the Chinese remainder theorem
inapplicable.  For positive moduli $m_i$, a finite system
$N\equiv a_i\pmod{m_i}$ is solvable precisely when
\[
 a_i\equiv a_j\pmod{\gcd(m_i,m_j)}\qquad\text{for every }i,j.
\]
Necessity follows by subtraction.  For sufficiency, keep the strongest
prime-power modulus for each prime; compatibility makes all weaker
prescriptions agree with it, and the ordinary CRT combines the
resulting coprime prime powers~\cite[Theorem~3.1, pp.~2--3]{conradcrt}.
For example, $N\equiv1\pmod4$ and $N\equiv1\pmod6$ give
$N\equiv1\pmod{12}$, whereas $4\mid N$ and $2\nmid N$ are
incompatible.  A block-certificate construction must satisfy these
compatibility conditions and its separate tail budget.  Neither
M\"obius inversion nor CRT alone provides both inputs.

\subsubsection*{The specific obstacle, proved}

Here the honest answer is stronger than an obstacle: the dense branch as stated is \emph{false},
and the falsifier is an explicit, natural, divisor-dense support. This also corrects the reduction
recorded earlier in this document, that universal \#257 passes ``with no loss'' to
$\mathrm{O1}' := \forall A\ \text{infinite},\ \mathrm{Cert}(A)$.

\begin{prop}[Squarefree support: exact engine ceiling, not an open value]
\label{prop:squarefree}
Let $A=\{n\ge2:n\text{ is squarefree}\}$.  Then
\[
  c_{A}(n)=2^{\omega(n)}-1,
\]
which is odd for every $n\ge2$.  Consequently neither the digitwise nor the
carry-aware divisibility-first block-certificate schema has an instance at
any even base.  The two exact no-go theorems are
\lean{ErdosProblems.Erdos257.SquarefreeSupportIncidence.not_exists_carry_certificates_squarefreeSupport}{ErdosProblems/Erdos257/SquarefreeSupportIncidence.lean:274}
and
\lean{ErdosProblems.Erdos257.SquarefreeSupportIncidence.not_exists_digitwise_certificates_squarefreeSupport}{ErdosProblems/Erdos257/SquarefreeSupportIncidence.lean:292}.
\end{prop}
\leannote{Lean: \lproof{ErdosProblems/Erdos257/PaperCompleteR21/SquarefreeSupportEngineCeiling.lean}{20}{paper_squarefree_support_engine_ceiling}.}
\ev{Lean} \scale{uniform} \coord{binary digits}.

This does not leave the squarefree value open.  Duverney--Tachiya's
Corollary~1.2 and Example~1.1 prove that, for every $h\ge1$,
\[
  1,\quad
  \sum_{n\ {\rm squarefree}}\frac1{2^n-1},\quad\ldots,\quad
  \sum_{n\ {\rm squarefree}}\frac1{2^{hn}-1}
\]
are linearly independent over $\Q$.  After deleting the rational $n=1$
term, squarefree support is therefore irrational at every base $2^j$.
Among the integer bases, the cited theorem leaves the non-powers of two
uncovered.  This is cited mathematics, not formalised here.  Proposition~\ref{prop:squarefree}
is a failure of two proof schemas on a value already known to be irrational.

The coordinate dependence is exact.  Adjoin $1$ to the support.  The series
changes by the rational number $(b-1)^{-1}$, so its irrationality is
unchanged, while the incidence becomes $2^{\omega(n)}$:
\lean{ErdosProblems.Erdos257.SquarefreeSupportIncidence.supportCoeff_squarefreeShiftedSupport}{ErdosProblems/Erdos257/SquarefreeSupportIncidence.lean:314}
and
\lean{ErdosProblems.Erdos257.SquarefreeSupportIncidence.irrational_erdosSupportSeries_squarefreeSupport_iff_shifted}{ErdosProblems/Erdos257/SquarefreeSupportIncidence.lean:335}.
At base $2$, CRT then supplies arbitrarily long opening blocks with
$\omega(N+r)\ge r$:
\lean{ErdosProblems.Erdos257.SquarefreeSupportIncidence.exists_omega_ge_block}{ErdosProblems/Erdos257/SquarefreeSupportIncidence.lean:441}
and
\lean{ErdosProblems.Erdos257.SquarefreeSupportIncidence.exists_digitwise_block_squarefreeShiftedSupport}{ErdosProblems/Erdos257/SquarefreeSupportIncidence.lean:485}.
The middle-window and height conditions remain; no certificate or new
irrationality proof follows.  The durable conclusion is methodological:
test a no-go statement under rational finite changes before attributing it
to the underlying value.

\subsection{What the experiments and obstructions do not imply}
\label{sec:invent-shape}

Everything in this subsection is inference from the accumulated failures, not theorem. It is
labelled that way throughout, and where the evidence does not determine the answer it says so
rather than manufacturing a story.

\subsubsection*{Finite observations and their limits}
The recorded obstructions concern specified coordinate systems.  Fixed
precision, a bounded carry state and a prescribed suffix can lose
information needed by the actual-selector condition.  They do not show
that the underlying deterministic sequence has no arithmetic structure.
The relevant exact statements remain
\lean{fixedPrecisionTropicalNoGo}{TropicalCurvatureCarry.lean:137},
\lean{balancedPulse_no_autonomous_decoder}{GenericTailOrbitRigidity.lean:247}
and
\lean{affineBinaryOrbit_mod_twoPow_eq}{GenericTailOrbitRigidity.lean:307}.
Their hypotheses, rather than the fact that an attempted argument failed,
determine their scope.

The inherited experiment log reports $200{,}000$ seam rows with counts
$R:M:U=100197:49899:49898$, a roughly geometric run-length histogram, and a
longest equal run of $15$ in $40{,}000$ binary digits.  It also reports
agreement with OEIS A211706 and a near-integer phase constant of about
$2.5$.  These computations and the external digit comparison were not rerun
in this revision.  They are descriptive reports, not certificates of
normality, independence, a random sampling model, or an infinite
anti-concentration estimate.  A statement about ``sampling noise'' would
require a specified probabilistic model that is not supplied here.

The previously reported exceptional row $r=7$ and crossing $(s,d)=(10,7)$
are likewise statements about the tested range only.  The exact range and
reproduction artefacts must accompany them before they are used as
numerical evidence outside this historical record.  No claim that this is
the only exceptional row of the infinite orbit is made.

These coordinates all descend from the same greedy-carry representation;
they are not independent tests of a common random hypothesis.  Strict
strict tail domination gives an injective coding, but says neither that its
digits are random nor that averaging over indices is impossible.  The
measure-one and nowhere-dense properties of the achievement set do not
decide membership of a specified rational.  The correct research question
is whether a precise arithmetic estimate can be proved for the actual
selector, not whether finite data resemble a random model.

\subsubsection*{Why the support and rational-target problems remain distinct}

Proposition~\ref{prop:squarefree} is a reason to separate the two halves
rather than assimilate them.  On the universal side the present
certificate engine first meets a coordinate-dependent parity obstruction.
The squarefree value is nevertheless already known to be irrational at
power-of-two bases, and adjoining $1$ removes the opening-block obstruction
without changing irrationality.  Thus the no-go diagnoses this engine; it
does not diagnose the value or the universal problem.

Short-window near-integer anti-concentration is one possible missing
input for the particular mechanisms analysed here.  No equivalence shows
that it is necessary for either problem, and no argument identifies it as
a common necessary obstruction to both.  The squarefree example shows why
finite rational changes of support must be checked before an opening-block
failure is interpreted as a property of the value.

\subsubsection*{What the evidence does and does not determine}

It does not determine whether a short proof exists. Nothing in a catalogue of failed methods
bounds the length of a successful one, and this document should not pretend otherwise.

It records limitations of specified methods.  The fixed-precision
completion theorem and the balanced-pulse countermodel concern the
particular summaries defined above, not all fixed-precision arithmetic
or all finite-state methods with additional arithmetic input. An
irrationality-measure bound on $E$ does not classify a specified rational in $\mathcal{A}$.
Instantiating the present block-certificate engine at a general support is blocked, on the
universal side, by Proposition~\ref{prop:squarefree}. Survival through any tested finite depth does
not by itself establish membership; a fatal-gap certificate would establish non-membership. These
facts do not decide rational-target membership, and they do not yield an impossibility theorem for
all probabilistic, measure-theoretic or averaging methods. Each recorded obstruction is reusable
in a way that a failed attempt is not. That is the contribution: not a proof, not a reduction, and
not a map of every region in which a proof cannot live.
% ; - part a257_p5 ; -
\clearpage
\part*{V. Outstanding obligations and provenance}
\addcontentsline{toc}{part}{V. Outstanding obligations and provenance}
\section{What is open, stated exactly}

Erd\H{o}s \#257 asks whether $\sum_{n\in A}1/(2^n-1)$ is irrational for \emph{every} infinite
$A\subseteq\mathbb{N}_{\ge 1}$. Nothing in this paper decides this, in either the universal form
or the single-target half-value form below.  Parts I--IV do prove
unconditional theorems for specified support classes, as well as uniform
finite-denominator and achievement-set results.  The unresolved
membership conditions must be distinguished from those results: some
are equivalences, some are unproved sufficient conditions, and some
have only finite computational checks. This section states the exact open targets, in display mathematics, with their Lean sites
where formalised and their precise logical relationships to one another. \textbf{Read the
implication table in \S\ref{sec:v-table} before attacking any single target below}: several
apparent waypoints are \emph{equivalent} to the full problem, not easier approaches to it.

\subsection{Notation fixed for this section}

The functional $\mathtt{erdosSupportSeries}\ 2\ (-)$ denotes the base-$2$ Erd\H{o}s support series
$X_{A}(2) := \sum_{n\in A} 1/(2^n-1)$
(\lean{erdosSupportSeries}{CertificateKernel.lean:8897}). $c_{A}(n) :=
|\{d\mid n : d\in A\}|$ (\lean{supportCoeff}{CertificateKernel.lean:8856}) is the divisor-indicator
coefficient whose generating tail feeds every certificate below. $\mathcal{A} :=
\ensuremath{\mathcal A} = \{x\mid \exists A,\ 0\notin A\wedge x =
\mathtt{positiveMersenneSupportValue}\ A\}$ (\lean{mersenneAchievementSet}{GreedyAchievementSet.lean:571}).
$G := \mathtt{greedyMersenneSupport}(1/2)$ is the canonical greedy support for target $1/2$;
$\mathrm{rem}(s)$ is its integer seam
remainder at row $s$, and a \emph{reset} at row $r$ is a row where the greedy orbit returns to the
near-$2^r$ regime after a run of takes or skips.

\subsection{The universal statement}

\begin{defn}[Universal \#257]
\[
\mathrm{U257} \;:=\; \forall\, A\subseteq\mathbb{N}_{\ge 1},\ A.\mathrm{Infinite} \;\Longrightarrow\;
\mathrm{Irrational}\!\left(X_{A}(2)\right).
\]
\end{defn}
This is the literal statement of Erd\H{o}s \#257. \ev{Open}. The corpus's own sufficient-condition
engine is
\[
\mathrm{Cert}(A)\ :=\ \forall q>0,\ \exists N,K,L,C,\ K\le L,\
\left(\forall r\in[1,K],\ 2^r\mid c_{A}(N+r)\right)\wedge{}
\]
\[
\sum_{r=K+1}^{L}c_{A}(N+r)\cdot 2^{L-r}\le C\ \wedge\
\exists t,\ 0<c_{A}(N+L+1+t)\ \wedge\ q(C+N+L+2)<2^L,
\]
with $\mathrm{Cert}(A)\Rightarrow \mathrm{Irrational}(X_{A}(2))$
already proved unconditionally for \emph{every} $A$
(\lean{irrational_erdosSupportSeries_of_weighted_coeff_certificates}{CertificateKernel.lean:9031},
\ev{Lean}, \scale{n/a}; the theorem is universally quantified over $A$; only its hypothesis is
unsupplied).  Therefore the stronger statement
\[
\mathrm{O1'}\ :=\ \forall A,\ A.\mathrm{Infinite}\Longrightarrow \mathrm{Cert}(A).
\]
would imply $\mathrm{U257}$, but the reduction is \emph{not} lossless and
$\mathrm{O1'}$ is actually false.  For
$A=\{n\ge2:n\text{ squarefree}\}$ the incidence is
$2^{\omega(n)}-1$, hence odd for every $n\ge2$, and the digitwise
certificate schema has no instance at any even base
(\lean{ErdosProblems.Erdos257.SquarefreeSupportIncidence.not_exists_digitwise_certificates_squarefreeSupport}{ErdosProblems/Erdos257/SquarefreeSupportIncidence.lean:292},
\ev{Lean}).  This is a method ceiling, not a counterexample to irrationality:
Duverney--Tachiya prove the squarefree value irrational at every
power-of-two base.  Likewise the pairwise-coprime required input does not cover the
primes because $\sum_p1/p$ diverges, even though Tao--Ter\"av\"ainen prove
the base-$2$ prime-support value irrational.  The honest universal target is
$\mathrm{U257}$ itself; $\mathrm{O1'}$ is a rejected sufficient
strengthening, not an open equivalent normal form.

\subsection{Whether one half is represented}

\begin{defn}[Half membership]
\[
\mathrm{HALF}\ :=\ \left(\tfrac12:\mathbb{R}\right)\in\mathcal{A}.
\]
\end{defn}
$\mathrm{HALF}$ is a \emph{single instance} of the negation of $\mathrm{U257}$: any witness
$A$ with $0\notin A$, $X_{A}(2)=1/2$ is infinite
(\lean{finite_boolSupport_ne_half}{HalfCarryReachability.lean:589}, \ev{Lean}, \scale{uniform} ; 
finite supports have odd-denominator value, so can never equal a dyadic rational), hence a
genuine counterexample to $\mathrm{U257}$. Thus $\mathrm{HALF}\Rightarrow\neg\mathrm{U257}$.  The reverse
implication is not established: failure of half-membership would not decide $\mathrm{U257}$ if false, since $\mathrm{U257}$ could still fail at a
different target $t\in\mathbb{Q}$. $\mathrm{HALF}$ is \ev{Open}. It is logically equivalent, by an
unconditional Lean iff, to a purely combinatorial statement about one fixed orbit:
\[
\mathrm{HALF}\ \Longleftrightarrow\ \left(\mathtt{greedyMersenneSkippedSupport}(1/2)\right).\mathrm{Infinite}
\]
(\lean{half_mem_mersenneAchievementSet_iff_greedySkippedSupport_infinite}{GreedyAchievementSet.lean:2583},
\ev{Lean}, \scale{uniform}, coordinate \coord{divisor counts and finite sums}, target-generic in the reverse
direction; the forward direction uses only that the target is rational, so the same theorem
holds verbatim for every rational $t\ge 0$ in place of $1/2$).

\subsubsection{Conditions equivalent to half-membership}

The following three conditions are equivalent to half-membership.
They express the same unresolved assertion in different terms.  The
proofs of equivalence identify the content still to be established;
they do not compare the difficulty of possible proofs.

\begin{prop}[Infinitely many positive skips are equivalent to half-membership]
\label{prop:cpgs-equiv}
Define
\[
\mathrm{CPGS}\ :=\ \forall N,\ \exists c\ge N,\ c\ \text{skipped by the rational half-greedy
orbit}\ \wedge\ 0<\ensuremath{r}(1/2)(c-1)
\]
(\lean{CofinalPositiveHalfGreedySkips}{BooleanMobiusSkipRowCofinal.lean:22}). The positivity
conjunct is unconditionally true for every skipped $c$
(\lean{localMersennePrefixValue_halfGreedy_lt_half}{BooleanMobiusCriticalCapacityCofinal.lean:138},
\ev{Lean}, \scale{uniform}; an odd-denominator parity fact), hence deletable without changing
the Prop's meaning. What survives, $\forall N\,\exists c\ge N$ skipped, is exactly
$(\mathtt{greedyMersenneSkippedSupport}(1/2)).\mathrm{Infinite}$, so
\[
\mathrm{CPGS}\ \Longleftrightarrow\ \mathrm{HALF}.
\]
\end{prop}
\leannote{Lean: \lproof{ErdosProblems/Erdos257/PaperCompleteR21/SquareDepthAndHalfMembershipEquivalences.lean}{44}{paper_cpgs_equiv}.}
\ev{Math}. \scale{cofinal}.
\coord{divisor counts and finite sums, greedy orbit}. The cited implication
\lean{cofinalExactLocalMersenneHalfRows_of_positiveHalfGreedySkips}{BooleanMobiusSkipRowCofinal.lean:84}
therefore gives an equivalent formulation, not a proof of membership.

\begin{prop}[Terminal carry bounds are equivalent to half-membership]
\label{prop:strip-equiv}
Consider finite sets $D\subseteq\{2,\ldots,M\}$ at arbitrarily large
depths $M$ with
\[
 |\operatorname{ihc}(D,M-1)|\le2\lfloor\sqrt M\rfloor+4.
\]
This is the exact terminal-strip condition in
\lean{HalfCarryCofinalTerminalOnlyStrip}{TerminalOnlyCofinal.lean:34}.
It is equivalent to $1/2\in\mathcal A$: the forward direction follows
from the displayed finite-approximation estimate; for the converse,
truncate an achieving support at $M=k^2$ and apply
Lemma~\ref{lem:sqrt-witness}.  Thus constant $4$ already suffices at
cofinally many depths.  The relaxed constant $6$ gives a bound at every
depth, using $\sqrt M\le\lfloor\sqrt M\rfloor+1$, but that relaxation is
not needed for the cofinal statement.
\end{prop}
\leannote{Lean: \lproof{ErdosProblems/Erdos257/PaperCompleteR21/SquareDepthAndHalfMembershipEquivalences.lean}{130}{paper_terminal_strip_equiv}, \lproof{ErdosProblems/Erdos257/PaperCompleteR21/SquareDepthAndHalfMembershipEquivalences.lean}{112}{paper_relaxed_constant_six_every_depth}.}
\ev{Math}. \scale{cofinal}. \coord{support
rigidity, half-carry recurrence}. This corrects the claim that the existence condition is logically weaker.
An equivalent formulation may still be useful for finding a proof.

\begin{obs}
Propositions~\ref{prop:cpgs-equiv} and~\ref{prop:strip-equiv} do not decide
half-membership.  They show that the proposed existence conditions are
equivalent to it, not that those formulations have no possible use.  The
upper bound on the centred carry in
\lean{greedy_half_infinite_of_mobiusCenteredHalfCarry_sqrtBound}{HalfCarryReachability.lean:834}
has the same equivalence for the greedy support $A=G$.  For arbitrary
$A$ with $1\notin A$, an upper bound alone yields only
$x_A\ge1/2$, as the identity
$\operatorname{ihc}(A,N)=2^{N+1}(1/2-x_A)+\mathrm T(N+1)$ shows.
The greedy inequality $x_G\le1/2$ is what turns that conclusion into
equality.
\end{obs}

\subsection{Exact finite sums at unbounded depths}

\begin{defn}
For every integer $N$, require some $n\ge\max\{N,1\}$ and
$D\subseteq\{2,\ldots,n\}$ with
\[
 \sum_{d\in D}\left\lfloor\frac{2^n}{2^d-1}\right\rfloor
       =2^{n-1}-1.
\]
(\lean{CofinalExactLocalMersenneHalfRows}{BooleanMobiusCofinalExactRows.lean:38}, \ev{Open},
\scale{cofinal}). Supports at different depths need not agree.
This is the condition in Definition~\ref{record:257bm-d4};
no strict weakening of half-membership is asserted.
\end{defn}
$\mathrm{CofinalExactLocalMersenneHalfRows}\Rightarrow\mathrm{HALF}$ is proved unconditionally
via compactness of $\mathcal{A}$ (continuous image of $(\mathbb{N}\to\mathrm{Fin}\,2)$ under the
digit-coding map, hence closed)
(\lean{half_mem_mersenneAchievementSet_of_cofinalExactLocalRows}{BooleanMobiusCofinalExactRows.lean:71},
\ev{Lean}). The cited theorem already proves that a representing support exists,
because that is what membership of $\mathcal A$ means.  It need not
choose a support by a specified algorithm from the row witnesses.
An ordinary diagonal subsequence argument gives a compatible limiting
support: successively fix each coordinate along an infinite
subsequence, then use the uniform tail bound to pass to its value.
This uses continuity of the digit map, also recorded in
\lean{continuous_positiveMersenneDigitValue}{GreedyAchievementSet.lean:628}.  The diagonal argument preserves $2\in A$, since all sufficiently deep exact rows
contain exponent two
(\lean{two_mem_of_exact_localMersenneQuotient}{BooleanMobiusExactRowRankTwo.lean:23}).
This describes an ordinary existence proof (\ev{Math}); no additional
formal extraction theorem or effective algorithm from purely
existential row data is claimed.

Proposition~\ref{prop:collapsed-list}(c) already identifies
$\mathrm{CofinalExactLocalMersenneHalfRows}$ as logically equivalent to $1/2\in\mathcal{A}$.
Equivalence of a membership condition, constructive extraction of a coherent support, and a
rejected sufficient strengthening are not interchangeable descriptions. Absence of a Lean
extraction theorem does not make O3 a strictly weaker substitute for $\mathrm{HALF}$. The remaining
work on O3 is to supply the cofinal exact rows, or any equivalent membership condition.

\subsubsection{A capacity condition for exact sums at unbounded depths}
The finite construction uses a $(c-2)$-bit capacity test. Write
\[
\ensuremath{S}\ D\ k\ M\ :=\ 2^{M-k}-\ensuremath{Q}\ D\ M - 1
\]
(\lean{localBinarySuffix}{BooleanMobiusLocalRepair.lean:102}), with
natural-number subtraction understood as truncated at zero.  The loose
bound below requires a critical crossing: $c\ge4$,
$D\subseteq\{2,\ldots,c-1\}$, and
$0<1/2-V(D)<w_c$.  A merely below-half core is not enough:
$D=\varnothing$ gives $S(D,1,2c-2)=2^{2c-3}-1$.
For a critical-crossing core the two bounds are
\begin{align}
&\text{(loose, proved)}\quad \ensuremath{S}\ D\ 1\ (2c-2)\ <\ 2^{c-1}
\tag{$\ast$}\\
&\text{(sharp, needed)}\quad \ensuremath{S}\ D\ 1\ (2c-2)\ <\ 2^{c-2}
\tag{$\ast\ast$}
\end{align}
The implication $(\ast)$ from the critical-crossing hypotheses is
proved uniformly in $c$
(\lean{localBinarySuffix_two_mul_sub_two_lt_upperHalfCapacity}{BooleanMobiusSkippedCoreExactRow.lean:123},
\ev{Lean}, \scale{uniform}).

$(\ast\ast)$ is exactly one bit sharper and, via
\lean{localBinarySuffix_two_mul_sub_two_lt_criticalCapacity_iff}{BooleanMobiusSkippedCoreCriticalCapacity.lean:22}
(under the stated core hypotheses), is equivalent to
$Q(D\cup\{c\},2c-2)\ge2^{2c-3}$.
Here $c$ is the added exponent; the quotient is evaluated at depth
$2c-2$, not at depth $c$.
The proof of $(\ast)$ gives the more informative additive bound
$S<2^{c-2}+|D|$ before using $|D|\le c-2$.
One possible cofinal supply is therefore the following exclusion
of a specified integer interval:
\[
\begin{gathered}
\text{(O3-supply)}\qquad \forall N\ \exists c\ge\max(N,4)\ \exists D\subseteq\{2,\ldots,c-1\},\\
0<1/2-V(D)<w_c,\qquad
S(D,1,2c-2)\notin[\,2^{c-2},\,2^{c-2}+c-3\,].
\end{gathered}
\]
The excluded band contains $c-2$ integers.  The sharp capacity
bound itself allows $2^{c-2}$ nonnegative integer values.
The numerical widths differ from those in Section~\ref{sec:o4}, but the
conditions concern different quantities.  Width alone does not prove a
logical implication between them.

The sharp-fill implication
\lean{exists_exactRowStrictUpperFill_of_skippedCoreSharpCapacity}{BooleanMobiusSkipRow.lean:238}
(\ev{Lean}) assumes $(\ast\ast)$ for a below-half core
$D\subseteq\{2,\ldots,c-1\}$, but does not require a critical
crossing or identify that core with a real greedy prefix.  This gives
another formulation in which to seek the required finite rows.
It does not establish a strict logical weakening of half-membership,
nor does it supply cores satisfying the bound at unbounded ranks.
Moreover, the integer gap separation fixes the resulting exact row
whenever it exists; the formulation does not create multiple exact
rows at a fixed depth.

The corpus's complete inductive skeleton from the endpoint-six seed to O3 is already assembled
(\lean{cofinalExactLocalMersenneHalfRows_of_criticalQuotientSupply}{BooleanMobiusCriticalCapacityCofinal.lean:1309},
\ev{Lean}); it consumes the sharp capacity input, in its reduced canonical form
\[
 c\ge4,\ c\notin G\quad\Longrightarrow\quad
 Q\bigl((G\cap\{2,\ldots,c-1\})\cup\{c\},2c-2\bigr)
 \ge2^{2c-3}.
\]
(\lean{HalfGreedySkippedCriticalQuotientSupply}{BooleanMobiusCriticalCapacityCofinal.lean:178}),
at essentially every step (\ev{Open}, \scale{cofinal}). This condition is required along the specified greedy prefixes.  That
restriction on the witnesses does not by itself prove a strict logical
ordering with the cofinal condition above.

\subsection{Reset bounds and runs of right transitions}
\label{sec:o4}

\begin{defn}[Reset $\sqrt{\text{-escape}}$ from row 10]
\[
\mathrm{SQRTESC}\ :=\ \forall r\ge10\ \text{that are resets},\quad
\big|\mathrm{rem}(r+1)-2^{r+1}\big|>2^{(r+5)/2}
\qquad\Big(\text{equivalently } \Delta_{r+1}^2 > 2^{r+5}\Big).
\]
\end{defn}
The conditional implication uses the integer bound
$B(r)=2^{\lfloor(r+4)/2\rfloor}+2r+3$ described in Theorem~A:
$B(r)^2\le2^{r+5}$ for $r\ge10$, while a right branch at the
critical crossing would force $|\Delta_{r+1}|\le B(r)$.
Thus the actual implication chain is
\[
 \mathrm{SQRTESC}\ \Longrightarrow\
 \mathrm{LargestSkipLateStepSocket}\ \Longrightarrow\ \mathrm{HALF}.
\]
The last implication is
\lean{half_mem_mersenneAchievementSet_of_largestSkipLateStepSocket}{HalfCylinderLargestSkipInduction.lean:167};
the supplied release proof of the first is identified in Theorem~A.
These are sufficient implications, not an equivalence with
half-membership.

For comparison one may ask for
\[
 \mathrm{RUNBOUND}:\qquad
 \forall\,\text{reset }r\ge31,\quad L_r<(r-3)/2,
\]
where $L_r$ counts the consecutive right branches after reset $r$.
Along $L$ such steps the exact relation is
\[
 w_{r+1+L}=4^Lw_{r+1}
   -\sum_{j=0}^{L-1}4^{L-1-j}(4+p_{r+1+j}).
\]
The pulse sum and the threshold ranges cannot be dropped to claim
$\mathrm{RUNBOUND}\Leftrightarrow\mathrm{SQRTESC}$.
The pointwise run/deviation equivalence stated in earlier source
notes is not established by the bound used in the conditional proof.
The observation at $r=7$ (deviation $9$, square-root threshold $64$)
is outside both proposed late ranges; that fact alone does not
justify replacing the initial threshold $10$ by $31$.

\emph{Certified evidence, not a proof.} Exact integer arithmetic to row $200{,}000$: TARGET holds
at every reset row except $r=7$; the maximum observed R-run is $19$, at row $158{,}096$, against a
required threshold there of $\approx79{,}046$; a margin of $\sim4000\times$
(\ev{Cert}, \scale{fixed}). This is a finite verified list, \emph{not} a cofinal supply, and must
never be read as one: it is the exact instance of the difference between a finite and an unbounded parameter range this programme's own
doctrine warns against.

\begin{obs}[The sign of the reset deviation]
At all $1209$ observed resets, the sign of the deviation is perfectly predicted by branch type
(M$\Rightarrow$dev$>0$, U$\Rightarrow$dev$<0$, zero exceptions), and the margin over the threshold
grows from $1.11$ bits at $r=14$ to $\sim1246$ bits by row $2500$ (\ev{Cert}). Proving the sign
law reduces to the middle inequality
$4\mathrm{rem}(r)>p_r$.  For an upper reset the sign already follows
from $w_{r+1}=-(4o_r+p_r^+)<0$, with positive overshoot and
nonnegative pulse.  The middle inequality remains \ev{Open}.
Even after both signs are known, a separate lower bound for their
magnitudes is required for $\mathrm{SQRTESC}$.
\end{obs}

\begin{obs}[Dangerous-reset rigidity; argument B]
\label{obs:dangerous-reset-rigidity}
If reset $r$ is \emph{dangerous} ($|\mathrm{rem}(r+1)-2^{r+1}|\le2^{(r+5)/2}$) and the preceding
reset $r_0$ has a pure R-run of length $L'=r-r_0-1>(r+5)/4$, then $w_{r_0+1}$ is confined to at
most two adjacent integers, determined by the divisor-pulse stream alone. A chain of $n$
consecutive dangerous resets with long runs forces $n-1$ nested exact congruence towers. This
shows the failure mode is a tower of exact integer coincidences, not drift; it is \ev{Math}, and
it does \emph{not} by itself exclude the two pinned values; an isolated dangerous reset after a
short run is untouched by this theorem.
\end{obs}

\subsection{The remaining
\texorpdfstring{\(-2,-1\)}{negative-two, negative-one} middle-cell exclusion}

\begin{defn}
$\mathtt{SeamTwoSidedDyadicCellEscape}$ excludes three integer values of the middle coordinate,
$4\cdot\mathrm{rem}(s)-\mathrm{belowPulse}(s)-4\notin\{-3,-2,-1\}$, plus one right-pulse-leak
bound $4\cdot\mathrm{overshoot}(s)+\mathrm{abovePulse}(s)\le2^{s+2}$ under
$\mathrm{overshoot}(s)\le2^s$ (\lean{SeamTwoSidedDyadicCellEscape}{HalfCylinderMiddleCarryLowerBound.lean:4374}).
\end{defn}
Granted this, induction from a \texttt{decide}-checked row-5 base propagates
$\forall s\ge5,\ \mathrm{rem}(s)\le2^s\vee\mathrm{overshoot}(s)\le2^s$
(\lean{SeamTwoSidedDyadicCellEscape.twoSided}{HalfCylinderMiddleCarryLowerBound.lean:4448}, \ev{Lean}
given the hypothesis, \scale{uniform in the consequence}). One of the three cells, $-3$, is already
closed unconditionally in the all-right-tail configuration
(\lean{finalMiddleCell_neg_three_not_last}{HalfCylinderFinalMiddleCellEscape.lean:587}, \ev{Lean}).
The two remaining cells, $-2$ and $-1$, are \ev{Open}: what would close O5 is their exclusion for the
actual quotient-greedy orbit (not a hypothetical one), together with the right-pulse-leak inequality. This excludes specified integer values in one coordinate.  The
square-root reset condition instead bounds the magnitude of a
different quantity by a rank-dependent threshold.  Their numerical
widths do not order the two hypotheses, and no implication between
them is established here.

A distinct sufficient condition concerns upper resets only.  For
every $d\ge13$ at which the upper carry occurs, set
$E_d=4\mathrm{overshoot}(d)+\mathrm{abovePulse}(d)$ and require,
for every $0\le j\le d$,
\[
 2^{d-j+1}<E_d\quad\text{or}\quad
 E_d+2(d+j)\le2^{d-j+1}.
\]
This is the condition named
$\mathrm{SeamUpperResetDyadicBandEscape}$
(\lean{SeamUpperResetDyadicBandEscape}{HalfCylinderMiddleCarryLowerBound.lean:4566}, \ev{Open},
proved \ev{Cert} for rows $13\le d\le30$ by the cited finite computation
(\lean{seamUpperResetDyadicBandEscape_through_thirty}{HalfCylinderUpperResetBandCertificates.lean:78})).
This band has linear width $O(d)$, whereas the other bound uses a width
of order $2^{d/2}$.  The quantities being bounded are different, so this
comparison of widths does not prove a logical implication or strict
weakening.  The linear-width condition suffices for half-membership by
$\lean{half_mem_mersenneAchievementSet_of_upperResetDyadicBandEscape}{HalfCylinderMiddleCarryLowerBound.lean:4790}$.
It is proved logically equivalent to a single-critical-index check per row
(\lean{seamUpperResetCriticalBandEscape_iff}{HalfUpperResetCriticalBand.lean:883}, \ev{Lean}), which is a genuine
quantifier-complexity reduction, not a strength reduction.

\subsection{Whether one twenty-first is represented}

For the second rational target, the recorded equivalence is
\[
 \frac1{21}\in\mathcal A
 \quad\Longleftrightarrow\quad
 \neg\,\mathtt{TwentyOneFatalAlignedBranch}.
\]
\lean{one\_div\_twenty\_one\_mem\_iff\_not\_fatalAlignedBranch}{TwentyOneQuotientGreedy.lean:3507}
\quad\ev{Lean}.  The fatal branch means a finite fatal greedy gap followed
by cofinite exact selection, eventual quotient/rational alignment, and
eventual visits to every doubling block.  If it occurs, the canonical
quotient remainder is eventually strictly larger than $2^R$ and follows one
explicit affine recurrence:
\lean{twentyOneFatalAlignedBranch\_eventually\_strict\_supercapacity}{TwentyOneQuotientGreedy.lean:5625}
and
\lean{twentyOneFatalAlignedBranch\_eventually\_affine\_supercapacity}{TwentyOneQuotientGreedy.lean:5658}.
Any unbounded sequence of closed returns instead proves membership.

No finite support on ranks at least $2$ sums to $1/21$
(\lean{finiteErdosSum\_ne\_one\_div\_twenty\_one}{ErdosProblems/Erdos257/HalfCounterexampleFrontier.lean:61}).
Thus membership would already be an infinite-support rational counterexample
to $\mathrm{U257}$.  The missing step is to exclude the fatal aligned
affine-supercapacity branch (or force unbounded closed returns); neither is
proved. \ev{Open}

\subsection{Implication table}
\label{sec:v-table}

Read left to right: $\Rightarrow$ records a proved implication
with the hypotheses stated in the relevant section, and
$\Leftrightarrow$ records an equivalence.  Source proofs and ordinary
arguments retain the evidence distinctions explained above; the table
is not a fresh kernel replay.  ``No implication stated'' means that
no comparison is established here.  Quantitative bounds are not ranked
solely by their widths.

{\small
\begin{longtable}{@{}p{0.25\linewidth}p{0.43\linewidth}p{0.23\linewidth}@{}}
\toprule
\papertablehead
Target & Relation to its endpoint & Status \\
\midrule
Universal irrationality & Would exclude both rational targets; the stronger universal certificate condition is false & \ev{Open} \\
Half-membership & Would give a counterexample to universal irrationality & \ev{Open} \\
$1/21$ membership & Equivalent to excluding the specified failure branch; would give a counterexample & \ev{Open} \\
Infinitely many greedy skips (Proposition~\ref{prop:cpgs-equiv}) & $\Leftrightarrow$ O2 & \ev{Open}, equivalent condition \\
Terminal-only strip (Prop.~\ref{prop:strip-equiv}) & $\Leftrightarrow$ O2 & \ev{Open}, equivalent condition \\
The centred square-root bound for $G$ & $\Leftrightarrow$ O2 & \ev{Open}, equivalent condition \\
Exact finite sums at unbounded depths & $\Leftrightarrow$ O2 (Prop.~\ref{prop:collapsed-list}(c)); Lean lands $\Rightarrow$ & \ev{Open}, equivalent condition \\
(O3-supply) linear-width capacity band & $\Rightarrow$ O3 (Lean, given lemma) & \ev{Open} \\
The quotient condition on skipped greedy prefixes & $\Rightarrow$ O3 (Lean) & \ev{Open}; no strict comparison asserted \\
Reset square-root bound & $\Rightarrow$ O2 (release-source proof); no run-length equivalence asserted & \ev{Open} \\
The reset sign condition & no implication to half-membership stated & \ev{Open} \\
Rigidity after long runs & partial rigidity, no exclusion & \ev{Open} \\
The remaining middle-cell exclusions and right-branch bound & $\Rightarrow$ two-sided bound (Lean, given hyp.) & \ev{Open} \\
The upper-reset band condition & $\Rightarrow$ O2 (Lean), linear-width & \ev{Open}, cited check for $13\le r\le30$ \\
\bottomrule
\end{longtable}}

\subsection{Possible inputs and their scope}

The following comparisons identify missing hypotheses in the recorded
arguments.  They neither predict the length of a proof nor classify all
possible methods.

\textbf{O1 (universal).} The question is $\mathrm{U257}$ itself. The stronger schema
$\mathrm{O1'}=\forall A\,\mathrm{Cert}(A)$ is a rejected sufficient strengthening, not future work:
Proposition~\ref{prop:squarefree} already exhibits a squarefree support on which neither the
digitwise nor the carry-aware block-certificate required input exists, while Duverney--Tachiya still
prove that squarefree value irrational. Supplying $\mathrm{Cert}(A)$ for arbitrary infinite
supports is therefore not an open equivalent normal form of $\mathrm{U257}$. Remaining work on O1
is the universal irrationality statement, or a different method that covers supports the
certificate schemas cannot.

\textbf{O2/O3/O5 (Mersenne-specific).} The sharp-capacity gap $(\ast)\to(\ast\ast)$ and the
$-2,-1$ middle-cell exclusion are both, at bottom, requests for anti-concentration of an explicit
divisor-count quantity ($\ensuremath{S}$, or $4\cdot\mathrm{rem}-\mathrm{belowPulse}-4$)
away from specified short integer windows.  A congruence argument,
a quantitative distribution estimate, or another method would have to
supply the stated inequalities.  The reformulations do not decide
which kind of argument can do so.

\textbf{O4 (reset separation).} The specified quotient-row argument
needs $|\mathrm{rem}(r+1)-2^{r+1}|>2^{(r+5)/2}$ at its reset rows.
A weighted divisor window is one proposed source of information.
After division by $2^{r+1}$, the displayed condition becomes
$|\lambda_{r+1}|>2^{(3-r)/2}$.  A bound for
$\operatorname{dist}(\Theta_L(M),\mathbb Z)$ concerns a different
quantity; Section~\ref{sec:invent-R1} displays the additional
skip-set term that a transfer argument would need to control.
No such implication or uniform estimate is supplied.
Part~IV compares this with a totient-window hypothesis for a particular
argument about Problem~249.  The analogy does not identify the two
statements, transfer a theorem between them, or show that either
hypothesis is necessary for its original problem.

\subsection{Summary}

The results here do not decide the universal statement
(O1${}={}$U257) or the two rational membership targets $1/2$
(O2) and $1/21$ (O2b).  The conditions surveyed here have different
logical roles.
Those equivalent to O2, including O3 by
Proposition~\ref{prop:collapsed-list}(c), do not establish membership
merely by reformulating it.  The sufficient conditions O4, O5, and
O3-supply have not been supplied at the required cofinal scale, and no
converse is proved for them.  The proposed universal certificate
strengthening $\mathrm{O1'}$ is a separate, rejected statement: it is
already false on squarefree support.  Finite computational checks,
however large their margins, are distinct from all these quantified
assertions. The sufficient support theorems in the first part are proved for their
stated classes.  The identities, finite checks and equivalences in this
section organise the remaining questions but do not resolve either
rational target or the universal problem.

\clearpage
\section{Logarithmic cost under arithmetic sampling}\label{sec:257-logarithmic-sampling}
Positive fractional covers retain their geometric-kernel estimate, but their
optimised cost cannot be bounded uniformly by a positive logarithmic divisor
cost.  The obstruction occurs on complete arithmetic periods, rather than
only in incomplete-period errors.  We give the construction and the surviving
initial-interval estimate here; the short note records their consequences
beside its cube-family calculation.

The proofs in this section are ordinary arguments from the 17 September
research note~\cite[Theorem~1, Corollary~3 and Proposition~4]{endpoint2026}.
They have not had independent review or a fresh Lean verification.
For the underlying Lambert subseries, Kova\v{c}--Tao treat the same fixed-base
object in~\cite[Section~2.1.2 and Remark~4.1]{kovactao}; their strict-tail
argument gives distinct support coding and Cantor topology.
Van Doorn--Kova\v{c} treat a freely chosen denominator sequence.
Their convention is $n_{i+1}/n_i\ge\lambda$ at every index.
For each $1<\lambda<2$, they construct such a sequence whose finite
reciprocal sums include every rational in $[0,2]$; no sequence with
$n_{i+1}/n_i\ge2$ can do this on any nonempty open
interval~\cite[Theorem~1, Lemma~4 and Corollary~5]{vandoornkovac}.
This is a statement about the rational points of an interval,
not an interval consisting entirely of finite sums; there are
only countably many finite subsets.
Those results provide context, not an arithmetic logarithmic estimate or a
solution of fixed-base rational membership.  The proof dependencies below
are positive divisor majorisation, elementary prime selection, the Chinese
remainder theorem, and the original kernel estimate
\eqref{eq:257-mixed-finite-kernel}.  The short note's logarithmic lower bound
and cube-family calculation concern the same fractional-cover cost; their
structured upper bound is not a uniform logarithmic converse.

\subsection{The two finite costs}
Let $F$ be a finite nonempty subset of $\mathbb N_{>0}$, and write
\[
 Q=\operatorname{lcm}(F),\qquad f_F(n)=\#\{a\in F:a\mid n\},\qquad
 U_F(N)=\sum_{r\ge1}2^{-r}f_F(N+r)
       =\sum_{a\in F}\frac{2^{N\bmod a}}{2^a-1}.
\]
For $0<t\le1$ define, with natural logarithms,
\begin{equation}
 \kappa_1(F;t)=\min\left\{\sum_{d\mid Q}\frac{c_d}{d}:
 c_d\ge0,\quad
 \log(1+f_F(s)/t)\le\sum_{d\mid s}c_d\quad(s\mid Q)\right\}.
 \label{eq:257-logarithmic-cost}
\end{equation}
The constraints also hold at every positive integer $n$, by replacing $n$
with $\gcd(n,Q)$.  The programme is finite and feasible; a bounded cost
sublevel bounds each coefficient, so its minimum is attained.
For $L\ge1$, let $\mathbb P_L$ be uniform sampling of $N=Lm$ over a
period of $U_F(Lm)$, for which $Q/\gcd(Q,L)$ is a permissible length.
This notation does not assert that the displayed period is least.

We use $K_*(F)$ for the infimum of the existing fractional-cover cost
\begin{equation}
 K=\sum_j\frac{C_j\eta_j^{-\alpha_j}}{2^{\alpha_j}-1},\qquad
 C_j=\sum_{d\ge1}\frac{c_{j,d}}d,
 \label{eq:257-optimised-cover-cost}
\end{equation}
over finite or countable families of finite sets $F_j$ of positive integers covering
$F$, exponents $0<\alpha_j\le1$, weights $\eta_j>0$ with
$\sum_j\eta_j=1$, and coefficients $c_{j,d}\ge0$ satisfying
$f_{F_j}(n)^{\alpha_j}\le\sum_{d\mid n}c_{j,d}$ for every $n\ge1$.
The covering sets need not be disjoint or contained in $F$.  Covers of infinite
cost do not lower the infimum, and a finite-cost cover exists, for example
the single set $F$ with exponent one.  This is the manuscript's fractional cost,
not a redefinition of the logarithmic cost.

\begin{thm}[arithmetic logarithmic counterexample]\label{thm:257-logarithmic-counterexample}
For every integer $H\ge2$ and every real $A_0\ge0$, there are a
squarefree positive integer $L$ and a finite nonempty set $F$ of distinct
squarefree positive integers such that
\[
 \min F>\max\{L,A_0\},\qquad
 \kappa_1(F;1)\le\frac{30\log2}{H},\qquad
 \mathbb P_L(U_F>1)\ge1-e^{-1}.
\]
For these $F,L$, some $R\ge0$ satisfies
$\mathscr D_{L;R,L}\mathbf1_{\{U_F>1\}}>1/2$.
Consequently, no absolute constant $C$ satisfies
\begin{equation*}
 \mathscr D_{L;R,M}\mathbf1_{\{U_F>t\}}
 \le C(1+L/M)\kappa_1(F;t)
 \tag{E}
\end{equation*}
uniformly in finite $F$, positive integers $L,M$, integers $R\ge0$,
and $0<t\le1$.
\end{thm}
\leannote{Lean: \lproof{ErdosProblems/Erdos257/PaperCompleteR21/ArithmeticCounterexampleAssembly.lean}{1645}{arithmetic_logarithmic_counterexample}, \lproof{ErdosProblems/Erdos257/PaperCompleteR21/ArithmeticCounterexampleAssembly.lean}{1859}{no_absolute_dyadic_kappaOne_constant}.}

\begin{proof}
The modulus makes many small divisor conditions automatic.  Fresh tag
primes then make the remaining incidences independent, and taking all
least common multiples turns their number into an exponential incidence.
A logarithm pays only for the individual tags, while the shifted potential
can exceed one at many different offsets.

Choose a finite set of primes $P$ such that, with
\[
 L=\prod_{p\in P}p,\qquad D=\tau(L),\qquad
 B=\prod_{p\in P}(1+1/p),
\]
we have $B\ge H^2$.  Such a set exists because the sum of reciprocal primes
diverges and $\log(1+1/p)\ge1/(2p)$.
Here $\tau$ counts positive divisors and $\sigma$ denotes their sum.
Complementary divisors satisfy
\[
 \frac{\sigma(r)}r\,
 \frac{\sigma(L/r)}{L/r}=B\qquad(r\mid L).
\]
Thus the set
$G=\{r\mid L:\sigma(r)/r\ge H\}$ has at least $D/2$ elements.
Also $D\ge B\ge H^2$.

Choose a finite set $\mathcal B$ of parent primes, all exceeding
$\max\{L,H,D,A_0\}$, with
\[
 \frac4D\le\sum_{q\in\mathcal B}\frac1q\le\frac5D.
\]
Divergence of the prime reciprocal sum after any finite cutoff permits
stopping at the first crossing of $4/D$; the overshoot is less than $1/D$.
For each $q\in\mathcal B$ and $d\mid L$, choose a finite set $T_{q,d}$
of tag primes, all greater than $H$, disjoint from $P$, from $\mathcal B$,
and from every other tag set, such that
\[
 \frac{4d}{H}\le S_{q,d}:=\sum_{p\in T_{q,d}}\frac1p\le\frac{5d}{H}.
\]
The same stopping argument works successively: each new reciprocal is
less than $1/H\le d/H$, and only finitely many primes have been excluded.
No bound on the size of these prime sets is needed.

Let $\mathcal I_q=\{(d,p):d\mid L,\ p\in T_{q,d}\}$, and define
\[
 F_q=\left\{q\operatorname{lcm}\{dp:(d,p)\in I\}:
                         I\subseteq\mathcal I_q\right\},
 \qquad F=\bigcup_{q\in\mathcal B}F_q,
\]
where the least common multiple of the empty set is one.
Every tag prime occurs in just one labelled pair.  Its presence in the
least common multiple therefore recovers whether that pair was selected.
Hence different subsets give distinct exponents.  All exponents are
squarefree; different parent primes give disjoint sets, since no parent
prime occurs in any other set.  Their minimum is greater than
$\max\{L,A_0\}$.  The nonempty tag sets with $d=L$ also show $L\mid Q$.

For an integer $n\ge1$ put
$Z_q(n)=\sum_{(d,p)\in\mathcal I_q}\mathbf1_{\{dp\mid n\}}$.
The least-common-multiple encoding gives the exact identity
\[
 f_{F_q}(n)=\mathbf1_{\{q\mid n\}}2^{Z_q(n)}.
\]
Indeed, when $q\mid n$, a selected least common multiple divides $n$
exactly when every selected pair does.  Consequently
\[
 \log(1+f_{F_q}(n))
 \le(\log2)\left(\mathbf1_{\{q\mid n\}}
          +\sum_{(d,p)\in\mathcal I_q}\mathbf1_{\{qdp\mid n\}}\right).
\]
All moduli in this majorant divide $\operatorname{lcm}(F_q)$.
Its cost is at most
\begin{equation}
 \frac{\log2}{q}\left(1+\sum_{d\mid L}\frac{S_{q,d}}d\right)
 \le\frac{\log2}{q}(1+5D/H).
 \label{eq:257-parent-log-cost}
\end{equation}

Now sample $N$ uniformly modulo $Q$, subject to $L\mid N$.
For $r\in G$, the event $q\mid N+r$ has probability $1/q$.
For a fixed parent these events are disjoint as $1\le r\le L<q$.
Conditional on one such event, the tag-prime events remain independent
by the Chinese remainder theorem.  Since $d\mid N+r$ is equivalent to
$d\mid r$, the mean of $Z_q(N+r)$ is
\[
 \mu_r=\sum_{d\mid r}S_{q,d}
 \ge\frac4H\sigma(r)\ge4r,
\]
and its variance is at most $\mu_r$.
Chebyshev's inequality gives
\[
 \mathbb P_L\bigl(Z_q(N+r)<r\mid q\mid N+r\bigr)
 \le\frac{\mu_r}{(\mu_r-r)^2}
 \le\frac4{9r}\le\frac49.
\]
Thus success has conditional probability at least $1/2$.
On success the $r$th term of $U_{F_q}(N)$ is at least
$2^{-r}2^r=1$.  Later positive terms, for example those supplied by the
exponent $q\in F_q$, make $U_{F_q}(N)>1$ strictly.
The event $E_q$ of success for some $r\in G$ consequently satisfies
\[
 \mathbb P_L(E_q)\ge\frac{|G|}{2q}\ge\frac{D}{4q}.
\]
Different parent groups depend on disjoint prime coordinates, so the
$E_q$ are independent under this conditional uniform measure.  It follows
that
\[
 \mathbb P_L(U_F>1)\ge
 1-\prod_{q\in\mathcal B}(1-\mathbb P_L(E_q))
 \ge1-\exp\left(-\frac D4\sum_{q\in\mathcal B}\frac1q\right)
 \ge1-e^{-1}.
\]
Since $f_F=\sum_qf_{F_q}$ and
$\log(1+\sum_qx_q)\le\sum_q\log(1+x_q)$ for $x_q\ge0$,
we may sum the positive majorants in~\eqref{eq:257-parent-log-cost}:
\[
 \kappa_1(F;1)\le\frac{5\log2}{D}(1+5D/H)
 \le\frac{30\log2}{H}.
\]

Since $L\mid Q$, write $P_0=Q/L$ and $p=\mathbb P_L(U_F>1)$.  The bounded periodic
sequence $\mathbf1_{\{U_F(Lm)>1\}}$ satisfies, for every integer $T\ge1$,
\[
 \left|\frac1T\sum_{m=1}^{T}\mathbf1_{\{U_F(Lm)>1\}}-p\right|
 \le\frac{P_0}{T}.
\]
Summing over the dyadic window gives
\[
 \left|\mathscr D_{L;R,M}\mathbf1_{\{U_F>1\}}-p\right|
 \le\frac{2P_0}{M2^R}.
\]
Take $M=L$ and $2^R\ge16P_0$.  The error is at most $1/8$, whereas
$p\ge1-e^{-1}>5/8$, so the mean exceeds $1/2$.
A putative bound (E) would instead give at most $60C\log2/H$.
Choosing an integer $H>120C\log2$ contradicts it.
\end{proof}

\subsection{The arithmetic lower bound for every fractional cover}
\begin{cor}[finite-functional separation]\label{cor:257-logarithmic-separation}
For every finite nonempty $F$, with $Q=\operatorname{lcm}(F)$,
\begin{equation}
 K_*(F)\ge\max_{\ell\mid Q}
   \mathbb P\bigl(U_F(N)>1\mid\ell\mid N\bigr),
 \label{eq:257-arithmetic-cover-lower}
\end{equation}
where $N$ is uniform modulo $Q$.
The supports of Theorem~\ref{thm:257-logarithmic-counterexample} satisfy
$K_*(F)\ge1-e^{-1}$ and $\kappa_1(F;1)\le30\log2/H$.
In particular, no absolute $C$ gives $K_*(F)\le C\kappa_1(F;1)$
for all finite nonempty $F$, even after all covering sets, exponents, coefficients
and positive weights have been optimised.
\end{cor}
\leannote{Lean: \lproof{ErdosProblems/Erdos257/PaperCompleteR21/ArithmeticCoverLowerBoundPaperForm.lean}{40}{sup_condExceedProb_le_paperCoverCost}, \lproof{ErdosProblems/Erdos257/PaperCompleteR21/ArithmeticCounterexampleAssembly.lean}{1900}{exists_support_paperCoverCost_ge}, \lproof{ErdosProblems/Erdos257/PaperCompleteR21/ArithmeticCounterexampleAssembly.lean}{1912}{no_absolute_paperCoverCost_kappaOne_constant}.}
\begin{proof}
Fix one finite-cost cover in~\eqref{eq:257-optimised-cover-cost}, and write
$U_j=U_{F_j}$ and $B_j=2^{\alpha_j}$.
Coverage gives $U_F\le\sum_jU_j$.  If $U_F>1$, then $U_j>\eta_j$
for at least one index.  Subadditivity and the divisor majorants therefore
give the pointwise estimate
\[
 \mathbf1_{\{U_F>1\}}
 \le\sum_j\eta_j^{-\alpha_j}U_j^{\alpha_j}
 \le\sum_{j,d}\eta_j^{-\alpha_j}c_{j,d}w_{B_j,d}.
\]
Nonnegative interchange and the original geometric-kernel estimate
\eqref{eq:257-mixed-finite-kernel}, with its factor $1+4L/M$, imply
\[
 \mathscr D_{L;R,M}\mathbf1_{\{U_F>1\}}\le(1+4L/M)K.
\]
For fixed $L,M$, let $R$ tend to infinity.  Periodicity of the left-hand
indicator gives $\mathbb P_L(U_F>1)\le(1+4L/M)K$.
Then let $M$ tend to infinity, with $L$ fixed, to obtain
$\mathbb P_L(U_F>1)\le K$.
For $L=\ell\mid Q$ this is the conditional probability in the statement.
It holds for every admissible cover, so take their infimum and the maximum
over the finitely many divisors $\ell$.  Finally use the constructed
$L\mid Q$ and let $H$ tend to infinity.
\end{proof}

\subsection{What survives without arithmetic restriction}
\begin{prop}[ordinary initial intervals]\label{prop:257-logarithmic-initial-interval}
For every finite nonempty $F$, integer $X\ge1$, and $0<t\le1$,
\[
 \frac1X\#\{1\le N\le X:U_F(N)>t\}
 \le\frac{2}{\log(4/3)}\kappa_1(F;t).
\]
\end{prop}
\leannote{Lean: \lproof{ErdosProblems/Erdos257/PaperCompleteR21/LogarithmicInitialInterval.lean}{491}{logarithmic_initial_interval}.}
\begin{proof}
Choose a positive logarithmic divisor majorant
$g(n)=\sum_{d\mid n}c_d$ of cost $K=\sum_{d\mid Q}c_d/d$, putting
$c_d=0$ when $d\nmid Q$, and let $a=\log(4/3)$.
For positive integers $n,Y$,
\[
 g(n)\le Kn,\qquad
 \sum_{n=1}^{Y}g(n)=\sum_{d\mid Q}c_d\lfloor Y/d\rfloor\le KY.
\]
If $U_F(N)>t$, there is an $r\ge1$ with
$\log(1+f_F(N+r)/t)>ar$.
Otherwise every $f_F(N+r)/t\le(4/3)^r-1$, and
\[
 \frac{U_F(N)}t\le\sum_{r\ge1}2^{-r}\bigl((4/3)^r-1\bigr)=1,
\]
a contradiction.
If $K\ge a/2$ the required counting bound is trivial.
For $K<a/2$, a witnessing integer $n=N+r$ satisfies
$a(n-N)<g(n)\le Kn$, whence $n<aN/(a-K)\le2X$ when $N\le X$.
A fixed $n$ can witness at most $g(n)/a$ integer values of
$N$, since $1\le n-N<g(n)/a$.
The union bound thus gives
\[
 \#\{1\le N\le X:U_F(N)>t\}
 \le\frac1a\sum_{n=1}^{2X}g(n)\le\frac{2KX}{a}.
\]
Take the infimum over the admissible majorants.
\end{proof}

The proposition also bounds every $L=1$ dyadic observation window, since
such a window averages ordinary initial-interval means.  It does not bound
conditional arithmetic sampling uniformly in $L$.  Freezing a finite
initial support in the irrationality argument requires exactly that
arithmetic restriction.  The proposed logarithmic-tail irrationality
corollary therefore loses this proof; its conclusion is not disproved.
The corollary above separates finite functionals, not infinite-support
irrationality classes.  The weighted, fractional-cover, mixed and
finite-period noncollapse results are unchanged.  No conclusion here
decides universal \#257 or membership of $1/2$ or $1/21$.

\clearpage
\section*{Two proof details and source comparators}
The elementary atom comparison can be sharpened to
$\Delta_{b,A}(N)\le\Delta_{2,A}(N)$ for all real $b\ge2$.
After cancelling $b-1$, differentiate the ratio
$\sum_{i<r}b^i/\sum_{i<d}b^i$: its numerator is
$\sum_{i<r\le j<d}(i-j)b^{i+j-1}\le0$.
Only the rational-lattice step needs an integer base.

For the scalar positive-cover cost, put $z=\log_2t$.  Direct differentiation
of $y^z/(y-1)$ on $1<y\le2$ gives
\[
 \Psi(t)=t\quad(1\le t\le4),\qquad
 \Psi(t)=\frac{z^z}{(z-1)^{z-1}}\quad(t>4).
\]
These are ordinary proof details, not additional replayed declarations.

The original periodic and linear-independence papers of Luca and Tachiya
are now distinguished from their later overview
\cite{lucatachiya2014periodic,lucatachiya2014independence,lucatachiya2017}.
Their theorem formulations used in this revision were read in the
supplied author-written overview; the two original publisher PDFs were
not successfully retrieved in this revision.  Van Assche's nonvanishing and
small-linear-form argument for the full Lambert series
\cite[Lemma~1 and (34)--(35)]{vanassche2001} is a comparator for the missing
upper-height control, not a result for arbitrary thinnings.

For achievement-set background, Hornich's original bibliographic record
\cite{hornich1941} is separated from the exposition actually used here
\cite{nitecki2013}.  G\l\k{a}b and Prus-Wi\'sniowski's survey
\cite{glabprus2025} provides a modern map of geometry and cardinality
questions, not a rational-membership theorem for the present set.

\section*{Acknowledgements}
I thank Wouter van Doorn for advice on mathematical exposition, especially
on making restrictive hypotheses intelligible, avoiding unnecessary
notation, and writing for a first-time reader.  His comments concerned a
note on Problem~243; this acknowledgement does not imply review of the
mathematics of the present paper.

\section*{Statements and declarations}

\paragraph{Artefact and data availability.} The linked Lean declarations,
fixed toolchain, and library manifest are published in the companion
repository. This manuscript provides navigation and exposition rather than
proof authority.

\paragraph{Authority boundary.} Kernel checking establishes that a proposition
was proved; it does not authorise the exposition, literature judgements, or any
claim that Problem 257 is solved.

\paragraph{Funding and competing interests.} This work received no external
funding. The author declares no competing interests.

\begin{thebibliography}{99}
\small
\raggedright
\setlength{\itemsep}{0pt}
\bibitem{erdos1948} P.~Erd\H{o}s,
  \href{https://www.renyi.hu/~p_erdos/1948-04.pdf}{\emph{On arithmetical
  properties of Lambert series}}, J.~Indian Math.\ Soc.\ 12 (1948), 63--66.
\bibitem{erdos1968} P.~Erd\H{o}s,
  \href{https://users.renyi.hu/~p_erdos/1969-09.pdf}{\emph{On the
  irrationality of certain series}}, Math.\ Student 36 (1968), 222--226
  (issued 1969); \href{https://www.renyi.hu/~p_erdos/1969-09.pdf}{five-page scan}.  The theorem on p.~222 treats pairwise-coprime support with
  convergent reciprocal sum at every integer base $b\ge2$; the claimed
  removal of pairwise coprimality is stated without proof.  The same page
  says the reciprocal-sum condition could be replaced by a weaker but more
  complicated condition, and p.~226 suggests
  $\sum_{n_i<x}1/n_i=o(\log\log x)$ for pairwise coprime supports.
\bibitem{lucatachiya2017} F.~Luca and Y.~Tachiya,
  \href{https://www.kurims.kyoto-u.ac.jp/~kyodo/kokyuroku/contents/pdf/2014-14.pdf}
  {\emph{Linear independence results for the values of divisor functions
  series}}, RIMS K\^oky\^uroku No.~2014 (2017), 138--150.  Theorem~A (p.~139)
  restates Theorem~1.1 of their paper \emph{Irrationality of Lambert series
  associated with a periodic sequence}, Int.\ J.~Number Theory 10 (2014),
  no.~3, 623--636,
  \href{https://doi.org/10.1142/S1793042113501121}{DOI}: for a purely periodic
  integer sequence $a(n)$, not identically zero,
  $\sum_{n\ge1}a(n)/(q^n-1)$ is irrational for every integer $q$ with $|q|>1$.
  Example~2 (p.~140) treats the odd support.
\bibitem{nitecki2013} Z.~Nitecki,
  \href{https://arxiv.org/abs/1106.3779v2}{\emph{Subsum sets: intervals, Cantor
  sets, and Cantorvals}}, arXiv:1106.3779v2 (2013).  Theorem~4 (p.~9) shows
  that when every term exceeds its tail the subsum set is a Cantor set of
  Lebesgue measure $\lim_n2^nX_n$, where $X_n$ is the $n$th tail, and credits
  this to H.~Hornich (1941).
\bibitem{kanekosuzukitachiya} H.~Kaneko, Y.~Suzuki, and Y.~Tachiya,
  \href{https://arxiv.org/abs/2601.20743v1}{\emph{Refinements of Erd\H{o}s's
  irrationality criterion for certain sparse infinite series}},
  arXiv:2601.20743v1 (2026).  The averaged tail $R_c(q,x,z)$, a sum over
  offsets $j\ge z$, is~(1.7) and Theorem~1 is on p.~3; its rational-integer
  form, Theorem~3, is on p.~5.  Lemma~1 (pp.~6--7) derives irrationality from
  arbitrarily small nonzero scaled tails, and Lemma~2 (pp.~7--8) shows that
  most scaled tails in a range are small under a counting condition on the
  support.  Their criteria assume sparse coefficient supports.
\bibitem{duverneytachiya} D.~Duverney and Y.~Tachiya,
  \href{https://danielduverney.fr/documents/theorie-des-nombres/DuverneyTachiya190522.pdf}
  {\emph{Refinement of the Chowla--Erd\H{o}s method and linear independence
  of certain Lambert series}},
  Forum Math.\ 31 (2019), no.~6, 1557--1566,
  \href{https://doi.org/10.1515/forum-2018-0299}{DOI}.  Page numbers refer to
  the linked author preprint.  Theorem~1.1 (pp.~2--3) is the refined
  Chowla--Erd\H{o}s criterion; its proof in Section~2 (pp.~5--7) selects an
  index by minimising a local coefficient mass along an arithmetic
  progression, (2.3)--(2.9).  Corollary~1.2 gives the general $F_s(E)$
  theorem and its monomial images under $|q|\operatorname{lcm}(1,\dots,\ell)\le s$, and at
  every integer base when $s=\infty$; Example~1.1 gives the joint squarefree
  family at all bases \(2^j\), \(j\ge1\).  Both are on p.~4, which also
  relates Corollary~1.2 to the conjecture of Erd\H{o}s and Graham; the proof
  of Corollary~1.2 is on pp.~10--11.
\bibitem{taoteravainen2025} T.~Tao and J.~Ter\"av\"ainen,
  \href{https://arxiv.org/abs/2512.01739v2}{\emph{Quantitative correlations and
  some problems on prime factors of consecutive integers}}, arXiv:2512.01739v2
  (submitted December 2025, revised April 2026).  Theorem~1.3 (p.~4) proves
  $\sum_{n\ge1}\omega(n)/2^n=\sum_p(2^p-1)^{-1}$ irrational, settling the
  prime-support case of \#257 at base~$2$; the paragraph after it states that
  the prime-support result extends to every integer base $b\ge2$, and
  separately that full prime-power support is irrational at base $2$.
  The modifications for those two assertions are left to the reader;
  arbitrary subsets of the primes or prime powers are not treated there.  Theorem~3.1 is on p.~24, and the
  proof of Theorem~1.3 is Section~5, pp.~44--56.
\bibitem{zudilin2004} W.~Zudilin, \emph{Heine's basic transform and a
  permutation group for $q$-harmonic series}, Acta Arith.\ 111 (2004), no.~2,
  153--164, \href{https://doi.org/10.4064/aa111-2-4}{DOI}.  Theorem~1 (p.~154)
  gives $\mu(h_p(1))\le2.46497868\ldots$ for $h_p(1)=\sum_{n\ge1}1/(p^n-1)$,
  uniformly in the integer $p$ with $|p|\ge2$; the remark at the end of
  Section~3 (p.~159) corrects the computation of the author's earlier paper.
\bibitem{campbell2026} J.~M.~Campbell,
  \href{https://arxiv.org/abs/2605.24160v1}{\emph{On the binary digits of the
  Erd\H{o}s--Borwein constant}}, arXiv:2605.24160v1 (2026).  Theorem~1 (p.~12)
  proves that the block \texttt{11} occurs infinitely often in the binary
  expansion of $\sum_{n\ge1}1/(2^n-1)$, answering a question of Crandall
  recalled in the introduction.
\bibitem{matomakiradziwill} K.~Matom\"aki and M.~Radziwi\l\l,
  \emph{Multiplicative functions in short intervals}, Ann.\ of Math.\ (2) 183
  (2016), no.~3, 1015--1056,
  \href{https://doi.org/10.4007/annals.2016.183.3.6}{DOI}.  Page numbers
  refer to arXiv:1501.04585v4.
\bibitem{mangerel2021} A.~P.~Mangerel, \emph{Divisor-bounded multiplicative
  functions in short intervals}, Res.\ Math.\ Sci.\ 10 (2023), no.~1, Paper
  No.~12, \href{https://doi.org/10.1007/s40687-023-00376-0}{DOI}.  Page numbers
  refer to arXiv:2108.11401v2.
\bibitem{sun2026} Y.-C.~Sun,
  \href{https://arxiv.org/abs/2401.08432v3}{\emph{The critical-window profile
  for the $k$-fold divisor function in almost all short intervals}},
  arXiv:2401.08432v3 (2026).
\bibitem{lucatachiya2014periodic}
F.~Luca and Y.~Tachiya,
\emph{Irrationality of Lambert series associated with a periodic sequence},
International Journal of Number Theory \textbf{10} (2014), no.~3, 623--636.
\href{https://doi.org/10.1142/S1793042113501121}{doi:10.1142/S1793042113501121}.

\bibitem{lucatachiya2014independence}
F.~Luca and Y.~Tachiya,
\emph{Linear independence of certain Lambert series},
Proceedings of the American Mathematical Society \textbf{142} (2014), no.~10, 3411--3419.
\href{https://doi.org/10.1090/S0002-9939-2014-12102-2}{doi:10.1090/S0002-9939-2014-12102-2}.

\bibitem{vanassche2001}
W.~Van Assche,
\emph{Little $q$-Legendre polynomials and irrationality of certain Lambert series},
The Ramanujan Journal \textbf{5} (2001), 295--310.
\href{https://doi.org/10.1023/A:1012930828917}{doi:10.1023/A:1012930828917};
\href{https://arxiv.org/abs/math/0101187v1}{arXiv:math/0101187v1}.

\bibitem{hornich1941}
H.~Hornich,
\emph{\"Uber beliebige Teilsummen absolut konvergenter Reihen},
Monatshefte f\"ur Mathematik und Physik \textbf{49} (1941), 316--320.
\href{https://doi.org/10.1007/BF01707309}{doi:10.1007/BF01707309}.

\bibitem{glabprus2025}
S.~G\l\k{a}b and F.~Prus-Wi\'sniowski,
\emph{Achievement sets -- current results and open problems},
arXiv:2512.17285v1 (2025), preprint.
\url{https://arxiv.org/abs/2512.17285v1}.

\bibitem{kovactao} V.~Kova\v{c} and T.~Tao, \emph{On several irrationality
  problems for Ahmes series}, Acta Math.\ Hungar.\ 175 (2025), no.~2, 572--608,
  \href{https://doi.org/10.1007/s10474-025-01528-0}{DOI}.  Page numbers refer
  to arXiv:2406.17593v4.  Remark~4.1 (p.~13) records the strict tail
  inequality and the Cantor structure.  Theorem~2.3 (p.~5; proof pp.~13--14)
  constructs rational merged sums from several bases under its mass
  hypothesis; the construction merges several bases and leaves fixed-base \#257 untouched.
\bibitem{endpoint2026} \emph{The logarithmic endpoint fails under arithmetic
  sampling}, AI-assisted ordinary proof note, 17 September 2026.
  Theorem~1, Corollary~3 and Proposition~4.  Unpublished working note;
  independent review and fresh Lean verification are outstanding.
\bibitem{vandoornkovac} W.~van Doorn and V.~Kova\v{c},
  \emph{Lacunary sequences whose reciprocal sums represent all rational
  numbers in an interval}, Acta Arith.\ 223 (2026), 275--295,
  \href{https://doi.org/10.4064/aa251001-13-1}{DOI}.
  The theorem and page references used here refer to
  arXiv:2509.24971v3 (3 December 2025),
  \url{https://arxiv.org/abs/2509.24971v3}.
  Theorem~1 is on p.~2; Lemma~4 and Corollary~5 are on pp.~5--7.

\bibitem{erdos1975} P.~Erd\H{o}s, \emph{Some problems and results on the
  irrationality of the sum of infinite series}, J.~Math.\ Sci.\ 10 (1975),
  1--7.

\bibitem{bkkkz2026} K.~Barreto, J.~Kang, S.-H.~Kim, V.~Kova\v{c}, and
  S.~Zhang, \href{https://arxiv.org/abs/2601.21442v3}{\emph{Irrationality of
  rapidly converging series: a problem of Erd\H{o}s and Graham}},
  arXiv:2601.21442v3 (2026), to appear in Bull.\ London Math.\ Soc.

\bibitem{conradcrt} K.~Conrad,
  \href{https://kconrad.math.uconn.edu/blurbs/ugradnumthy/crt.pdf}
  {\emph{The Chinese remainder theorem}}, expository notes,
  Theorem~3.1 (pp.~2--3).  The compatibility criterion for noncoprime
  moduli used above is obtained by the displayed prime-power reduction.

\bibitem{erdos1932bertrand} P.~Erd\H{o}s,
  \href{https://www.renyi.hu/~p_erdos/1932-01.pdf}
  {Beweis eines Satzes von Tschebyschef}, 1932, pp.~194--198,
  especially \S\S4--5.

\end{thebibliography}
